\documentclass[10pt, a4paper, oneside]{memoir}
%%%%%
\DisemulatePackage{tabularx} % This is needed as memoir, bidi (for bidirectional text e.g., Arabic/Persian), polyglossia, and hyperref have a conflict. This solution works.
\usepackage{tabularx}
\usepackage{xltabular} %For tabularx across multiple pages  
%%%%%

\usepackage[all]{nowidow}
\usepackage[no-math]{fontspec}
\usepackage[dvipsnames]{xcolor}


\usepackage[hang,flushmargin, bottom, multiple]{footmisc}
\setlength{\footnotesep}{\baselineskip}% Keeping the footnote mark and text properly aligned, especially for longer paragraph-style footnotes
\makeatletter
\renewcommand\@makefntext[1]{\leftskip=0em\hskip0em\@makefnmark#1}
\makeatother
\setlength{\footparindent}{2em}

%Wrapping footnote numbers around edited text for Persian critical edition
\makeatletter
\newcommand{\wrapfootnote}[1]{\stepcounter\@mpfn%
% marks in the text
\protected@xdef\@thefnmark{\thempfn}%
\@footnotemark$^-$#1$^-$\@footnotemark%
% marks at bottom
\protected@xdef\@thefnmark{\thempfn\thinspace---\thinspace\thempfn}%
\@footnotetext}
\makeatother


\usepackage{ragged2e}
\usepackage{marginnote}
\usepackage{rotating} % For sideways figures
\usepackage{subcaption} % For subfigures
\usepackage{sidecap} % For captions on the side of figures
\usepackage{cancel}% For using \cancel{} command in the critical apparatus
\usepackage[shortlabels]{enumitem}
\usepackage{xpatch}
\usepackage{array, multicol, paracol, multirow} % For multiple columns to be merged in table environment. 
%paracol helps multicol-like layout across page breaks automatically without messing with the text across columns. 
\newcolumntype{C}[1]{>{\centering\arraybackslash}X}

\usepackage{booktabs}
% \usepackage{amsmath}
\usepackage{amscd, amsmath, amssymb, amsthm, yhmath, MnSymbol} % American Mathematical Society packages for metamathematical symbols.
% \usepackage{xfrac} % For \sfrac in math mode (in line with text))
\usepackage{marvosym, wasysym} % For Astronomical symbols
% \allowdisplaybreaks % for automatically splitting equations in align environment across pages
\usepackage{fourier-orns} % Decorative florets
% \newcommand{\decolinebreak}{\plainbreak{1}\fancybreak{\floweroneleft\floweroneright}\plainbreak{1}}
\usepackage{makecell} % For using \makecell{ \\ } in the tabular environment to case line breaks in a cell
\usepackage{wrapfig} % For wrapping figure
\usepackage{pst-3dplot, multido, pst-intersect, pst-tools, pst-poly} % PSTricks images


%%%%%%%%%%%%%%%%%%%%%%%%%%%%%%%%%%%%%%%
\usepackage{polyglossia}
\defaultfontfeatures{Ligatures=TeX}
\setmainfont[Ligatures={Common}, Numbers=OldStyle]{Linux Libertine O}
\setdefaultlanguage{english}
\setotherlanguage{sanskrit}
\setotherlanguage{hindi}
\setotherlanguage{latin}

\newfontfamily\critfont{Linux Libertine O}
\newfontfamily\sanskritfont[Script=Devanagari, Scale=1]{Shobhika-Regular}
\newfontfamily\latinfont[Script=Latin, Scale=1]{Linux Libertine O}
\newfontfamily\greekfont[Script=Greek, Mapping = tex-text]{Linux Libertine O}
\newfontfamily\metersymbol[Path=Fonts/]{new_athena_unicode.ttf}% Font supporting metrical short and long symbols for scansion (UNICODE), and dotted circle for diacritics
%~~~~~~~~~~~~~~~~~~~~~~~~~~~~~~~~~~~~~~~~~~~~~~~~~~~~~~~~~~~~~~~%
\newfontfamily\hindifont[Script=Devanagari,Mapping=velthuis-sanskrit,Scale=1]{Shobhika-Regular} % for typing (using velthuis) Sanskrit text in critical apparatus 
%~~~~~~~~~~~~~~~~~~~~~~~~~~~~~~~~~~~~~~~~~~~~~~~~~~~~~~~~~~~~~~~%
% Font declaration for displaying only Sanskrit vowel-signs (matras), especially for critical apparatus when certain vowel signs need to be crossed out without affecting the consonant. The Shobhika font (along with a direct Unicode declaration of the ligature and ``Script=Devanagari'' turned OFF) helps produce the dependent vowel-mark WITHOUT the dotted circle.
\newfontfamily\sanskritfontvowel [Mapping=velthuis-sanskrit, Scale=1]{Shobhika-Regular}
\newcommand{\tsv}[1]{\sanskritfontvowel{#1}} 
%~~~~~~~~~~~~~~~~~~~~~~~~~~~~~~~~~~~~~~~~~~~~~~~~~~~~~~~~~~~~~~~%
%Alternative formulation for when using Velthuis scheme for Sanskrit Devanagari within the thesis, except for Critical edition--   %
% Make the tilde into a normal letter of the alphabet %
\def\maketildeletter{\catcode`\~=11}
\def\dev{\maketildeletter \texthindi}
% The usage for this prescription is "{{\dev{j~n\={a}na}}" --  notice the double braces and "\dev{}'' taking the argument to be typeset
%NOTE: This does not seem to work for the lemma or variant. Within the \Lemma{} or \Variant{} environments, just type in devanagari in \tsans{}. e.g, \Lemma{@\tsans{नाञ्जा}@}
%~~~~~~~~~~~~~~~~~~~~~~~~~~~~~~~~~~~~~~~~~~~~~~~~~~~~~~~~~~~~~~~%
\newcommand{\ts}[1]{\textsanskrit{\bfseries #1}} % Sanskrit in Devanagari (in bold)
\newcommand{\tsnb}[1]{\textsanskrit{#1}} % Sanskrit in Devanagari (not bold)
\newcommand{\tms}[1]{{\small \textsanskrit{#1}}} % Sanskrit in Devanagari in the margin
\newcommand{\rs}[1]{\textit{\textlatin{#1}}} % Sanskrit (transliterated) in Latin
\newcommand{\rms}[1]{\small \textit{\textenglish{#1}}} % Sanskrit (transliterated) in Latin in the margin
\newcommand{\msf}[1]{\metersymbol{#1}} % For short (23D1), long (2013), and short-over-long (23D3) metrical symbols, and dotted circle for diacritics
\newcommand{\te}[1]{\textenglish{#1}} % Text in English
\newcommand{\tsans}[1]{\texthindi{#1}} % Text in Ordinary Sanskrit (typed in velthuis)



%%%%%%%%%%%%%%%%%%%%%%%%%%%%%%%%%%%%%%%
% Sanskrit and Persian words
% Sanskrit authors and word (unitalicised)

\newcommand{\Nityananda}{Nityā\-nanda}
\newcommand{\Nityanandafull}{Nityā\-nanda Miśra}
\newcommand{\Devadatta}{Deva\-datta}
\newcommand{\Laksmana}{Lakṣmaṇa}
\newcommand{\IcchaDulinahatta}{Icchā Ḍulīnahaṭṭa}
\newcommand{\Indrapuri}{Indrapurī}
\newcommand{\MahendraSuri}{Mahen\-dra Sūri}
\newcommand{\MalayenduSuri}{Malayendu Sūri}
\newcommand{\MallikarjunaSuri}{Malli\-kārjuna Sūri}
\newcommand{\Munisvara}{Munīśvara}
\newcommand{\Munisvarafull}{Munīśvara Viśva\-rūpa}
\newcommand{\Kamalakara}{Kamalā\-kara}
\newcommand{\KamalakaraBhatta}{Kamalākara Bhaṭṭa}
\newcommand{\Aryabhata}{Ārya\-bhaṭa}
\newcommand{\AryabhataII}{Ārya\-bhaṭa~II}
\newcommand{\Brahmagupta}{Brahma\-gupta}
\newcommand{\BhaskaraI}{Bhāskara~I}
\newcommand{\BhaskaraII}{Bhāskara~II}
\newcommand{\Bhaskaracarya}{Bhāskarācārya}
\newcommand{\Somesvara}{Some\-śvara}
\newcommand{\Paramesvara}{Para\-meśvara}
\newcommand{\RaghunathaRaja}{Raghunātha Rāja}
\newcommand{\Suryadeva}{Sūrya\-deva}
\newcommand{\Nilakantha}{Nīla\-kaṇṭha}
\newcommand{\Ranganatha}{Raṅganātha}
\newcommand{\Jagannatha}{Jagan\-nātha}
\newcommand{\Krsnadasa}{Bihārī Kṛṣṇa\-dāsa Miśra}
\newcommand{\VedangarayaMalajit}{Mālajit Vedāṅgarāya}
\newcommand{\Srimalaji}{Śrīmālajī}
\newcommand{\Jagannathafull}{Jagan\-nātha Samrāṭ}
\newcommand{\Kevlaramafull}{Kevala\-rāma Jyotiṣa\-rāya}
\newcommand{\Kevalarama}{Kevala\-rāma}
\newcommand{\Krsnadaivajna}{Kṛṣṇa Daivajña}
\newcommand{\Nayanasukhopadhyaya}{Nayana\-sukho\-pādhyāya}
\newcommand{\Nayanasukha}{Nayana\-sukha}
\newcommand{\Lalla}{Lalla}
\newcommand{\SudhakaraDvivedi}{Sudhākara Dvivedi}
\newcommand{\NrsimhaDaivajna}{Nṛsiṃha Daivajña}
\newcommand{\Sripati}{Śrī\-pati}
\newcommand{\Makkibhatta}{Makkibhaṭṭa}
\newcommand{\Madhava}{Mādhava}
\newcommand{\Sangamagrama}{Saṅgamagrāma}
\newcommand{\NilakanthaSomayaji}{Nīla\-kaṇṭha Somayājī}
\newcommand{\SankaraVariyar}{śaṅkara Vāriyar}
\newcommand{\Bhattotpala}{Bhaṭṭotpala}
\newcommand{\Vatesvara}{Vaṭeśvara}
\newcommand{\PutumanaSomayaji}{Putamana Soma\-yājī}
\newcommand{\AcyutaPisarati}{Acyuta Piṣāratī}
\newcommand{\Kashi}{Kāśī}
\newcommand{\Lanka}{Laṅkā}
\newcommand{\Panini}{Pāṇini}
\newcommand{\Siva}{śiva}
\newcommand{\Romaka}{Romaka}
\newcommand{\Romasa}{Romaśa}
\newcommand{\Romakacarya}{Romakācārya}
\newcommand{\Romasamuni}{Romaśa\-muni}
\newcommand{\Narada}{Nārada}
\newcommand{\Narayana}{Nārā\-yaṇa}
\newcommand{\Dhumraputra}{Dhūmra\-putra}
\newcommand{\Surya}{Sūrya}
\newcommand{\Brahma}{Brahmā}
\newcommand{\Caturveda}{Caturveda}
\newcommand{\Prthudakaswami}{Pṛthūdaka\-svāmī}
\newcommand{\Varahamihira}{Varāha\-mihira}
\newcommand{\Srisena}{Śrīṣeṇa}
\newcommand{\Latadeva}{Lāṭa\-deva}
\newcommand{\Jnanaraja}{Jñāna\-rāja}
\newcommand{\Vasistha}{Vasiṣṭha}
\newcommand{\Pulastya}{Pulastya}
\newcommand{\Garga}{Garga}
\newcommand{\Srisavayana}{śrīṣavāyana}
\newcommand{\Jisnuja}{Jiṣṇuja}
\newcommand{\Visrama}{Viśrāma}
\newcommand{\Jambusara}{Jambūsara}
\newcommand{\Karttika}{Kārttika}
\newcommand{\Samvat}{Saṃvat}
\newcommand{\Siddhanta}{Siddhānta}
\newcommand{\siddhanta}{siddhānta}
\newcommand{\Siddhantic}{Siddhāntic}
\newcommand{\siddhantic}{siddhāntic}
\newcommand{\Samhita}{Saṃhitā}
\newcommand{\Saurapaksa}{Saura\-pakṣa}
\newcommand{\Brahmana}{Brāhmaṇa}
\newcommand{\Brahmapaksa}{Brāhma\-pakṣa}
\newcommand{\Samarasimha}{Samara\-siṃha}
\newcommand{\Rama}{Rāma}
\newcommand{\Suryadasa}{Sūryadāsa}
\newcommand{\Pandita}{Paṇḍita}
\newcommand{\Nagari}{Nāgarī}
\newcommand{\pada}{pāda}
\newcommand{\JagannathaPanditaraja}{Jagan\-nātha Paṇḍita\-rāja}
\newcommand{\KavindracaryaSaravati}{Kavīndrā\-cārya Sara\-svatī Vidyā\-nidhāna}
\newcommand{\Kavindracarya}{Kavīndrā\-cārya}
\newcommand{\Balabhadra}{Bala\-bhadra}
% \newcommand{\ParamanandaJyotisaraya}{Paramā\-nanda Jyotiṣarāya}
\newcommand{\Jyotihsastra}{Jyotiḥ\-śāstra}

%Sanskrit texts and chapters

\newcommand{\spastakrantyadhikara}{\textit{spaṣṭa\-krāntyādhi\-kāra}}
\newcommand{\triprasnadhikara}{\textit{tripraśnādhi\-kāra}}
\newcommand{\mimamsadhikara}{\textit{mīmāṃsā\-dhikāra}}
\newcommand{\udayastadhikara}{\textit{udayāstādhi\-kāra}}
\newcommand{\Sarvasiddhantaraja}{\textit{Sarva\-siddhānta\-rāja}}
\newcommand{\Shahajahamganita}{\textit{Śāha\-jahāṃ\-gaṇita}}
\newcommand{\Romakasiddhanta}{\textit{Romaka\-siddhānta}}
\newcommand{\Romasasiddhanta}{\textit{Romaśa\-siddhānta}}
\newcommand{\Siddhantasindhu}{\textit{Siddhānta\-sindhu}}
\newcommand{\Siddhantaraja}{\textit{Siddhāntarāja}}
\newcommand{\Ganitadhyaya}{\textit{Gaṇitā\-dhyāya}}
\newcommand{\ganitadhyaya}{\textit{gaṇitādhyāya}}
\newcommand{\goladhyaya}{\textit{golā\-dhyāya}}
\newcommand{\Goladhyaya}{\textit{Golā\-dhyāya}}
\newcommand{\Bijaganita}{\textit{Bīja\-gaṇita}}
\newcommand{\Lilavati}{\textit{Līlāvatī}}
\newcommand{\golapada}{\textit{golapāda}}
\newcommand{\yantradhyaya}{\textit{yantrā\-dhyāya}}
\newcommand{\Yantraraja}{\textit{Yantrarāja}}
\newcommand{\Yantrasiromani}{\textit{Yantra\-śiromaṇi}}
\newcommand{\Siddhantatattvaviveka}{\textit{Siddhānta\-tattva\-viveka}}
\newcommand{\Siddhantasarvabhauma}{\textit{Siddhānta\-sārva\-bhauma}}
\newcommand{\Samratsiddhanta}{\textit{Samrāṭ\-siddhānta}}
\newcommand{\Samratsiddhantakaustubha}{\textit{Samrāṭ\-siddhānta\-kaustubha}}
\newcommand{\Siddhantasarakaustubha}{\textit{Siddhānta\-sāra\-kaustubha}}
\newcommand{\Purvardha}{\textit{Pūrvārdha}}
\newcommand{\Asayaprakasini}{\textit{Āśaya\-prakāśinī}}
\newcommand{\Siddhantatattvartha}{\textit{Siddhānta\-tattvārtha}}
\newcommand{\bhagrahayuti}{\textit{bhagrahayuti}}
\newcommand{\Astadhyayi}{\textit{Aṣṭā\-dhyāyī}}
\newcommand{\Aryabhatiya}{\textit{Ārya\-bhaṭīya}}
\newcommand{\Aryabhatiyabhasya}{\textit{Ārya\-bhaṭīya\-bhāṣya}}
\newcommand{\Mahasiddhanta}{\textit{Mahā\-siddhānta}}
\newcommand{\Brahmasphutasiddhanta}{\textit{Brāhma\-sphuṭa\-siddhānta}}
\newcommand{\bhagrahayutyadhikara}{\textit{bhagrahayutyadhikāra}}
\newcommand{\grahayutyadhikara}{\textit{grahayutyadhikāra}}
\newcommand{\Suryasiddhanta}{\textit{Sūrya\-siddhānta}}
\newcommand{\Jatakapaddhati}{\textit{Jātaka\-paddhati}}
\newcommand{\spastadhikara}{\textit{spaṣṭādhikāra}}
\newcommand{\Mahabhaskariya}{\textit{Mahā\-bhāskarīya}}
\newcommand{\sastodhyaya}{\textit{ṣaṣṭo\avagraha dhyāya}}
\newcommand{\Laghubhaskariya}{\textit{Laghu\-bhāskarīya}}
\newcommand{\Siddhantasiromani}{\textit{Siddhānta\-śiromaṇi}}
\newcommand{\Grahaganitadhyaya}{\textit{Graha\-gaṇitā\-dhyāya}}
\newcommand{\Gola}{\textit{Gola}}
\newcommand{\udayastavasana}{\textit{udayāstavāsanā}}
\newcommand{\grahachayadhikara}{\textit{gra\-hacchāyādhikāra}}
\newcommand{\Gudarthaprakasika}{\textit{Gūḍārthaprakāśikā}}
\newcommand{\Sisyadhivrddhidatantra}{\textit{śiṣyadhī\-vṛddhida\-tantra}}
\newcommand{\candrasrngonnatyadhikara}{\textit{candra\-śṛṅgon\-natyādhi\-kāra}}
\newcommand{\Nutanatilaka}{\textit{Nūtana\-tilaka}}
\newcommand{\Varttika}{\textit{Vārttika}}
\newcommand{\Vasanabhasya}{\textit{Vāsanā\-bhāṣya}}
\newcommand{\Vasanavarttika}{\textit{Vāsanā\-vārttika}}
\newcommand{\Siddhantasekhara}{\textit{Siddhānta\-śekhara}}
\newcommand{\candradhyaya}{\textit{candrādhyāya}}
\newcommand{\Tantrasamgraha}{\textit{Tantra\-saṅgraha}}
\newcommand{\LaghuVivrti}{\textit{Laghu\-vivṛti}}
\newcommand{\Khandakhadyaka}{\textit{Khaṇḍa\-khādyaka}}
\newcommand{\Purvakhanda}{\textit{Pūrva\-khaṇḍa}}
\newcommand{\Vivrtibhasya}{\textit{Vivṛti\-bhāṣya}}
\newcommand{\Vatesvarasiddhanta}{\textit{Vaṭeśvara\-siddhānta}}
\newcommand{\udayastamayadhikara}{\textit{udayā\-sta\-mayādhi\-kāra}}
\newcommand{\Karanapaddhati}{\textit{Karaṇapaddhati}}
\newcommand{\Goladipika}{\textit{Goladīpikā}}
\newcommand{\Sphutanirnaya}{\textit{Sphuṭa\-nirṇaya}}
\newcommand{\Rasigolasphutaniti}{\textit{Rāśi\-gola\-sphuṭā\-nīti}}
\newcommand{\JicaUlugbegi}{\textit{Jīca Ulugbegī}}
\newcommand{\Ulakabegijica}{\textit{Ulakā\-begī\-jīca}}
\newcommand{\Pancasiddhantika}{\textit{Pañca\-siddhāntikā}}
\newcommand{\Siddhantasundara}{\textit{Siddhānta\-sundara}}
\newcommand{\bhuvanakosadhikara}{\textit{bhuvana\-kośādhikāra}}
\newcommand{\bhutasamkhya}{\textit{bhūta\-saṃkhyā}}
\newcommand{\Hayatagrantha}{\textit{Hayata\-grantha}}
\newcommand{\Rekhaganita}{\textit{Rekhā\-gaṇita}}
\newcommand{\Ukara}{\textit{Ukarā}}
\newcommand{\Sarahatajakiravirajandi}{\textit{Śarahatajkirā Virajandī}}
\newcommand{\Yantrarajasyarasala}{\textit{Yantrarājasya Rasāla}}
\newcommand{\Visavava}{\textit{Vīsavāva}}
\newcommand{\Yantrarajavicaravimsadhyayi}{\textit{Yantrarāja\-vicāra\-vimśādhyāyī}}
\newcommand{\Bijapallava}{\textit{Bīja\-pallava}}
\newcommand{\Karmaprakasa}{\textit{Karma\-prakāśa}}
\newcommand{\Parasiprakasa}{\textit{Pārasī\-prakāśa}}
\newcommand{\SamskrtaParasikaPadaPrakasa}{\textit{Saṃskṛta\-pārasīka\-padaprakāśa}}
\newcommand{\tajika}{\textit{tājika}}
\newcommand{\ramala}{\textit{ramala}}
\newcommand{\Ramalavaicitrya}{\textit{Ramalavaicitrya}}
\newcommand{\MlecchaMataNirupana}{\textit{Mleccha\-mata\-nirūpaṇa}}
\newcommand{\Siddhantasamhitasarasamuccaya}{\textit{Siddhānta\-saṃhitā\-sāra\-samuccaya}}
\newcommand{\TajikaNilakanthi}{\textit{Tājika\-nīla\-kaṇṭhī}}
% \newcommand{\Jahangirvondaratnakara}{\textit{Jahāṃgīr\-vinoda\-ratnā\-kara}}
\newcommand{\Phirangicandracchedyopayogika}{\textit{Phiraṅgi\-candra\-cchedyo\-payogika}}
\newcommand{\Drkpaksasarani}{\textit{Dṛk\-pakṣa\-sāraṇī}}
\newcommand{\Hayanaratna}{\textit{Hāyana\-ratna}}
\newcommand{\Giridharananda}{\textit{Giri\-dharānanda}}
\newcommand{\Khetakautuka}{Kheṭa\-kautuka}

%Islamicate authors and place names

\newcommand{\Maragha}{Marāgha}
\newcommand{\Samarqand}{Samarqand}
\newcommand{\Shahjahanabad}{Shāh\-jahānā\-bād}
\newcommand{\alBiruni}{al\=/Bīrūnī}
\newcommand{\alBirunifull}{Abu\Alif l\=/Rayḥān al\=/Bīrūnī}
\newcommand{\UlughBegfull}{Sulṭān Ulugh Beg}
\newcommand{\UlughBeg}{Ulugh Beg}
\newcommand{\alAmilifull}{Bahā\Alif\ l\=/Dīn al\=/\Ayn Āmilī}
\newcommand{\Babur}{Bābur}
\newcommand{\Baburfull}{Ẓahīr al\=/Dīn Muhammad Bābur}
\newcommand{\ShahIsmaili}{Shāh Ismā\Ayn īl\=/I}
\newcommand{\ShahbinIsmail}{Shāh Tahmāsp\-b.\@ Ismā\Ayn īl}
\newcommand{\Humayun}{Humāyūn}
\newcommand{\Akbar}{Akbar}
\newcommand{\Jahangir}{Jahāngīr}
\newcommand{\Shahjahan}{Shāh Jahān}
\newcommand{\Awrangzeb}{Awrangzīb \AynCap{Ā}lamgīr}
\newcommand{\ShahAbbasI}{Shāh \AynCap{A}bbās–I}
\newcommand{\ShahAbbasII}{Shāh \AynCap{A}bbās–II}
\newcommand{\KibnLabban}{Kūshyār b.\@ Labbān}
\newcommand{\alTusifull}{Naṣīr al\=/Dīn al\=/Ṭūsī}
\newcommand{\alTusi}{al\=/Ṭūsī}
\newcommand{\alKashi}{al\=/Kāshī}
\newcommand{\alKashifull}{Jamshīd Mas\Ayn ūd al\=/Kāshī}
\newcommand{\alKhwarizmi}{al\=/Khwārizmī}
\newcommand{\alKhwarizmifull}{Muḥammad b.\@ Mūsā al\=/Khwārizmī}
\newcommand{\IbnYunis}{Ibn~Yūnis}
\newcommand{\ibnYunisfull}{Abu\Alif l‐Ḥasan \AynCap{A}lī b.\@ Yūnus al‐Ṣadafī}
\newcommand{\Gilani}{Kūshyār Gīlānī}
\newcommand{\Gilanifull}{Kūshyār b.\@ Labbān al\=/Jīlī}
\newcommand{\abuNasr}{Abū Naṣr Manṣūr b.\@ \AynCap{A}lī b.\@ \Ayn Irāq}
\newcommand{\abuWafa}{Abu\Alif l\=/Wafā Būzhjānī}
\newcommand{\alHasib}{Ḥabash al\=/Ḥāsib}
\newcommand{\alMaghribi}{Muḥyī l\=/Dīn al\=/Maghribī}
\newcommand{\MullaFarid}{Mullā Farīd}
\newcommand{\MullaFaridfull}{Mullā Farīd al\=/Dīn Mas\Ayn ūd b.\@ Ḥāfiẓ Ibrāhīm Dihlavī}
\newcommand{\MullaTayyib}{Mullā Ṭayyib}
\newcommand{\AsifKhan}{Aṣif Khān}
\newcommand{\AsafKhan}{Abu\Alif l\=/Ḥasan Āṣaf Khān}
\newcommand{\AsafKhanshort}{Āṣaf Khān}
\newcommand{\Timur}{Sulṭān Amīr Tīmūr}
\newcommand{\ibmHaytham}{Ḥasan b.\@ al\=/Haytham}
\newcommand{\alShirazifull}{Qutb al\=/Dīn al\=/Shirāzī}
\newcommand{\alJaghmini}{Maḥmūd al\=/Jaghmīnī}
\newcommand{\alRumifull}{Qāḍī Zāda al\=/Rūmī}
\newcommand{\alRumi}{al\=/Rūmī}
\newcommand{\aliQushji}{\AynCap{A}lī Qushjī}
\newcommand{\aliQushjifull}{\AynCap{A}lā\Alif\ l\=/Dīn \AynCap{A}lī Qushjī}
\newcommand{\alBirjandi}{Niẓām al\=/Dīn al\=/Bīrjandī}
\newcommand{\AbulFadlAllami}{Abu\Alif l\=/Faḍl \AynCap{A}llāmī}
\newcommand{\ShahFathullahShirazi}{Shāh Fatḥallāh Shīrāzī}
\newcommand{\AmirFathullahofShiraz}{Amīr Fatḥallāh of Shīrāz}
\newcommand{\KhanKhana}{Mirzā \AynCap{A}bd\Alif l\=/Raḥīm Khān\=/i Khānān}
\newcommand{\KhanKhanashort}{Khān\=/i Khānān}
\newcommand{\Rahim}{Raḥīm}
\newcommand{\MalikShah}{Malik Shāh}
\newcommand{\InayatKhan}{\Ayn{I}nāyat Khān}
\newcommand{\TahirKhan}{Muḥammad Ṭāhīr Khān}
\newcommand{\FiruzShah}{Sulṭān Fīrūz Shāh Ṭughlāq}
\newcommand{\Nastaliq}{Nasta\Ayn līq}
\newcommand{\AbdalHamidLahori}{\AynCap{A}bd al\=/Hamīd Lāhūrī}
\newcommand{\MuhammadAbidda}{Muḥammad Ābidda}
\newcommand{\AbualFaydi}{Abu\Alif l\=/Fayḍ Fayḍī}
\newcommand{\AtaAllahRushdi}{\AynCap{A}ṭa\Alif\thinspace Allāh Rushdī}
\newcommand{\Badaunifull}{\AynCap{A}bd al\=/Qādir b.\@ Mulūk Shāh Badā\Alif ūnī}
\newcommand{\Badauni}{Badā\Alif ūnī}
\newcommand{\MirzaKhanbFakhralDinMuhammad}{Mīrzā Khān b.\@ Fakhr al-Dīn Muḥammad}
\newcommand{\Rumipoet}{Jalāl al\=/Dīn Rūmī}
\newcommand{\ShahShuja}{Shāh Shujā\Alif}
\newcommand{\AzamKhan}{\AynCap{Ā}ẓam Khān}
\newcommand{\VazirKhan}{Vazīr Khān}
\newcommand{\HakimShaykhIlmalDinAnsari}{Ḥakīm Shaykh \Ayn{I}lm al-Dīn Ansārī}
\newcommand{\AbdallahKhanFiruzJung}{\AynCap{A}bdallāh Khān Fīrūz Jung}
\newcommand{\Subadar}{Ṣūbadār}
\newcommand{\MahabatKhanKhaniKhanan}{Mahābat Khān Khān\=/i Khānān}
\newcommand{\ZafarKhanAhsan}{Ẓafar Khān Aḥsan}
\newcommand{\KhwajaSabirKhanDauran}{Khwāja Ṣābir Khān\=/i Daurān}
\newcommand{\NasiriKhan}{Naṣīrī Khān}
\newcommand{\ItiqadKhan}{I\Ayn tiqād Khān}
\newcommand{\MirzaShapur}{Mirzā Shāpūr}
\newcommand{\MehmetII}{Sulṭān Meḥmed~II}
\newcommand{\MuslihalDinMuhammadLari}{Muṣliḥ al\=/Dīn Muḥammad al\=/Lārī}
%Arabic and Persian works

\newcommand{\ZijNasiri}{\textit{Zīj\=/i Nāṣirī}}
\newcommand{\alZijalJami}{\textit{al\=/Zīj al\=/Jāmi}\Ayn}
\newcommand{\alZijalBaligh}{\textit{al\=/Zīj al\=/Bāligh}}
\newcommand{\HakimiZij}{\textit{al\=/Zīj al\=/Kabīr al\=/Ḥākimī}}
\newcommand{\ZijIlkhani}{\textit{Zīj\=/i Ilkhānī}}
\newcommand{\ZijSultani}{\textit{Zīj\=/i Jadīd\=/i Sulṭānī}}
\newcommand{\ZijMirzai}{\textit{Zīj\=/i Jadīd\=/i Mīrzā}\Ayn\textit{ī}}
\newcommand{\ZijUlughBeg}{\textit{Zīj\=/i Ulugh Beg}}
\newcommand{\KhaqaniZij}{\textit{Zīj al\=/Khāqānī}}
\newcommand{\KhaqaniZijfull}{\textit{Zīj al\=/Khāqānī fī Takmīl al\=/Zīj al\=/Ilkhānī}}
\newcommand{\Almagest}{\textit{Alma\-gest}}
\newcommand{\Elements}{\textit{Elements}}
\newcommand{\Sphaerica}{\textit{Sphae\-rica}}
\newcommand{\KitabalMijisti}{\textit{Kitāb al\=/Majisṭī)}}
\newcommand{\TahriralMijisti}{\textit{Taḥrīr al\=/Majisṭī}}
\newcommand{\RasaelalBiruni}{\textit{Rasā}\Ayn \textit{il}\Alif \textit{l\=/Bīrūnī}}
\newcommand{\JadalalTaqwim}{\textit{Jadwal al\=/Taqwīm}}
\newcommand{\Tajalazyaj}{\textit{Tāj al\=/Azyāj}}
\newcommand{\QanunalMasudi}{\textit{al\=/Qānūn al\=/Mas}\Ayn\textit{ūdī}}
\newcommand{\ZijiShahJahani}{\textit{Zīj\=/i Shāh Jahānī}}
\newcommand{\ZijiShahJahanifull}{\textit{Kārnamāh-i Ṣaḥib Qirān-i Thānī, Zīj\=/i Shāh Jahānī}}
\newcommand{\AiniAkbari}{\textit{Ā\Alif in\=/i Akbarī}}
\newcommand{\ZijJadidiGurani}{\textit{Zīj\=/i  Sa}\Ayn\textit{īd-i Jadīd-i Gūrgānī}}
\newcommand{\ZijiRahimi}{\textit{Zīj\=/i Raḥīmī}}
\newcommand{\zijiHisabi}{\textit{zīj\=/i ḥiṣābī}}
\newcommand{\zijiRasadi}{\textit{zīj\=/i raṣadī}}
\newcommand{\MulakhkhasiShahJahanNama}{\textit{Mulakhkhaṣ-i Shāhjahān Nāma}}
\newcommand{\Padshahnama}{\textit{Pādshāh\-nāma}}
\newcommand{\Risaladarilmalhaya}{\textit{Risāla dar} \AynCapI\textit{Ilm al\=/Hay}\Alif \textit{a}}
\newcommand{\KitabTahirUsulliUqlidus}{\textit{Kitāb Taḥrīr Uṣūl li\=/Uqlīdus}}
\newcommand{\Tahrialukarr}{\textit{Taḥrīr al\=/Ukarr}}
\newcommand{\SharhalTadhkirah}{\textit{Sharḥ al\=/Tadhkira}}
\newcommand{\alTadhkirafiilmalhaya}{\textit{al\=/Tadhkira fī }\AynCapI\textit{Ilm al\=/Hay}\Alif \textit{a}}
\newcommand{\RisaladarBistBabMarifatUsturlab}{\textit{Risāla\=/yi Bīst Bāb dar ma}\Ayn\textit{Rifat\=/i Usṭurlāb}}
% \newcommand{\KutubalMatawassitat}{\textit{Kutub al\=/Mutawassiṭāt}}
\newcommand{\MuntakhabalTawarikh}{\textit{Muntakhab al\=/Tawārīkh}}
\newcommand{\TarikhiBadani}{\textit{Tārīkh\-/i Badā}\Alif \textit{ūnī}}
\newcommand{\TuhfatalHind}{\textit{Tuḥfat al\=/Hind}}
\newcommand{\TarjumayiLilawati}{\textit{Tarjuma\=/yi Līlāwatī}}
\newcommand{\TarjumayiBijganit}{\textit{Tarjuma\=/yi Bījganit}}
\newcommand{\AkhlaqiNasiri}{\textit{Akhlāq\=/i Nāsirī}}
\newcommand{\MasnaviyiManavi}{\textit{Masnavī\=/yi Ma}\Ayn\textit{navī}}
\newcommand{\SharhilmalHaya}{\textit{Sharḥ\=/i Risāla dar} \AynCapI\textit{Ilm al\=/Hay}\Alif \textit{a}}
%Arabic and Persian words

\newcommand{\aan}{\Alif\textit{ān}}
\newcommand{\ajza}{\textit{ajzā}\Alif}
\newcommand{\aks}{\Ayn\textit{aks}}
\newcommand{\alawwal}{\textit{al\=/awwal}}
\newcommand{\alburuj}{\textit{al\=/burūj}}
\newcommand{\aljayb}{\textit{al\=/jayb}}
\newcommand{\aljuzi}{\textit{al\=/juz}\Alif \textit{ī}}
\newcommand{\alkabah}{\textit{al\=/ka}\Alif \textit{bah}}
\newcommand{\alislah}{\textit{al\=/iṣlāḥ}}
\newcommand{\allayl}{\textit{al\=/layl}}
\newcommand{\almail}{\textit{al\=/mā}\Ayn\textit{il}}
\newcommand{\almakus}{\textit{al\=/ma}\Ayn\textit{kūs}}
\newcommand{\almayl}{\textit{al\=/māyl}}
\newcommand{\alnahar}{\textit{al\=/nahār}}
\newcommand{\alqaws}{\textit{al\=/qaws}} 
\newcommand{\alraml}{\textit{al\=/raml}}
\newcommand{\alsatrayn}{\textit{al\=/saṭrayn}}
\newcommand{\althani}{\textit{al\=/thānī}}
\newcommand{\amal}{\Ayn\textit{amal}}
\newcommand{\aqrab}{\textit{aqrab}}
\newcommand{\aqsam}{\textit{aqsām}}
\newcommand{\arbai}{\textit{arba\Ayn}\textit{i}}
\newcommand{\arbaia}{\Alif\textit{arba}\Ayn\textit{a}}
\newcommand{\ard}{\Ayn \textit{arḍ}}
\newcommand{\avqat}{\textit{avqāt}}
\newcommand{\avval}{\textit{avval}}
\newcommand{\az}{\textit{az}}
\newcommand{\azam}{\textit{a}\Ayn \textit{ẓam}}
\newcommand{\bab}{\textit{bāb}}
\newcommand{\balad}{\textit{balad}}
\newcommand{\bayn}{\textit{bayn}}
\newcommand{\biaqtab}{\textit{bi aqṭāb}}
\newcommand{\biban}{\textit{bībān}}
\newcommand{\bud}{\textit{bu}\Ayn\textit{d}}
\newcommand{\daneshmand}{\textit{dāneshmand}}
\newcommand{\daqaiq}{\textit{daqā}\Alif \textit{iq}}
\newcommand{\daqiqa}{\textit{daqīqa}}
\newcommand{\dar}{\textit{dar}}
\newcommand{\daraj}{\textit{daraj}}
\newcommand{\daraji}{\textit{daraji}}
\newcommand{\daravardan}{\textit{dar\thinspace āvardan}}
\newcommand{\darardan}{\textit{dar\thinspace ārdan}}
\newcommand{\darb}{\textit{ḍarb}}
\newcommand{\dawr}{\textit{dawr}}
\newcommand{\dayiri}{\textit{dāyiri}}
\newcommand{\duvum}{\textit{duvum}}
\newcommand{\fadla}{\textit{faḍla}}
\newcommand{\falak}{\textit{falak}}
\newcommand{\farman}{\textit{farmān}}
\newcommand{\fasl}{\textit{faṣl}}
\newcommand{\fatammalha}{\textit{fata}\Alif\textit{mmalhā}}
\newcommand{\fudala}{\textit{fuḍalā}\Alif}
\newcommand{\fusul}{\textit{fuṣūl}}
\newcommand{\ghurub}{\textit{ghurūb}}
\newcommand{\har}{\textit{har}}
\newcommand{\hasil}{\textit{ḥāṣil}}
\newcommand{\haya}{\textit{hay}\Alif \textit{a}}
\newcommand{\hisabaljummal}{\textit{ḥisāb al\=/jummal}}
\newcommand{\hissi}{\textit{ḥiṣṣi}}
\newcommand{\hissat}{\textit{ḥiṣṣat}}
\newcommand{\hissatalbud}{\textit{ḥiṣṣat al\=/bu}\Ayn \textit{d}}
\newcommand{\hissayiard}{\textit{ḥiṣṣa-yi} \Ayn \textit{arḍ}}
\newcommand{\ghayat}{\textit{ghāyat}}
\newcommand{\GurkaniAlam}{\textit{Gūrkāni Ālam}}
\newcommand{\ijam}{\textit{i}\Alif \textit{jām}}
\newcommand{\Ilahi}{\textit{Ta}\Alif \textit{rīkh-i Ilāhī}}
\newcommand{\ilmalhaya}{\Ayn \textit{ilm al\=/hay}\Alif \textit{a}}
\newcommand{\inhiraf}{\textit{inḥirāf}}
\newcommand{\inhirafalqibla}{\textit{inḥirāf al\=/qibla}}
\newcommand{\inkhifad}{\textit{inkhifāḍ}}
\newcommand{\inqilab}{\textit{inqilāb}}
\newcommand{\intiha}{\Alif\textit{intihā}}
\newcommand{\iqlim}{\textit{iqlīm}}
\newcommand{\irtifa}{\Alif\textit{irtifā}\Ayn}
\newcommand{\istikhraj}{\textit{istikhrāj}}
\newcommand{\istiva}{\textit{istiva}}
\newcommand{\itidal}{\textit{i}\Ayn \textit{tidāl}}
\newcommand{\jadr}{\textit{jaḏr}}
\newcommand{\jadval}{\textit{jadval}}
\newcommand{\jahat}{\textit{jahat}}
\newcommand{\Jalali}{\textit{Jalālī}}
\newcommand{\jam}{\textit{jam}\Ayn}
\newcommand{\jayb}{\textit{jayb}}
\newcommand{\jins}{\textit{jins}}
\newcommand{\kardan}{\textit{kardan}}
\newcommand{\kawkab}{\textit{kawkab}}
\newcommand{\kawakib}{\textit{kawākib}}
\newcommand{\kharij}{\textit{khārij}}
\newcommand{\khasmohor}{\textit{khās muhr}}
\newcommand{\kitabkhana}{\textit{kitāb-khāna}}
\newcommand{\khatt}{\textit{khaṭṭ}}
\newcommand{\kulli}{\textit{kullī}}
\newcommand{\ma}{\textit{mā}}
\newcommand{\madrasa}{\textit{madrasa}}
\newcommand{\majmu}{\textit{majmū}\Ayn}
\newcommand{\makkah}{\textit{makkah}}
\newcommand{\maktabkhana}{\textit{maktab\-khāna}}
\newcommand{\mamarr}{\textit{mamarr}}
\newcommand{\mankus}{\textit{mankūs}}
\newcommand{\maqala}{\textit{maqāla}}
\newcommand{\maqalat}{\textit{maqā\-lāt}}
\newcommand{\maratib}{\textit{marātib}}
\newcommand{\marfu}{\textit{marfū}\Ayn}
\newcommand{\marfuat}{\textit{marfū}\Ayn\textit{āt}}
\newcommand{\marifat}{\textit{ma\Ayn rifat}}
\newcommand{\marri}{\textit{mārri}}
\newcommand{\martaba}{\textit{martaba}}
\newcommand{\matla}{\textit{maṭla}\Ayn}
\newcommand{\matali}{\textit{maṭāli}\Ayn}
\newcommand{\mayl}{\textit{mayl}}
\newcommand{\miyan}{\textit{miyān}}
\newcommand{\muabbiran}{\textit{mu}\Ayn\textit{abbirān}}
\newcommand{\muarrab}{\textit{mu}\Ayn\textit{arrab}}
\newcommand{\mukhtalif}{\textit{mukhtalif}}
\newcommand{\muharrir}{\textit{muḥarrir}}
\newcommand{\munajjim}{\textit{munajjim}}
\newcommand{\munhatt}{\textit{munḥaṭṭ}}
\newcommand{\munshi}{\textit{munshi}}
\newcommand{\muaddil}{\textit{mu}\Ayn \textit{addil}}
\newcommand{\muqaddima}{\textit{muqaddima}}
\newcommand{\mustaqim}{\textit{mustaqīm}}
\newcommand{\mutarjiman}{\textit{mutarjimān}}
\newcommand{\mutamarris}{\textit{muta\-marris}}
\newcommand{\nisf}{\textit{niṣf}}
\newcommand{\panjum}{\textit{panjum}}
\newcommand{\qibla}{\textit{qibla}}
\newcommand{\qif}{\textit{qīf}}
\newcommand{\qism}{\textit{qism}}
\newcommand{\qismat}{\textit{qismat}}
\newcommand{\qiran}{\textit{qirān}}
\newcommand{\qaws}{\textit{qaws}}
\newcommand{\rabi}{\textit{rābi}\Ayn}
\newcommand{\rasm}{\textit{rasm}}
\newcommand{\rub}{\textit{rub}\Ayn}
\newcommand{\ruyat}{\textit{ru}\Alif \textit{yat}}
\newcommand{\sad}{\textit{ṣād}}
\newcommand{\saat}{\textit{sā}\Ayn \textit{āt}}
\newcommand{\sahm}{\textit{sahm}}
\newcommand{\samt}{\textit{samt}}
\newcommand{\Shahishani}{\textit{Tā}\Alif \textit{rīkh\=/i Ilāhī Shāhishānī}}
\newcommand{\shudan}{\textit{shudan}}
\newcommand{\tadil}{\textit{ta}\Ayn \textit{dīl}}
\newcommand{\tafadul}{\textit{tafāḍul}}
\newcommand{\tali}{\textit{ṭāli}\Ayn}
\newcommand{\taqvim}{\textit{taqvīm}}
\newcommand{\tamam}{\textit{tamām}}
\newcommand{\tashil}{\textit{tashīl}}
\newcommand{\tashkil}{\textit{tashkīl}}
\newcommand{\tawali}{\textit{ṭawali}\Ayn}
\newcommand{\tazig}{\textit{tāzīg}}
\newcommand{\thani}{\textit{thānī}}
\newcommand{\thaniya}{\textit{thāniya}}
\newcommand{\thawani}{\textit{thawānī}}
\newcommand{\tul}{\textit{ṭūl}}
\newcommand{\ulama}{\Ayn \textit{ulamā}\Alif}
\newcommand{\Umara}{\Ayn\textit{Umarā}\Alif}
\newcommand{\ulumalaqliyah}{\Ayn \textit{ulūm} \mbox{\textit{al\=/}\Ayn \textit{aqlīyah}}}
\newcommand{\UlumiAqliyah}{\Ayn \textit{Ulūm}\=/i \Ayn \textit{Aqlīyah}}
\newcommand{\ustadh}{\textit{ustādh}}
\newcommand{\uy}{\textit{u}}
\newcommand{\va}{\textit{va}}
\newcommand{\vaqt}{\textit{vaqt}}
\newcommand{\vazir}{\textit{vazīr}}
\newcommand{\ya}{\textit{yā}}
\newcommand{\yik}{\textit{yik}}
\newcommand{\zabanifarsi}{\textit{zabān\=/i fārsī}}
\newcommand{\zabanihind}{\textit{zabān\=/i hindī}}
\newcommand{\ziyadi}{\textit{ziyādi}}
\newcommand{\zij}{\textit{zīj}}
\newcommand{\Zij}{\textit{Zīj}}
\newcommand{\zijat}{\textit{zījāt}}
\newcommand{\zill}{\textit{ẓill}}


%Sanskrit words

\newcommand{\aadhara}{\textit{ādhara}}
\newcommand{\aayana}{\textit{āyana}}
\newcommand{\acarya}{\textit{ācārya}}
\newcommand{\adhikara}{\textit{adhikāra}}
\newcommand{\adhyaya}{\textit{adhyāya}}
\newcommand{\adyaa}{\textit{ādyā}}
\newcommand{\amsaka}{\textit{aṃśaka}}
\newcommand{\anka}{\textit{aṅka}}
\newcommand{\anga}{\textit{aṅga}}
\newcommand{\antara}{\textit{antara}}
\newcommand{\anya}{\textit{anya}}
\newcommand{\anayana}{\textit{ānayana}}
\newcommand{\anumana}{\textit{anumāna}}
\newcommand{\anyatama}{\textit{anyatama}}
\newcommand{\anyatara}{\textit{anyatara}}
\newcommand{\apama}{\textit{apama}}
\newcommand{\apakrama}{\textit{apakrama}}
\newcommand{\apara}{\textit{apara}}
\newcommand{\asus}{\textit{asus}}
\newcommand{\ayana}{\textit{ayana}}
\newcommand{\bahu}{\textit{bāhu}}
\newcommand{\bana}{\textit{bāṇa}}
\newcommand{\bimba}{\textit{bimba}}
\newcommand{\bha}{\textit{bha}}
\newcommand{\bhaga}{\textit{bhāga}}
\newcommand{\bhagrahanam}{\textit{bhagrahāṇām}}
\newcommand{\bhatraya}{\textit{bhatraya}}
\newcommand{\bhava}{\textit{bhāva}}
\newcommand{\bhavana}{\textit{bhavana}}
\newcommand{\bhuja}{\textit{bhuja}}
\newcommand{\bhujaa}{\textit{bhujā}}
\newcommand{\bhuvanakosa}{\textit{bhuvanakośa}}
\newcommand{\cakra}{\textit{cakra}}
\newcommand{\cala}{\textit{cala}}
\newcommand{\candra}{\textit{candra}}
\newcommand{\capa}{\textit{cāpa}}
\newcommand{\Devanagari}{Deva\-nāgarī}
\newcommand{\dhanus}{\textit{dhanus}}
\newcommand{\dhruva}{\textit{dhruva}}
\newcommand{\dina}{\textit{dina}}
\newcommand{\dorjya}{\textit{dorjyā}}
\newcommand{\dos}{\textit{dos}}
\newcommand{\drk}{\textit{dṛk}}
\newcommand{\drkkarma}{\textit{dṛkkarma}}
\newcommand{\drz}{\textit{dṛś}}
\newcommand{\drsti}{\textit{dṛṣṭi}}
\newcommand{\dvaya}{\textit{dvaya}}
\newcommand{\dvitiya}{\textit{dvitīya}}
\newcommand{\dvitiyaa}{\textit{dvitīyā}}
\newcommand{\dyu}{\textit{dyu}}
\newcommand{\dyucara}{\textit{dyucara}}
\newcommand{\dyujiva}{\textit{dyujīvā}}
\newcommand{\dyujya}{\textit{dyujyā}}
\newcommand{\dyuguna}{\textit{dyuguṇa}}
\newcommand{\ganita}{\textit{gaṇita}}
\newcommand{\gola}{\textit{gola}}
\newcommand{\golabandha}{\textit{golabandha}}
\newcommand{\golavid}{\textit{golavid}}
\newcommand{\gotra}{\textit{gotra}}
\newcommand{\graha}{\textit{gra\-ha}}
\newcommand{\grahaja}{\textit{grahaja}}
\newcommand{\grahajya}{\textit{grahajya}}
\newcommand{\grahana}{\textit{grahaṇa}}
\newcommand{\ista}{\textit{iṣṭa}}
\newcommand{\isu}{\textit{iṣu}}
\newcommand{\jiva}{\textit{jīvā}}
\newcommand{\jnana}{\textit{jñāna}}
\newcommand{\jya}{\textit{jyā}}
\newcommand{\jyaka}{\textit{jyakā}}
\newcommand{\jyardha}{\textit{jyārdha}}
\newcommand{\jyotihsastra}{\textit{jyotiḥ\-śāstra}}
\newcommand{\jyotirvid}{\textit{jyotir\-vid}}
\newcommand{\jyotisa}{\mbox{\textit{jyotiṣa}}}
\newcommand{\Jyotisaraja}{\textit{Jyotiṣa\-rāja}}
\newcommand{\jyotisika}{\textit{jyoti\-ṣika}}
\newcommand{\kadamba}{\textit{kadamba}}
\newcommand{\kanda}{\textit{kāṇḍa}}
\newcommand{\karana}{\textit{karaṇa}}
\newcommand{\karna}{\textit{karṇa}}
\newcommand{\keli}{\textit{keli}}
\newcommand{\kendra}{\textit{kendra}}
\newcommand{\khaga}{\textit{khaga}}
\newcommand{\kheta}{\textit{kheṭa}}
\newcommand{\kodanda}{\textit{kodaṇḍa}}
\newcommand{\kosthaka}{\textit{koṣṭhaka}}
\newcommand{\koti}{\textit{koṭi}}
\newcommand{\kotijiva}{\textit{koṭijīvā}}
\newcommand{\kotijya}{\textit{koṭijyā}}
\newcommand{\kranti}{\textit{krānti}}
\newcommand{\ksetra}{\textit{kṣetra}}
\newcommand{\lagna}{\textit{lagna}}
\newcommand{\lava}{\textit{lava}}
\newcommand{\madhyama}{\textit{madhyama}}
\newcommand{\mandala}{\textit{maṇḍala}}
\newcommand{\mangalacarana}{\textit{maṅgalā\-caraṇa}}
\newcommand{\mata}{\textit{mata}}
\newcommand{\maurvi}{\textit{maurvī}}
\newcommand{\mimamsa}{\textit{mīmāṃsā}}
\newcommand{\mleccha}{\textit{mleccha}}
\newcommand{\muni}{\textit{muni}}
\newcommand{\nabhaga}{\textit{nabhaga}}
\newcommand{\nabhoga}{\textit{nabhoga}}
\newcommand{\naksatra}{\textit{nakṣatra}}
\newcommand{\paksa}{\textit{pakṣa}}
\newcommand{\pandita}{\textit{paṇḍita}}
\newcommand{\para}{\textit{para}}
\newcommand{\paraa}{\textit{parā}}
\newcommand{\parama}{\textit{parama}}
\newcommand{\pata}{\textit{pāta}}
\newcommand{\phala}{\textit{phala}}
\newcommand{\purva}{\textit{pūrva}}
\newcommand{\prakara}{\textit{prakāra}}
\newcommand{\prakarana}{\textit{prakaraṇa}}
\newcommand{\prama}{\textit{pramā}}
\newcommand{\pramana}{\textit{pramāṇa}}
\newcommand{\prathama}{\textit{prathama}}
\newcommand{\prathamaa}{\textit{prathamā}}
\newcommand{\pratibimba}{\textit{pratibimba}}
\newcommand{\pratita}{\textit{pratita}}
\newcommand{\pratyaksa}{\textit{pratyakṣa}}
\newcommand{\prokta}{\textit{prokta}}
\newcommand{\rasa}{\textit{rasa}}
\newcommand{\rasi}{\textit{rāśi}}
\newcommand{\rata}{\textit{rata}}
\newcommand{\rekha}{\textit{rekhā}}
\newcommand{\rsi}{\textit{ṛṣi}}
\newcommand{\sabda}{\textit{śabda}}
\newcommand{\sadrksa}{\textit{sadṛkṣa}}
\newcommand{\sadrz}{\textit{sadṛś}}
\newcommand{\sadrza}{\textit{sadṛśa}}
\newcommand{\samyuti}{\textit{saṃyuti}}
\newcommand{\samskara}{\textit{saṃskāra}}
\newcommand{\sara}{\textit{śara}}
\newcommand{\sarani}{\textit{sāraṇī}}
\newcommand{\sasta}{\textit{ṣaṣṭha}}
\newcommand{\sastra}{\textit{śā\-stra}}
\newcommand{\sastrin}{\textit{śā\-strin}}
\newcommand{\sastric}{\textit{śā\-stric}}
\newcommand{\siddha}{\textit{siddha}}
\newcommand{\sinjini}{\textit{siñjinī}}
\newcommand{\spasta}{\textit{spaṣṭa}}
\newcommand{\sphuta}{\textit{sphuṭa}}
\newcommand{\sphutaa}{\textit{sphuṭā}}
\newcommand{\srngonnati}{\textit{śṛṅgon\-nati}}
\newcommand{\sruti}{\textit{śruti}}
\newcommand{\sthana}{\textit{sthāna}}
\newcommand{\sthula}{\textit{sthūla}}
\newcommand{\su}{\textit{su}}
\newcommand{\sukla}{\textit{śukla}}
\newcommand{\suksmatara}{\textit{sūkṣma\-tara}}
\newcommand{\surya}{\textit{sūrya}}
\newcommand{\sutra}{\textit{sūtra}}
\newcommand{\tika}{\textit{ṭīkā}}
\newcommand{\trtiya}{\textit{tṛtīya}}
\newcommand{\tribhajyaka}{\textit{tribha\-jyakā}}
\newcommand{\trijya}{\textit{trijyā}}
\newcommand{\triguna}{\textit{triguṇa}}
\newcommand{\triprasna}{\textit{tripraśna}}
\newcommand{\tulaa}{\textit{tulā}}
\newcommand{\udaya}{\textit{udaya}}
\newcommand{\udu}{\textit{uḍu}}
\newcommand{\unnatamsa}{\textit{unnatāṃśa}}
\newcommand{\upalabdhi}{\textit{upalabdhi}}
\newcommand{\upapatti}{\textit{upapatti}}
\newcommand{\vakya}{\textit{vākya}}
\newcommand{\valana}{\textit{valana}}
\newcommand{\valaya}{\textit{valaya}}
\newcommand{\vastu}{\textit{vastu}}
\newcommand{\Vedangaraja}{\textit{Vedāṅga\-rāja}}
\newcommand{\vidus}{\textit{vidus}}
\newcommand{\viloma}{\textit{viloma}}
\newcommand{\vilomaa}{\textit{vilomā}}
\newcommand{\viksepa}{\textit{vikṣepa}}
\newcommand{\visikha}{\textit{viśikha}}
\newcommand{\visuva}{\textit{viṣuva}}
\newcommand{\visuvat}{\textit{viṣuvat}}
\newcommand{\vrti}{\textit{vṛti}}
\newcommand{\vrtta}{\textit{vṛtta}}
\newcommand{\vyaksa}{\textit{vyakṣa}}
\newcommand{\yantra}{\textit{yantra}}
\newcommand{\yavana}{\textit{yavana}}
\newcommand{\Yavana}{\textit{Yavana}}
\newcommand{\yuktartha}{\textit{yuktārtha}}
\newcommand{\yukta}{\textit{yukta}}
\newcommand{\yukti}{\textit{yukti}}
\newcommand{\yuga}{\textit{yuga}}
\newcommand{\yugala}{\textit{yugala}}
\newcommand{\zabda}{\textit{śabda}}
\newcommand{\zinjini}{\textit{śiñjinī}}

\newcommand{\sardulavikridita}{\textit{śārdūla\-vikrīḍita}}

\newcommand{\Sin}{\text{Sin}}
\newcommand{\Cos}{\text{Cos}}
\newcommand{\arcSin}{\text{arcSin}}
\newcommand{\arcCos}{\text{arcCos}}
\newcommand{\VerSin}{\text{VerSin}}

% Arabic Diacritical marks for transliteration, taken from semtrans.sty. The entire package not implemented as conflict with other packages
\newcommand*\Alif{\char"02BE} %right
\newcommand*\Ayn{\char"02BF} %left
\newcommand*\AynCapI{\char"02BF} % Capital Italicised
\newcommand*\AynCap{\char"02BF\kern-0.25ex} % Capital Non Italicised
% Uses in front of capital A, with the words enclosed in braces, eg. \AynA{U}, otherwise there otherwise the glyph is unrecognisable. The kerning is to prevent the extra space when not italicised.


%Persian genitive
\newcommand{\idafaconsonant}{\textit{\=/i}}
\newcommand{\idafavowel}{\textit{\=/yi}}

\newcommand{\avagraha}{\texttt{'}\kern-.2ex}% Sanskrit Avagraha character
\newcommand{\selip}{\raisebox{-.75ex}{\tsans{@}}} % Sanskrit elipsis
%%%%%%%%%%%%%%%%%%%%%%%%%%%%%%%%%%%%%%%

% Commands for latin words with non-sentence breaking punctuation
\newcommand{\vid}{see}
\newcommand{\Vid}{See}
\newcommand{\ie}{i.e.\@ }
\newcommand{\eg}{e.g.\@ }
\newcommand{\fl}{\emph{fl}.\@ }
\newcommand{\viz}{\emph{viz}.\@ }
\newcommand{\scl}{\textit{sc}.\@ }
\newcommand{\etc}{etc.\@ } 
\newcommand{\etcp}{etc.} 
\newcommand{\cf}{\emph{cf}.\@ }
\newcommand{\Cf}{\emph{Cf}.\@ }
\newcommand{\ce}{\textsc{ce}}
\newcommand{\bce}{\textsc{bce}}
\newcommand{\ah}{\textsc{ah}} % Anno Hegirae (Hijri date)
\newcommand{\circa}{\emph{c}.\thinspace\nobreak }
\newcommand{\ante}{\emph{ante}}
\newcommand{\post}{\emph{post}}
\newcommand{\lit}{\emph{lit}.\@ }
\newcommand{\Ibid}{Ibid.\@ }
\newcommand{\opcit}{\textit{op.\thinspace cit.\@ }}
\newcommand{\sic}{\textit{sic}}
\newcommand{\alias}{\textit{alias}}



% MSS Sigla
\newcommand{\SjA}{Sj\kern0.15ex\textsubscript{A}}
\newcommand{\SjB}{Sj\kern0.15ex\textsubscript{B}}


%%%%%%%%%%%%%%%%%%%%%%%%%%%%%%%%%%%%%%%
% Bibliography

\usepackage[autostyle]{csquotes}
\usepackage[backend=biber, style=ext-authoryear-icomp, uniquename=false, natbib=true, language=autobib, autolang=langname, ibidtracker=false, doi=true]{biblatex} % You can change the style as per your requirement as given in the biblatex manual
% backref=true option produces "(cit. on p./pp. #)." at the end of the entries. Use if needed. 

% \DeclareOuterCiteDelims{parencite}{\bibopenbracket}{\bibclosebracket} % Using square brackets with parencite 

\addbibresource{bibliography.bib} % Place the .bib file in your project directory
\setlength\bibitemsep{0.5\baselineskip} % Sets the space between bibliography entries
\DeclareBibliographyAlias{webpage}{online}  % For the error ``No driver for entry type 'webpage'.''
\DeclareBibliographyAlias{url}{online}  % For the error ``No driver for entry type 'url'.''


%%%%%%%%%%%%%%%%%%%%%%%%%%%%%%%%%%%%%%%
% Extra trimmings


\setlength{\parindent}{0pt}
\setlength{\parskip}{0.75em}
% Ligature for fancy Ampersand
\chardef\&="E050
\newcommand{\degree}{\text{°}} % Degree symbol Unicode
\newcommand{\signederenvoi}{$\wedge\kern-.4ex\wedge$} % interlinear correction mark (Persian ms SjA)
\hbadness=99999

% For redefining \wideparen{} for arcs in mathmode (needs yhmath packge)
\DeclareFontFamily{OMX}{yhex}{}
\DeclareFontShape{OMX}{yhex}{m}{n}{<->yhcmex10}{}
\DeclareSymbolFont{yhlargesymbols}{OMX}{yhex}{m}{n}
\DeclareMathAccent{\wideparen}{\mathord}{yhlargesymbols}{"F3}

\renewcommand\labelitemi{--} % Using dashes instead of bullets in itemized lists
% \usepackage[all]{nowidow}



%%%%%%%%%%%%%%%%%%%%%%%%%%%%%%%%%%%%%%%
\usepackage[bookmarksopen=true, bookmarksdepth=4, hyperfootnotes=true, plainpages=false, bookmarks, bookmarksnumbered, colorlinks, linktocpage=true, linkcolor=blue, citecolor=blue, filecolor=black, urlcolor=blue]{hyperref}
\usepackage{bookmark}

%%%%%%%%%%%%%%%%%%%%%%%%%%%%%%%%%%%%%%%
% Glossary (LOAD AFTER HYPERREF FOR HYPERLINKED ENTRIES BUT BEFORE THE POLYGLOSSIA PERSIAN)
\usepackage[automake, xindy, acronym, stylemods, entrycounter, subentrycounter=true,numberedsection=nameref]{glossaries-extra} % For glossary

% entrycounter puts a counter on all parent entries in the glossary. An entries number can then be referenced by \glsrefentry{"label"}--- not used here, but good to note.

 \renewcommand*{\glsentrycounterlabel}{} % Used with option entrycounter to prevent numbered glossary from showing up, but the pagenumber of a glossary enrty can then be referenced using \pageref{glsentry-"label"} e.g., \pageref{glsentry-distance}
\usepackage{glossary-mcols} % For multicolumn (abbreviations)
\usepackage{glossary-tree}
\makeglossaries
% Glossary preamble

\setglossarypreamble[main]{%
This glossary lists Persian and Sanskrit technical expressions from (i) the chapter-titles of \ZijiShahJahani\ Discourse~II and \Siddhantasindhu\ Part~II in \S~\ref{chapter_title_comparision_persian_sanskrit}, (ii) the Persian text of \ZijiShahJahani\ Discourse II.6 in \S~\ref{zijshahjahan_persian_english}, and (iii) the Sanskrit text of \Siddhantasindhu\ Part II.6 in \S~\ref{siddhantasindhu_sanskrit_english}. Individual entries are grouped together under their common English translation. At the end of each entry, appropriate \textbf{chapter-numbers} and/or \textbf{passage-markers} in square brackets indicate its location in \S\S~\ref{chapter_title_comparision_persian_sanskrit}, \ref{zijshahjahan_persian_english}, or \ref{siddhantasindhu_sanskrit_english}. \Vid\ \S~\ref{glossary_appendix_format} for a more detailed description of the format of the glossary.\bigskip
}


\newglossaryentry{knowledge}
{
        name={knowledge},
        description={\tfarsi{معرفت} (\marifat); \tsans{{\tsnb{ज्ञा}}na} (\textit{jñāna})\newline \textit{passim}, seen in almost all chapter-titles of Discourse~II and Part~II}
}

\newglossaryentry{distance_ce_true_declination}
{
        name={distance from the celestial equator---},
        description={\phantom{x}\nopagebreak}
}

\newglossaryentry{distance_celestial_object}
{
        name={distance of a celestial object from the celestial equator},     
        description={\tfarsi{بعد كوكب از معدّل النها} (\bud\idafaconsonant\ \kawkab\ \az\ \muaddil\ \alnahar) \hyperlink{Pii6}{II.6}},
        plural={distance of celestial objects from the celestial equator},
        parent={distance_ce_true_declination}
}
\newglossaryentry{distance}
{
        name={distance},
        description={\tfarsi{بعد}  (\bud) [\hyperlink{PEpass5}{5}, \hyperlink{PEpass7}{7}]},
        parent={distance_ce_true_declination}
}

\newglossaryentry{true_declination}
{
        name={true declination},
        description={\tsans{spa.s.ta-kraanti} (\textit{spaṣṭa-krānti}) \hyperlink{Sii6}{II.6}, [\hyperlink{SEpass5}{5}], \tsans{sphu.ta-kraanti} (\textit{sphuṭa-krānti}) [\hyperlink{SEpass7}{7}], \tsans{sphu.ta-apama} (\textit{sphuṭa-apama}) [\hyperlink{SEpass11}{11}]},
        parent={distance_ce_true_declination}
}

\newglossaryentry{celestial_equator}
{
        name={celestial equator},     
        description={\tfarsi{معدّل النها}  (\muaddil\ \alnahar) \hyperlink{Pii5}{II.5}; \tsans{vi.suva-v.rtta} (\textit{viṣuva-vṛtta})\footnote{\,Kh attests \tsans{vi\underline{.sa}va-v.rtta} (\textit{vi\underline{ṣa}va-vṛtta}), an irregular form of \tsans{vi.suva-v.rtta} (\textit{viṣuva-vṛtta}), \vid\ \S\thinspace\ref{siddhantasindhu_sanskrit_english}:~footnote~\ref{vishava_deviant_3}.} [\hyperlink{SEpassB}{β}]}
}

\newglossaryentry{ecliptic}
{
        name={ecliptic},     
        description={\tsans{bhavana-cakra} (\textit{bhavana-cakra}) [\hyperlink{SEpassC}{γ}]} 
}
% \tfarsi{فلك البروج} (\falak\ \alburuj); UNATTESTED (=ecliptic)

\newglossaryentry{latitude_parent}
{
        name={latitude---},
        description={\phantom{x}\nopagebreak}
}

\newglossaryentry{latitude_celestial_object}
{
        name={latitude of a celestial object},
        description={\tfarsi{عرض کوکب} (\ard\idafaconsonant\ \kawkab) [\hyperlink{PEpass1}{1}]; \tsans{khagasya baa.na} (\textit{khagasya bāṇa}) [\hyperlink{SEpass1}{1}, \hyperlink{SEpass5}{5}]},
        parent={latitude_parent}
}

\newglossaryentry{latitude}
{
        name={latitude},
        description={\tfarsi{عرض}  (\ard) [\hyperlink{PEpass5}{5}, \hyperlink{PEpass6}{6}, \hyperlink{PEpass10}{10}]; \tsans{"sara} (\textit{śara}) [\hyperlink{SEpass6}{6}], \tsans{baa.na} (\textit{bāṇa}) [\hyperlink{SEpass10}{10}]},
        parent={latitude_parent}
}
 
\newglossaryentry{celestial_object}
{
        name={celestial object},
        description={\tfarsi{کوکب} (\kawkab), pl.\thinspace \tfarsi{كواكب} (\kawakib)   \hyperlink{Pii13}{II.13}, \hyperlink{Pii22}{22}, [\hyperlink{PEpass3}{3}, \hyperlink{PEpass4}{4}, \hyperlink{PEpass5}{5}];
       \tsans{graha} (\textit{graha}) \hyperlink{Sii7}{II.7}, \tsans{nak.satra} (\textit{nakṣatra}) \hyperlink{Sii12}{II.12}, \hyperlink{Sii13}{13}, \tsans{khaga} (\textit{khaga}) \hyperlink{Sii22}{II.22}, \tsans{nabhoga} (\textit{nabhoga}) [\hyperlink{SEpassA}{α}, \hyperlink{SEpassD}{δ}], \tsans{dyucara} (\textit{dyucara}) [\hyperlink{SEpassB}{β}], \tsans{bha} (\textit{bha}) [\hyperlink{SEpassD}{δ}]},
        plural={celestial objects}
}

\newglossaryentry{second_declination_parent}
{
        name={second declination---},
        description={\phantom{x}\nopagebreak}
}

\newglossaryentry{second_declination_degree}
{
        name={second declination of its degree},
        description={\tfarsi{میل ثاني درجه او} (\mayl\idafaconsonant\ \thani\idafavowel\ \daraji\idafavowel\ \uy) [\hyperlink{PEpass1}{1}]},
        parent={second_declination_parent}
}

\newglossaryentry{other_declination}
{
        name={other declination},
        description={\tsans{anyatara-apama} (\textit{anyatara-apama}) [\hyperlink{SEpass1}{1}]\newline
        synonymous with \tsans{anya-kraanti} (\textit{anya-krānti}) and identified with the \tsans{dvitiiya-kraanti} (\textit{dvitīya-krānti}) `second declination'},
        parent={second_declination_parent}
}
% \tsans{anya-kraanti} (\textit{anya-krānti})} UNATTESTED (=other declination)

% UNATTESTED
% \newglossaryentry{second_declination}
% {
%         name={second declination},
%         description={\tsans{dvitiiya-kraanti} (\textit{dvitīya-krānti})},
%         parent={second_declination_parent}
% }

\newglossaryentry{Cosine_second_declination_parent}
{
        name={Cosine of the second declination---},
        description={\phantom{x}\nopagebreak}
}


\newglossaryentry{Cosine_second_declination_degree}
{
        name={Cosine of the second declination of the degree},
        description={\tfarsi{جيب تمام میل ثاني درجه} (\jayb\idafaconsonant\ \tamam\idafaconsonant\ \mayl\idafaconsonant\ \thani\idafavowel\ \daraji) [\hyperlink{PEpass3}{3}, \hyperlink{PEpass4}{4}]},
        parent={Cosine_second_declination_parent}
}        

\newglossaryentry{Cosine_other_declination}
{
        name={Cosine of the other declination},
        description={\tsans{anya-kraanti-ko.tijyaa} (\textit{anya-krānti-koṭijyā}) [\hyperlink{SEpass3}{3}]},
        parent={Cosine_second_declination_parent}
}
 
\newglossaryentry{Cosine_second_declination}
{
        name={Cosine of the second declination},
        description={\tsans{dvitiiya-kraanti-ko.tijyaa} (\textit{dvitīya-krānti-koṭijyā}) [\hyperlink{SEpass4}{4}]},
        parent={Cosine_second_declination_parent}
}       

\newglossaryentry{direction_sum_parent}
{
        name={direction of the sum---},
        description={\phantom{x}\nopagebreak}
}

\newglossaryentry{direction_sum}
{
        name={direction of the sum},
        description={\tfarsi{جهت مجموع}  (\jahat\idafaconsonant\ \majmu) [\hyperlink{PEpass1}{1}, \hyperlink{PEpass10}{10}]},
        parent={direction_sum_parent}
}

\newglossaryentry{one_direction}
{
        name={one direction},
        description={\tfarsi{یک جهت}  (\yik\ \jahat) [\hyperlink{PEpass1}{1}, \hyperlink{PEpass10}{10}]; \tsans{eka\=/di"s} (\textit{eka\=/diś})} [\hyperlink{SEpass1}{1}],
        parent={direction_sum_parent} 
}


\newglossaryentry{sum}
{
        name={sum},
        description= {\newline to \textbf{sum}\quad  \tfarsi{جمع کردن} (\jam\ \kardan) [\hyperlink{PEpass1}{1}, \hyperlink{PEpass10}{10}]\\[.2\baselineskip]
        \textbf{sum}\quad \tsans{sa.myuti} (\textit{saṃyuti}) [\hyperlink{SEpass1}{1}, \hyperlink{SEpass10}{10}]},
}


\newglossaryentry{difference_parent}
{
        name={difference---},
        description={\phantom{x}\nopagebreak}
}

\newglossaryentry{difference}
{
        name={difference},
        description={\tfarsi{تفاضل} (\tafadul) [\hyperlink{PEpass1}{1}, \hyperlink{PEpass10}{10}]; \tsans{antara}  (\textit{antara}) [\hyperlink{SEpass1}{1}], \tsans{vivara} (\textit{vivara}) [\hyperlink{SEpassB}{β}, \hyperlink{SEpassC}{γ}], \tsans{viyuti} (\textit{viyuti})} [\hyperlink{SEpass10}{10}],
        parent={difference_parent}
}

\newglossaryentry{subtraction}
{
        name={made to be subtracted},
        description={\tsans{vi"sodhita} (\textit{viśodhita}) [\hyperlink{SEpass10}{10}]},
        parent={difference_parent}
}

\newglossaryentry{argument_of_distance}
{
        name={argument of the distance---},
        description={\phantom{x}\nopagebreak}
}

\newglossaryentry{share_of_distance}
{
        name={share of the distance},
        description={\tfarsi{حصّهٔ بعد} (\hissi\idafavowel\ \bud) [\hyperlink{PEpass1}{1}], \tfarsi{حصّة البعد} (\hissatalbud) [\hyperlink{PEpass7}{7}]},
        parent={argument_of_distance}
}

\newglossaryentry{curve_true_declination}
{
        name={curve of true declination},
        description={\tsans{spa.s.ta-kraanti-a"nka} (\textit{spaṣṭa-krānti-aṅka}) [\hyperlink{SEpass7}{7}]},
        parent={argument_of_distance}
}
% \tsans{sphu.ta-apama-a"nka} (\textit{sphuṭa-apama-aṅka}) NOT ATTESTED ANYWHERE BY ITSELF

\newglossaryentry{share_true_declination}
{
        name={share of the true declination},
        description={\tsans{sphu.ta-apama-a.m"sa} (\textit{sphuṭa- apama-aṃśa}) [\hyperlink{SEpass1}{1}]},
        parent={argument_of_distance}
}


\newglossaryentry{direction_share_of_distance}
{
        name={direction of the share of the distance},
        description={\tfarsi{جهت حصّهٔ بعد} (\jahat\idafaconsonant\ \hissi\idafavowel\ \bud) [\hyperlink{PEpass1}{1}, \hyperlink{PEpass3}{3}]}
}

\newglossaryentry{direction}
{
        name={direction},
        description={\tfarsi{جهت} (\jahat) [\hyperlink{PEpass3}{3}, \hyperlink{PEpass6}{6}, \hyperlink{PEpass9}{9}--\hyperlink{PEpass11}{11}]; \tsans{di"s} (\textit{diś}) [\hyperlink{SEpass11}{11}]}
}

\newglossaryentry{direction_difference_parent}
{
        name={direction of the difference---},
        description={\phantom{x}\nopagebreak}
}

\newglossaryentry{direction_residue}
{
        name={direction of the residue},
        description={\tfarsi{جهت فضل} (\jahat\idafaconsonant\ \fadla) [\hyperlink{PEpass1}{1}, \hyperlink{PEpass10}{10}]},
        parent={direction_difference_parent}
}

\newglossaryentry{different_directions}
{
        name={different directions},
        description={\tfarsi{جهت مختلف} (\jahat\idafaconsonant\ \mukhtalif) [\hyperlink{PEpass10}{10}]},
        parent={direction_difference_parent}
}

\newglossaryentry{lowering}
{       
        name={lowering---},
        description={\phantom{x}\nopagebreak}
}        

\newglossaryentry{having_been_lowered}
{
        name={having been lowered},
        description={\tsans{adharii-k.rta} (\textit{adharī-kṛta}) [\hyperlink{SEpass2}{2}, \hyperlink{SEpass8}{8}, \hyperlink{SEpass11}{11}]},
        parent={lowering}
}

\newglossaryentry{should_be_lowered}
{
        name={should [be made] lower},
        description={\tsans{adha.h kuryaat} (\textit{adhaḥ kuryāt}) [\hyperlink{SEpass6}{6}]},
        parent={lowering}
}


\newglossaryentry{low_multiplication}
{
        name={low-multiplication},
        description={\newline to \textbf{low-multiply} \quad \tfarsi{منحطّ ضرب کردن} (\munhatt\idafaconsonant\ \darb\ \mbox{\kardan}) [\hyperlink{PEpass2}{2}, \hyperlink{PEpass6}{6}, \hyperlink{PEpass8}{8}, \hyperlink{PEpass11}{11}]}
}

\newglossaryentry{low_division}
{
        name={low-division},
        description={to \textbf{low-divide}\quad \tfarsi{منحطّ قسمت کردن } (\munhatt\idafaconsonant\ \qismat\ \kardan) [\hyperlink{PEpass9}{9}]}
}

\newglossaryentry{Sine_arg_dist}
{
        name={Sine of the argument of the distance---},
        description={\phantom{x}\nopagebreak}
}

\newglossaryentry{Sine_share_of_distance}
{
        name={Sine of the share of the distance},
        description={\tfarsi{جيب حصّهٔ بعد} (\jayb\idafaconsonant\ \hissi\idafavowel\ \bud) [\hyperlink{PEpass2}{2}, \hyperlink{PEpass3}{3}, \hyperlink{PEpass4}{4}]},
        parent={Sine_arg_dist}
}

\newglossaryentry{Sine_curve_true_declination}
{
        name={Sine of the curve of true declination},
        description={\tsans{sphu.ta-apama-a"nka-si{\tsnb{ञ्जि}}nii}  (\textit{sphuṭa-apama-aṅka-siñjinī}) [\hyperlink{SEpass2}{2}], \tsans{sphu.ta-kraanti-a"nka-jivaa} (\textit{sphuṭa-krānti-aṅka-jīvā}) [\hyperlink{SEpass3}{3}], \tsans{sphu.ta-kraanti-a"nka-jyaa} (\textit{sphuṭa-krānti-aṅka-jyā}) [\hyperlink{SEpass4}{4}]},
        parent={Sine_arg_dist}
}

\newglossaryentry{sine_declination}
{
        name={Sine of the declination---},
        description={\phantom{x}\nopagebreak}
}

\newglossaryentry{Sine_distance}
{
        name={Sine of the distance},
        description={\tfarsi{جيب بعد} (\jayb\idafaconsonant\ \bud) [\hyperlink{PEpass2}{2}, \hyperlink{PEpass3}{3}, \hyperlink{PEpass4}{4}, \hyperlink{PEpass6}{6}],},
        parent={sine_declination}
}

\newglossaryentry{Sine_distance_celestial_object}
{
        name={Sine of the distance of the celestial object},
        description={\tfarsi{جيب بعد کوکب} (\jayb\idafaconsonant\ \bud\idafaconsonant\ \kawkab) [\hyperlink{PEpass11}{11}],},
        parent={sine_declination}
}

\newglossaryentry{Sine_true_declination}
{
        name={Sine of the true declination},
        description={\tsans{sphu.ta-apama-jyakaa} (\textit{sphuṭa-apama-jyakā}) [\hyperlink{SEpass2}{2}], 
                    \tsans{spa.s.ta-apama-jyakaa} (\textit{spaṣṭa-apama-jyakā}) [\hyperlink{SEpass3}{3}], 
                    \tsans{spa.s.ta-kraanti-jyaa} (\textit{spaṣṭa-krānti-jyā}) [\hyperlink{SEpass4}{4}, \hyperlink{SEpass6}{6}]},
        parent={sine_declination}
}

\newglossaryentry{Cosine_maximum_declination}
{
        name={Cosine of the maximum declination},
        description={\tfarsi{جیب تمام میل کلّی} (\jayb\idafaconsonant\ \tamam\idafaconsonant\ \mayl\idafaconsonant\ \kulli) [\hyperlink{PEpass3}{3}, \hyperlink{PEpass6}{6}],
        \tsans{parama-kraanti-ko.tijyaa} (\textit{parama-krānti-koṭijyā}) [\hyperlink{SEpass3}{3}, \hyperlink{SEpass6}{6}]}
} 

\newglossaryentry{table_Cosine_maximum_declination}
{
        name={table of the Cosine of the maximum declination},
        description={\mbox{\tfarsi{جدول جيب تمام میل کلّی}}  (\jadval\idafaconsonant\ \jayb\idafaconsonant\ \tamam\idafaconsonant\ \mayl\idafaconsonant\ \kulli) [\hyperlink{PEpass4}{4}];
        \tsans{parama-kraanti-ko.tijyaa-ko.s.thika} (\textit{parama-krānti-koṭijyā-koṣṭhika}) [\hyperlink{SEpass4}{4}], \tsans{parama-kraanti-ko.tijyaa-ko.s.thaka} (\textit{parama-krānti-koṭijyā-koṣṭhaka}) [\hyperlink{SEpass6}{6}]}
} 


\newglossaryentry{Cosine_longitude_increased_ninety}
{
        name={Cosine of the ecliptic longitude increased by ninety degrees---},
        description={\phantom{x}\nopagebreak}
}

\newglossaryentry{Cosine_inverse_declination_degree_celestial_object}
{
        name={Cosine of the inverse declination of the degree of a celestial object},
        description={\newline\tfarsi{جیب تمام میل منکوس درجه كوكب}  (\jayb\idafaconsonant\ \tamam\idafaconsonant\ \mayl\idafaconsonant\ \mankus\idafaconsonant\ \daraji\idafavowel\ \kawkab) [\hyperlink{PEpass2}{2}]},
        parent={Cosine_longitude_increased_ninety}
}


\newglossaryentry{day_Sine_increased_by_three_signs}
{
        name= {day-Sine [of the longitude] increased by three zodiacal signs},
        description={\tsans{sa-bha-traya-dyujiivaa} (\textit{sa-bha-traya-dyujīvā}) [\hyperlink{SEpass2}{2}]},
        parent={Cosine_longitude_increased_ninety}   
}

% UNATTESTED
% \newglossaryentry{inverse_declination}
% {
%         name={inverse declination of the degree of a celestial object},
%         description={\tfarsi{میل منکوس درجه كوكب} (\mayl\idafaconsonant\ \mankus\idafaconsonant\ \daraji\idafavowel\ \kawkab)}
% }

\newglossaryentry{result}
{
        name={result},
        description={\tfarsi{حاصل} (\hasil) [\hyperlink{PEpass2}{2}--\hyperlink{PEpass4}{4}, \hyperlink{PEpass6}{6}, \hyperlink{PEpass8}{8}, \hyperlink{PEpass10}{10}]; \tsans{phala} (\textit{phala}) \hyperlink{Sii2}{II.2}}
}

\newglossaryentry{multiplication}
{
        name ={multiplication},
        description ={\newline
           to \textbf{multiply}\quad \tfarsi{ضرب کردن} (\darb\ \kardan) [\hyperlink{PEpass3}{3}]\\[0.2\baselineskip]
        \textbf{is multiplied}\quad \tsans{ni-hanyate} (\textit{ni-hanyate}) [\hyperlink{SEpass2}{2}]\\[0.2\baselineskip]
         \textbf{having been multiplied}\quad \tsans{hata} (\textit{hata}) [\hyperlink{SEpass3}{3}, \hyperlink{SEpass8}{8}, \hyperlink{SEpass11}{11}]\\[0.2\baselineskip]
         \textbf{to/must be multiplied}\quad \tsans{sa.mgu.nya} (\textit{saṃguṇya}) [\hyperlink{SEpass6}{6}]}
}

\newglossaryentry{maximum_declination}
{
        name={maximum declination},
        description={\tfarsi{میل کلّی} (\mayl\idafaconsonant\ \kulli) [\hyperlink{PEpass7}{7}, \hyperlink{PEpass10}{10}]; \tsans{parama-kraanti} (\textit{parama-krānti}) [\hyperlink{SEpass7}{7}], \tsans{parama-apama} (\textit{parama-apama}) [\hyperlink{SEpass10}{10}]}
}        
        
\newglossaryentry{division}
{
        name={division},
        description={\newline to \textbf{divide}\quad\tfarsi{قسمت کردن} (\qismat\ \kardan) [\hyperlink{PEpass3}{3}, \hyperlink{PEpass4}{4}]\\[.2\baselineskip]
        \textbf{having been divided}\quad \tsans{aapta} (\textit{āpta}) [\hyperlink{SEpass3}{3}] \\[.2\baselineskip]
        {[one]} \textbf{should divide}\quad \tsans{bhajet} (\textit{bhajet}) [\hyperlink{SEpass4}{4}] \\[.2\baselineskip]
        \textbf{having been divided}\quad \tsans{bhaajita} (\textit{bhājita}) [\hyperlink{SEpass9}{9}]}
}   

\newglossaryentry{quotient}
{
        name={quotient---},
        description={\phantom{x}\nopagebreak}
}

\newglossaryentry{quotient_division}
{
        name={quotient of division},
        description={\tfarsi{خارج قسمت} (\kharij\idafaconsonant\ \qismat) \hyperlink{Pii1}{II.1}, [\hyperlink{PEpass3}{3}, \hyperlink{PEpass4}{4}, \hyperlink{PEpass9}{9}]},
        parent={quotient}
}

\newglossaryentry{obtained}
{
        name={obtained},
        description={\tsans{labdha} (\textit{labdha}) [\hyperlink{SEpass4}{4}]},
        parent={quotient}
}

\newglossaryentry{table}
{
        name={table},
        description={\tfarsi{جدول} (\jadval) [\hyperlink{PEpass6}{6}]; \tsans{ko.s.thaka} (\textit{koṣṭhaka}) [\textit{varia lectio}: \tsans{ko.s.thika} (\textit{koṣṭhika})]  \hyperlink{Sii11}{II.11}},
        plural={tables}
}        

\newglossaryentry{declination_parent}
{
        name={declination---},
        description={\phantom{x}\nopagebreak}
}

\newglossaryentry{declination_degree_celestial_object}
{
        name={declination of the degree of a celestial object},
        description={\tfarsi{میل درجه كوكب} (\mayl\idafaconsonant\ \daraji\idafavowel\ \kawkab) [\hyperlink{PEpass10}{10}]},
        parent={declination_parent}
}  

\newglossaryentry{declination_celestial_object}
{
        name={declination of a celestial object},
        description={\tsans{khagasya kraanti} (\textit{khagasya krānti}) [\hyperlink{SEpass7}{7}]},
        parent={declination_parent}
}

\newglossaryentry{declination}
{
        name={declination},
        description={\tfarsi{میل} (\mayl) [\hyperlink{PEpass6}{6}]; \tsans{kraanti} (\textit{krānti}) \hyperlink{Sii5}{II.5}, [\hyperlink{SEpass5}{5}, \hyperlink{SEpass6}{6}]},
        parent={declination_parent}
}

\newglossaryentry{declination_degree}
{
        name={declination of its degree},
        description={\tfarsi{میل درجه او} (\mayl\idafaconsonant\ \daraji\idafavowel\ \uy) [\hyperlink{PEpass5}{5}, \hyperlink{PEpass7}{7}]},
        parent={declination_parent}
}        

\newglossaryentry{ecliptic_distance_equinox}
{
        name={ecliptic longitude---},
        description={\phantom{x}\nopagebreak}
}

% UNATTESTED
% \newglossaryentry{distance_degree_celestial_object_nearest_equinox}
% {   
%         name={distance of the degree of a celestial object from the nearest equinox},
%         description={\newline\tfarsi{بعد درجه کوکب از اعتدال اقرب} (\bud\idafaconsonant\ \daraji\idafavowel\ \kawkab\ \az\ \itidal\idafaconsonant\ \aqrab)},
%         parent={ecliptic_distance_equinox}
% }

% UNATTESTED
% \newglossaryentry{arc_ecliptic_longitude_celestial_object}
% {
%         name={arc of ecliptic longitude of a celestial object},
%         description={\tsans{khagasya bhujaa} (\textit{khagasya bhujā})},
%         parent={ecliptic_distance_equinox}
% } 

% UNATTESTED
% \newglossaryentry{degree_celestial_object}
% {
%         name={degree of a celestial object},
%         description={\tfarsi{درجه کوکب} (\daraji\idafavowel\ \kawkab)},
%         parent={ecliptic_distance_equinox}
% }

\newglossaryentry{its_degree}
{
        name={its degree},
        description={\tfarsi{درجه او} (\daraji\idafavowel\ \uy) [\hyperlink{PEpass6}{6}]\newline
        in other words, \tfarsi{درجه کوکب} (\daraji\idafavowel\ \kawkab) `degree of a celestial object'},
        parent={ecliptic_distance_equinox}
} 

\newglossaryentry{Sine_latitude_parent}
{
        name={Sine of the latitude---},
        description={\phantom{x}\nopagebreak}
}
    
\newglossaryentry{Sine_latitude_celestial_object}
{
        name={Sine of the latitude of a celestial object},
        description={\tfarsi{جیب عرض کوکب} (\jayb\idafaconsonant\ \ard\idafaconsonant\ \kawkab) [\hyperlink{PEpass9}{9}]; \tsans{nabhoga-vi"sikhasya si}\tsnb{ञ्जि}\tsans{nii} (\textit{nabhoga-viśikhasya siñjinii}) [\hyperlink{SEpass9}{9}]},
        parent={Sine_latitude_parent}
}

\newglossaryentry{Sine_latitude}
{
        name={Sine of its latitude},
        description={\tfarsi{جيب عرض او} (\jayb\idafaconsonant\ \ard\idafaconsonant\ \uy) [\hyperlink{PEpass6}{6}]},    
        parent={Sine_latitude_parent}
}

\newglossaryentry{Sine_the_latitude}
{
        name={Sine of the latitude},
        description={\tsans{baa.na-jyaa}  (\textit{bāṇa-jyā}) [\hyperlink{SEpass6}{6}]},
        parent={Sine_latitude_parent}
}

\newglossaryentry{direction_latitude_parent}
{       name={direction of the latitude---},
        description={\phantom{x}\nopagebreak}
}

\newglossaryentry{direction_latitude_celestial_object}
{
        name={direction of the latitude of a celestial object},
        description={\tfarsi{جهت عرض کوکب} (\jahat\idafaconsonant\ \ard\idafaconsonant\ \kawkab) [\hyperlink{PEpass9}{9}]},
        parent={direction_latitude_parent}
}

\newglossaryentry{direction_latitude}
{
        name={direction of the latitude},
        description={\tfarsi{جهت عرض} (\jahat\idafaconsonant\ \ard) [\hyperlink{PEpass6}{6}]; \tsans{baa.na-di"s} (\textit{bāṇa-diś}) [\hyperlink{SEpass6}{6}]},
        parent={direction_latitude_parent}
} 

\newglossaryentry{Sine_complement_ecliptic_longitude}
{
        name={Sine of the complement of the ecliptic longitude---},
        description={\phantom{x}\nopagebreak}
}

\newglossaryentry{Sine_distance_celestial_object_nearest_solstice}
{
        name={Sine of the distance of the degree of a celestial object from the nearest \mbox{solstice}},
        description={\tfarsi{جيب بعد درجه کوکب از انقلاب اقرب} (\jayb\idafaconsonant\ \bud\idafaconsonant\ \daraji\idafavowel\ \kawkab\ \az\ \inqilab\idafaconsonant\ \aqrab) [\hyperlink{PEpass8}{8}]},
        parent={Sine_complement_ecliptic_longitude}
}

\newglossaryentry{Sine_koti_celestial_object}
{
        name={Sine of the complement of the arc of ecliptic longitude of a celestial object},
        description={\newline\tsans{khagasya ko.ti-si}\tsnb{ञ्जि}\tsans{nii} (\textit{khagasya koṭi-siñjini}) [\hyperlink{SEpass8}{8}]},
        parent={Sine_complement_ecliptic_longitude}
}        

\newglossaryentry{Cosine_latitude_parent}
{
        name={Cosine of the latitude---},
        description={\phantom{x}\nopagebreak}
}
        
\newglossaryentry{Cosine_latitude_celestial_object}
{
        name={Cosine of the latitude of a celestial object},
        description={\tfarsi{جیب تمام عرض کوکب} (\jayb\idafaconsonant\ \tamam\idafaconsonant\ \ard\idafaconsonant\ \kawkab) [\hyperlink{PEpass8}{8}]},
        user1={Cosine of the latitude of the celestial object},
        parent={Cosine_latitude_parent}
}

\newglossaryentry{Cosine_latitude}
{
        name={Cosine of its latitude},
        description={\tsans{sva-baa.na-ko.tijiivaa} (\textit{sva-bāṇa-koṭijīvā}) [\hyperlink{SEpass8}{8}]},
        parent={Cosine_latitude_parent}
}

\newglossaryentry{Sine_distance_solstitial_colure_parent}
{
        name={Sine of the distance along the `circle congruent to the ecliptic' from the\linebreak solstice---},
        description={\phantom{x}\nopagebreak}
} 

\newglossaryentry{Sine_distance_celestial_object_solstitial_colure}
{
        name={Sine of the distance of a celestial object from the `circle passing through the four poles'},
        description={\tfarsi{جیب بعد کوکب از
        \tfarsib{دایرهٔ ماره باقطاب اربعه}} (\jayb\idafaconsonant\ \bud\idafaconsonant\ \kawkab\ \az\ \guillemotleft\dayiri\idafavowel\ \marri\ \biaqtab\idafaconsonant\ \arbai\guillemotright) [\hyperlink{PEpass8}{8}]},
        parent={Sine_distance_solstitial_colure_parent}
}        

\newglossaryentry{Sine_congruent_koti}
{
        name={Sine of the congruent complementary arc},
        description={\tsans{sad.rk.sa-ko.ti-si}\tsnb{ञ्जि}\tsans{nii} (\textit{sadṛkṣa-koti-siñjinī}) [\hyperlink{SEpass8}{8}]},
        parent={Sine_distance_solstitial_colure_parent}
}


\newglossaryentry{Cosine_distance_solstitial_colure}
{
        name={Cosine of the distance of a celestial object from the `circle passing through the four poles'},
        description={\tfarsi{جیب تمام بعد از
        \tfarsib{دایرهٔ ماره باقطاب اربعه}} (\jayb\idafaconsonant\ \tamam\idafaconsonant\ \bud\ \az\ \guillemotleft\dayiri\idafavowel\ \marri\ \biaqtab\idafaconsonant\ \arbai\guillemotright) [\hyperlink{PEpass9}{9}, \hyperlink{PEpass11}{11}]}
}

\newglossaryentry{Sine_distance_cong_circle_equinox}
{
        name={Sine of the distance along the `circle congruent to the ecliptic' from the \linebreak equinox---},
        description={\phantom{x}\nopagebreak}
}

\newglossaryentry{lowered_Sine_congruent_base}
{
        name={lowered Sine of the congruent arc},
        description={\tsans{adhara-sad.rk.sa-dorjyaa} (\textit{adhara-sadrkṣa-dorjyā}) [\hyperlink{SEpass9}{9}]},
        parent={Sine_distance_cong_circle_equinox}
}

\newglossaryentry{Sine_congruent_base}
{
        name={Sine of the congruent arc},
        description={\tsans{sad.rk.sa-baahu-jyakaa} (\textit{sadrkṣa-bāhu-jyakā}) [\hyperlink{SEpass11}{11}]},
        parent={Sine_distance_cong_circle_equinox}
}

\newglossaryentry{arc}
{
        name={arc},
        description={\tfarsi{قوس} (\qaws) [\hyperlink{PEpass9}{9}]; \tsans{dhanus} (\textit{dhanus}) [\hyperlink{SEpassB}{β}, \hyperlink{SEpassC}{γ}, \hyperlink{SEpass9}{9}], \tsans{koda.n.da} (\textit{kodaṇḍa}) [\hyperlink{SEpassD}{δ}], \tsans{caapa} (\textit{cāpa}) [\hyperlink{SEpass11}{11}]}
}

\newglossaryentry{table_of_Sine}
{
        name={table of Sine},
        description={\tfarsi{جدول جيب} (\jadval\idafaconsonant\ \jayb) [\hyperlink{PEpass9}{9}]}
}

\newglossaryentry{arc_max_latitude_parent}
{
        name={arc of maximum latitude---},
        description={\phantom{x}\nopagebreak}
}

\newglossaryentry{first_arc}
{
        name={first arc},
        description={\tfarsi{قوس اوّل} (\qaws\idafaconsonant\ \avval) [\hyperlink{PEpass9}{9}, \hyperlink{PEpass10}{10}]},
        parent={arc_max_latitude_parent}
}        

\newglossaryentry{maximum_latitude}
{
        name={maximum latitude},
        description={\tsans{para-i.su} (\textit{para-iṣu}) [\hyperlink{SEpassC}{γ}, \hyperlink{SEpass10}{10}], \tsans{para-"sara} (\textit{para-śara}) [\hyperlink{SEpass9}{9}]},
        parent={arc_max_latitude_parent}
}        

\newglossaryentry{ninety_parent}
{
        name={ninety degrees---},
        description={\phantom{x}\nopagebreak}
}
  
\newglossaryentry{one_quarter}
{
        name={one-quarter},
        description={\tfarsi{ربع} (\rub) [\hyperlink{PEpass10}{10}]},
        parent={ninety_parent}
}

\newglossaryentry{ninety}
{
        name={ninety},
        description={\tsans{abhra-nava} (\textit{abhra-nava})\footnote{\,\bhutasamkhya\ word-numerals: \textit{abhra} `0'and  \textit{nava} `9' forming `90'} [\hyperlink{SEpass10}{10}]},
        parent={ninety_parent}
}

\newglossaryentry{exceeds_parent}
{
        name={exceeds---},
        description={\phantom{x}\nopagebreak}
}

\newglossaryentry{exceeds}
{
        name={exceeds},
        description={to \textbf{exceed}\quad \tfarsi{زیادی شدن} (\ziyadi\ \shudan) [\hyperlink{PEpass10}{10}]},
        parent={exceeds_parent}            
}           

\newglossaryentry{greater}
{
        name={greater},
        description={\tsans{adhika} (\textit{adhika}) [\hyperlink{SEpass10}{10}]},
        parent={exceeds_parent}
}
 

\newglossaryentry{whole_sum}
{
        name={whole sum},
        description={\tfarsi{تمام مجموع} (\tamam\idafaconsonant\ \majmu) [\hyperlink{PEpass10}{10}]}
}

\newglossaryentry{one_eighty}
{
        name={one hundred and eighty degrees---},
        description={\phantom{x}\nopagebreak}
}

\newglossaryentry{one_half}
{
        name={one-half},
        description={\tfarsi{نصف} (\nisf) [\hyperlink{PEpass10}{10}]},
        parent={one_eighty}
}

\newglossaryentry{one_hundred_eighty}
{
        name={one hundred and eighty},
        description={\tsans{kha-a.s.ta-bhuu} (\textit{akha-aṣṭa-bhū})\footnote{\,\bhutasamkhya\ word-numerals: \textit{kha} `0', \textit{aṣṭa} `8', and  \textit{bhū} `1' forming `180'} [\hyperlink{SEpass10}{10}]},
        parent={one_eighty}
}

\newglossaryentry{arc_max_declination}
{
        name={arc of maximum argument of the distance---},
        description={\phantom{x}\nopagebreak}
}

\newglossaryentry{second_arc}
{
        name={second arc},
        description={\tfarsi{قوس دوم} (\qaws\idafaconsonant\ \duvum) [\hyperlink{PEpass10}{10}]},
        parent={arc_max_declination}
}

\newglossaryentry{maximum_true_declination}
{
        name={maximum true declination},
        description={\tsans{para-sphu.ta-apama} (\textit{para-sphuṭa-apama}) [\hyperlink{SEpassB}{β}]},
        parent={arc_max_declination}
}    

\newglossaryentry{maximum_true_declination_celestial_object}
{
        name={maximum true declination of a celestial object},
        description={\newline\tsans{grahasya para-sphu.ta-apama} (\textit{grahasya para-sphuṭa-apama}) [\hyperlink{SEpass10}{10}]},
        parent={arc_max_declination}
}        

\newglossaryentry{Sine_max_arg_dist}
{
        name={Sine of the maximum argument of the distance---},
        description={\phantom{x}\nopagebreak}
}

\newglossaryentry{Sine_second_arc}
{
        name={Sine of the second arc},
        description={\tfarsi{جیب قوس دوم} (\jayb\idafaconsonant\ \qaws\idafaconsonant\ \duvum) [\hyperlink{PEpass11}{11}]},
        parent={Sine_max_arg_dist}
}

\newglossaryentry{Sine_maximum_true_declination}
{
        name={Sine of the maximum true declination},
        description={\tsans{para-sphu.ta-kraanti-bhava-jyakaa} (\textit{para-sphuṭa-krānti-bhava-jyakā}) [\hyperlink{SEpass11}{11}]},
        parent={Sine_max_arg_dist}
}        


\newglossaryentry{direction_second_arc}
{
        name={direction of the second arc},
        description={\tfarsi{جهت قوس دوم} (\jahat\idafaconsonant\ \qaws\idafaconsonant\ \duvum) [\hyperlink{PEpass11}{11}]}
}

\newglossaryentry{direction_same_different_parent}
{
        name={direction of the sum or the difference---},
        description={\phantom{x}\nopagebreak}
}

\newglossaryentry{conjunction_disjunction_direction}
{
        name={direction of the conjunction or the disjunction},
        description={\tsans{yuti-viyoga\=/di"s} (\textit{yuti-viyoga\=/diś}) [\hyperlink{SEpass10}{10}], \tsans{sa.myoga-viyoga\=/di"s} (\textit{saṃyoga-viyoga\=/diś}) [\hyperlink{SEpass11}{11}]},
        parent={direction_same_different_parent}
}

\newglossaryentry{own_direction}
{
        name={own direction},
        description={\tsans{sva-di"s} (\textit{sva-diś}) [\hyperlink{SEpass1}{1}]},
        parent={direction_same_different_parent}
}

\newglossaryentry{same_different_directions}
{
        name={same or different directions},
        description={\tsans{sama-bhinna-di"s} (\textit{sama-bhinna-diś}) [\hyperlink{SEpass10}{10}]},
        parent={direction_same_different_parent}
}

\newglossaryentry{product_result_multiplication}
{
        name={result/product of multiplication---},
        description={\phantom{x}\nopagebreak}
}

\newglossaryentry{result_multiplication}
{
        name={result of multiplication},
        description={\tfarsi{حاصل ضرب} (\hasil\idafaconsonant\ \darb) \hyperlink{Pii1}{II.1}},
        parent={product_result_multiplication}
} 

\newglossaryentry{product_multiplication}
{
        name={product of multiplication},
        description={\tsans{gu.nita-phala} (\textit{guṇita-phala}) [\hyperlink{SEpass4}{4}, \hyperlink{SEpass6}{6}]},
        parent={product_result_multiplication}
}
       

\newglossaryentry{result_multiplication_division}
{
        name={result of multiplication and division},
        description={ \tsans{gu.nana-bhajana-phala} (\textit{guṇana-bhajana-phala}) \hyperlink{Sii1}{II.1}}
}

\newglossaryentry{extract}
{
        name={extract},
        description={\newline 
        to \textbf{extract}/\textbf{bring out}\quad \tfarsi{درآردن} (\darardan) [\hyperlink{PEpass4}{4}, \hyperlink{PEpass6}{6}]\\[.2\baselineskip]
        \textbf{having extracted}\quad \tsans{utthaaya} (\textit{utthāya}) [\hyperlink{SEpass4}{4}]\\[0.2\baselineskip]
        \textbf{may extract}\quad \tsans{utthaapayet} (\textit{utthāpayet}) [\hyperlink{SEpass6}{6}]}
}


\newglossaryentry{circle}
{
        name={circle},
        description={\tsans{v.rtta} (\textit{vṛtta}) [\hyperlink{SEpassA}{α}]}
}
% \tfarsi{دایره} (\textit{dāyiri}); UNATTESTED (=CIRCLE)

\newglossaryentry{ecliptic_pole}
{
        name={ecliptic pole},
        description={\tsans{kadamba} (\textit{kadamba}) [\hyperlink{SEpassA}{α}]}
}

\newglossaryentry{celestial_pole}
{
        name={celestial pole},
        description={\tsans{vi.suvat-dhruva} (\textit{viṣuvat-dhruva})\footnote{\,Kh attests \tsans{vi\underline{.sa}va-dhruva} (\textit{vi\underline{ṣa}va-dhruva}), an irregular form
        of \tsans{vi.suvat-dhruva} (\textit{viṣuvat-dhruva}), \vid\ \S\thinspace\ref{siddhantasindhu_sanskrit_english}:~footnote~\ref{vishava_deviant}. Another \acrshort{regular}-form is \tsans{vi.suva-dhruva} (\textit{viṣuva-dhruva}).}  [\hyperlink{SEpassA}{α}]}
}

\newglossaryentry{solstitial_colure}
{
        name={solstitial colure---},
        description={\phantom{x}\nopagebreak}
}


\newglossaryentry{circle_four_poles}
{
        name={circle passing through the four poles},
        description={\tsans{dhruva-catu.ska-yaata-v.rtta} (\textit{dhruva-catuṣka-yāta-vṛtta}) [\hyperlink{SEpassA}{α}]\newline
        identified with the \tfarsi{دایرهٔ ماره باقطاب اربعه} (\textit{\dayiri\idafavowel\ \marri\ \biaqtab\idafaconsonant\ \arbai}) `circle passing through the four poles'},
        parent={solstitial_colure}
}
 
\newglossaryentry{circle_solstice}
{       
        name={solstitial colure},
        description={\tsans{aayana-v.rtta} (\textit{āyana-vṛtta}) [\hyperlink{SEpassA}{α}, \hyperlink{SEpassB}{β}, \hyperlink{SEpassC}{γ}, \hyperlink{SEpassD}{δ}]},        
        parent={solstitial_colure}
}       

\newglossaryentry{well_rounded}
{
        name={well rounded},
        description={\tsans{su-v.rtta} (\textit{su-vṛtta}) [\hyperlink{SEpassA}{α}]}
}

\newglossaryentry{pair_equinoctial_points}
{
        name={pair of equinoctial points},
        description={\tsans{vi.suvat-dvaya} (\textit{viṣuvat-dvaya})\footnote{\,Kh attests \tsans{vi\underline{.sa}va-dvaya} (\textit{vi\underline{ṣa}va-dvaya}), an irregular form of \tsans{vi.suvat-dvaya} (\textit{viṣuvat-dvaya}), \vid\ \S\thinspace\ref{siddhantasindhu_sanskrit_english}:~footnote~\ref{vishava_deviant_2}.} [\hyperlink{SEpassA}{α}]}
}


\newglossaryentry{circle_congruent_ecliptic_parent}
{
    name={circle passing through the equinoctial points and the celestial object---},
    description={\phantom{x}\nopagebreak}
}

\newglossaryentry{circle_congruent_ecliptic}
{
        name={circle congruent to the ecliptic},
        description={\tsans{bhacakra-sad.r"sa-v.rtta} (\textit{bhacakra-sadṛśa-vṛtta}) [\hyperlink{SEpassA}{α}, \hyperlink{SEpassD}{δ}], \tsans{bhav.rtta-sad.r"sa-v.rtta} (\textit{bhavṛtta-sadṛśa-vṛtta}) [\hyperlink{SEpassB}{β}], \tsans{bhacakra-sad.rk.sa-v.rtta} (\textit{bhacakra-sadṛkṣa-vṛtta}) [\hyperlink{SEpassC}{γ}]},
        parent={circle_congruent_ecliptic_parent}
}

\newglossaryentry{knower_spheres}
{
        name={knower of spheres},
        description={\tsans{gola-vid} (\textit{gola-vid}) [\hyperlink{SEpassA}{α}]}
}
     

\newglossaryentry{conjunction_equinox_node}
{
        name={conjunction of the equinoctial point and the node of the orbit of a celestial object},
        description={\tsans{vi.suva-paata-yuga} (\textit{viṣuva-pāta-yuga})\footnote{\,Kh attests \tsans{vi\underline{.sa}va-paata-yuga} (\textit{vi\underline{ṣa}va-pāta-yuga}), an irregular of \tsans{vi.suva-paata-yuga} (\textit{viṣuva-pāta-yuga}), \vid\ \S~\ref{siddhantasindhu_sanskrit_english}:~footnote~\ref{vishava_deviant_4}.} [\hyperlink{SEpassC}{γ}]}
}

\newglossaryentry{equinoctial_point}
{
        name={equinoctial point},
        description={\tsans{vi.suvat} (\textit{viṣuvat})\footnote{\,\,Kh attests \tsans{vi\underline{.sa}vat} (\textit{vi\underline{ṣa}vat}), an irregular form of \tsans{vi.suvat} (\textit{viṣuvat}), \vid\ \S\thinspace\ref{siddhantasindhu_sanskrit_english}:~footnote~\ref{vishava_deviant_5}.} [\hyperlink{SEpassD}{δ}]}
}

\newglossaryentry{distance_equinox}
{
        name={distance along the `circle congruent to the ecliptic' from the equinox---},
        description={\phantom{x}\nopagebreak}
}

\newglossaryentry{congruent_bhuja}
{
        name={congruent arc},
        description={\tsans{sad.r"s-bhujaa} (\textit{sadṛś-bhujā}) [\hyperlink{SEpassD}{δ}], \tsans{sad.rk.sa-baahu} (\textit{sadṛkṣa-bāhu}) [\hyperlink{SEpass9}{9}]},
        parent={distance_equinox}
}      


\newglossaryentry{distance_solstitial_colure}
{
        name={distance along the `circle congruent to the ecliptic' from the solstice---},
        description={\phantom{x}\nopagebreak}
}

% UNATTESTED
% \newglossaryentry{distance_celestial_object_solstice}
% {
%         name={distance of a celestial object from the `circle passing through the four poles'},
%         description={\newline\tfarsi{بعد کوکب از
%         \tfarsib{دایرهٔ ماره باقطاب اربعه}} (\textit{\bud\idafaconsonant\ \kawkab\ \az\ \guillemotleft\dayiri\idafavowel\ \marri\ \biaqtab\idafaconsonant\ \arbai\guillemotright})},
%         parent={distance_solstitial_colure}
% }

\newglossaryentry{congruent_koti}
{
        name={congruent complementary arc},
        description={\tsans{sad.r"s-ko.ti} (\textit{sadṛś-koṭi}) [\hyperlink{SEpassD}{δ}]\newline
        identified with the \tfarsi{بعد کوکب از
        \tfarsib{دایرهٔ ماره باقطاب اربعه}} (\textit{\bud\idafaconsonant\ \kawkab\ \az\ \guillemotleft\dayiri\idafavowel\ \marri\ \biaqtab\idafaconsonant\ \arbai\guillemotright}) `distance of a celestial object from the ``circle passing through the four poles"'},
        parent={distance_solstitial_colure}
}

% UNATTESTED
% \newglossaryentry{ecliptic_distance_solstitial_colure}
% {
%         name={complement of ecliptic longitude---},
%         description={\phantom{x}\nopagebreak}
% }

% \newglossaryentry{distance_degree_celestial_object_nearest_solstice}
% {
%         name={distance of the degree of a celestial object from the nearest solstice},
%         description={\newline\tfarsi{بعد درجه کوکب از انقلاب اقرب}  (\textit{\bud\idafaconsonant\ \daraji\idafavowel\ \kawkab\ \az\ \inqilab\idafaconsonant\ \aqrab})},
%         parent={ecliptic_distance_solstitial_colure}
% }

% \newglossaryentry{arc_complementary_ecliptic_longitude_celestial_object}
% {
%         name={complement of the arc of ecliptic longitude of a celestial object},
%         description={\tsans{khagasya ko.ti} (\textit{khagasya koṭi})},
%         parent={ecliptic_distance_solstitial_colure}
% }


\newglossaryentry{reducing_from_ninety}
{       
        name={having been reduced from ninety},
        description={\tsans{navatita"scyuta} (\textit{navatitaś-cyuta}) [\hyperlink{SEpass9}{9}]}
}

\newglossaryentry{celestial_hemisphere}
{
        name={celestial hemisphere},
        description={\tsans{gola} (\textit{gola}) [\hyperlink{SEpass10}{10}]}
}

\newglossaryentry{same}
{
        name={same},
        description={\tsans{sama} (\textit{sama}) [\hyperlink{SEpass11}{11}]}
}

%%%%%%%%%%%%%%%%%%%%%%

\newglossaryentry{shadow}
{
        name={shadow of a gnomon},
        description={\tfarsi{ظلّ} (\zill) \hyperlink{Pii4}{II.4};  \tsans{chaayaa} (\textit{chāyā}) \hyperlink{Sii4}{II.4}}
}

\newglossaryentry{definition}
{
        name={definition},
        description={\tsans{lak.sa.na} (\textit{lakṣaṇa}) \hyperlink{Sii2}{II.2}, \hyperlink{Sii8}{8}--\hyperlink{Sii13}{13}, \hyperlink{Sii17}{17}--\hyperlink{Sii19}{19}, \hyperlink{Sii22}{22}},
        plural={definitions}
}


\newglossaryentry{Sine}
{
        name={Sine},
        description={\tfarsi{جيب} (\jayb) \hyperlink{Pii3}{II.3}; \tsans{jyaa} (\textit{jyā}) \hyperlink{Sii3}{II.3}}
}
% \tsans{jyakaa} (\textit{jyakā}), \tsans{si}\tsnb{ञ्जि}\tsans{nii} (\textit{siñjinī })UNATTESTED by itself

\newglossaryentry{versd_sine_parent}
{
        name={Versed Sine---},
        description={\phantom{x}\nopagebreak}
}

\newglossaryentry{Sagitta}
{
        name={Sagitta},
        description={\tfarsi{سهم} (\sahm) \hyperlink{Pii3}{II.3}},
        parent={versd_sine_parent}
} 

\newglossaryentry{Versed_Sine}
{
        name={Versed Sine},
        description={\tsans{"sara} (\textit{śara}) \hyperlink{Sii3}{II.3}},
        parent={versd_sine_parent}
}

% UNATTESTED
% \newglossaryentry{Cosine}
% {
%         name={Cosine},
%         description={\tfarsi{جيب تمام} (\jayb\idafaconsonant\ \tamam); \tsans{ko.tijiivaa} (\textit{koṭijīvā}), \tsans{ko.tijyaa} (\textit{koṭijyā}), \tsans{dyu\=/jiivaa} (\textit{dyu-jīvā}) [\lit `day-Sine']}
% }

\newglossaryentry{successive_difference}
{
        name={difference between successive rows---},
        description={\phantom{x}\nopagebreak}
}
  
\newglossaryentry{between_two_lines}
{
        name={between two lines},
        description={\tfarsi{ما بين السطرين} (\ma\ \bayn\ \alsatrayn) \hyperlink{Pii2}{II.2}},
        parent={successive_difference}
}

\newglossaryentry{difference_between_two_cells}
{
        name={difference between two cells},
        description={\tsans{dvi-ko.s.tha-antara} (\textit{dvi-koṣṭha-antara}) \hyperlink{Sii2}{II.2}},
        parent={successive_difference}
}        

\newglossaryentry{demonstration}
{
        name={demonstration},
        description={\tsans{saadhana} (\textit{sādhana}) \hyperlink{Sii2}{II.2}, \tsans{upapatti} (\textit{upapatti}) \hyperlink{Sii15}{II.15}}
}

\newglossaryentry{method_of_identity}
{
        name={method of identity},
        description={\tsans{ananyatva-prakaara} (\textit{ananyatva-prakāra}) \hyperlink{Sii15}{II.15}\newline \textit{tacitly}, an argument following a method identical to one previously stated.}
}

\newglossaryentry{method_of_interpolation}
{
        name={method of interpolation},
        description={\tfarsi{عمل تعدیل} (\amal\idafaconsonant\ \tadil), \lit `operation of adjustment' \hyperlink{Pii2}{II.2}}
}

\newglossaryentry{time_parent}
{
        name={time---},
        description={\phantom{x}\nopagebreak}
}

\newglossaryentry{desired_time}
{
        name={desired time},
        description={\tsans{abhimata-samaya} (\textit{abhimata-samaya}) \hyperlink{SiiInc}{II\thinspace\textit{incipit}}},
        parent={time_parent}
}

\newglossaryentry{time}
{
        name={time},
        description={\tfarsi{وقت} (\vaqt), pl.\thinspace \tfarsi{اوقات} (\avqat) \hyperlink{PiiInc}{II\thinspace\textit{incipit}}},
        plural={times},
        parent={time_parent}
}

\newglossaryentry{circle_of_declination}
{
        name={circle of declination},
        description={\tsans{kraanti-suutra} (\textit{krānti-sūtra}) \hyperlink{Sii5}{II.5}}
}        
       

\newglossaryentry{declination_parts_ecliptic}
{
        name={declination of parts of the ecliptic},
        description={\tfarsi{ميل اجزاء فلك البروج}  (\mayl\idafaconsonant\ \ajza\idafaconsonant\ \falak\ \alburuj) \hyperlink{Pii5}{II.5}}
}

\newglossaryentry{maximum_elevation_depression_celestial_object}
{
        name={maximum elevation and depression of a celestial object},
        description={\newline\tfarsi{غایت ارتفاع و انخفاض كوكب} (\ghayat\idafaconsonant\ \irtifa\ \va\ \inkhifad\idafaconsonant\ \kawkab) \hyperlink{Pii7}{II.7}},
        plural={maximum elevation and depression of celestial objects}
}

\newglossaryentry{degrees_maximum_elevation}
{
        name={degrees of the maximum elevation},
        description={\tsans{parama-unnata-a.m"sa} (\textit{parama-unnata-aṃśa}) \hyperlink{Sii7}{II.7}},
}        

\newglossaryentry{degrees_maximum_depression}
{
        name={degrees of the maximum depression},
        description={\tsans{adha.hstha-parama-bhaaga} (\textit{adhaḥstha-parama-bhāga}) \hyperlink{Sii7}{II.7}}
}

\newglossaryentry{elevation_parent}
{
        name={elevation---},
        description={\phantom{x}\nopagebreak}
}

\newglossaryentry{elevation}
{
        name={elevation},
        description={\tfarsi{ارتفاع} (\irtifa) \hyperlink{Pii14}{II.14}, \hyperlink{Pii15}{15}, \hyperlink{Pii21}{21}, \hyperlink{Pii22}{22}},
        parent={elevation_parent}
}

\newglossaryentry{desired_degree_elevation}
{
        name={desired degrees of elevation},
        description={\tsans{abhiipsita-unnata-a.m"sa} (\textit{abhīpsita-unnata-aṃśa}) \hyperlink{Sii14}{II.14},
         \tsans{abhii.s.ta-unnata-a.m"sa} (\textit{abhīṣṭa-unnata-aṃśa}) \hyperlink{Sii15}{II.15}, \hyperlink{Sii21}{21}, \hyperlink{Sii22}{22}},
        parent={elevation_parent}
}

\newglossaryentry{depression_parent}
{
        name={depression---},
        description={\phantom{x}\nopagebreak}
}

\newglossaryentry{depression}
{
        name={depression},
        description={\tfarsi{انخفاض} (\inkhifad) \hyperlink{Pii14}{II.14}, \hyperlink{Pii22}{22}},
        parent={depression_parent}
}

\newglossaryentry{degrees_depression}
{
        name={degrees of depression},
        description={\tsans{adhara-a.m"sa} (\textit{adhara-aṃśa}) \hyperlink{Sii14}{II.14}, \hyperlink{Sii15}{15}, \tsans{adhara-a.m"saka} (\textit{adhara-aṃśaka}) \hyperlink{Sii22}{II.22}},
        parent={depression_parent}
}

\newglossaryentry{right_ascension}
{
        name={right ascension of zodiacal signs in degrees---},
        description={\phantom{x}\nopagebreak}
}

\newglossaryentry{ascension_line_terrestrial_observer}
{
        name={ascensions [of the ecliptic] at the line of the terrestrial equator},
        description={\newline \tfarsi{مطالع خطّ استوا} (\matali\ \khatt\idafaconsonant\ \istiva)
        \hyperlink{Pii8}{II.8}},
        parent={right_ascension}
}

\newglossaryentry{ascension_right_sphere}
{
        name={ascensions [of the ecliptic] in the right sphere},
        description={\tfarsi{مطالع فلك مستقیم} (\matali\ \falak\idafaconsonant\ \mustaqim) \hyperlink{Pii8}{II.8}},
        parent={right_ascension}
}

\newglossaryentry{rising_terrestrial_equator_degrees}
{
        name={rising [of the zodiacal signs] at the terrestrial equator in degrees},
        description={\tsans{vyak.sa-udaya-a.m"sa} (\textit{vyakṣa-udaya-aṃśa}) \hyperlink{Sii8}{II.8}},
        parent={right_ascension}
}

\newglossaryentry{rising_Lanka_degrees}
{
        name={rising [of the zodiacal signs] at Laṅkā in degrees},
        description={\tsans{la"nkaa-udaya-a.m"sa} (\textit{laṅkā-udaya-aṃśa}) \hyperlink{Sii8}{II.8}},
        parent={right_ascension}
}

\newglossaryentry{ascendant_rising_point}
{
        name={ascendant or rising zodiacal sign---},
        description={\phantom{x}\nopagebreak}
}

\newglossaryentry{ascendant}
{
        name={ascendant},
        description={\tfarsi{طالع} (\tali), pl.\thinspace \tfarsi{طوالع} (\tawali) \hyperlink{Pii11}{II.11}, \hyperlink{Pii21}{21}, \hyperlink{Pii22}{22}; \tsans{lagna} (\textit{lagna}) \hyperlink{Sii21}{II.21}},
        plural={ascendants},
        parent={ascendant_rising_point}
}

\newglossaryentry{ascendant_each_time}
{
        name={ascendant at each time},
        description={\tfarsi{طالع هر وقت} (\tali\ \har\ \vaqt) \hyperlink{PiiInc}{II\thinspace\textit{incipit}}},
        parent={ascendant_rising_point}
}

\newglossaryentry{ascendant_that_time}
{
        name={ascendant at that time},
        description={\tsans{taatkaalika-lagna} (\textit{tātkālika-lagna}) \hyperlink{SiiInc}{II\thinspace\textit{incipit}}},
        parent={ascendant_rising_point}
}

\newglossaryentry{ascendant_ecliptic_degrees}
{
        name={ecliptic degrees of the ascendants},
        description={\tsans{vilagna-a.m"saka} (\textit{vilagna-aṃśaka}) \hyperlink{Sii11}{II.11}},
        parent={ascendant_rising_point}
}


\newglossaryentry{rising_parent}
{
        name={ascension or rising ---},
        description={\phantom{x}\nopagebreak}
}

\newglossaryentry{ascension_measure}
{
        name={[measure of] ascension},
        description={\tfarsi{مطلع} (\matla), pl.\thinspace \tfarsi{مطالع} (\matali) \hyperlink{Pii11}{II.11}, \hyperlink{Pii13}{13}},
        user1={ascensions},
        parent={rising_parent}
}

\newglossaryentry{rising}
{
        name={rising},
        description={\tfarsi{طالع} (\tali), pl.\thinspace \tfarsi{طوالع} (\tawali) \hyperlink{Pii13}{II.13}; \tsans{udaya} (\textit{udaya}) \hyperlink{Sii12}{II.12}},
        parent={rising_parent}
}

\newglossaryentry{time_rising}
{
        name={time of rising},
        description={\tsans{udaya-samaya} (\textit{udaya-samaya}) \hyperlink{Sii13}{II.13}},
        parent={rising_parent}
}
\newglossaryentry{setting_parent}
{
        name={descension or setting---},
        description={\phantom{x}\nopagebreak}
}

\newglossaryentry{setting}
{
        name={setting},
        description={\tfarsi{غروب} (\ghurub) \hyperlink{Pii13}{II.13}},
        parent={setting_parent}
}
%  \tsans{asta} (\textit{asta}) UNATTESTED (=setting)

\newglossaryentry{time_setting}
{
        name={time of setting},
        description={\tsans{asta-samaya} (\textit{asta-samaya}) \hyperlink{Sii13}{II.13}},
        parent={setting_parent}
}

\newglossaryentry{oblique_ascension}
{
        name={oblique ascension zodiacal signs in degrees---},
        description={\phantom{x}\nopagebreak},
}

\newglossaryentry{ascension_ecliptic_locality}
{
        name={ascensions [of the ecliptic] of a locality},
        description={\tfarsi{مطالع بلد} (\matali\idafaconsonant\ \balad) \hyperlink{Pii10}{II.10}},
        parent={oblique_ascension}
}

\newglossaryentry{rising_location_degree}
{
        name={rising [of the zodiacal signs] in one's own location in degrees},
        description={\tsans{nija-udaya-a.m"sa} (\textit{nija-udaya-aṃśa}) \hyperlink{Sii10}{II.10}, \tsans{sva\=/udaya-a.m"sa} (\textit{sva\=/udaya-aṃśa}) \hyperlink{Sii11}{II.11}, \hyperlink{Sii22}{22}, \tsans{nija-udaya-a.m"saka} (\textit{nija-udaya-aṃśaka}) \hyperlink{Sii13}{II.13}},
        parent={oblique_ascension}
}


\newglossaryentry{equation_of_daylight_parent}
{
        name={equation of daylight---},
        description={\phantom{x}\nopagebreak}
}

\newglossaryentry{equation_of_daylight}
{
        name={equation of daylight},
        description={\tfarsi{تعديل النها} (\tadil\ \alnahar) \hyperlink{Pii9}{II.9}},
        parent={equation_of_daylight_parent}
}

\newglossaryentry{ascensional_difference}
{
        name={ascensional difference},
        description={\tsans{cara} (\textit{cara}) \hyperlink{Sii9}{II.9}},
        parent={equation_of_daylight_parent}
}
        
\newglossaryentry{arc_of_daylight}
{
        name={arc of daylight},
        description={\tfarsi{قوس النهار} (\qaws\ \alnahar) \hyperlink{Pii9}{II.9}}
}

\newglossaryentry{arc_of_night}
{
        name={arc of night},
        description={\tfarsi{قوس الليل} (\qaws\ \allayl) \hyperlink{Pii9}{II.9}}
}

\newglossaryentry{hours_of_daylight}
{
        name={hours of daylight},
        description={\tfarsi{ساعات النهار} (\saat\ \alnahar) \hyperlink{Pii9}{II.9}}
}

\newglossaryentry{hours_of_night}
{
        name={hours of night},
        description={\tfarsi{ساعات الليل} (\saat\ \allayl) \hyperlink{Pii9}{II.9}}
}

\newglossaryentry{hours_of_day_and_night}
{
        name={hours of day and night},
        description={\tsans{dina-raatri-horaa} (\textit{dina-rātri-horā}) \hyperlink{Sii9}{II.9}}
}

\newglossaryentry{oblique_diurnal_circle}
{
        name={oblique diurnal circle},
        description={\tsans{dina-raatri-vaama-v.rtta} (\textit{dina-rātri-vāma-vṛtta}) \hyperlink{Sii9}{II.9}}
}


\newglossaryentry{inverse_procedure}
{
        name={inverse procedure},
        description={\tsans{viloma-kriyaa} (\textit{viloma-kriyā}) \hyperlink{Sii11}{II.11}}
}

\newglossaryentry{inverse_method}
{
        name={inverse method},
        description={\tfarsi{عمل عکس} (\amal\idafaconsonant\ \aks) \hyperlink{Pii11}{II.11}}
}

\newglossaryentry{degree_transit_parent}
{
        name={ecliptic longitude of the meridian ecliptic point---},
        description={\phantom{x}\nopagebreak}
}

\newglossaryentry{degree_transit_celestial_object}
{       name={[ecliptic] degree of the [meridian] transit [at the time of rising] of a celestial object},
        description={\tfarsi{درجهٔ ممرّ كوكب} (\daraji\idafavowel\ \mamarr\idafaconsonant\ \kawkab) \hyperlink{Pii12}{II.12}},
        sort={ecliptic degree of meridian transit at the time of rising of a celestial object},
        parent={degree_transit_parent}
}

\newglossaryentry{degree_meridian_ecliptic_rising_time}
{       name={degrees of the [meridian] ecliptic point at the [time of] rising of a celestial object},
        description={\tsans{bha-udaya-lagna-a.m"sa} (\textit{bha-udaya-lagna-aṃśa}) \hyperlink{Sii12}{II.12}},
        parent={degree_transit_parent}
}        

\newglossaryentry{co_ascension}
{
        name={right ascension of the meridian ecliptic point---},
        description={\phantom{x}\nopagebreak}
}

\newglossaryentry{ascension_transit}
{
        name={ascensions of [the degrees] of [meridian] transit},
        description={\tfarsi{مطالع  ممرّ} (\matali\idafaconsonant\ \mamarr) \hyperlink{Pii12}{II.12}},
        parent={co_ascension}
}

\newglossaryentry{degree_equatorial_ascension}
{
        name={degrees of equatorial ascension of the [meridian] ecliptic point at the [time of] rising of a celestial object},
        description={\tsans{bha-udaya-lagna-vyak.sa-udaya-a.m"sa} (\textit{bha-udaya-lagna-vyakṣa-udaya-aṃśa}) \hyperlink{Sii12}{II.12}},
        parent={co_ascension}
}


\newglossaryentry{rule_of_three}
{
        name={rule of three},
        description={\tsans{trai-raa"sika} (\textit{trai-rāśika}) \hyperlink{Sii2}{II.2}}
}

\newglossaryentry{correlated_numbers}
{
        name={correlated numbers},
        description={\tsans{paraspara-sambandhi-raa"si} (\textit{paraspara-sambandhi-rāśi}) \hyperlink{Sii2}{II.2}}
}

\newglossaryentry{azimuth_degrees}
{
        name={azimuth---},
        description={\phantom{x}\nopagebreak}
}

\newglossaryentry{azimuth}
{
        name={azimuth},
        description={\tfarsi{سمت} (\samt) \hyperlink{Pii14}{II.14}, \hyperlink{Pii15}{15}},
        parent={azimuth_degrees}
}

\newglossaryentry{azimuth_own_location}
{
        name={degrees of azimuth in one's own location},
        description={\tsans{sva-di"s-a.m"sa} (\textit{sva-diś-aṃśa}) \hyperlink{Sii14}{II.14}, \hyperlink{Sii15}{15}},
        parent={azimuth_degrees}
}


\newglossaryentry{degree_azimuth}
{
        name={degree of azimuth},
        description={\tsans{di"s-a.m"sa} (\textit{diś-aṃśa}) \hyperlink{Sii20}{II.20}},
        parent={azimuth_degrees}
}

\newglossaryentry{meridian_parent}
{
        name={local meridian---},
        description={\phantom{x}\nopagebreak}
}

\newglossaryentry{line_midday}
{
        name={line of midday},
        description={\tfarsi{خطّ نصف النهار} (\khatt\ \nisf\ \alnahar) \hyperlink{Pii16}{II.16}; \tsans{madhyaahna-rekhaa} (\textit{madhyāhna-rekhā}) \hyperlink{Sii16}{II.16}},
        parent={meridian_parent}
}

\newglossaryentry{meridian_line}
{
        name={line of the meridian},
        description={\tsans{yaamya-uttara-rekhaa} (\textit{yāmya-uttara-rekhā}) \hyperlink{Sii16}{II.16}},
        parent={meridian_parent}
}

\newglossaryentry{north_south_direction}
{
        name={north-south direction},
        description={\tsans{saumya-yaamya-di"s} (\textit{saumya-yāmya-diś}) \hyperlink{Sii20}{II.20}},
        user1={northern and southern directions}
}

\newglossaryentry{local_long_lat_parent}
{
        name={local terrestrial coordinates---},
        description={\phantom{x}\nopagebreak}
}

\newglossaryentry{longitude_and_latitude_terrestrial}
{
        name={longitude and latitude of a locality},
        description={\tfarsi{طول و عرض بلد} (\tul\ \va\ \ard\idafaconsonant\ \balad) \hyperlink{Pii17}{II.17}},
        parent={local_long_lat_parent}
}

\newglossaryentry{degrees_terrestrial_long_latitude}
{
        name={degrees of [terrestrial] longitude and latitude},
        description={\tsans{de"saantara-ak.sa-a.m"sa} (\textit{deśāntara-akṣa-aṃśa}) \hyperlink{Sii17}{II.17}},
        parent={local_long_lat_parent}
}        

\newglossaryentry{altitude_ecliptic_pole}
{
        name={altitude of the ecliptic pole or nonagesimal point---},
        description={\phantom{x}\nopagebreak}
}

\newglossaryentry{latitude_visible_climate}
{
        name={latitude of the visible climate},
        description={\tfarsi{عرض اقلیم رؤیت}  (\ard\idafaconsonant\ \iqlim\idafaconsonant\ \ruyat) \hyperlink{Pii18}{II.18}},
        parent={altitude_ecliptic_pole}
}

\newglossaryentry{ecliptic_zenith_distance}
{
        name={zenith distance of the nonagesimal point},
        description={\tsans{d.rkk.sepa} (\textit{dṛkkṣepa}) \hyperlink{Sii18}{II.18}},
        parent={altitude_ecliptic_pole}
}


\newglossaryentry{ecliptic_pole_zenith_distance}
{
        name={zenith distance of the ecliptic pole},
        description={\tsans{d.rggati} (\textit{dṛggati}) \hyperlink{Sii18}{II.18}},
        parent={altitude_ecliptic_pole}
}

\newglossaryentry{distance_celestial_object_parent}
{
        name={distance between celestial objects---},
        description={\phantom{x}\nopagebreak}
}

\newglossaryentry{distance_between_celestial_object}
{
        name={distance between two celestial objects},
        description={\tfarsi{بعد میان دو کوکب} (\bud\idafaconsonant\ \miyan\idafaconsonant\ \mbox{\duvum\idafaconsonant} \kawkab) \hyperlink{Pii19}{II.19}},
        parent={distance_celestial_object_parent}
}

\newglossaryentry{degrees_separation_two_celestial_objects}
{
        name={degrees [of separation] between two celestial objects},
        description={\tsans{dvi-nak.satra-antara-a.m"saka} (\textit{dvi-nakṣatra-antara-aṃśaka}) \hyperlink{Sii19}{II.19}},
        parent={distance_celestial_object_parent}
}

\newglossaryentry{determination}
{
        name={determination},
        description={\tfarsi{استخراج} (\istikhraj), \lit `bringing out or extraction' \hyperlink{Pii19}{II.19}}
}

\newglossaryentry{azimuth_qibla}
{
        name={azimuth of \qibla},
        description={\tfarsi{سمت قبله} (\samt\idafaconsonant\ \qibla) \hyperlink{Pii20}{II.20}}
}

% \newglossaryentry{inclination_parent}
% {
%         name={inclination---},
%         description={\phantom{x}\nopagebreak}
% }

\newglossaryentry{inclination_azimuth_qibla}
{
        name={inclination of the azimuth of \qibla},
        description={\tfarsi{انحراف سمت قبله} (\inhiraf\idafaconsonant\ \samt\idafaconsonant\ \qibla) \hyperlink{Pii20}{II.20}}
        % parent={inclination_parent}
}

% \newglossaryentry{inclination}
% {
%         name={inclination},
%         description={\tfarsi{انحراف} (\inhiraf)},
%         parent={inclination_parent}
% }

\newglossaryentry{genus_parent}
{
        name={place-value of sexagesimal digits---},
        description={\phantom{x}\nopagebreak}
}

\newglossaryentry{genus}
{
        name={genus},
        description={\tfarsi{جنس} (\jins) \hyperlink{Pii1}{II.1}},
        parent={genus_parent}
}

\newglossaryentry{belonging_genus}
{
        name={belong to a particular genus},
        description={\tsans{jaatiiya} (\textit{jātīya}) \hyperlink{Sii1}{II.1}},
        parent={genus_parent}
}


\newglossaryentry{square_root}
{
        name={square root},
        description={\tfarsi{جذر} (\jadr) \hyperlink{Pii1}{II.1}; \tsans{muula} (\textit{mūla}) \hyperlink{Sii1}{II.1}}
}

\newglossaryentry{position}
{
        name={position},
        description={\tfarsi{مرتبه} (\martaba), pl.\thinspace \tfarsi{مراتب} (\maratib) \hyperlink{Pii1}{II.1};  \tsans{sthaana} (\textit{sthāna}) \hyperlink{Sii1}{II.1}},
        plural={positions}        
}

\newglossaryentry{raising_up_parent}
{
        name={integer-numbers of revolution---},
        description={\phantom{x}\nopagebreak}
}

\newglossaryentry{elevation_sexagesimal}
{
        name={elevated [rank]},
        description={\tfarsi{مرفوع} (\marfu), pl.\thinspace \tfarsi{مرفوعات} (\marfuat), \lit `raised up' \hyperlink{Pii1}{II.1}},
        plural={elevated [ranks]},
        parent={raising_up_parent}
}

\newglossaryentry{revolution}
{
        name={revolution},
        description={\tsans{parivarta} (\textit{parivarta}) \hyperlink{Sii1}{II.1}},
        parent={raising_up_parent}
}

\newglossaryentry{degrees_parent}
{
        name={degree---},
        description={\phantom{x}\nopagebreak}
}

\newglossaryentry{degree}
{
        name={degree},
        description={\tfarsi{درج} (\daraj) \hyperlink{Pii1}{II.1}},
        parent={degrees_parent}
}

\newglossaryentry{parts_of_degree}
{
        name={fractional parts of a degree},
        description={\tfarsi{اجزاء درج} (\ajza\idafaconsonant\ \daraj) \hyperlink{Pii1}{II.1}},
        parent={degrees_parent}
}

\newglossaryentry{minute}
{
        name={minute},
        description={\tfarsi{دقيقه} (\daqiqa), pl.\thinspace \tfarsi{دقایق} (\daqaiq) \hyperlink{Pii1}{II.1}; \tsans{kalaa} (\textit{kalā}) \hyperlink{Sii1}{II.1}},
        plural={minutes}
}

\newglossaryentry{second}
{
        name={second},
        description={\tfarsi{ثانيه} (\thaniya), pl.\thinspace \tfarsi{ثواني} (\thawani) \hyperlink{Pii1}{II.1}},
        plural={seconds}
}


\newglossaryentry{digit}
{
        name={digit},
        description={\tsans{a"nka} (\textit{aṅka}) \hyperlink{Sii1}{II.1}},
        plural={digits}
}
\newacronym{predicative}{PRED}{predicative}
\newacronym{prepositional_phrase}{PP}{prepositional phrase}
\newacronym{compound}{COMP}{compound}
\newacronym{modifier}{MOD}{modifier}
\newacronym{nominative}{NOM}{nominative case}
\newacronym{accusative}{ACC}{accusative case}
\newacronym{genitive}{GEN}{genitive case}
\newacronym{dative}{DAT}{dative case}
\newacronym{ablative}{ABL}{ablative case}
\newacronym{locative}{LOC}{locative case}
\newacronym{optative}{OPT}{optative mood}
\newacronym{imperative}{IMP}{imperative mood}
\newacronym{noun_phrase}{NP}{noun phrase}
\newacronym{noun}{NOUN}{noun}
\newacronym{singular}{SING}{singular}
\newacronym{plural}{PL}{plural}
\newacronym{subject_fronting}{SBJ-FRNT}{subject-fronting}
\newacronym{infinitive}{INF}{infinitive}
\newacronym{imperfect}{IMPERF}{past imperfective}
\newacronym{gerund}{GER}{gerund}
\newacronym{gerundive}{GDV}{gerundive}
\newacronym{participle}{PTCP}{participle}
\newacronym{indicative}{IND}{indicative mood}
\newacronym{present}{PRES}{present tense}
\newacronym{past}{PAST}{past tense}
\newacronym{future}{FUT}{future tense}
\newacronym{passive}{PASS}{passive voice}
\newacronym{active}{ACT}{active voice}
\newacronym{middle}{MID}{middle voice}
\newacronym{verb}{VB}{verb}
\newacronym{class}{CL}{class}
\newacronym{third}{3\textsuperscript{rd}}{third person}
\newacronym{first}{1\textsuperscript{st}}{first person}
\newacronym{masculine}{MASC}{masculine gender}
\newacronym{feminine}{FEM}{feminine gender}
\newacronym{neuter}{NEUT}{neuter gender}
\newacronym{subjunctive}{SBJV}{subjunctive mood}
\newacronym{stem}{STEM}{stem}
\newacronym{preverb}{PVERB}{preverb}
\newacronym{negation}{NEG}{negative form (negation)}
\newacronym{causative}{CAUS}{causative verb}
\newacronym{denominative}{CAUS}{denominative verb}
\newacronym{indeclinable}{INDECL}{indeclinable}
\newacronym{irregular}{IRREG}{irregular form}
\newacronym{regular}{REG}{regular form}
\newacronym{primary_nominal_derivative}{NOM-D$_\text{1}$}{primary nominal derivative}
\newacronym{secondary_nominal_derivative}{NOM-D$_\text{2}$}{secondary nominal derivative}
\newacronym{adjective}{ADJ}{adjective}
\newacronym{action}{ACTION}{action}

\renewcommand{\glsnamefont}[1]{\textbf{#1}}
\renewcommand*{\glstextformat}[1]{\textbf{#1}}


\renewcommand*{\glssubentrycounterlabel}{$\hookrightarrow$\space} % Allows the symbol (hooked right arrow) for dependent/similar entries in to appear when using index-style glossary 
% \let\glsalias\glsuseri
\glsaddall[types={main}] % Adds all main (glossary_shortened) entries but not all acronyms (only those acronyms used in the text are added)

%%%%%%%%%%%%%%%%%%%%%%%%%%%%%%%%%%%%%%%

% This is a part of polyglossia but as it text written in the reverse direction, it needs to be loaded after hyperref else there are conflicts

\setotherlanguage[numerals=eastern]{persian}
% ~~~~~~~~~~~~~~~~~~~~~~~~~~~~~~~~~~~~~~~~~~~~~~~~~~~~~~~~~~~~~~~%
\newfontfamily\persianfont[Path=Fonts/,Script=Arabic,Scale=1.1,WordSpace=2]{XB_Niloofar.ttf}% For typing Persian
\newcommand{\tfarsi}[1]{\textpersian{#1}} % Text in Farsi/Persian

\newfontfamily\persianbold[Path=Fonts/,Script=Arabic,Scale=1.1,WordSpace=2]{XB_NiloofarBd.ttf}% For typing Persian (Bold)
\newcommand{\tfarsib}[1]{{\persianbold{#1}}} % Text in Farsi/Persian Bold

\newfontfamily\persiannastlq[Path=Fonts/,Script=Arabic,Scale=1.1,WordSpace=2, Contextuals=Swash]{IranNastaliq.ttf}%
% For Persian Nasta'līq font
\newcommand{\tfarsinastlq}[1]{{\persiannastlq{#1}}} 

% ~~~~~~~~~~~~~~~~~~~~~~~~~~~~~~~~~~~~~~~~~~~~~~~~~~~~~~~~~~~~~~~%
%For preventing line break at hyphen used in Romanised Arabic & Sanskrit words (Used in sanskrit_persian_words.tex
\usepackage[shortcuts]{extdash} % Load last, it freaks out when other things follow it.
\pdfstringdefDisableCommands{\def\-{-}} % Fixes the "\EXD@break" hyperref warnings from popping up
% Standard LaTeX dashes
% - Standard LaTeX hyphen

% extdash breakable dashes
% Words hyphened with these dashes can also be broken at other positions than the dash, i.e., in the pre word and the post word
% \-/ hyphen


% extdash unbreakable dashes
% No line breaks possible at the hyphen, i.e., only breaks in the post word allowed, no pre-hyphen word break allowed: Use \=/ for words like al\=/Birūnī which prevents an orphan "al-" at the end of the line
% \=/ hyphen

% ~~~~~~~~~~~~~~~~~~~~~~~~~~~~~~~~~~~~~~~~~~~~~~~~~~~~~~~~~~~~~~~%


\addtocounter{secnumdepth}{2} % For numbered subsections 
\setcounter{tocdepth}{2}
\usepackage{chngcntr}
\counterwithout{figure}{subsection}
\counterwithout{table}{subsection}
\counterwithout{section}{chapter} % Resets section numbering to 1, 2, etc.
\hangsecnum



\title{Persian Astronomy in Sanskrit\\[5pt]
A comparative study of \MullaFarid's \ZijiShahJahani, Discourse II (\circa 1629/30 \ce) and its Sanskrit translation in \Nityananda's \Siddhantasindhu\ (\circa early 1630s \ce)}
% \MullaFarid's \ZijiShahJahani, Discourse II, Chapter VI (1628 \ce) and its translation in \Nityananda's \Siddhantasindhu\ (\circa early 1630 \ce)}

\author{Anuj Misra\\[10pt] 

{\small Department of Cross-Cultural and Regional Studies,}\\{\small University of Copenhagen, Denmark.}\\{\small \url{anuj.misra@hum.ku.dk}}}
\date{}

\OnehalfSpacing
\setaftersecskip{\baselineskip}% For slightly increased space after section title 
\pagestyle{plain}

\begin{document}
\sloppy
\sloppybottom
\maketitle
\thispagestyle{empty}
\vfill
This study is a part of my research project on `Early Modern Exchanges in Sanskrit Astral Sciences' (EMESAS). This project has received funding from the European Union’s Horizon 2020 research and innovation programme under the Marie Skłodowska-Curie grant agreement \textnumero\ 836055.
\clearpage


\thispagestyle{empty}

\begin{abstract}
Starting from the late medieval period of Indian history, Islamicate and Sanskrit astral sciences exchanged ideas in complex discourses shaped by the power struggles of language, culture, and identity. The practice of translation played a vital role in transporting science across the physical and mental realms of an ever-changing society. The present study begins by looking at the culture of translating astronomy in late-medieval and early-modern India. This provides the historical context to then examine the language with which \Nityananda, a seventeenth-century Hindu astronomer at the Mughal court of Emperor \Shahjahan, translated into Sanskrit the Persian astronomical text of his Muslim colleague \MullaFarid. \Nityananda's work is an example of how secular innovation and sacred tradition expressed themselves in Sanskrit astral sciences.

This article includes a comparative description of the contents in the second discourse of \MullaFarid's \ZijiShahJahani\ (\circa 1629/30) and the second part of \Nityananda's \Siddhantasindhu\ (\circa early 1630s), along with a critical examination of the sixth chapter from both these works. The chapter-titles and the contents of the sixth chapter in Persian and Sanskrit are edited and translated into English for the very first time. The focus of this study is to highlight the linguistic (syntactic, semantic, and communicative) aspects in \Nityananda's Sanskrit translation of \MullaFarid's Persian text. The mathematics of the chapter is discussed in a forthcoming publication. An indexed glossary of technical terms from the edited Persian and Sanskrit text is appended at the end of the work.
\end{abstract}
\vfill
 \paragraph{Keywords} Sanskrit astral sciences, Persian astronomy, true declination, science in Mughal India, culture of translating science, \ganita\ \jyotisa, \Nityananda, \MullaFarid, \Sarvasiddhantaraja, \Siddhantasindhu, \ZijiShahJahani, \ZijUlughBeg. 
\clearpage
\tableofcontents*
\clearpage

% Main matter 
\section{Introduction}\label{introduction}
\renewcommand{\thefootnote}{\arabic{footnote}} % Hindu-Arabic footnote numbering

Over the course of the history of Sanskrit mathematical astronomy (\ganita\ \jyotisa), foreign ideas have evoked a full range of emotions that extend from affinity to apathy, going all the way to antipathy. These reactions are a reflection of the intellectual diversity of Indian astral sciences (\jyotihsastra). Historical actors may have chosen to accept, reject, or ignore foreign ideas based on their scientific convictions; however, those choices could only be expressed under the aegis of the political and sociocultural institutions of the times. 

With the turn of the seventeenth century \ce, Sanskrit astronomers/astrologers (\jyotisa s) and their Persianate counterparts (\munajjim s) were working under the common patronage of the imperial court of Mughal India.\footnote{\,The \GurkaniAlam\ or the Mughal Empire was an early-modern Muslim empire in South Asia led by monarchs of the Timurid dynasty. From 1526 to 1857 \ce, the successors of \Baburfull, the first Mughal Emperor, extended their dominion over large swathes of the Indian subcontinent, and in doing so, helped create a highly complex cosmopolitan society extending beyond the imperial borders. I refer to this cultural sphere of influence of the Mughal rule as Mughal India.} At the court of Emperor \Shahjahan\ (r.\thinspace 1628--58), we find the Gauḍa Brahmin \Pandita\ \Nityanandafull\ (\fl 1630/50) working alongside \MullaFaridfull\ (d.\thinspace \circa 1629/32; henceforth identified as \MullaFarid) to translate into Sanskrit the latter's Persian \zij\ (set of astronomical tables) the  \ZijiShahJahani\ (\circa 1629/30). \Nityananda's Sanskrit translation, the \Siddhantasindhu\ (\circa early 1630), was his first attempt at explaining Islamicate\footnote{\,I use the word Islamicate (instead of Islamic) to indicate the cultural outputs (\eg artistic, literary, or scientific works) of Muslim societies educated in the Arabic and Persian language traditions but not directly connected to the Islamic faith or any particular geographic region \parencite[\vid\ recent discussions on Islamicate Secularities in][]{Dressleretal}.\label{islamicate_defintion}} computations and astronomical tables to his fellow Sanskrit \jyotisa s. 

By the end of the decade, he included several of these Islamicate ideas in his canonical treatise the \Sarvasiddhantaraja\ `The King of all \Siddhanta s' (1639). The \Sarvasiddhantaraja\ is composed in the style of a traditional Sanskrit \textit{siddhānta} (a canonical treatise in astronomy) and has a tripartite structure: the \ganitadhyaya\ `chapter on computations', the \goladhyaya\ `chapter on spheres', and the \yantradhyaya\ `chapter on instruments'.\footnote{\,\textcite[Sections~1.1 and 1.2 on pp.\thinspace 1--20]{Misrathesis} offers a fuller discussion on Sanskrit astronomy in early-modern India, in particular, the contribution of \Nityananda\ and his \Sarvasiddhantaraja. Also, contemporary studies like \textcite{PingreeSSR}, \textcite{MonetelleSine1}, and \textcite{MonetelleSine2} discuss Islamicate influences in the mathematical computations described in the \Sarvasiddhantaraja.} In contrast, the \Siddhantasindhu\ mimics the structure and content of the Persian \ZijiShahJahani\ quite intimately. 

In this study, I compare the general structure of the \ZijiShahJahani\ and the \Siddhantasindhu\ in parallel, and subsequently focus on a chapter from each of these works that discusses the same topic, \viz the declination of a celestial object. My aim is to highlight the semantic and communicative aspects in \Nityananda's Sanskrit translation of \MullaFarid's Persian text. I defer all remarks on the mathematics in \Nityananda's text to \textcite[forthcoming]{MisraTD}. Instead, I first begin by discussing the practice of translating Sanskrit, Arabic, or Persian astronomical texts during the late-medieval and early-modern periods of Indian history. This overview, built from separate studies on the history, philosophy, and language of astral sciences in India,\footnote{\,For example, \vid\ \textcites{PingreeIslamAstSkt, PingreeIndianReception, Ansaritransmission, Ansarihindu, Ohasi} for historical accounts of Sanskrit and Persian astronomy in India; \textcites{Choudhuri, Minkowskilearnedbrahmin} for surveys of Sanskritic (Hindu and Jain) scholars under Muslim patronage; and \textcite{Truschke, Minkowski, MinkowskiSuryadasa, Nair} for linguistic and philosophical reforms affecting Sanskrit \jyotihsastra\ in early-modern India.} gives us the context to situate \Nityananda's works in the world of seventeenth-century Mughal India. His writings can then be seen as an ongoing dialogue between different scientific traditions in a changing society, instead of simply being judged as a `failure' and an `elaborate apology for using Muslim astronomy' \parencite[270]{PingreeSSR}. 


%%%%%%%%%%%%%%%%%%%%%%%%%%%%%%%%%%%%%%%%%%%%%%%%%%%%%%%%%%%%%%%%%%%%%%%%%%%%%%%%%%%%%%%%%
\subsection[Translating the astral sciences in Persianate India]{Translating the astral sciences in Persianate India\footnote{\,The word Persianate refers to a sociocultural association with the Persian language (Fārsī) extending beyond the ethnic identity and geographical boundaries of Persia (much like the word Islamicate distinguishes itself from Islamic, \vid\ footnote~\ref{islamicate_defintion}). I use the expression `Persianate India' to refer the geographical regions of late-medieval and early-modern India where Persian culture (expressed in its art, language, literature, and science) directly influenced society \parencite[\vid][]{Eaton2019}.}}\label{translating_the_astral_sciences_in_persianate_india}
%%%%%%%%%%%%%%%%%%%%%%%%%%%%%%%%%%%%%%%%%%%%%%%%%%%%%%%%%%%%%%%%%%%%%%%%%%%%%%%%%%%%%%%%%

\subsubsection{Before the Mughal court} \label{pre_Mughal_court}
Sanskrit texts on astronomical instruments (\yantra) written in the late fourteenth century offer some of the earliest extant evidence of a relationship between Islamicate and Sanskrit mathematical astronomy.\footnote{\,Sanskrit \tajika\ texts are Indian adaptations of Islamicate astrology that were composed from the thirteenth century \ce. \textcite{PingreeAstralOmens} provides a historical summary of the \tajika\ literature in Sanskrit, while \textcite{Gansten} studies the transmission of Perso-Arabic \tazig-astrology in the \Karmaprakasa\ (\circa 1274) of \Samarasimha, the earliest preserved Sanskrit \tajika\ work.\label{tajika_footnote}} \textcite{SarmaAstrolabeinSanskrit} provides an excellent overview of Sanskrit texts on astrolabes, many of which include lengthy discussion on Islamicate mathematical astronomy. \MahendraSuri's \Yantraraja\ (1370), along with his student \MalayenduSuri's commentary on it (in 1382), is the earliest and most recognised of such works \parencite{PlofkerAstrolabe}. 

\MahendraSuri\ was a Jain monk-astronomer at the court of \FiruzShah\ (r.\thinspace 1351--88), a Turko-Indian ruler of the pre-Mughal Sultanate of Delhi, and is thought to have worked in close association with `unnamed Muslim astronomers at Fīrūz's court' \parencite[148]{SarmaAstrolabeinSanskrit}. \citeauthor{SarmaAstrolabeinSanskrit} qualifies him as a `mediator between the Islamic and Sanskritic tradition of learning' (p.\thinspace 149). This is perhaps justly so, as three hundred years later,  the language and structure of \MahendraSuri's \Yantraraja\ continues to echo in the works of several seventeenth-century authors. The \Yantrasiromani\ (\circa 1612/15) of \Visrama\ of \Jambusara, the \Vasanavarttika\ (1621) of \NrsimhaDaivajna\ of \Kashi, and the \Sarvasiddhantaraja\ (1639) of \Nityananda\ are three such examples.\footnote{\,\MahendraSuri's \Yantraraja, along with \Visrama's \Yantrasiromani\ is edited by \textcite{Raikva}. \NrsimhaDaivajna's \Vasanavarttika, a commentary on \BhaskaraII's \Siddhantasiromani\ (1150), is edited by \textcite{Chaturvedi}. There are no known editions or translations of \Nityananda's \Sarvasiddhantaraja\ in its entirety. \textcite[149]{SarmaAstrolabeinSanskrit} describes how the structure of the \yantradhyaya\ `chapter on instruments' from the \Sarvasiddhantaraja\ mimics that of \MahendraSuri's \Yantraraja\ on the basis of MS~264 from the Asiatic Society of Bombay.} 

%--------------------%

\subsubsection{At the Mughal court} \label{Mughal_court}
During the Mughal rule of India, the practice of translation became an administrative activity under the patronage of the Mughal emperors. Various kinds of literary, historical, religious and scientific texts in Sanskrit were chosen to be translated into Persian. As \textcites{AlamSubramanyam, Truschke} have astutely observed, these translations served, more than anything else, to help the Mughal crown conceive and consolidate its self-identity as a ruling establishment harmonious with locally existing notions of kingship.\footnote{\,\textcite{Haider} offers an excellent study on the role of language and translations in the context of intercultural communication and Mughal state-building. \textcite{Israel} builds on this to examine the complex ways in which translation processes and political discourses are mobilised to shape cultural and national identities.} %
% \footnote{\,In talking about the scientific activities at the Mughal court of Emperor \Humayun\ (r.\thinspace 1530--56)  \textcite{Anooshahr} notes that `the court’s network of patronage reflected the cosmopolitan (and cosmocratic) ambitions of the emperor, extending to intellectuals from Shiraz, Herat, Istanbul, Gwalior and Samarqand. The changes continued and intensified as the remaining decades of the century unfolded' (p.\thinspace 315).} %
For a small group of professionals, however, these translation projects offered more immediate opportunities for employment at the Mughal court and with it, a chance for social recognition.

We learn from the sixteenth-century Mughal historian \Badaunifull\ that Emperor \Akbar\ (r.\thinspace 1556--1605) established a scriptorium  (\maktabkhana) where secretaries, scholars, and scribes worked collaboratively to produce Persian editions of Sanskrit texts.\footnote{\,\Vid\ \Badauni's \MuntakhabalTawarikh, \textcite[Vol.\thinspace II, p.\thinspace 344]{Lees} for the Persian text; its English translation can be found in \textcite[Vol.\thinspace II, p.\thinspace 356]{Lowe}. Also \vid\ \textcite[Chapter~6, pp.\thinspace 203--222]{Rizvi} for a descriptive list of the Sanskrit works translated at \Akbar's \textit{maktabkhāna}, including \AbualFaydi's Persian translation \TarjumayiLilawati\ of \BhaskaraII's \Lilavati\ (\circa mid-twelfth century \ce) from 1587.} According to \Badauni, Sanskrit interpreters (\muabbiran) and Persian translators (\mutarjiman) worked separately at different stages of the translation process.\footnote{\,\MuntakhabalTawarikh: Persian text in \textcite[Vol~II, pp.\thinspace 320--321]{Lees} and its English translation in \textcite[Vol.\thinspace II, pp.\thinspace 329--330]{Lowe}.} Starting with a vernacular paraphrasing of the Sanskrit text by Hindu/Jain scholars (\pandita s or \sastrin s), perhaps in a colloquial dialect of Hindavī, Khaṛībolī, or Braja\-bhāṣā,\footnote{\,As \textcite{AlamPersian1998} observes, `[h]indavī was recognized as a semi-official language by the Sūr Sultāns (1540--55) and their chancellery rescripts bore transcriptions in the Devanāgarī script of the Persian contents. The practice is said to have been introduced by the Lodīs (1451--1526)' (p.\thinspace 319). \Vid\ \textcite{Behl2012} for a study of the Hindavī literary traditions in pre-Mughal India; \textcite{Imre} for the emergence of Khaṛībolī literature in Northern India; and \textcite{Busch} for poetry in Braja\-bhāṣā at the Mughal courts.} a preliminary Persian translation was prepared by Muslim clerks/secretaries (\muharrir s). This was then refined by more accomplished Persian scholars (\ustadh s or \mutamarris es) into its final form over several revisions \parencite[564--566]{Hodivala}. 

It is reasonable to think that this process also occurred in reverse as Persian texts were translated into Sanskrit. In his \AiniAkbari, \Akbar's chronicler \AbulFadlAllami\ mentions at least one instance where a Persian astronomical text was translated into Sanskrit: the \ZijMirzai\ (alias \ZijUlughBeg) was translated into Sanskrit (\JicaUlugbegi) under the superintendence of \AmirFathullahofShiraz\ with the assistance of Kishan Joshī, Gaṅgādhar, and Mahesh Mahānand %
% \footnote{\,\textcite[\P~2.4 on p.\thinspace 367]{Sarmajyotisaraja} identifies Kishan Joshī as \Krsnadaivajna\ (\fl \circa 1600/25 \ce), a favourite of Emperor \Jahangir\ (r.\@ 1605-- 1627 \ce). \Krsnadaivajna\ belonged to a prominent family of Devarāta Brahmins who had emigrated to \Kashi\ \parencite[pp.\thinspace 53a--55b in Volume A2]{PingreeCESS}. In the seventeenth century, their academic rivalry with a family of Bhāradvāja Brahmins, also immigrants in \Kashi, ensured many Sanskrit texts debating Islamicate astronomy were authored by both sides \parencite[\vid][122-127]{Minkowskilearnedbrahmin}. \Munisvarafull\ (b.\thinspace 1603 \ce, Devarātagotra, the nephew of \Krsnadaivajna) and \Kamalakara\ (b.\thinspace 1610 \ce, Bhāradvājagotra, nephew of \NrsimhaDaivajna\ of \Kashi), two of \Nityananda's contemporaries, were noted scholars from these rival families. Along with \Nityananda's \Sarvasiddhantaraja\ (1639 \ce), \Munisvara's \Siddhantasarvabhauma\ (1646 \ce) and \Kamalakara's \Siddhantatattvaviveka\ (1658 \ce) form the corpus of Sanskrit \siddhanta s composed in seventeenth-century Mughal India that discuss Islamicate mathematical astronomy.}
\parencite[110]{BlochmannAiniAkbari}.\footnote{\,H.\thinspace Blochmann translated the first two books of the \AiniAkbari\ into English in 1873 (published by the Asiatic Society of Bengal). \citeauthor{BlochmannAiniAkbari} edited and revised the second edition in 1927, which was then reprinted in 1977. As \textcite[footnote~20 on p.\thinspace 367]{Sarmajyotisaraja} points out, Blochmann's statement on translating the \ZijUlughBeg\ by a consortium of Sanskrit scholars is indeed `hopelessly garbled' in his English translation in all three editions.} % 

In \Shahjahan's reign, beginning in 1628, we find a Persian translation \TarjumayiBijganit\ (\circa 1634/35) of \BhaskaraII's  \Bijaganita, a celebrated twelfth-century Sanskrit treatise on Algebra, written by the Mughal architect \AtaAllahRushdi\ and dedicated to Emperor \Shahjahan\ \parencite[384--386]{Ansaribhaskara}. This is also around the same time when \Nityananda\ translated \MullaFarid's Persian \ZijiShahJahani\ into his Sanskrit \Siddhantasindhu. Although there are no intermediaries (interpreters/translators) that are explicitly named in either of these works, there are historical precedents from the literary traditions, particularly those patronised by the Muslim nobility of early-modern India, to suppose the presence of bilingual interlocutors.\footnote{\,For instance, based on her study of the vernacular literary culture of early-modern North India, \textcite{Orsini} remarks that `it is better to understand the literary culture in fifteenth-century north-India as a multilingual and multilocation literary culture---with a trend towards Persian-Hindavi bilinguality in the domains of politics and literature of the various regional Sultans and in the Sufi religious and literary practices' (pp.\thinspace 238--239).} 

Beyond these literary traditions, we learn from \textcite[327--328]{AlamPersian1998} that
by the middle of the seventeenth century, most administrative positions in the Mughal chancellery were occupied by Persian-speaking Hindu \munshi s, many of whom made significant contributions to Persian literature.\footnote{\,Many Hindu \munshi s at the Mughal courts wrote epistolary prose (\textit{inshā}) and composed poetry in Persian. The story of Chandar Bhān Brahman (d.\thinspace \circa 1666--70), a Brahmin \munshi\ who lived through the reign of four Mughal emperors, is a fascinating tale of how a Hindu secretary came to be regarded as one of the great Persian prose stylists and poets of his era. \textcite{Kinra} offers a particularly compelling account of the literary, social, and political worlds of \Shahjahan's Mughal India through the life and works of Chandar Bhān Brahman.} %
The power and prestige associated with being literate in Persian extended beyond the circles of Hindu imperial administrators---including, of course, Hindu nobility like the Rājapūta kings---and even reached low-ranking officials in smaller towns and villages. By the time of \Shahjahan's reign, Persian classics like \AkhlaqiNasiri\ of \alTusifull\ or \MasnaviyiManavi\ of \Rumipoet\ became regular reading material even among the less-prominent Hindus associated with the Mughal state \parencite[328]{AlamPersian1998}. 

\citeauthor{AlamPersian1998}'s observations allow us to see how Persian became a tool of socioeconomic mobility for the professional classes in seventeenth-century Mughal India. The Sanskrit \jyotisa s (and perhaps, even the Muslim \munajjim s) served as astrologers for various high-ranking Hindus in the Mughal realm.\footnote{\,For example, \VedangarayaMalajit\ (\fl 1643, also known as \Srimalaji) was a \jyotisa\ at \Shahjahan's court. His admittance to the imperial court was presumably mediated by his immediate patron Rāja Giridhara Dāsa, the Rājapūta King of Ajmer, to whom, \Srimalaji\ dedicated his \Giridharananda\ `The joy of Giridhara' \parencite[121--122]{Minkowskilearnedbrahmin}. As \citeauthor{Minkowskilearnedbrahmin} remarks, `the presence of a \jyotisa\ at a particular court appears in some cases to have been rather notional. The Banarsī paṇḍits, in particular, received gifts, honors, or patronage simultaneously from several courts, large and small' (p.\thinspace 116).\label{srimalaji_footnote}} These Hindus, as \textcite{Truschke} describes them, `joined the Mughal administration and became absorbed into Persian-speaking communities' (p.\thinspace 8). Essentially, they were now a part of the Mughal Persianate elite. As their consultant astrologer, the ability to be reasonably bilingual (fluent in vernacular Hindi and conversant in Persian) would have been a competitive advantage and social distinction for any Sanskrit \jyotisa.

From the seventeenth century, the linguistic hegemony of Persian that served the political ambitions of the Mughal crown was met with the rising popularity (and patronage) of vernacular literature among the Persianate elite, \parencite[\eg \vid][chapters~3--4]{AllisonHindiPoetry}. The prominence of Hindavī/Brajabhāṣā literature, coupled with a politico-cultural shift towards the vernaculars (in other words, treating vernacular texts as sources of cultural history instead of those written in Sanskrit) led to renewed ways in which Persian writers engaged with the vernacular cultures.\footnote{\,For instance, \MirzaKhanbFakhralDinMuhammad\ wrote his encyclopedic Persian digest \TuhfatalHind\ (\circa 1674/75) `Gift from India' on the `current Indian sciences' (\Ayn \textit{ulūm\=/i mutadāwila\-/yi hindiya}) during the reign of Mughal Emperor \Awrangzeb\ (r.\thinspace 1658--1707). His book includes discussions on various topics of ordinary and academic interests peculiar to the people of who spoke \textit{Braj Bhākhā} (Brajabhāṣā). \Vid\ \textcite{Ziauddin} for an English translation of Mirzā Khān's elaborate exposition of the grammar of \textit{Braj Bhākhā}; and more generally, \vid\ \textcite[342--348]{AlamPersian1998} for a historical summary of the relationship between Persian and Hindavī at the Mughal courts.} 

% Sanskrit and vernacular accounts also reveal how select Hindu astrologers were conferred the title  (Skt.\@ \Jyotisaraja) (Skt.\@ \Vedangaraja), both implying `Royal Astrologer', by successive Mughal emperors beginning with Akbar \parencite[\eg \vid][58--61]{Truschke}.\footnote{\,The Sanskrit compositions of these title-holders, often including encomiums to their generous patrons, also reflect the Islamicate heritage of the Muslim  benefactors: for instance, the \TajikaNilakanthi\ (1587 \ce) of \Akbar's Jotik Rāi \Nilakantha\ (\fl 1569/87 \ce) was based on Islamic \tajika\ astrology \parencite[pp.\thinspace 177b--189a in Volume A3]{PingreeCESS}; or the \SamskrtaParasikaPadaPrakasa\ (1643 \ce) of \Shahjahan's \VedangarayaMalajit\ (\fl 1643\ce) was a Sanskrit text on Persian grammar \parencite[pp.\thinspace 421a--422a in Volume A4]{PingreeCESS}.} %

Sanskrit poets also learned to adapt to this shift towards the vernaculars. Many scholars maintain that the literary eminence of Sanskrit at the Mughal court began to wane in the reign of \Shahjahan\
\parencites[\eg][]{PollockDeathofSanskrit, Truschke}.%
% (\eg \cite[404--412]{PollockDeathofSanskrit}; \cite[pp.\thinspace 13--16]{Truschke}).
\footnote{\,In contrast, Sanskrit poetry (\textit{kāvya}) composed outside the central Mughal court played a critical role in elaborating the vernacular cultures and identities. As \textcite{BronnerShulman} elaborate in their study, Sanskrit was employed to articulate regional distinctiveness instead of occluding it.} The accounts of two Hindi-speaking Brahmin poets at \Shahjahan's court, \KavindracaryaSaravati\ (\fl \circa 1600/75) and \JagannathaPanditaraja\ (\fl \circa\ 1620/60), describe how two eminent Sanskrit scholars ingratiated themselves with the emperor and his retinue by composing panegyrics in Braja\-bhāṣā and singing Hindustānī \textit{dhrupad} songs at the Mughal court (\cite[50--53]{Truschke}).

For a lesser-known Sanskrit \jyotisa\ like \Nityananda, however, one can imagine that the changing tides of patronage and the competition to find patrons, would have presented very different challenges to those faced by courtly bards singing encomiums. \Nityananda's name appears in the annals of Sanskrit \Jyotihsastra\ as the author of \Siddhantasindhu---a Sanskrit translation of a Persian original sponsored by \AsafKhanshort, the prime minister (\vazir\idafaconsonant\ \azam) of \Shahjahan\ and a highly influential Mughal elite. It is, therefore, not inconceivable that \Nityananda\ might have had some basic level of Persian literacy to begin with, or at the very least, developed it through his interactions with \MullaFarid\ (in vernacular Hindi) over the course of his commission. 

\paragraph{Sanskrit manuals on learning Persian} 
Between the fourteenth and eighteenth centuries, several Sanskrit compendiums were authored to teach Persian to Sanskrit-speaking audiences \parencites[\eg \vid][]{SarmaSanskritmanuals, Truschkepersiansanskritlexica}. Typically, these compendiums comprised of two sections composed in metrical Sanskrit verses: namely, the \textit{kośa prakaraṇa}, a bilingual Persian-Sanskrit lexicon, and the \textit{vyākaraṇa prakaraṇa}, a section on the rules of Persian grammar. The \Parasiprakasa\ (\circa 1575) of \Krsnadasa\ dedicated to \Akbar\ and the \SamskrtaParasikaPadaPrakasa\ (1643) of \VedangarayaMalajit\ sponsored by \Shahjahan\ are two prominent exemplars \parencite[\vid][]{SarmaSanskritPersianLexica}. The former contains a general collection of Persian words, whereas, the latter includes a specialised lexicon on technical terms in Islamicate astrology/astronomy. It is doubtful if either of these manuals were ever sufficient to learn Persian. However, their value in promoting Persian as a language of sociopolitical influence in Mughal India is certainly conceivable.\footnote{A statement in support of this idea is found in the words of the Sanskrit scholar \Pandita\ \Suryadasa\ Daivajña (b.\thinspace 1508). \Suryadasa\ wrote a versified glossary of Perso-Arabic astrological terms as a section of the fifth chapter in his \Siddhantasamhitasarasamuccaya\ (1583). He begins the section by claiming (in v.\thinspace 56) that the knowledge of the `technical terms stated in the science of the foreigners'  (\textit{yavana-śāstra-uktā saṃjñā}) will be `useful in the royal court' (\textit{narapati-sabhā-upayogya}) and will also be `beneficial to astrologers' (\textit{upakāra-artha daivavidām}); \vid\ \textcite[p.\thinspace 329--330]{MinkowskiSuryadas} for the Sanskrit text of this verse, and also an overview of \Suryadasa's contributions in promoting Islamicate astrology in Sanskrit. More generally, \vid\ \textcite{AlamPersian2003} for an excellent study on the cultural and political role of Persian in the polity of Mughal India.}  

% More recently, \textcite{Nair} has studied how Sanskrit philosophical and religious texts were translated into Persian by collectives of Sanskrit and Persian scholars at the Mughal courts. \citeauthor{Nair} interprets philosophical translations as cross-cultural intellectual `dialogues' expressed through the creation of novel vocabulary. This interpretation can also be extended to technical translations, particularly those that help bridge language barriers in communicating foreign science. A hundred years after \Nityananda, his words continued to echo in the conversations between Sanskrit and Islamicate astronomy beyond the Mughal court.
%--------------------%

\subsubsection{Away from the Mughal court} \label{away_Mughal_court}
By the turn of the eighteenth century, the locus of Sanskrit patronage shifted from the Mughal court to the courts of the vassal states under Mughal suzerainty. Among these subimperial sponsors, the royal patronage of Mahārāja Savāī Jayasiṃha of Jayapura (Jaipur) is particularly pertinent to the history of Sanskrit astronomy. Savāī Jayasiṃha II (r.\thinspace 1699--1743) was the Kachvāha Raj\-put King of Āmera (and later Jayapura) who invested in Sanskrit astronomy both academically and economically. He not only paid for the construction of five astronomical observatories in India but also instituted an ambitious project to translate Islamicate scientific works into Sanskrit; in particular, Arabic and Persian version of Greco-Islamicate mathematics and astronomy---and to a lesser extent, even the European astronomical tables brought to him by the Jesuits \parencites{SarmaJaiSinghSanskrit, PingreeAstronomersprogress}.\footnote{\,In her doctoral dissertation, \textcite{Johnson_Roehr} describes the sociopolitical impact of Savāī Jayasiṃha's urban observatories, in particular, the emplacement of ancillary knowledge-systems (like accounting, masonry, \etcp) within the local landscape of his newly built city of Jayapura (Jaipur). Her observations locate these subsidiary activities within Jayasiṃha's programme of assimilating Islamicate and European astronomy, and in that capacity, offer an interesting parallel to the patronage of professional interpreters, scribes, accountants, and clerks in early-modern society of Mughal India \parencite[\eg  \vid][Chapter~7 `The Making of a Munshī']{AlamSubramanyam}.} %

\textcite{SharmaSawaiJaiSingh} provides a descriptive account of the Hindu astronomers, astrologers, observers, and scribes recruited under Savāī Jayasiṃha's programme. Among these names, \Jagannathafull\ (\fl \circa 1720/40), \Nayanasukhopadhyaya\ (\fl 1729), and \Kevlaramafull\ (\fl \circa 1730/80) are three notable Hindu \jyotisa s who translated the science of the \yavana s (foreigners) into Sanskrit.\footnote{\,\Kevalarama\ authored several works at Jayasiṃha's court; one of them is believed to be the \Drkpaksasarani\ (\circa 1733), a Sanskrit adaptation of Philippe de La Hire's \textit{Tabulae Astronomicae} (1702) based on its 1727 Latin edition 
\parencite{PingreeLaHireSanskrit}. There is at least one other Sanskrit text, the \Phirangicandracchedyopayogika\ (\circa 1732/24), also inspired by La Hire's work but collectively authored by the \jyotisa s at Jayasiṃha's court (\cite[248--249]{SAT_Montelle_Plofker}; \cite{PingreeLaHireJaysimha}).\label{kevalarama_works_foreign}} %
Table~\ref{sanskrit_translations_of_islamic_texts} lists some of the more prominent Sanskrit translations of  Arabic and Persian works, particularly, those that were commissioned by Savāī Jayasiṃha in the early eighteenth century \parencite[131--151]{PingreeJaipur}.\footnote{\,These works are translations in an \textit{explicit} sense; there are other works, mostly \siddhanta s composed in the late seventeenth- and early eighteenth-century Mughal India, that implicitly engage and discuss Islamicate astronomy. \Nityananda's \Sarvasiddhantaraja\ (1639), \Munisvara's \Siddhantasarvabhauma\ (1646), and \Kamalakara's \Siddhantatattvaviveka\ (1658) are three canonical examples of Sanskrit texts that discuss Islamicate mathematical astronomy. More on this in \S~\ref{sanskrit_translation_of_persian_zij} \label{explicit_implicit_translation_sanskrit}.}

%%%%%%%%%%%%%%%
\begin{table}[!hbp]
    \centering
    \renewcommand{\arraystretch}{1.25}
    \renewcommand{\baselinestretch}{1.25}\selectfont
    \begin{tabularx}{\textwidth}{cX}
    \hline
    Date of composition & Sanskrit text\\
    \hline
    \textit{ante} 1694 &
    the \Hayatagrantha, an anonymously authored Sanskrit translation of \aliQushji's Persian text \Risaladarilmalhaya\ (1458) `Treatise on Astronomy (\ilmalhaya)', edited by  \textcite{BhattacaryaHayata}\\
    1726 & the \Rekhaganita\ of \Jagannathafull\ (1652--1744), a Sanskrit translation of \alTusifull's Arabic text
    \KitabTahirUsulliUqlidus\ (\circa 1248) `The recension of Euclid's \Elements', edited by \textcite{Trivedi}\\
    \circa 1726--1732 & the \Samratsiddhanta\ of \Jagannathafull, a Sanskrit translation of \alTusifull's  Arabic recension \TahriralMijisti\ (1247) `Commentary on [Ptolemy's] \Almagest'; three versions of this text are attested: the earliest, called \Samratsiddhantakaustubha, is from 1726, while the two later expanded versions are from 1730 and 1732 respectively; the text dated 1732 is edited by \textcite{SharmaSamratSiddhanta}\\
    1729 & the \Ukara\ of \Nayanasukhopadhyaya\ (with the assistance of \MuhammadAbidda), a Sanskrit translation of \alTusifull's Arabic recension \Tahrialukarr\ (1253) `Commentary on [Theodosius'] \Sphaerica', edited by \textcite{BhattacaryaUkara}\\
    1729 & the \Sarahatajakiravirajandi\ of \Nayanasukhopadhyaya\ (with the assistance of \MuhammadAbidda), a Sanskrit translation of Chapter~11 from Book~II of \alTusifull's \alTadhkirafiilmalhaya\ `Memoirs on Astronomy' (1261--1274) with \alBirjandi's \SharhalTadhkirah\ (1507) `Commentary on the \textit{Tadhkira}', edited by \textcite{KusubaPingree}\\
    \circa 1730 & the \Yantrarajasyarasala\ aliases \Visavava, \Yantrarajavicaravimsadhyayi\ of \Nayanasukhopadhyaya\ (suspected), a Sanskrit translation of \alTusifull's Persian text \RisaladarBistBabMarifatUsturlab\ `Treatise in Twenty Chapters on the Knowledge of the Astrolabe' (\circa 1240), edited by \textcite{BhattacaryaYantrarajavicara}\\
    \hline
    \end{tabularx}
    \captionof{table}{Major Sanskrit translations of Arabic and Persian astronomical texts}
    \label{sanskrit_translations_of_islamic_texts}
\end{table}
%%%%%%%%%%%%%%%
% \pagebreak %%%%%%%%%%%%%%%%%%%%% PAGEBREAK FOR FORMATTING ON PDF, CAN BE REMOVED LATER 

It is worth noting that Savāī Jayasiṃha possessed a copy of \Nityananda's \Siddhantasindhu\ (\circa early 1630s).\footnote{\,MS~Museum~23 (444 folia) of the \Siddhantasindhu\ held at the Maharaja Sawai Man Singh II Museum at the City Palace of Jaipur attests that it was copied by Gaṅgārāma of Kaśmīra for Mahārāja Jayasiṃha on Thursday 6 April 1727 \ce\ (stated on f.\thinspace 443). A second note (in Hindi) appears on f.\thinspace 444v indicating it copied (from an earlier copy?) by Gaṇgāyaratna on \circa 24 May 1726 \ce\ \parencite[142--143]{PingreeJaipur}.\label{ss_museum_23_description}} It is very likely he also possessed a copy of \Nityananda's \Sarvasiddhantaraja\ (1639). As \textcite[79]{PingreeAstronomersprogress} notes, the earliest version of \Jagannathafull's \Samratsiddhanta\ (\ie the \Samratsiddhantakaustubha\ from 1726; the third entry in Table~\ref{sanskrit_translations_of_islamic_texts}) includes the astronomical parameters of \UlughBeg\ derived from \Nityananda's \Sarvasiddhantaraja. \citeauthor{PingreeAstronomersprogress} continues on to say: `From \Jagannatha's use of [the astronomical parameters] we come to realize what has long been suspected, that \Nityananda's arguments, originally advanced in 1639, finally found a receptive audience, nearly a century later, at Jayasiṃha's court' (p.\thinspace 79). In fact, \Nityananda's technical vocabulary also provides some of the terminology with which \Kevalarama\ (\vid\ footnote~\ref{kevalarama_works_foreign}) translates European astronomy into Sanskrit around the mid eighteenth-century \ce\ \parencite[283]{PingreeSSR}.

The method of translation at Savāī Jayasiṃha's court becomes evident with one of his astronomers' own statement on the process. \Nayanasukhopadhyaya, in his \Ukara\ and his \Sarahatajakiravirajandi\ (the fourth and fifth entry in Table~\ref{sanskrit_translations_of_islamic_texts}), expressly mention \MuhammadAbidda\ dictating the Arabic passages while he composes them into Sanskrit \parencite[73-74]{SarmaJaiSinghSanskrit}. As \textcite{KusubaPingree} remark, \Nayanasukhopadhyaya\ `did not simply render the Arabic commentary together with the original into Sanskrit literally, but expanded those passages that he found particularly difficult' (p.\thinspace 7). This suggests that between \MuhammadAbidda's  dictation of the Arabic passages and \Nayanasukhopadhyaya's translation of those passages into Sanskrit, they communicated directly or through an intermediary in a common link  language (perhaps, a colloquial dialect of Hindavī, Brajabhāṣā, or Rājasthānī). 

Bidirectional translations of texts between Sanskrit and other languages (\eg vernaculars like Hindustānī or Bāṅglā, or even European languages like English or German) continued beyond the eighteenth-century, with different methods and motivations \parencites[\eg \vid][]{Dodsontranslation, Raina, Gallien}. Texts in the exact sciences were included in many translation projects, and undoubtedly, they were repurposed to serve the ambitions of the benefactor and the beneficiary alike.  

\subsubsection{Why \Nityananda?} \label{nityananda_bio}

It is uncertain why \Nityananda\ was chosen to translate \MullaFarid's Persian \ZijiShahJahani\ into Sanskrit. Based on what we know, \Nityananda\ was not a decorated astronomer: he did not hold any titles like \textit{Jotik Rāi} or \textit{Vedāṅgarāya}, even though  \Shahjahan\ conferred such a title on \VedangarayaMalajit\ (\fl 1643) (\vid\ footnote~\ref{srimalaji_footnote}). \Nityananda\ identifies himself in the colophon of his \Sarvasiddhantaraja\ (1639) as a Gauḍa Brahmin of Mudgala \gotra\ (patronymic) from Indrapurī (Old Delhi), and provides a genealogy of his Brahmin ancestors beginning with his father: \Nityananda, son of \Devadatta, son of \Narayana, son of \Laksmana\ son of \IcchaDulinahatta\ \parencite[\eg \vid][228]{PetersonCatalogue}. Beyond this register of names, we have no reliable information on who these other Brahmins were, where they came from, or what works they wrote (if any).

Until any new evidence suggests otherwise, we believe \Nityananda's association with the Mughal court begins with \AsafKhanshort\ employing him in \circa \post\ 1628.\footnote{\,In contrast, \MullaFarid\ first joined the service of \KhanKhana---a prominent Mughal nobility
during the reigns of the Mughal emperors \Akbar\ and \Jahangir, `a man of the sword as well as of the pen' \parencite[75]{Lefevre}---in 1597 and remained in his service till he joined the court of \Shahjahan\ in 1628 \parencite[34]{Ghori}.} As I describe below (in  \S~\ref{zijshahjahani_mulla_farid}), a royal decree (\farman) was issued to bring Muslim and Hindu astronomers together to prepare the \ZijiShahJahani\ under \AsafKhanshort's supervision. This may have been the ticket for \Nityananda's entry to the Mughal court. His fluency in vernacular Hindi (as a resident of Delhi) and competency in Sanskrit astronomy (presumably, attested through testimony) might have brought him to \Shahjahan's court seeking patronage as a Sanskrit \jyotisa.
 
We do not know whether \Nityananda\ continued to remain at the Mughal court after composing the \Siddhantasindhu, or even if his association was ever exclusive to begin with. His second book, the \Sarvasiddhantaraja\ (1639), is a complex syncretism of Sanskrit \siddhantic\ astronomy and Islamicate theories. To my knowledge, there are no explicit references to any patrons in this work; however, judging by the scale and scope of the work, it seems very likely that he had continued access to intellectual and financial resources throughout its production \parencite[more on this in][]{MisraTD}.\footnote{\,A single extant manuscript of a text called the \Shahajahamganita, allegedly authored by \Nityananda, is currently held at the Anup Sanskrit Library in Bikaner (MS~5291, Serial \textnumero 787, 12 folia, injured, \vid \cite[p.\thinspace 174a in Volume A3]{PingreeCESS}; \cite[393]{AnupSanskritLibrary}). I have not been able to consult this manuscript to verify its professed authorship.} 

 

% \KhanKhanashort\ (1556--1627 \ce), better known as \Rahim\ in the world of Hindi-poetry and Nabbāba Khānakhānā in Sanskrit, was a polyglot man of letters and a patron  of Hindi, Sanskrit, and Persian literature. He composed the \Khetakautuka\ (\circa late sixteenth century \ce), a Sanskrit text of 124 verses explaining Persian \tajika\ astrology with transliterated Arabic/Persian technical terms \parencite[pp.\thinspace 79b--80a in Volume A2]{PingreeCESS}.}

% \Shahjahan's \jyotisa s like of \Nityananda, \Balabhadra\footnote{\,\Balabhadra\ (\fl 1629/53 \ce) composed his \Hayanaratna\ for \Shahjahan's second son \ShahShuja\ at Rāja\-mahala in 1629 \ce. The \Hayanaratna\ is a well-renowned Sanskrit digest (\textit{nabandha}) of excerpts from earlier Sanskrit writings on Arabic \tajika\ genethlialogy; edited by \textcite[forthcoming]{Ganstenhayanaratna}.}

% The Arabic divination method of geomancy (\khatt\ \alraml) appears as Sanskrit \ramala-texts from the beginning of the seventeenth-century, \eg \Ramalavaicitrya\ (\circa \ante\ 1657 \ce) of \Rama\ \parencite[pp.\thinspace 424a--425bin Volume A5]{PingreeCESS}. According to \textcite[79--80]{PingreeJyotihsastra}, these texts were introduced from Persian sources in the early Mughal period.

% A comprehensive analytical survey of Persian texts (originals and translations) authored between the thirteenth and nineteenth centuries in India has been undertaken by the \href{http://www.perso-indica.net/}{Perso-Indica Project}, an online peer-reviewed publication led by Fabrizio Speziale (EHESS, Paris) and Carl W.\@ Ernst (University of North Carolina, Chapel Hill) as editors-in-chief \parencite{PersoIndica}.


%%%%%%%%%%%%%%%%%%%%%%%%%%%%%%%%%%%%%%%%%%%%%%%%%%%%%%%%%%%%%%%%%%%%%%%%%%%%%%%%%%%%%%%%%
\subsection{Islamicate \zij es in Mughal India}\label{islamicate_zijes_mughal_india}
%%%%%%%%%%%%%%%%%%%%%%%%%%%%%%%%%%%%%%%%%%%%%%%%%%%%%%%%%%%%%%%%%%%%%%%%%%%%%%%%%%%%%%%%%

By the seventeenth century, Arabic and Persian astronomical texts were regularly studied at Islamic institutions of higher learning (\madrasa) in Mughal India; particularly, at those institutions that focused on teaching the rational sciences (\ulumalaqliyah).\footnote{\textcite[1--88]{Sufi} provides an extensive chronological study of the evolution of curriculum in Indian \madrasa s through the dynastic rule of the Ghaznavids, the Ghurids, the Delhi Sultanates, and eventually the Mughals.} \textcite[Table~1 on p.\thinspace 278]{Ansaritransmission} lists the names of prominent Islamicate scholars whose works are extant in several copies in Indian libraries. It includes the works of \abuNasr\ (d.\thinspace 1036), \Gilanifull\ (d.\thinspace 1029), \ibmHaytham\ (965--1041), \alBirunifull\ (973--1048), \alJaghmini\ (\circa early thirteenth century), \alTusifull\ (1201--1274), \alShirazifull\ (1236--1311), \alKashifull\ (d.\thinspace 1436), \UlughBegfull\ (1394--1499), and many others; also \vid\ \textcite[Section~III on pp.\thinspace 279--281 and Appendix~I on pp.\thinspace 288--294]{Ansaritransmission}. The large numbers of manuscript witnesses suggest the prevalence of these works in the repertoire of Muslim scholars (\ulama, \fudala) in early-modern India. 
% Of notable significance are the works of  \abuNasr\ (d.\thinspace 1036 \ce), \ibmHaytham\ (965--1041 \ce), \alBirunifull\ (973--1048 \ce), \alJaghmini\ (\circa early thirteenth century \ce), \alTusifull\ (1201--1274 \ce), \alShirazifull\ (1236--1311 \ce), \alKashifull\ (d.\thinspace 1436 \ce), \UlughBegfull\ (1394--1499 \ce), \alRumifull\ (\circa 1440 \ce), \aliQushjifull\ (d.\thinspace 1474 \ce), \alBirjandi\ (d.\thinspace 1525 \ce), and \alAmilifull\ (1547--1622 \ce), \vid\ \textcite[Section~III on pp.\thinspace 279--281 and Appendix~I on pp.\thinspace 288--294]{Ansaritransmission}.
% \footnote{\,\textcite[Section~III on pp.\thinspace 279--281]{Ansaritransmission} describes various astronomical works of Islamicate scholars that were studied in Indian \madrasah s. Also, \vid\  Appendix~I (\Ibid, pp.\thinspace 288--294) for a catalogue-list of astronomical works in India.}

The lists of \zij es enumerated in the \AiniAkbari\ of \AbulFadlAllami\ (1551--1602), the chronicler of \Akbar, and the \ZijiShahJahani\ of \MullaFarid\ composed during the rule of \Shahjahan\ provide further information on the astronomical tables in the imperial library (\kitabkhana) of early seventeenth-century Mughal India \parencite[Appendices~A and B on pp.\thinspace 45--48]{Ghori}. 
% \footnote{\,\textcite{Seyller} provides a comprehensive description and valuation of the manuscripts in the collections of the Mughal library. The paper begins by observing that `[s]hortly after the death of Emperor Akbar in 1605, an inventory of the vast holdings of the imperial Mughal library recorded a total of 24,000 volumes with a value of 6,463,731 rupees.' (p.\thinspace 243).}

An important \zij\ in these accounts is the \ZijSultani\ (alias \ZijUlughBeg\ or \ZijJadidiGurani) of \UlughBegfull\ composed by a collaborative team of astronomers (\alRumi, \alKashi, \UlughBeg, and \aliQushji) at the Observatory of Ulugh Beg in Samarqand in 1438/39.\footnote{\,\Vid\ \textcite[277--279]{RosenfeldandIhsanoglu}, \textcite[54]{KingSamso}, and \textcite[pp.\thinspace 127b--128a and pp.\thinspace 166b--167b]{Kennedysurvey} for a descriptive survey of the contents and manuscripts of the \ZijUlughBeg.} As \textcite[p.\thinspace 581a]{Ansarisurvey} notes, the \ZijUlughBeg\ was translated into Sanskrit (\JicaUlugbegi\thinspace\footnote{\,The \JicaUlugbegi\ (or \Ulakabegijica) is extant in a few fragmentary manuscripts containing only tables and star catalogues. The largest (and most complete) manuscript appears to be MS~Museum 45 held at the Maharaja Sawai Man Singh II Museum at the City Palace of Jaipur. According to \textcite[135]{PingreeJaipur} it contains 100 folia measuring  17 $\times$  28\textonehalf\ cm, contains only tables written in \Nagari\ numerals, and was acquired from S\={u}rata by Nandar\={a}ma Jo\'{s}\={i} for 20\textonehalf\ rupees. \label{jicaulughbeg}}) by a consortium of Muslim and Hindu scholars led by \ShahFathullahShirazi\ (d.\thinspace 1589) during the reign of \Akbar. From the sixteenth century, the preeminence of the \ZijUlughBeg\ in Mughal India made its structure the standard with which subsequent \zij es were composed. The \ZijiRahimi\ and the \ZijiShahJahani, both composed by \MullaFarid, are two such examples.\footnote{\,\textcite[Section~3.2 on pp.\thinspace 582--583]{Ansarisurvey} reviews the structure and contents of \MullaFarid's \ZijiRahimi\ (\circa 1615/17) dedicated to his patron \KhanKhana, a prominent Mughal nobility
during the reigns of the Mughal emperors \Akbar\ and \Jahangir.}

%%%%%%%%%%%%%%%%%%%%%%%%%%%%%%%%%%%%%%%%%%%%%%%%%%%%%%%%%%%%%%%%%%%%%%%%%%%%%%%%%%%%%%%%%
\subsection{The \ZijiShahJahani\ (\circa 1629/30) of \MullaFarid}\label{zijshahjahani_mulla_farid}
%%%%%%%%%%%%%%%%%%%%%%%%%%%%%%%%%%%%%%%%%%%%%%%%%%%%%%%%%%%%%%%%%%%%%%%%%%%%%%%%%%%%%%%%%

\MullaFarid's \ZijiShahJahanifull\thinspace\footnote{\,The \ZijiShahJahanifull\ `Grand Accomplishment of the Second Lord of the Conjunction, the \Zij\ of \Shahjahan' uses the royal epithet of \Shahjahan\ as the `Second Lord of the Conjunction' born on the auspicious conjunction (\qiran) of Jupiter and Venus at his natal hour on 5 January 1592 \ce. The use of `second' in the title is to establish a direct descent from the first Lord of Auspicious Conjunction, \Timur\ \parencite[1105--1106]{Chann}.} (\ZijiShahJahani\ for short) is a set of astronomical tables written at the behest of \AsafKhan\ (d.\thinspace 1641), the prime minister (\vazir\idafaconsonant\ \azam) and father-in-law to \Shahjahan. It was commissioned to institute a new calendar of \Shahjahan, the \Shahishani, beginning on the first day of Farvardin in his ascensional year 1037 \ah\ (21 March 1628). This was in keeping with previous regnal year calendar like \Jalali\ era (epoch~21 March 1079) of the Seljuk Sultan \MalikShah\ or the \Ilahi\ `Divine Era' (epoch~20 March 1555/56) of \Shahjahan's grand father \Akbar. \textcite[Section~3.2 on pp.583--585]{Ansarisurvey}, \textcite[34--36]{Ghori}, \textcite[357--358]{RosenfeldandIhsanoglu}, and \textcite[307]{STMI} survey the context, the structure, and manuscripts of the \ZijiShahJahani. For our present purpose, we note the following points from these surveys:
\begin{enumerate}[topsep=0pt]
    \item The \ZijiShahJahani\ (like the \ZijUlughBeg) consist of a detailed prolegomenon (\muqaddima) consisting of five sections (\qism, pl.\thinspace \aqsam) followed by four discourses (\maqala, pl.\thinspace\maqalat) on four different subjects, each containing several chapters (\bab, pl.\thinspace \biban) that are further divided into sections (\fasl, pl.\thinspace\fusul).
    \item \MullaFarid\ classifies the \ZijiShahJahani\ as a \zijiHisabi\ or a `computational table that revises and updates the parameters' of the \ZijUlughBeg\ (and different from a \zijiRasadi\ or an `observational table based on findings from direct observations').
    \item A lack of time to conduct newer observations, in part due to the advancement of age and the ailing health of \MullaFarid, meant that large parts of the \ZijiShahJahani\ were reproductions of corresponding parts of the \ZijUlughBeg. However, as \textcite[585]{Ansarisurvey} notes, the tables in the \ZijiShahJahani\ outnumber those in the \ZijUlughBeg. \MullaFarid\ includes the auxiliary tables for simplification (\tashil) so that the true longitudes (\taqvim) of celestial objects can be computed directly (without any interpolation). 
    \item By a royal decree of \Shahjahan, the \ZijiShahJahani\ was to be translated into `the language of Hindustan by Indian astronomers in consultation with Persian astronomers, for the sake of public utility'.\footnote{\,Excerpted from the \MulakhkhasiShahJahanNama, an abridged history of \Shahjahan\  written by his seventeenth-century court chronicler \TahirKhan, alias \InayatKhan, \vid\ \textcite{Ansarisurvey} for \TahirKhan's Persian text (Appendix~II.\thinspace A4 on p.\thinspace 597) and its English translation (p.\thinspace 584a)} \textcite[34]{Ghori} also observes that the \ZijiShahJahani\ was prepared by \MullaFarid\ in `collaboration of his brother \MullaTayyib\ and other scholars of Muslim and Hindu astronomy under the over-all supervision of the Vazir \AsifKhan\ [\sic]'.     
\end{enumerate}

\subsubsection{Manuscripts of the \ZijUlughBeg}\label{manuscripts_zij_shahjahani}
In this study, I outline the twenty-two chapters (\biban) in the second discourse (\maqala\idafaconsonant\ \duvum) of the \ZijiShahJahani, and among these, I examine the sixth chapter on the methods to determine the \gls{distance_celestial_object} (\bud\idafaconsonant\ \kawkab\ \az\ \muaddil\ \alnahar), \ie its declination, in closer detail. To this end, I have consulted (parts of) the two manuscripts of the \ZijiShahJahani\ that were available to me. Table~\ref{mss_description_zij_shah_jahani} provides a description of these manuscripts and their assigned sigla.

\begin{table}[!htbp]
\centering
\renewcommand{\arraystretch}{1.5}
\renewcommand{\baselinestretch}{1.25}\selectfont
\begin{tabularx}{\textwidth}{c X}
\hline
Siglum & Manuscript description\\
\hline
\SjA & MS~2735 (MS Ind.\@ Inst.\@ Pers.\@ 12) from the Bodleian Library Oxford, entitled \textbf{Kâranâma i ṣâḥibkirân thânî zîj i shâhjahânî}, 380 folia (incomplete) with 25 lines per folio, 13\,¼ $\times$ 9\,⅜ inches, Persian \Nastaliq, written with red and black ink, \circa 17\textsuperscript{th} century \ce\ \parencite[p.\thinspace 61b]{Beeston}. \\
& Mr Alasdair Watson, the Bahari Curator of Persian Collections at the Bodleian Library, very kindly provided me with photographs of folia that include the sixth chapter of the second discourse. The folio numbers do not appear on the images; Mr Watson identified them as ff.\thinspace 21b--22a (personal communication). I have not had the opportunity to inspect the other folia of this manuscript. The summary of the Persian chapter-titles in \S~\ref{chapter_title_comparision_persian_sanskrit} rely entirely on the reading in MS~\SjB.\\
\SjB & MS~Or.\thinspace 372 (Sch 7909) from the British Library London, labelled \textbf{Farìd Ibráhím Zíj E Sháhjaháni Persian} (on the microfilm cover), entitled \tfarsi{کارنامه صاحبقران ثانی زیج شاه جهانی} %
\ZijiShahJahanifull, 419 folia with 31 lines per folio, 13\,¾ $\times$ 8\,½ inches, Persian \Nastaliq, treble-ruled text frame, \circa 17\textsuperscript{th} century \ce\ \parencite[pp.\thinspace 459b--460b]{Rieu}. \\
& Folio numbers are written at the top left corner of \textit{folium rectum} (b-side) in western Arabic numerals, (possibly) by a European owner/cataloguer/librarian.\\
& I am grateful to Dr Benno Van Dalen (from the \textit{Ptolemaeus Arabus et Latinus} project at the Bayerische Akademie der Wissenschaften München) for providing a digitised black-and-white photocopy of this manuscript to me. \\
\hline
\end{tabularx}
\caption{Description of the manuscripts of the \ZijiShahJahani}
\label{mss_description_zij_shah_jahani}
\end{table}

To my knowledge, the \ZijiShahJahani\ (or any part of it) has never been edited or translated into any major modern European or Indian (vernacular) language.


%%%%%%%%%%%%%%%%%%%%%%%%%%%%%%%%%%%%%%%%%%%%%%%%%%%%%%%%%%%%%%%%%%%%%%%%%%%%%%%%%%%%%%%%%
\subsection{The \Siddhantasindhu\ (\circa early 1630) of \Nityananda}\label{siddhantasindhu_nityananda} 
%%%%%%%%%%%%%%%%%%%%%%%%%%%%%%%%%%%%%%%%%%%%%%%%%%%%%%%%%%%%%%%%%%%%%%%%%%%%%%%%%%%%%%%%%

\Pandita\ \Nityanandafull\ completed his translation of \MullaFarid's \ZijiShahJahani\ in the early 1630s and named it \Siddhantasindhu\ `the Ocean of \Siddhanta s'. This enormous work occupied around 440 folia measuring 45 $\times$ 33 cm (approximately). At the time, ten copies of this work were made and distributed among the Muslim nobility of northern Mughal India (more on this in \S~\ref{circualtion_siddhantasindhu}). Today, however, only a handful of near-complete manuscripts of this work survive. And among these, the four manuscripts at the Maharaja Sawai Man Singh II Museum at the City Palace of Jaipur (India) are the best preserved copies. One of these manuscripts (Khasmohor 4960, part of the \khasmohor\ or `special seal' collection) bears the royal insignia of \Shahjahan\ himself. \textcite[138--143]{PingreeJaipur} describes the four manuscripts at the City Palace Museum Library in Jaipur, while \textcites[p.\thinspace 173b in Volume A3 and p.\thinspace 184a in Volume A5]{PingreeCESS} provides the catalogue references of the other (fragmentary) manuscripts located elsewhere.

\subsubsection{Manuscript of the \Siddhantasindhu} \label{manuscript_siddhantasindhu} 
Parallel to the selection from the \ZijiShahJahani, I outline the twenty-two chapters (\adhyaya s) of the second part (\dvitiya\ \kanda) of the \Siddhantasindhu, and among these, I focus on the sixth chapter on the methods to determine the \gls{true_declination} (\spasta-\kranti) of a celestial object. I have consulted the only copy of \Siddhantasindhu\ that has been available to me. A description of the manuscript, and its assigned siglum, is given in Table~\ref{mss_description_siddhantasindhu}.

\begin{table}[!htbp]
\centering
\renewcommand{\arraystretch}{1.5}
\renewcommand{\baselinestretch}{1.25}\selectfont
\begin{tabularx}{\textwidth}{c X}
\hline
Siglum & Manuscript description\\
\hline
Kh & MS~4962 from the Khasmohor Collection at the City Palace Library of Jaipur, entitled (in Hindustānī) \tsans{pothi siddhaa.mtasi.mdhu kii} `Book of \Siddhantasindhu', 436 folia (incomplete: missing ff.\thinspace 1 and 3; tears and damages on f.\thinspace 2) with 21--30 lines per folio, 37 $\times$ 25 cm, Sanskrit \Nagari, written with red and black ink parallel to the shorter edge, double-ruled text frame, left-binding with side-sewing stitches, red-and-blue striped cloth-covered boards and book flap, belonging to Jagannātha Jośī and acquired for 100 rupees, \circa early 18\textsuperscript{th} century \ce\ \parencite[143]{PingreeJaipur}. \\
& Folio numbers are written at the bottom left corner of \textit{folium versum} in \Nagari\ numerals by the same hand as the scribe.\\
&\textbf{Remark}\quad The metrical verses are often introduced by the word \tsans{cha.mda.h} (\textit{chaṃḍaḥ}) `meter', \eg f.\thinspace 4v:\thinspace 4, 15~Kh. Its abbreviated form \tsans{cha.m} (\textit{chaṃ}) also appears in several places, \eg on f.\thinspace 5v:\thinspace 7~Kh. Overall, Kh uses the double-\textit{daṇḍa} `\tsans{||}' to indicate half-stanza breaks, and frequently places the verse number, each time beginning with one, between two sets of double-\textit{daṇḍa}s, \eg \tsans{||~1~||}. The prose-passages are unnumbered.\\
& I gratefully acknowledge Dr Chandramani Singh, (retired) head curator of the Maharaja Sawai Man Singh II Museum Library, for her assistance in helping me get a digital copy of this rare and private manuscript.\\
\hline
\end{tabularx}
\caption{Description of the manuscript of the \Siddhantasindhu}
\label{mss_description_siddhantasindhu}
\end{table}

There are no published editions, translations, or studies of \Nityananda's \Siddhantasindhu\ to my knowledge.\footnote{\,\textcite[231--232]{PetersonCatalogue} provides an excerpt containing the first thirty-two verses (from the prolegomenon) and the colophon (from the end of the second part) of \Nityananda's \Siddhantasindhu. This is presumably transcribed from MS~RORI (Alwar) 2627 = MS~2014 Alwar, 441 folia, copied in 1855 \ce\ \parencite[p.\thinspace 184a in Volume A5]{PingreeCESS}; however, \citeauthor{PetersonCatalogue} does not identify the shelf mark of the manuscript. \textcite[128--129]{Minkowskilearnedbrahmin} makes a few remarks on \Nityananda's \Siddhantasindhu\ (based on \citeauthor{PetersonCatalogue}'s extract), while \textcite[269--27]{PingreeSSR} summarises (very briefly) his observations on the \Siddhantasindhu\ based on the manuscripts held at the Maharaja Sawai Man Singh II Museum at the City Palace of Jaipur.\label{peterson_catalogue_siddhantasindhu_reference}} At a very minimum, a comprehensive description of the structure and contents of the entire \Siddhantasindhu\ is certainly needed; however, such a task lies well beyond the scope of this study. Instead, I present below a few salient remarks from the \Siddhantasindhu\ that relate to the \ZijiShahJahani.

% \pagebreak %%%%%%%%%%%%%%%%%%%%% PAGEBREAK FOR FORMATTING ON PDF, CAN BE REMOVED LATER 

\begin{enumerate}[topsep=0pt]
    \item \textbf{On the content of MS~Kh}\quad \label{ms_kh_content_siddhantasindhu}The \Siddhantasindhu\ also begins with a detailed prolegomenon (\textit{granthārambha}, synonymous with \muqaddima) consisting of five sections (\prakara, synonymous with \qism) on ff.\thinspace 3r--11v~Kh\footnote{\,I cite manuscript references in the format $\langle$f.\thinspace folio$_\text{\#}$~Siglum$\rangle$ or $\langle$f.\thinspace folio$_\text{\#}$:\,line$_\text{\#}$~Siglum$\rangle$ throughout this paper. For instance, `f.\thinspace 22r: 9--11 Kh' indicates lines 9 to 11 on f.\thinspace 22r in Kh. \label{reference_folio_line_syntax}} (incomplete). Thereafter, it follows the structure of the \ZijiShahJahani\ with each part (\kanda, identified with \maqala) addressing a different subject and containing several chapters (\adhyaya, synonymous with \bab). However, 
    only the first three \maqalat\ of the \ZijiShahJahani\ appear to have been translated in the \Siddhantasindhu\ (as the three \kanda s); the fourth \maqala\ on miscellaneous astronomical calculations does not appear in Kh. A brief description of the content of Kh is as follows:
    \begin{itemize}
        \item The first part (\prathama-\kanda) with seven chapter describing the different calendrical eras (\textit{śāka}), \viz Arabic or Hijri (\textit{ārbīya}), \Shahjahan's (\textit{shāhjahām̐nīya}), Roman (\textit{raumīya}), Persian (\textit{phārasīya}), Malakī/Jalālī (\textit{ma\-lakīya}), Saṃvat (\textit{Hindukīya}), and Chinese-Uighūr Animal (\textit{khitāyīya-tur\-kīya}) on ff.\thinspace 12r--16v~Kh.
        \item The second part (\dvitiya\ \kanda) with twenty-two chapters describing various topics on finding the desired time (\textit{abhimata-samaya}) and the ascendant at that time (\textit{tātkālika-lagna}), as well as other topics related to it, on ff.\thinspace 17r--28v~Kh. More on this in \S~\ref{chapter_title_comparision_persian_sanskrit}.
        \item Tables (\kosthaka s) from ff.\thinspace 29r--97v~Kh.
        \item The third part (\trtiya-\kanda) with fifteen chapters describing the true (\sphuta) position and motion of celestial objects, and other topics related to it on ff.\thinspace 98v--111v~Kh.
        \item Tables (\kosthaka s) from ff.\thinspace 112r--436v~Kh.
    \end{itemize}
    %%%%
    \item \textbf{On \Shahjahan}\quad\label{shahjahan_epithet_siddhantasindhu} In the preamble, \Nityananda\ extols \Shahjahan\ with his encomiastic poetry \parencite[\eg \vid][129]{Minkowskilearnedbrahmin} and transliterates his Persian regnal epithet into \Nagari\ (on f.\thinspace 5v:\thinspace 15--16~Kh) as    
    \begin{multicols}{2}
    \tsans{abala-mujaphara-"saahibbadiina-\newline mahammada-saahiba-kiraana-saanii-\newline "saahajahaa-baadi"saaha-gaajii}
    \columnbreak
    
    \textit{abala-mujaphara-śāhibbadīna-mahammada-sāhiba-kirāna-sānī-śāhajahā-bādiśāha-gājī}
    \end{multicols}
        In his \Padshahnama, the seventeenth-century Mughal chronicler \AbdalHamidLahori\ states that upon ascending to the throne, Prince Shahāb al-Dīn Muḥammad Khurram (\Shahjahan) assumed the regnal name `Abū\Alif l-Muẓaffar Shahāb al-Dīn Muḥammad Ṣāhib-i Qirān-i Thānī' \parencite[6]{Elliot}. Among several other imperial epithets, \Shahjahan\ was called \tfarsi{صاحب قران ثاني} (\textit{ṣāhib-i qirān-i thānī}) `Second Lord of the Conjunction', \tfarsi{پادشاه غازی} (\textit{pādshāh-i ghāzī}) `Conqueror of Emperors' and \tfarsi{شاه جهان} (\textit{shāh-i jahān}) `King of the World'.
        
       \Nityananda\ traces the male ancestors of \Shahjahan\ from Tīmūr (\textit{taimūra}), Mīrān Shāh (\textit{mīrā-sāhā}), Sulṭān Muḥammad (\textit{sullā-mahaṃma}), Abū Sa\Ayn īd (\textit{abūsayīda}), Umar Shaykh (\textit{umara-śekha}), Bābur (\textit{bābara}), Humāyūn (\textit{humāū}), Akbar Shāh Jalāl al\=/Dīn (\textit{akabaraḥ śāha\-jallāladīna}), Jahāngīr (\textit{śrī-jahāṃgīra}), and finally, Shāh Jahān (\textit{śrīmān-sāhajahāṃ}) \parencite[vv.\thinspace 5--11 on p.\thinspace 230]{PetersonCatalogue}.\footnote{\,MS~Kh begins on f.\thinspace 2r in the middle of describing the genealogy of \Shahjahan. With the first folio missing and heavy damage to the second, I have relied on \citeauthor{PetersonCatalogue}'s transcription (\vid\ footnote~\ref{peterson_catalogue_siddhantasindhu_reference}) to fill in the missing parts.}
    %%%%    
     \item \textbf{On \AsafKhanshort}\label{asaf_khan_siddhantasindhu}\quad \AsafKhan, the prime minister (\vazir\idafaconsonant\ \azam) of \Shahjahan, is mentioned by name in the \Siddhantasindhu. On f.\thinspace 2v:\thinspace 10--11~Kh, we find the name \tsans{ma.mtrii vaasaaphakhaa/} `Minster \AsafKhanshort\ (\textit{vāsapha-khām̐})'. Through vv.\thinspace 21--23 on the same folio (lines 8--18), \Nityananda\ generously praises him as
     \tsans{yo raajyaahvayama.n.dapasya sud.r.dha.h stambha.h} (line 15) `he who is the steadfast pillar of [this] pavilion called the Empire' and \tsans{var.naa"sramapaalayat} (lines 17--18) `[he who is] protecting the [Hindu social system of] \textit{varṇāśrama}'.  
    %%%%
    \item \textbf{On \MullaFarid}\quad \label{mulla_farid_sanskrit_sidhantasindhu}\Nityananda\ identifies \MullaFarid\ by his name. On f.\thinspace 6r:\thinspace 2~Kh, we find the statement
    \tsans{mullaaphariida.m ibaraahimaputra.m .dhilliinivaasina.m} `\MullaFarid\ (\textit{mullā-pharīda}), the son of Ibrāhīm (\textit{ibarāhima-putra}) [and] resident of Delhi (\textit{ḍhillī-nivasin})'. This description agrees with \textit{ibn} [\textit{Ḥāfiẓ}] \textit{Ibrāhīm} `son of Ibrāhīm, [a man who has memorised the Qur\Alif ān]' and \textit{Dehlavī} `resident of Delhi' in \MullaFarid's full name. 
    %%%%
     \item \textbf{On the \ZijiShahJahani}\quad \label{zij_shahjahani_sanskrit_siddhantasindhu}
    Further along, on f.\thinspace 6v:\thinspace 19--20~Kh, we find the name of \MullaFarid's text, transliterated into \Nagari\ from the Persian \ZijiShahJahanifull, as
    \begin{multicols}{2}
    \tsans{kaaranaamai-saahiba-kiraana-saanii \newline jiica-"saaha-jahaanii}
     \columnbreak
    
     \textit{kāranāmai-sāhiba-kirāna-sānī \newline jīca-śāha-jahānī}
    \end{multicols}
    %%%% 
    \item \textbf{On \zij es and their types}\label{zij_types_defintion_siddhantasindhu}\quad    \Nityananda\ first distinguishes between the terms \zij, \tashil, and \taqvim\ on f.\thinspace 7r:\thinspace 9--10~Kh. According to him,
    \begin{multicols}{2}
    \tsans{jiica iti siddhaanta.h |\\
    tasahiila iti saara.nii |\\
    takaviima iti grahasphu.tatvam |}
     \columnbreak

    \zij\ (\textit{jīca}) is \Siddhanta\ (canon);\\
    \tashil\ (\textit{tasahīla}) is \sarani\ (table);\\ \taqvim\ (\textit{takavīma}) is the true position of a celestial object (\textit{graha-sphuṭatva}); in other words, an ephemeris.
    \end{multicols}
    On f.\thinspace 8v:\thinspace 2--3~Kh, he defines a \zij\ as \tsans{yasmingranthe sthuulasuuk.smaga.nitaani bhavanti tasya naama jiica iti~|} `The book (\textit{grantha}) in which both gross (\textit{sthūla}) and subtle (\textit{sūkṣma}) computations (\textit{gaṇita}) are found, that is called a \zij\ (\textit{jīca})'.   
    
    \Nityananda, like \MullaFarid, also classifies \zij es into the two categories of \zijiRasadi\ and \zijiHisabi: 
    \begin{itemize}
            \item On f.\thinspace 7v:\thinspace4--5~Kh, he defines \tsans{jiica-rasadii} (\textit{jīca-rasadī}) as that work (\textit{tantra}) which is well-established (\textit{dṛḍhī-kṛtya}) by the rules of observations (\textit{rasada-vidhāna}) and state the motion of celestial objects (\textit{graha-bhukti}) with tables (\kosthaka s); and
            \item On f.\thinspace 7v:\thinspace 27--19~Kh, he defines \tsans{jiica-hisaabii} (\textit{jīca-hisābī}) as the work containing tables (\kosthaka s) that correct (\textit{śodhyate}) previous tables in the table-writing tradition (\textit{koṣṭhaka-lekhaka-paramparā}) or those that bring out the genuine result (\textit{vāstava-phala}) by simple procedures (\textit{sugama-prakāra}) of computations (\textit{gaṇita}). 
        \end{itemize}
     %%%%
% \pagebreak %%%%%%%%%%%%%%%%%%%%% PAGEBREAK FOR FORMATTING ON PDF, CAN BE REMOVED LATER 

     \item \textbf{On \UlughBeg\ and other Islamicate astronomers, and their \zij es}\label{ulugh_beg_others_zijes_siddhantasindhu}\quad   \Nityananda\ identifies \UlughBeg\ by name in several places in the preamble. For instance, on f.\thinspace 7r:\thinspace 26~Kh, he refers to \UlughBeg\ as \tsans{mirajaa ulaga-bega} (\textit{mirajā ulaga-bega}), and on the very next line, informs us of \UlughBeg's demise with \tsans{parame"svarastasya svargavaasa.m karotu} `May God (\textit{parameśvara}) grant him residence in heaven (\textit{svargavāsa})'. The \ZijUlughBeg\ is translated as \tsans{jiica-ulaga-begii} (\textit{jīca-ulaga-begi}) on f.\thinspace 7v:\thinspace5~Kh.
     
    Along with Ulugh Beg, \Nityananda\ also described the names and works of several Islamicate astronomers (on ff.\thinspace 7v--8v~Kh). For example:
     \begin{itemize}
         \item the \tsans{jiica-jaame} and \tsans{jiica-baaliga} of \tsans{go"siyaara} \\
         \alZijalJami\ (\textit{jīca-jāme}) and \alZijalBaligh\ (\textit{jīca-vāliga}) of \Gilanifull\ (\textit{gośiyāra}), or
         \item the \tsans{jiica-khaakaanii-takamiila-jiica-yiila\-khaanii} 
    of \tsans{maulaanaa jama"seda kaa"sii} \\
    \KhaqaniZijfull\ (\textit{jīca-khākānī-takamīla-jīca-yīlakhānī}) of \alKashifull\ (\textit{maulāna-jamaśeda-kāśī}).
     \end{itemize}
     The names of these \zij es, and the order in which they are listed, appear to be identical in both the \ZijiShahJahani\thinspace\footnote{\,\textcite[48]{Ghori} provides the list of \zij es in \MullaFarid's \ZijiShahJahani, \vid\ discussion in \S~\ref{islamicate_zijes_mughal_india}.}  and the \Siddhantasindhu.
    \end{enumerate}

\paragraph{Motivation and purpose of composition}\label{motivation_purpose_of_siddhantasindhu} 
The twenty-fourth verse from the prolegomenon of \Nityananda's \Siddhantasindhu\ is barely legible on the damaged second folio (verso) of MS~Kh; however, it is attested in \textcite[231]{PetersonCatalogue} from the Alwar manuscript. The verse (in the \sardulavikridita\ meter) describes how \AsafKhanshort, having derived this inspiration (\textit{preraṇā}) from \Shahjahan, ordered \Nityananda\ to compose a proper treatise (\textit{su-tantra}) for the benefit of people (\textit{loka-upakāra}).\footnote{\,\Nityananda\ eulogises \Shahjahan\ as the `crest jewel of the ornamental crown of kings' (\textit{nṛpāla-mukuṭa-ālaṅkāra-cūḍāmaṇi})---an epithet he repeats in the colophons of every part of the \Siddhantasindhu, \eg \vid\ Part~II colophon on p.\thinspace \pageref{partII_colophon_sanskrit}.} Therefore, as \Nityananda\ says, he endeavoured (\textit{samīhe}) to compose the \Siddhantasindhu, a pure (\textit{amala}) and clear (\textit{sphuṭa}) Ocean of \Siddhanta s (\textit{siddhānta-sindhu}), resembling the illustrious \ZijiShahJahani\ (\textit{śrīmat-sahājahāṃ-prakāśam}). 


\paragraph{Style of composition} \label{prose_poetry_mixed_form}   
\Nityananda's \Siddhantasindhu\ is a mixture of prose-sentences (\textit{gadya}) and metrical verses (\textit{padya}). The transition from prose to poetry is ubiquitous throughout the text. It is worth noting that most \textit{explicit} Sanskrit translations of Islamicate texts (\eg the texts listed in Table~\ref{sanskrit_translations_of_islamic_texts}) are entirely in prose. I discuss a few aspects of \Nityananda's choice of using prose and poetry in relation to Part~II.6 in \S~\ref{structure_chapter_six}.
    
\subsubsection{Circulation of the \Siddhantasindhu} \label{circualtion_siddhantasindhu}

According to two of the four manuscripts of \Nityananda's \Siddhantasindhu\ held at the Maharaja Sawai Man Singh II Museum at the City Palace of Jaipur, there were nine copies of the \Siddhantasindhu\ prepared for distributing among notable seventeenth-century Mughal elites (mostly, Mughal \Subadar s or Provincial Governors), and an author's copy for \Nityananda\ \parencite[142]{PingreeJaipur}. A note, in vernacular Hindi, on f.\thinspace 443v of MSS~Khasmohor 4960 and Museum 23,\footnote{\,MS~Khasmohor 4960 belonged to Pīthīnātha in 1717; the horoscope of his son (dated Thursday 10 October 1717) is written on f.\thinspace 443. It is said to have been previously purchased from Manasārāma in 1696 for 250\thinspace ½ rūpas (f.\thinspace 444). Both these notes appear in vernacular Hindi \parencite[138--142]{PingreeJaipur}. The provenance of MS~Museum 23 is previously described in footnote~\ref{ss_museum_23_description}.} begins by stating that the original Sanskrit text (\textit{jīcamūlakarī}) remains in the library (\textit{kitābkhāṃnā}) of the emperor (\textit{pātiśāha}, Pādishāh). Individual copies of the text were given to the following recipients:   
\begin{enumerate}[topsep=0pt]
    \item \AzamKhan\ (\textit{ājama khāṃ}) of Bengal (\textit{ba}[\textit{ṃ}]\textit{gāla}), the governor of Bengal from 1632 to 1635;
    \item \AbdallahKhanFiruzJung\ (\textit{avdullaha khaṃ}) of Patna (\textit{paṭaṇā}), the governor of Bihar from 1632 to 1639;
    \item Ṣāḥib Ṣuba[dār?] of Benaras (\textit{banārasī sahava} \textit{sūva}/\textit{mūva}), unidentified;
    \item \ItiqadKhan\ (\textit{itakada khāṃ}) of Delhi (\textit{dillī}), also known as \MirzaShapur\ (d.\thinspace 1650), the brother of \AsafKhanshort;
    \item \KhwajaSabirKhanDauran\ (\textit{khāṃ nadorā}) of Ujjain (\textit{ujjayaṇa}), also known as \NasiriKhan\ (\textit{navaśeri khāṃ}), governor of Malwa from
    1631 to 1638;
    \item \MahabatKhanKhaniKhanan\ (\textit{mahavata khā[ṃ] khāṃnakhāṃnā}) in Burhanpur (\textit{burahānapura}), \circa \post\ 1633;
     \item \VazirKhan\ (\textit{ujīra khāṃ}), also known as \HakimShaykhIlmalDinAnsari, of Lahore (\textit{lāhora}), governor of Lahore from 1628 to 1639;
     \item \ZafarKhanAhsan\ (\textit{japhara khāṃ}) of Kashmir (\textit{kaśmīra}), the governor of Kashmir from 1632 to 1639 and from 1642 to 1646; and
     \item an unnamed recipient in Multan (\textit{mulatāna}).
\end{enumerate}    



 \clearpage

\section{The Sanskrit translation of a Persian \zij} \label{sanskrit_translation_of_persian_zij}
%%%%%%%%%%%%%%%%%%%%%%%%%%%%%%%%%%%%%%%%%%%%%%%%%%%%%%%%%%%%%%%%%%%%%%%%%%%%%%%%%%%%%%%%%
\subsection{Chronology and influence}\label{chronology_influence}
%%%%%%%%%%%%%%%%%%%%%%%%%%%%%%%%%%%%%%%%%%%%%%%%%%%%%%%%%%%%%%%%%%%%%%%%%%%%%%%%%%%%%%%%%

\Nityananda's \Siddhantasindhu\ 
% is a translation of \MullaFarid's \ZijiShahJahani. It
was completed in the early 1630s, making it one of the earliest Sanskrit translations of a Persian \zij. The \JicaUlugbegi, a Sanskrit translation of \ZijUlughBeg\ commissioned during \Akbar's reign (\ante\ 1605), is perhaps the earliest example. However, the extant fragmentary manuscripts only contain tables and star catalogues (\vid\ footnote~\ref{jicaulughbeg}) and not the canon (the text associated with the tables, often describing the theory and use of the tables).\footnote{\,The \JicaUlugbegi\ was prepared by a consortium of astronomers led by \AmirFathullahofShiraz, with the assistance of Sanskrit interpreters (\vid\  \S~\ref{Mughal_court}). As \textcite[367]{Sarmajyotisaraja} has pointed out, some of these interpreters were Sanskrit \jyotisa s: \eg Kishan Joshī is identified as  \Krsnadaivajna\ (\fl \circa 1600/25). If the authors' version of \JicaUlugbegi\ contained the canon of the \ZijUlughBeg, the following two reasons make the loss of the Sanskrit canon a notable (and regrettable) event.
\begin{itemize}
    \item \MullaFarid\ copied large parts of the \ZijUlughBeg\ (almost) verbatim in his \ZijiShahJahani; these parts have been translated by \Nityananda\ in his \Siddhantasindhu. The Sanskrit canon in the \JicaUlugbegi\ may have served as \Nityananda's model text to recopy just as easily as the \ZijUlughBeg\ served \MullaFarid.
    \item \MullaFarid\ was a direct pupil of \AmirFathullahofShiraz, and is believed to have `learned the rational sciences (\UlumiAqliyah) including astronomy and astrology' from him \parencite[720]{Ansarimedievalindia}. This suggests that, at the very least, \MullaFarid\ (and by association, \Nityananda) would have been familiar with a version of the \JicaUlugbegi\ closer to the original than what is currently extant.
\end{itemize}} % 

In contrast, as \S~\ref{siddhantasindhu_nityananda} describes, the \Siddhantasindhu\ includes tables as well as three (out of the four) distinct portions (\kanda s `parts' in Sanskrit or \maqalat\ `discourse' in Persian) of the canon of \ZijiShahJahani. The full extent of similarity between the content of these two texts can only be determined by an extensive future study; for our present discussion, however, we note that the \Siddhantasindhu\ expressly declares itself to be a Sanskrit version of the Persian original. Other Sanskrit texts in mathematical astronomy, \eg the three \siddhanta s mentioned in footnote~\ref{explicit_implicit_translation_sanskrit}, discuss or dismiss Islamicate ideas but they are not translations of any particular text. 

The only other Sanskrit `translation' that could predate \Siddhantasindhu\ (early 1630s) is the anonymously authored \Hayatagrantha\ written sometime before 1694 (\vid\ Table~\ref{sanskrit_translations_of_islamic_texts}). \textcite[327]{PingreeIslamAstSkt} correctly identifies the \Hayatagrantha\ as a translation of \aliQushji's Persian \Risaladarilmalhaya\ (dedicated to the Ottoman \MehmetII, r.\@ 1451--1481). \aliQushji's work was known to the Mughal court of \Humayun\ (r.\thinspace 1530--56) via a commentary on it, the \SharhilmalHaya\ written by \MuslihalDinMuhammadLari\ (d.\thinspace 1572) and dedicated to his patron \Humayun\ \parencite[296]{Pourjavady}. 

Following the \SharhilmalHaya, the \Hayatagrantha\ also discusses an assortment of topics (\prakarana) on planetary motion (\textit{graha-gati-nirupaṇa}) and the Earth-sphere (\textit{bhūgola-varṇana}) in two chapters (\adhyaya). In each topic, it explains Islamicate astronomical terms and provides equivalent Sanskrit expressions for them. \textcite[p.\thinspace 3 of the preface]{BhattacaryaHayata} thinks the \Hayatagrantha\ is either a translation (\textit{anuvāda}) of the Arabic text (\textit{arabī-deśasya grantha}) or a paraphrase of its summary (\textit{sāraṃ gṛhitvā likhita}); \textcite{PingreeIndianReception} adds to that by suggesting the author was `helped by a collaborator who was versed in Persian and Islamic astronomy, at least at a level sufficient for understanding the \textit{Risālah}' (p.\thinspace 475).\footnote{\,According to \textcite[CESS A4, p.\thinspace 57ab]{PingreeCESS}, there are seven extant manuscripts of the \Hayatagrantha, the earliest of which was copied in Oudh (Uttar Pradesh) in 1694. \citeauthor{BhattacaryaHayata}'s editions is based on three manuscripts held at the Sarasvatī Bhavana Granthālaya in Varanasi. In his preface (p.\thinspace 12), he cites internal evidence from the manuscripts to suggest (erroneously) that the \Hayatagrantha\ was composed in Kāśī in the eighteenth century \parencite[\vid][326--327]{PingreeIslamAstSkt}.} 

Beyond these speculations on the nature of the \Hayatagrantha, its influence on \Nityananda's \Siddhantasindhu\ is what concerns us here, and this is a difficult thing to ascertain at the present time. At the very outset, the anonymous authorship and vague timeline of the \Hayatagrantha\ make it problematic to situate it in relation to the \Siddhantasindhu. The Sanskrit words for Islamicate astronomical terms in both these texts are quite similar; however, their presentations often differ. 
The \Hayatagrantha\ includes transliterations of Arabic/Persian technical words along with their equivalent expressions in Sanskrit, whereas the \Siddhantasindhu\ translates Islamicate terms into Sanskrit without (always) transliterating the original Arabic/Persian word.\footnote{\,It is worth noting that \Nityananda's \Siddhantasindhu\ is not sanitised of all Arabic or Persian words. There are several instances where Islamicate names of authors, works, calendrical elements, \etcp\ are transliterated in \Nagari\ (\eg the excerpts in \S~\ref{siddhantasindhu_nityananda}). The text from the second part (\dvitiya\ \kanda) indicates that \Nityananda\ translates Arabic/Persian astronomical terms into equivalent (or original) Sanskrit expressions without referring to the original words.} 

For example, the \Hayatagrantha\ glosses the Persian word \tsans{baaadkokiba} (\textit{bāad-kokiba}) and the Sanskrit word \tsans{spa.s.ta-kraanti} (\textit{spaṣṭa-krānti}) as the `true declination' of a celestial object \parencite[p.\thinspace 19, lines 2--4]{BhattacaryaHayata}. \Nityananda\ also 
refer to the `true declination' (of celestial objects) with the Sanskrit word \tsans{spa.s.ta-kraanti} (\textit{spaṣṭa-krānti}); however, he does not transliterate the corresponding Persian technical expression used by \MullaFarid, namely \tfarsi{بعد كواكب از معدّل النهار} (\bud\idafaconsonant\ \kawakib\ \az\ \muaddil\ \alnahar) the `distance of a celestial object from the celestial equator'. (\Vid\ the chapter-titles of Discourse \hyperlink{Pii6}{II.6} and Part \hyperlink{Sii6}{II.6} in \S~\ref{chapter_title_comparision_persian_sanskrit}, page~\pageref{discourse_part_6_chapter_title_example}).


Nevertheless, the use of the same Sanskrit words to translate Islamicate astronomy does not establish an interdependence between these texts per se. It could also indicate a common written source (\eg a bilingual technical lexicon) or point towards a more institutionalised setup (like at the scriptoriums or \maktabkhana s) where these translations were produced. The emergence of a common technical vocabulary is then another aspect of cross-traditional discourses that occurred in the seventeenth century, making this a topic of exploration for future studies.   

%%%%%%%%%%%%%%%%%%%%%%%%%%%%%%%%%%%%%%%%%%%%%%%%%%%%%%%%%%%%%%%%%%%%%%%%%%%%%%%%%%%%%%%%%
\subsection{The \Siddhantasindhu, Part II: content and context}
\label{comparative_overview_chapters_zij_sindhu}
%%%%%%%%%%%%%%%%%%%%%%%%%%%%%%%%%%%%%%%%%%%%%%%%%%%%%%%%%%%%%%%%%%%%%%%%%%%%%%%%%%%%%%%%%

The second part (\dvitiya\ \kanda) of \Nityananda's \Siddhantasindhu\ includes twenty-two chapters on various topics that help determine the time of rising (\textit{udaya-samaya}) and degrees of ascension (\textit{udaya-aṃśa}) of the zodiacal signs at one's local latitude. The arrangement of the chapters in the \Siddhantasindhu\ follows the order of the twenty-two chapters in the second discourse (\maqala\idafaconsonant\ \duvum) of the \MullaFarid's \ZijiShahJahani\ identically. Table~\ref{table_contents_part_two} provides a chapter-wise list of topics covered in the \ZijiShahJahani\ Discourse II and \Siddhantasindhu\ Part II. 

The Persian and Sanskrit chapter-titles from these two texts, along with my English translations of the titular text, are presented in parallel in \S~\ref{chapter_title_comparision_persian_sanskrit}.  A full list of Persian and Sanskrit technical expressions in the respective chapter-titles of \ZijiShahJahani\ Discourse~II and \Siddhantasindhu\ Part~II is included in the glossary (beginning on page~\pageref{main}) grouped under their common English translations. I offer below a few general remarks on the language and scope of these chapters.

%%%%%%%%%%%%%%%
\begin{table}[!htbp]
\centering
\renewcommand{\arraystretch}{1.25}
\renewcommand{\baselinestretch}{1.25}\selectfont
\begin{tabularx}{\linewidth}{cX}
    \hline
   Chapter & List of topics\\
    \hline   
    \hyperlink{Pii1}{II.1} & Sexagesimal place-values of digits\\
    \hyperlink{Pii2}{II.2} & Method of interpolation between successive entries in a table\\
    \hyperlink{Pii3}{II.3} & Calculating Sine and Versed Sine values\\
    \hyperlink{Pii4}{II.4} & Calculating the shadow of a gnomon (\ie Cotangent values)\\
    \hyperlink{Pii5}{II.5} & Declination of the zodiacal signs\\
    \hyperlink{Pii6}{II.6} & Calculating the true declination of a celestial object\\
    \hyperlink{Pii7}{II.7} & Calculating the maximum elevation and depression of a celestial object \\
    \hyperlink{Pii8}{II.8} & Right ascensions of the zodiacal signs at the terrestrial equator\\
    \hyperlink{Pii9}{II.9} & Calculating the equation and hours of daylights at a terrestrial location\\
    \hyperlink{Pii10}{II.10} & Calculating the oblique ascensions of the zodiacal signs\\
    \hyperlink{Pii11}{II.11} & Inverse calculation of the right ascensions of the zodiacal signs from their oblique ascensions\\
    \hyperlink{Pii12}{II.12} & Calculating the right ascension and ecliptic longitude of a zodiacal sign culminating at the time of rising of a celestial object\\
    \hyperlink{Pii13}{II.13} & Calculating the right ascension of celestial objects at the time of their rising and setting (at the local horizon)\\
    \hyperlink{Pii14}{II.14} & Calculating the azimuth from the altitude of a celestial object\\
    \hyperlink{Pii15}{II.15}& Calculating the altitude from the azimuth of a celestial object\\
    \hyperlink{Pii16}{II.16} & Determining the line of the local meridian\\
    \hyperlink{Pii17}{II.17} & Determining the latitude and longitude of a terrestrial location\\
    \hyperlink{Pii18}{II.18} & Calculating the zenith-distance of the nonagesimal point\\
    \hyperlink{Pii19}{II.19} & Calculating the distance (in degrees) between two celestial objects \\
    \hyperlink{Pii20}{II.20} & Determining the direction of Mecca/\Kashi\ from a terrestrial location\\
    \hyperlink{Pii21}{II.21} & Determining the ascendant zodiacal sign (at the local horizon) corresponding to the altitude of a celestial object \\
    \hyperlink{Pii22}{II.22} & Determining the altitude of a celestial object corresponding to an ascending zodiacal sign (at the local horizon)\\
\hline    
\end{tabularx}
\caption{List of topics commonly discussed in the twenty-chapters of the \ZijiShahJahani\ Discourse~II and the \Siddhantasindhu\ Part~II}
\label{table_contents_part_two}    
\end{table} 

%%%%%%%%%%%%%%%

\begin{enumerate}[topsep=0pt]
    \item \MullaFarid's Persian text describes the computations without necessarily defining the terms first. In contrast, \Nityananda's discussions often begin with definitions (\textit{lakṣaṇa}s) of technical terms before describing the computational methods.  
    The \ZijiShahJahani\ is written in the format of a traditional \zij, and accordingly, it assumes its readers are familiar with technical expressions in Arabic/Persian. \Nityananda's \Siddhantasindhu, however,  describes Islamicate astronomy in Sanskrit to readers largely unfamiliar with the form or the language of the text. The prefatory definitions, along with the use of typical Sanskrit deictic words (like \textit{atha} `now', \textit{tat} `its/their', \textit{tatra} `there', \etcp) to introduce them, suggest an emphasis on communicating ideas effectively rather than simply translating the Persian text.  
    %%
    \item The ability to translate Arabic/Persian technical terms into Sanskrit requires a conceptual understanding of both Islamicate and Sanskrit astronomy, as well as a linguistic competence in navigating between these languages. \Nityananda's translations, presumably mediated through vernacular Hindi, reflect his command on the language of Sanskrit astronomy.
    
    In some instances, his expressions are literal translations (\textit{śabda-anuvāda}) of Arabic/Persian words; for example, \gls{ascendant} \tfarsi{طالع} (\tali) as \tsans{lagna} (\textit{lagna}), 
    \gls{genus} [of digits] \tfarsi{جنس} (\jins) as \tsans{jaatii} (\textit{jātī}),  or \gls{line_midday} \tfarsi{خطّ نصف النهار} (\khatt\ \nisf\ \alnahar) as \tsans{madhyaahna-rekhaa} (\textit{madhyāhna-rekhā}). In other instances, they appear to be figurative translations (\textit{bhāva-anuvāda}) based on an implied equivalence of technical meaning; for example,
    \begin{itemize}
    \item \gls{distance_celestial_object} \tfarsi{بعد كواكب} (\bud\idafaconsonant\ \kawakib) as \gls{true_declination} \tsans{spa.s.ta-kraanti} (\textit{spaṣṭa-krānti}),
    \item \gls{equation_of_daylight} \tfarsi{تعديل النهار} (\tadil\ \alnahar) as \gls{ascensional_difference} \tsans{cara} (\textit{cara}), or 
    \item \gls{latitude_visible_climate} \tfarsi{عرض اقلیم رؤیت} (\ard\idafaconsonant\ \iqlim\idafaconsonant\ \ruyat) as \gls{ecliptic_zenith_distance} \tsans{d.rkk.sepa} (\textit{dṛkkṣepa}). 

    \end{itemize}
    

    %%
    \item A salient aspect of \Nityananda's translation is \textit{localisation}; in other words, adapting the foreign content to suit the local context. The passages in his Sanskrit translation (of the second part) explain technical terms in greater detail, while his Sanskrit vocabulary preserves the meaning of Arabic/Persian words without referring to them expressly. Beyond these  communicative and semantic measures, \Nityananda\ also changes the context in which these foreign computational methods are applied.
    
    For example, the twentieth chapter in the second discourse of the \ZijiShahJahani\ describes the method to determine the azimuth (and inclination) of \qibla; in essence, the direction of Mecca. \Nityananda\ translates this chapter as the method to determine the direction of \Kashi. Having first discussed the mathematics of finding the direction of \Kashi\ from one's own location, he then goes on to apply the method to find the directions of other cities like Agra (\textit{argalapura}) and Mecca (\textit{makkaapura}) as illustrative examples (\textit{udāharaṇa}). His translation not only captures the mathematical essence of the chapter from the \ZijiShahJahani\ but also translates the cultural context of locating sacred and imperial cities. 
    %%
    \item \label{triprasna_remark} In the colophon, \Nityananda\ states that the second part contains discussions `accompanied by many statements and rationales on the ``three questions"' (\textit{tripraśna-pracura-ukti-yukti-sahita}). The \triprasnadhikara\ is a separate chapter (\adhyaya\ or \adhikara) in Sanskrit \siddhanta s that discuss methods to find the cardinal directions (\textit{diś}), the local latitude (\textit{deśa}), and the times (\textit{kāla}) of various celestial and terrestrial phenomena. By referring to the \triprasna, \Nityananda\ again bring a familiar context to situate an otherwise curious collection of chapters. In fact, the practice of writing benediction verses (\mangalacarana) at the beginning of every part of the book and a closing colophon at the end of each part is a Sanskrit \siddhantic\ trait not seen in the \ZijiShahJahani. 
\end{enumerate}

%%%%%%%%%%%%%%%%%%%%%%%%%%%%%%%%%%%%%%%%%%%%%%%%%%%%%%%%%%%%%%%%%%%%%%%%%%%%%%%%%%%%%%%%%
\subsection{The \Siddhantasindhu, Part II.6: structure and language}
\label{chapter_six_zij_sindhu_comparision}
%%%%%%%%%%%%%%%%%%%%%%%%%%%%%%%%%%%%%%%%%%%%%%%%%%%%%%%%%%%%%%%%%%%%%%%%%%%%%%%%%%%%%%%%%

The sixth chapter (\textit{saṣṭhādhyāya}) from the second part of \Nityananda's \Siddhantasindhu\
describes three methods to compute the true declination of a celestial object. A celestial object is variously understood as a planet, a star, or an asterism that moves in the celestial sphere. Typically, this object posses a non-zero ecliptic latitude and hence its declination is different from the Sun that moves on the ecliptic. \Nityananda\ methods to compute the true declination of such an object are identical to those stated by \MullaFarid\ in the sixth chapter of the second discourse of his \ZijiShahJahani. An edition of the original text from \MullaFarid's \ZijiShahJahani\ Discourse~II.6 and \Nityananda's \Siddhantasindhu\ Part II.6, along with my English translations of corresponding Persian and Sanskrit passages, are included in \S\S~\ref{zijshahjahan_persian_english} and \ref{siddhantasindhu_sanskrit_english} separately. The technical glossary (beginning on page~\pageref{main}) also includes a list of technical expressions found in the Persian and Sanskrit passages of this chapter.

These methods rely on astronomical quantities that are distinctly Islamicate in their origin, \eg the \glslink{second_declination_parent}{second declination} of a celestial object or the \glslink{arc_max_declination}{arc of maximum argument of the distance}. However, \Nityananda's exposition of these Islamicate methods is uniquely original in its presentation and style. In the following paragraphs, I present a few remarks on the structure and the language of \Nityananda's Sanskrit text in comparison with \MullaFarid's Persian. The mathematical aspects of these methods is to appear in \textcite[forthcoming]{MisraTD}.


% The translated English expressions are emboldened for emphasis in the following paragraphs.     

\subsubsection{Structure of the text} \label{structure_content_zij_sindhu_chapter_six}
The methods to compute the \gls{distance_celestial_object} (\bud\idafaconsonant\ \kawkab\ \az\ \muaddil\ \alnahar) in the \ZijiShahJahani\ Discourse~II.6 are translated in the \Siddhantasindhu, Part II.6 as methods to compute the \gls{true_declination} (\spasta-\kranti). In \S\S~\ref{zijshahjahan_persian_english} and \ref{siddhantasindhu_sanskrit_english}, I have grouped the Persian and Sanskrit text of this chapter into comparable passages (numbered `[1]', `[2]', \etcp) to highlight their grammatical and mathematical likeness. (\Vid\ my editorial conventions in \S~\ref{chapter_vi_zij_sindhu}.) Table~\ref{comparision_passage_zij_sindhu} provides an outline of the passages in the two texts with a brief description of their content. 

\begin{table}[!htbp]
\centering
\renewcommand{\arraystretch}{1.2}
\renewcommand{\baselinestretch}{1.2}\selectfont
\begin{tabularx}{\textwidth}{ccX}
    \hline
   \multicolumn{2}{c}{Passages in II.6} & \multirow{2}{*}{Description}\\
    \cline{1-2}   
\ZijiShahJahani & \Siddhantasindhu &\\
\hline
\multicolumn{3}{c}{First Method}\\
{[\hyperlink{Ppass1}{1}]}$\,_\text{prose}$ & [\hyperlink{Spass1}{1}]$\,_\text{verse}$ & Definition of \glslink{argument_of_distance}{argument of the distance} \\
{[\hyperlink{Ppass2}{2}]}$\,_\text{prose}$ & [\hyperlink{Spass2}{2}]$\,_\text{verse}$ & First method to compute the \glslink{sine_true_declination}{Sine of the true declination}\\
\multicolumn{3}{c}{Second Method}\\
{[\hyperlink{Ppass3}{3}]}$\,_\text{prose}$ & [\hyperlink{Spass3}{3}]$\,_\text{verse}$ & Second method to compute the \glslink{sine_true_declination}{Sine of the true declination}\\
{[\hyperlink{Ppass4}{4}]}$\,_\text{prose}$ & [\hyperlink{Spass4}{4}]$\,_\text{prose}$ & Alternative second method to compute the \glslink{sine_true_declination}{Sine of the true declination} (using tables)\\
{[\hyperlink{Ppass5}{5}]}$\,_\text{prose}$ & [\hyperlink{Spass5}{5}]$\,_\text{prose}$ & Case one: no [ecliptic] \gls{latitude}\\
{[\hyperlink{Ppass6}{6}]}$\,_\text{prose}$ & [\hyperlink{Spass6}{6}]$\,_\text{prose}$ & Case two: [ecliptic] \gls{latitude} with no [first] \glslink{declination_parent}{declination}\\
{[\hyperlink{Ppass7}{7}]}$\,_\text{prose}$ & [\hyperlink{Spass7}{7}]$\,_\text{prose}$ & Case three: [first] \glslink{declination_parent}{declination} equals the \glslink{maximum_declination_parent}{obliquity of the ecliptic} \\
\multicolumn{3}{c}{Third Method}\\
-- & [\hyperlink{SpassA}{α}]$\,_\text{verse}$ & Definitions of the \glslink{solstitial_colure}{solstitial colure} and the \gls{circle_congruent_ecliptic}\\
-- & [\hyperlink{SpassB}{β}]$\,_\text{verse}$ & Definition of the arc of \gls{maximum_true_declination}\\
-- & [\hyperlink{SpassC}{γ}]$\,_\text{verse}$ & Definition of the arc of \gls{maximum_latitude}\\
-- & [\hyperlink{SpassD}{δ}]$\,_\text{verse}$ & Definitions of the \gls{congruent_bhuja} and the \gls{congruent_koti}\\
{[\hyperlink{Ppass8}{8}]}$\,_\text{prose}$ & [\hyperlink{Spass8}{8}]$\,_\text{verse}$ & Calculating the \glslink{Sine_distance_solstitial_colure_parent}{Sine of the distance along the `circle congruent to the ecliptic' from the solstice}\\
{[\hyperlink{Ppass9}{9}]}$\,_\text{prose}$ & [\hyperlink{Spass9}{9}]$\,_\text{verse}$ & Calculating the \glslink{arc_max_latitude_parent}{arc of maximum latitude}\\
{[\hyperlink{Ppass10}{10}]}$\,_\text{prose}$ & [\hyperlink{Spass10}{10}]$\,_\text{verse}$ & Calculating the \glslink{arc_max_declination}{arc of maximum argument of the distance}\\
{[\hyperlink{Ppass11}{11}]}$\,_\text{prose}$ & [\hyperlink{Spass11}{11}]$\,_\text{verse}$ & Third method to compute the \glslink{sine_true_declination}{Sine of the true declination}\\
\hline
\end{tabularx}
\caption{Description of the passages in \ZijiShahJahani\ Discourse~II.6 (in \S~\ref{zijshahjahan_persian_english}) and \Siddhantasindhu\ Part II.6 (in \S~\ref{siddhantasindhu_sanskrit_english}).}
\label{comparision_passage_zij_sindhu}
\end{table}

As Table~\ref{comparision_passage_zij_sindhu} shows, \Nityananda's third method includes four additional passages [\hyperlink{SpassA}{α}--\hyperlink{SpassD}{δ}] in metrical verses that are not found in \MullaFarid's Persian text. The mathematics of the third method requires a knowledge of several astronomical concepts that are not very commonly known in Sanskrit astronomy. \Nityananda\ defines these Sanskrit terms in the four passages before using them in the third method of computation. His style of interspersing the (\textit{mūla}-like) translations with (\textit{bhāṣya}-like) commentarial passages, in metrical Sanskrit, is as unique as the content of these passages themselves. These extra verses reveal his intentions (and ability) to communicate this new astronomy systematically and not slavishly translate the Persian into Sanskrit. \Nityananda's \Siddhantasindhu\ is among the earliest texts to Sanskritise Islamicate astronomy; and in doing so, conceive several original (or equivalent) Sanskrit expressions for Islamicate technical terms. 

\paragraph{Mixture of prose and poetry}\label{prose_poetry_siddhantasindhu_chapter_six}
\Nityananda's Sanskrit translation of the sixth chapter is a mixture of prose sentences and metrical verses. As Table~\ref{comparision_passage_zij_sindhu} indicates, \MullaFarid's Persian passages are written entirely in prose while the Sanskrit text includes eleven passages in metre [\hyperlink{Spass1}{1}--\hyperlink{Spass3}{3}, \hyperlink{SpassA}{α}--\hyperlink{SpassD}{δ}, and \hyperlink{Spass8}{8}--\hyperlink{Spass11}{11}] and four passages in prose [\hyperlink{Spass4}{4}--\hyperlink{Spass7}{7}]. The metrical verses are numbered while the prose passages are are not.\footnote{\,The numbering scheme restarts at one for every metrical part following a prose interlude, \vid\ Remark in Table~\ref{mss_description_siddhantasindhu}.} These poetic verses are composed in a range Sanskrit metres:
\begin{itemize}[topsep=0pt]
    \item Passages [\hyperlink{Spass2}{2}, \hyperlink{Spass8}{8}]: eight-syllabled \textit{pramāṇikā}
    \item Passage [\hyperlink{Spass3}{3}]: eight-syllabled \textit{anuṣṭubh}
    \item Passage [\hyperlink{Spass9}{9}, \hyperlink{Spass10}{10}]: eleven-syllabled \textit{rathoddhatā}
    \item Passages [\hyperlink{Spass1}{1}, \hyperlink{Spass11}{11}]: twelve-syllabled \textit{vaṃśasthavila}
    \item Passages [\hyperlink{SpassB}{β}, \hyperlink{SpassC}{γ}]: twelve-syllabled \textit{drutavilambita}
    \item Passage [\hyperlink{SpassA}{α}]: seventeen-syllabled \textit{pṛthvī}
    \item Passage [\hyperlink{SpassD}{δ}]: \textit{āryā} \textit{jāti}-metre
\end{itemize}

The \Siddhantasindhu\ is the first  explicit translation of an Islamicate astronomical text to use metrical Sanskrit (\vid\ remark on the style of composition on page~\pageref{prose_poetry_mixed_form}). Its intention to Sanskritise the content for local \jyotisa s may explain the use of metrical verses; however, it is not entirely clear why \Nityananda\ chooses to then translate certain Persian passages in prose. I list below my observations on these prose passages. 
\begin{enumerate}[topsep=0pt]
    \item All three methods described in the text are prescriptive: they outline the constituent terms and then suggest a computational formula using these terms. There are no mathematical derivation or explanations given. The prose passages include an interpolative method of computation (in passage [\hyperlink{Spass4}{4}]) and three special cases (in passages  [\hyperlink{Spass5}{5}--\hyperlink{Spass7}{7}]), both based on the formula of the second method (in passage [\hyperlink{Spass3}{3}]). They appear to illustrate the (use of the) formula rather than being ancillary to it.
    \item The interpolative method in passage [\hyperlink{Spass4}{4}], and again in passage [\hyperlink{Spass6}{6}], refers to the \gls{table_Cosine_maximum_declination}. MS~Kh of the \Siddhantasindhu\ includes the text of these passages but not the table itself. In comparison, MSS~\SjA\ (f.\thinspace 21b) and \SjB\ (f.\thinspace 16a) of the \ZijiShahJahani\ present the table alongside the Persian text of the chapter.
    \item The Sanskrit verses from the \Siddhantasindhu, Part II.6 are also found in the \Sarvasiddhantaraja, \ganitadhyaya\ `chapter on computations'. \Nityananda\ copies these metrical verses verbatim into the \spastakrantyadhikara\ `section on true declination' of his \Sarvasiddhantaraja\ but leaves out the four prose passages from the \Siddhantasindhu. A critical edition and technical translation of the \spastakrantyadhikara\ is to appear in \textcite[forthcoming]{MisraTD}. 
\end{enumerate}
 
\subsubsection{Language of the text} \label{language_content_zij_sindhu_chapter_six}
\Nityananda's Sanskrit passages of the sixth chapter follow \MullaFarid's  Persian text in more than their mathematical content. In the following remarks, I note some of the linguistic features of \Nityananda's compositions in comparison to \MullaFarid's Persian sentences. All grammatical terms are abbreviated in these remarks, with an expanded list of abbreviations included on page~\pageref{acronym}. 
\begin{enumerate}[topsep=0pt]
    \item \textbf{Grammatical similarity}
    \begin{enumerate}[topsep=0pt]
    \item  \label{subject_fronting_passage_1} \textbf{Subject-fronting}\quad  The inflected grammar of Sanskrit allows a flexible word-ordering in most prose sentences, and even more so, in metrical verses. \Nityananda\ utilises this syntactic freedom and composes some of his verses to resemble \MullaFarid's Persian statements quite closely. For example,
    the Persian text in passage~[\hyperlink{Ppass1}{1}] is a typical conditional sentence where the subjects precede the conditional proposition, \viz
    {\par\centering
    $\underbrace{\text{[Given]\enskip X \& Y}}_{\text{subject-fronted events}}$\quad:\quad $\underbrace{\text{if\enskip $\mathcal{C}$\,(X, Y)}}_{\text{if-clause}}$\quad $\longrightarrow$\quad $\underbrace{\text{then\enskip $\mathcal{A}_\text{1}$\,(X, Y),\enskip else\enskip $\mathcal{A}_\text{2}$\,(X, Y)}}_{\text{then-clause}}$\par}    
    where $\mathcal{C}$\,(X, Y) is the condition involving events `X' and `Y' in the if-clause, $\mathcal{A}_\text{1}$\,(X, Y) is the first action of the two events in the then-clause, and $\mathcal{A}_\text{2}$\,(X, Y) is the second alternative action of the two events in the then-clause. 
    
    Persian grammar specifies various kinds of conditional sentences, and accordingly, the verbs in the if-clause (protasis) are in different verbal moods, \eg if the condition is a proposition to be fulfilled, like in passage~[\hyperlink{Ppass1}{1}], the verb in the protasis is the subjunctive mood: \tfarsi{باشند} (\textit{bāshand})  \acrshort{present}-\acrshort{subjunctive} `should be'. The consequent then-clause (apodosis) can chose between indicative or subjunctive moods of verbs according to the context of the sentence. \Nityananda's Sanskrit verse in passage~[\hyperlink{Spass1}{1}] mirrors this subject-fronted conditional construction with the verb in the protasis in the optative mood: \tsans{bhavet} (\textit{bhavet}) \acrshort{optative}-\acrshort{active} `should be' (conditional possibility). The apodosis that follows states the consequent actions as statements to be understood as implicit instructions, \viz\ then [we take/do] $\mathcal{A}_\text{1}$\,(X, Y), otherwise [we take/do] $\mathcal{A}_\text{2}$\,(X, Y).
    
    \item \label{implied_modality}\textbf{Implied modality}\quad The realis mood of a present indicative statement in Persian also implies an irrealis future potential. For example, passage~[\hyperlink{Ppass2}{2}] includes two statements in the following form:
    {\par\centering
    \begin{tabular}{c} First (instructive)\\ statement \end{tabular}%
    \quad$\overset{\text{implied}}{\underset{\text{consequence}}{\longrightarrow}}$\quad%
    \begin{tabular}{c} Second (declarative)\\ statement. \end{tabular}\par}
    The second (declarative) statement uses the verb \tfarsi{بود} (\textit{buvad}) 
    \acrshort{present}-\acrshort{indicative} `is' to indicate `the result \textit{is} [something]'; this, in effect, also convey the meaning `the result \textit{will be} [something]'. \Nityananda's Sanskrit translation of this passage retains the form, and uses the verb \tsans{bhavet} (\textit{bhavet}) \acrshort{optative}-\acrshort{active} `will be' (future probability) to indicate a similar meaning: `the result \textit{will be} [something]'.
    \end{enumerate}

    \item \textbf{Semantic equivalents}\quad \Nityananda\ translates Islamicate astronomical terms using equivalent Sanskrit expressions, some of which, are literal translations; for example,
    \begin{itemize}
    \item  \gls{circle_four_poles} \tfarsi{دایرهٔ ماره باقطاب اربعه} (\textit{\dayiri\idafavowel\ \marri\ \biaqtab\idafaconsonant\ \arbai}) as \tsans{dhruva-catu.ska-yaata-v.rtta} (\textit{dhruva-catuṣka-yāta-vṛtta})
    \item \gls{latitude_celestial_object} \tfarsi{عرض کوکب} (\ard\idafaconsonant\ \kawkab) as \tsans{khagasya baa.na} (\textit{khagasya bāṇa}), or
    \item \gls{one_direction} \tfarsi{یک جهت}  (\yik\ \jahat) as  \tsans{eka\=/di"s} (\textit{eka\=/diś}).
        % \item \tfarsi{میل کلّی} (\mayl\idafaconsonant\ \kulli) `total/greatest declination' as \tsans{parama-kraanti} (\textit{parama-kraanti}),
    \end{itemize}
    In other instances, his translations use novel (seemingly, didactic) Sanskrit expressions to capture the implied mathematical meaning of Islamicate terms; for example,
    \begin{itemize}
        \item \gls{Cosine_inverse_declination_degree_celestial_object}\linebreak \tfarsi{جیب تمام میل منکوس درجه كوكب}  (\jayb\idafaconsonant\ \tamam\idafaconsonant\ \mayl\idafaconsonant\ \mankus\idafaconsonant\ \daraji\idafavowel\ \kawkab)\footnote{\,The use of the term `inverse declination' (\almayl\ \almakus) to indicate the [first] declination of the ecliptic longitude of a celestial object increased by ninety first appears in the works of thirteenth century \Maragha\ astronomers, \eg \ZijIlkhani\ of \alTusi\ \parencite[188]{HamadanialTusi} or \Tajalazyaj\ of \alMaghribi\ \parencite[196]{Dorce}.} as \gls{day_Sine_increased_by_three_signs} \tsans{sa-bha-traya-dyujiivaa} (\textit{sa-bha-traya-dyujīvā}),
        \item \gls{distance_celestial_object_solstice} \tfarsi{بعد کوکب از
        \tfarsib{دایرهٔ ماره باقطاب اربعه}} (\textit{\bud\idafaconsonant\ \kawkab\ \az\ \guillemotleft\dayiri\idafavowel\ \marri\ \biaqtab\idafaconsonant\ \arbai\guillemotright})  as \gls{congruent_koti}  \tsans{sad.r"s-ko.ti} (\textit{sadṛś-koṭi}),
        \item \gls{first_arc} \tfarsi{قوس اوّل} (\qaws\idafaconsonant\ \avval) as \gls{maximum_latitude} \tsans{para-i.su} (\textit{para-iṣu}), or
        \item \gls{second_arc} \tfarsi{قوس دوم} (\qaws\idafaconsonant\ \duvum) as \gls{maximum_true_declination} \tsans{para-sphu.ta-apama} (\textit{para-sphuṭa-apama}).
    \end{itemize}
    
    \item \textbf{Hybrid translations}\quad Certain Sanskrit words appear to be a mix of literal and figurative translations, \eg 
    the `argument of the distance', understood as the arc of the great circle passing through the two ecliptic poles and a celestial object, and lying between a celestial object and the celestial equator, is called \tfarsi{حصّهٔ بعد} (\hissi\idafavowel\ \bud) `share of the distance' in Persian.
    \Nityananda\ translates this Persian expression as \tsans{sphu.ta-apama-a.m"sa} (\textit{sphuṭa-apama-aṃśa}) `share of true declination' in passage~[\hyperlink{SEpass1}{1}], but then goes on to translate it as \tsans{sphu.ta-apama-a"nka} (\textit{sphuṭa-apama-aṅka}) `curve of the true declination'\footnote{\,Typically, the word \tsans{a"nka} (\textit{aṅka}) refers to a `digit/ number', or more literally, a `mark/sign'. \Nityananda\ uses the word to signify a geometrical `arc' or `curve' of a great circle. There are several geometrical explanations in the \Siddhantasindhu\ as well as the \Sarvasiddhantaraja\ that validate the interpretation of the word \textit{aṅka} as `arc/curve' \parencite[\eg \vid][279]{Misrathesis}.} in subsequent passages~[\hyperlink{SEpass2}{2}--\hyperlink{SEpass4}{4}]. The change from the semantic equivalent (\textit{aṃśa} `share') to the pragmatic (\textit{aṅka} `curve') can either be a simple (and serendipitous) artefact of the single manuscript witness, or it may be a more volitional change meant to emphasise the didactic meaning in using the word.  

    \item \textbf{Original expressions}\quad  As noted earlier, \Nityananda\ composes four passages~[\hyperlink{SpassA}{α}--\hyperlink{SpassD}{δ}] in his Sanskrit text to explain the
    terms involved in the third computational method. These terms are not directly expressed in the Persian text, and hence, \Nityananda's expressions are original in their language and their substance. For example, passage~[\hyperlink{SEpassA}{α}] defines the \gls{circle_congruent_ecliptic} \tsans{bhacakra-sad.r"sa-v.rtta} (\textit{bhacakra-sadṛśa-vṛtta}), a great circle passing through the two equinoctial points and a celestial object, and resembling the ecliptic. For planetary objects, this is the orbit of the planet with the longitude of its node being zero; or as \Nityananda\ later indicates in passage~[\hyperlink{SEpassC}{γ}], when the \gls{conjunction_equinox_node} \tsans{vi.suva-paata-yuga} (\textit{viṣuva-pāta-yuga}) is being assumed.
    
    \item \textbf{Synonymy and ambiguity}\quad In rendering \MullaFarid's Persian sentences into Sanskrit, \Nityananda\ employs a variety of Sanskrit synonyms (or near-synonyms) to translate Persian technical terms. For example, a celestial object is called \tfarsi{کوکب} (\kawkab) in Persian, whereas \Nityananda\ variously uses the words  \tsans{khaga} (\textit{khaga}), \tsans{graha} (\textit{graha}),  \tsans{dyucara} (\textit{dyucara}), \tsans{nak.satra} (\textit{nakṣatra}), \tsans{nabhoga} (\textit{nabhoga}), and \tsans{bha} (\textit{bha}) to describe such an object. The abundance of Sanskrit synonyms allows him to choose words that suits the metre of his verse; however, the metrical constraints also makes his translations a little vague in some instances. For example, the Persian text in passage~[\hyperlink{PEpass1}{1}] refers to the \tfarsi{جهت حصّهٔ بعد} (\jahat\idafaconsonant\ \hissi\idafavowel\ \bud) `direction of the share of the distance' as being in the \tfarsi{جهت مجموع}  (\jahat\idafaconsonant\ \majmu) `direction of the sum' or the \tfarsi{جهت فضل} (\jahat\idafaconsonant\ \fadla) `direction of the difference/residue'; in other words, dependant on the directions of the two quantities that make up the `share of the distance'. In his 
    Sanskrit translation in passage~[\hyperlink{SEpass1}{1}], \Nityananda\ simply states that the `share of the true declination' is in its \tsans{sva-di"s} (\textit{sva-diś}) `own direction'.
    
    \item \textbf{Lowering a sexagesimal number}\quad Of particular note is an Islamicate arithmetic operation that, to my knowledge, does appear in any Sanskrit astronomical or mathematical text till the \Siddhantasindhu.  
    
    Islamicate texts often include the operation \tfarsi{منحطّ} \munhatt\ `lowering' a sexagesimal number before multiplying (\gls{low_multiplication}) or before dividing (\gls{low_division}). In effect, shifting the fractional point leftwards to \textit{lower} the value of the sexagesimal number before operating on it; or in other words, dividing a number by the Radius (\textit{sinus totus}) of 60.\footnote{\,For instance, \alKashi\ uses the word \munhatt\ `to depress' a sexagesimal number, \ie divide it by the Radius of 60, in relation to his Sine computations in his \KhaqaniZij\ (\circa 1413/1414) \parencite[40]{Hamadani_Zadeh_Khaqani}.} \Nityananda\ uses the verb √\tsans{adharii-k.r} (√\textit{adharī-kṛ}) `to make [something] low' or `lower' to indicate a division by 60.\footnote{\,MS Benares (1963) 37079 of the \Sarvasiddhantaraja\ from the Sarasvatī Bhavana Granthālaya (Varanasi) parses the word \tsans{adhara} \textit{adhara} in relation to the divisor (in the last \pada\ of verse 6) as \begin{quote}\tsans{adhara-sa-.sa.s.ti-bhakta-bhaajaka-bhajana.m atraadhara-bhajana-sa.m{\tsnb{ज्ञं}} ucyate} (f.\thinspace 63r:~11)\newline
    \textit{adhara} is the sixtieth part [\lit with sixty divided] of the divisor of the division; here it is referred to as the \textit{adhara}-division by name.
    \end{quote}}
    His translation of the term is more literal compared to later authors like \Nayanasukhopadhyaya\ who, in his \Sarahatajakiravirajandi\ (1729), translates \munhatt\ from \alBirjandi's \SharhalTadhkirah\ (1507) as \textit{ṣaṣtyāpta}/\textit{ṣaṣṭibhakta} `divided by sixty'  \parencite[265]{KusubaPingree}.
\end{enumerate}

%%%%%%%%%%%%%%%%%%%%%%%%%%%%%%%%%%%%%%%%%%%%%%%%%%%%%%%%%%%%%%%%%%%%%%%%%%%%%%%%%%%%%%%%%
\subsection{Grammatical notes on translation} \label{grammatical_notes_translation}
%%%%%%%%%%%%%%%%%%%%%%%%%%%%%%%%%%%%%%%%%%%%%%%%%%%%%%%%%%%%%%%%%%%%%%%%%%%%%%%%%%%%%%%%%
In translating the Persian text from \MullaFarid's \ZijiShahJahani, Discourse~II.6 and \Nityananda's \Siddhantasindhu, Part~II.6 into English, I have interpreted the verbs (and verbal derivatives) based on their implied modality in the sentence. Hence, a particular form of a verb (\eg an optative active in Sanskrit) is sometimes translated differently in different sentences. I list below some of the main aspects of my translations. 

\paragraph{Persian to English} 
\begin{enumerate}[topsep=0pt]
    \item The Persian subjunctive mood expresses a variety of meanings based on the context of the sentence. \MullaFarid's Persian text of the sixth chapter uses the present subjunctive form of verbs quite commonly. Translating these verbs using English indicative forms does not fully capture the subjunctive mood of the original sentences. Hence, I translate the Persian subjunctive verbs in my English translations with the modal verb `should', \eg \tfarsi{باشد} (\textit{bāshad})  \acrshort{present}-\acrshort{subjunctive}-\acrshort{singular}·\acrshort{third} `[he/she/it] should be' instead of `[he/she/it] is' or `[he/she/it] is to be'. This helps distinguish between the Persian indicatives (realis) and subjunctives (irrealis) in the passages, particularly, in the case of conditional clauses.
                
    \item An impersonal passive sentence may be constructed in Persian with a third-person plural conjugation of the verb and a dismissive (or vague) subject. For example, the Persian text in passage~[\hyperlink{Ppass4}{4}] uses the verbs \tfarsi{درآرند} (\textit{dar ārand}) \acrshort{present}-\acrshort{indicative}-\acrshort{plural}·\acrshort{third} `[they] extract', along with a direct object (identified by the \tfarsi{را} (\textit{rā}) marker) and no specific subject. The only indication of the subject is found in the enclitic conjugation of the verb. A syntax-preserving translation of this sentence reads `they extract [the object]'; however, it can also be understood as the impersonal passive sentence `[the object] is extracted'. As grammatical opinions on passive constructions in Persian vary  \parencite[\eg \vid][261--264]{NematiPersianpassive}, I choose to retain the syntax-preserving form in my English translation in passage~[\hyperlink{PEpass4}{4}].
\end{enumerate}

\paragraph{Sanskrit to English} 
\begin{enumerate}[topsep=0pt]
\item The subjunctive mood is obsolete in classical Sanskrit and is replaced by the use of the  optative. Like the subjunctive, the optative mood also indicates various meanings depending on the context of the verb. \Nityananda's Sanskrit passages from the sixth chapter use several verbs in their optative active form. I translate these forms according to their syntactic location (\ie whether they occur in principal realis sentences or subordinating irrealis clauses) and their modal intention. For example, as I previously alluded in remarks~\ref{subject_fronting_passage_1} and \ref{implied_modality} in \S~\ref{language_content_zij_sindhu_chapter_six}, the verb \tsans{bhavet} (\textit{bhavet}) \acrshort{optative}-\acrshort{active}-\acrshort{singular}·\acrshort{third} implies both `[he/she/it] should be' (in conditional clauses) and `[he/she/it] will be' (in potential statements) depending on the context. 

\item More generally, I use the the English modal verb `should' to translate most Sanskrit optative forms to distinguish them from the indicative forms in English \eg \tsans{syaat} (\textit{syāt}) \acrshort{optative}-\acrshort{active}-\acrshort{singular}·\acrshort{third} `[he/she/it] should be/exist' instead of `[he/she/it] is/exists'. In such translations, the English modal verb `should' is more epistemic than deontic in conveying the irrealis mood (in other words, \textit{should} conveys the sentiment of possibility or inference and not a directive or exhortation). Also, the use `should' to mark the Persian subjunctive and the Sanskrit optative in the English translations reveal a syntactic similarity in the Persian and Sanskrit passages of the sixth chapter.
\end{enumerate}

A complete list of the Persian verbs in the \ZijiShahJahani, Discourse~II.6 and the Sanskrit verbs (and verbal derivatives) in \Siddhantasindhu, Part~II.6 can be found in Appendices~\ref{persian_verbs} and \ref{sanskrit_verbs} respectively.   
    

   
\clearpage

\section{Editorial Notes}\label{editoial_notes}
\renewcommand{\thefootnote}{\arabic{footnote}} % Hindu-Arabic footnote numbering

I discuss below the orthographic standards and transcription/transliteration schemes I have adopted in editing the Persian text (from \SjA\ \& \SjB) and the Sanskrit text (from Kh) in \S\S~\ref{persian_orthography}--\ref{transcription_transliteration_schemes}. A description of the typographic conventions in \S\S\ \ref{chapter_title_comparision_persian_sanskrit}, \ref{zijshahjahan_persian_english}, and \ref{siddhantasindhu_sanskrit_english}  of this study follows that in \S~\ref{typographic_conventions}. Therein, I explain with examples the various symbols and abbreviations used in the critical footnotes of these sections. Lastly, \S~\ref{glossary_appendix_format} describes the format of the Glossary (on page~\pageref{main}) and the Appendix (on page~\pageref{Appendix_verbs}).

%%%%%%%%%%%%%%%%%%%%%%%%%%%%%%%%%%%%%%%%%%%%%%%%%%%%%%%%%%%%%%%%%%%%%%%%%%%%%%%%
\subsection{Remarks on Persian orthography} \label{persian_orthography}
%%%%%%%%%%%%%%%%%%%%%%%%%%%%%%%%%%%%%%%%%%%%%%%%%%%%%%%%%%%%%%%%%%%%%%%%%%%%%%%%
The Persian text presented in this study follows the orthography of Classical Persian in which the manuscripts were written. It does, however, use modern punctuation marks for added clarity. 

Arabic loanwords are transcribed with their original spelling when they are attested as such, \eg the Arabic letter \tfarsi{ـي} (\textit{yā}\Ayn) is retained in the spelling of the word \tfarsi{ثاني} (\thani) following \SjA\ (f.\thinspace 21b:\thinspace21). In other instances, Persian spellings are used, \eg \tfarsi{درجه} (\daraji) instead of the Arabic \tfarsi{درجة} with the Arabic letter \tfarsi{ـة} (\textit{tā}\Ayn \textit{marbūṭah}). Persian words are presented in their unligated forms: \eg \tfarsi{آنرا} on f.21b:\thinspace 23~\SjA\ is transcribed as \tfarsi{آن را}. 

A few minor orthographic irregularities are seen in the Persian text of \SjA\ and \SjB. I note these below and emend them silently. However, when the reading of the text is affected by a grammatical ambiguity, I discuss my interpretation in corresponding footnotes. Scribal alterations, cancellations, copying errors, and marginalia are also noted in the footnotes.
\begin{enumerate}[topsep=0pt]
    \item Vocalisation marks are often omitted, in particular, on the syllable-initial \tfarsi{آ} (\textit{alif madda}). For example, \tfarsi{آن} (\textit{ān}) is simply written as \tfarsi{ان} on f.\thinspace 22a:\thinspace 2~\SjA\ and f.\thinspace 16b:\thinspace 1~\SjB.
    
    \item Arabic loan words are generally written without any diacritics. For example, the words \tfarsi{معدّل}  (\muaddil) (on f.\thinspace 21b:\thinspace 21~\SjA) and \tfarsi{اوّل} (\avval) (on f.\thinspace 14a:\thinspace 1~\SjB) are written without the over-letter diacritic \tfarsi{{\msf{◌}}\char"0651} (\textit{shadda}/\textit{tashdīd}).
    
    \SjA\ sometimes uses supplementary diacritics (\tashkil) to differentiate homographic words, \eg on f.\thinspace 22a:\thinspace8, the word \tfarsi{دور}, understood as `distant/remote', is explicitly written as \tfarsi{دُور}  (\textit{dūr}, \textsc{ipa} /duːɾ/) with the over-letter diacritic \tfarsi{{\msf{◌}}\char"064F} (\textit{ḍamma}/\textit{pish}) to differentiate it from \tfarsi{دَور} (\textit{dawr}, \textsc{ipa} /dawɾ/) `cycle/revolution'.   
    
    \item \SjA\ also has occasional diacritic points that are misplaced, for instance, the word \tfarsi{جهت} (\jahat) is spelt as \tfarsi{ـهـت}\.{\tfarsi{جـ}} on f.\thinspace 21b:\thinspace 23. The overdot over the word-initial letter \tfarsi {جـ} (\textit{jim}) is meaningless. 
    
    \item \SjB\ sometimes lacks diacritic points to indicate the phonetic distinction of consonants (\ijam). For example, the word-initial letter \tfarsi{بـ} (\textit{be}) in \tfarsi{بر} (\textit{bar}) appears without the underdot (on f.\thinspace 16a:\thinspace 30~\SjB). 
\end{enumerate}

%%%%%%%%%%%%%%%%%%%%%%%%%%%%%%%%%%%%%%%%%%%%%%%%%%%%%%%%%%%%%%%%%%%%%%%%%%%%%%%%
\subsection{Remarks on Sanskrit orthography}\label{sanskrit_orthography} 
%%%%%%%%%%%%%%%%%%%%%%%%%%%%%%%%%%%%%%%%%%%%%%%%%%%%%%%%%%%%%%%%%%%%%%%%%%%%%%%%
The Sanskrit text of Kh is fairly regular with occasional orthographic irregularities. Most of these are common scribal oversights seen in \Nagari\ palaeography, and hence, I emend these silently. They include 
\begin{itemize}[topsep=0pt]
    \item using the over-letter diacritic \tsans{\char"0902} (\textit{anusvāra}) for all conjoined nasal consonants;
    \item omitting the diacritical marks \tsans{\char"0903} (\textit{visarga}) for the terminal aspirate, \tsans{्} (\textit{virāma}) for inherent-vowel suppression, and \tsans{ऽ} (\textit{avagraha}) for prodelision of \textit{a/ā};
    \item using/omitting punctuation marks like the \tsans{||} (\textit{double-daṇḍa}) irregularly;
    \item using irregular (vernacular?) \Nagari\ letters (\eg {\tsv{\char"95F}} for \tsans{ya}) for Sanskrit, and retaining ill-formed vocalic signs (\eg \tsans{di\tsv{\char"940}});
    \item using doubled consonant irregularly, \eg  \tsans{ddha} in \tsans{arddha}  (after a vowel-suppressed \textit{r}-consonant) or across line (\textit{pāda}) breaks in a stanza;
    \item reversing conjunct-consonant pairs (\eg \tsans{dhba} for \tsans{bdha} or \tsans{nha} for \tsans{hna}); and
    \item confusing consonants like \tsans{ba} and \tsans{va}, \tsans{pa} and \tsans{ya}, \tsans{ma} and \tsans{sa}, \tsans{.sa} and \tsans{kha}, \etcp, and certain ligatures like \tsans{.s.ta} for \tsans{.s.tha}, \tsans{kta} for \tsans{tkta}, \etc
\end{itemize}

However, I discuss in footnotes the structure of those orthographic irregularities (mainly, morphosyntactic errors) that affect the reading of the text even as I emend the words accordingly. Other scribal errors like haplography (inadvertent omission) and dittography (inadvertent repetition) are also described in footnotes. 

%%%%%%%%%%%%%%%%%%%%%%%%%%%%%%%%%%%%%%%%%%%%%%%%%%%%%%%%%%%%%%%%%%%%%%%%%%%%%%%%
\subsection{Transcription and transliteration schemes}\label{transcription_transliteration_schemes}
%%%%%%%%%%%%%%%%%%%%%%%%%%%%%%%%%%%%%%%%%%%%%%%%%%%%%%%%%%%%%%%%%%%%%%%%%%%%%%%%
I adopt the following transcription/transliteration schemes to render Arabic, Persian, Sanskrit, and Hindi characters into the Roman (Latin) script.     
\begin{itemize}[topsep=0pt]
    \item Arabic and Persian texts are transcribed with the EI3 standard of phonetic transcription in Brill's \textit{Encyclopedia of Islam}, third edition \parencite[]{EIslam}.
    \item Sanskrit and Hindi texts are transliterated following the International Alphabet of Sanskrit Transliteration (IAST) scheme. For vernacular Hindi words, as well as \Devanagari-spellings of transliterated Persian words, I use the International Organisation for Standardisation (ISO) 15919 extension to transliterate certain characters, \eg {\tsnb{ड़ी}} is rendered as \textit{ṛī}, \tsans{khaa/} as \textit{khām̐}, \etc Commonly attested words of Indian origin (\eg Hindu, Brahmin, Mughal, Varanasi, \etcp) are presented without diacritics.
\end{itemize}

%%%%%%%%%%%%%%%%%%%%%%%%%%%%%%%%%%%%%%%%%%%%%%%%%%%%%%%%%%%%%%%%%%%%%%%%%%%%%%%%
\subsection{Typographic conventions} \label{typographic_conventions} 
%%%%%%%%%%%%%%%%%%%%%%%%%%%%%%%%%%%%%%%%%%%%%%%%%%%%%%%%%%%%%%%%%%%%%%%%%%%%%%%%
\subsubsection{Chapter-titles in \S~\ref{chapter_title_comparision_persian_sanskrit}}\label{chapter_titles_in_zij_sindhu}

\begin{enumerate}[topsep=0pt]
    \item \textbf{Layout}\quad The chapter-titles from Discourse II (\maqala\idafaconsonant\ \duvum) of the \ZijiShahJahani\ and Part II (\dvitiya\ \kanda) of the \Siddhantasindhu\ are placed in parallel columns in \S~\ref{chapter_title_comparision_persian_sanskrit}. The Persian title-text (to the left) and the Sanskrit title-text (to the right) can be identified by their corresponding chapter numbers in the left and right margins respectively, \eg `Discourse~II.2' (in the left margin) and `Book~II.2' (in the right margin) corresponding to the second chapter at the top of page~\pageref{chapter_number_example}. 
    I include the folia and line numbers of the 
    manuscripts at the end of the text in parentheses. 
    
    \item \textbf{Format of the translations}\quad My English translations of the chapter-titles are placed right below the corresponding Persian and Sanskrit text, parallel to each other. The technical terms in the translations are typeset in bold, and are accompanied by a Roman transcription/transliteration (in parentheses) of corresponding Persian and Sanskrit expressions. I indicate any additional words/expressions supplied for grammatical clarity by enclosing them in square brackets `[\,]' in my translations.
\end{enumerate}

\subsubsection{Chapter VI in \S\S~\ref{zijshahjahan_persian_english} \& \ref{siddhantasindhu_sanskrit_english}}\label{chapter_vi_zij_sindhu}

\begin{enumerate}[topsep=0pt]
     \item \textbf{Layout}\quad The Persian and Sanskrit text of the sixth chapter from Discourse II of the \ZijiShahJahani\ and from Part II of the \Siddhantasindhu\ are presented in \S~\ref{zijshahjahan_persian_english} and \S~\ref{siddhantasindhu_sanskrit_english} respectively. In both these sections, I place corresponding passages from the original text and their English translations on successive pages; \vid\ \S~\ref{structure_content_zij_sindhu_chapter_six} for my division of the text into comparable passages.
    \item \textbf{Ordering of the passages}\quad The passage-markers are enclosed in square brackets, \eg `[2]' or `[α]', and placed at the beginning of the passages. They appear in the right margin for the Persian text, and in the left margin for the Sanskrit text and English translations. 
    \item \textbf{Roman transcription/transliteration}\quad The technical expressions in the English translations are typeset in bold, and are accompanied by a Roman transcription/transliteration (in parentheses) of corresponding Persian and Sanskrit expressions. The Roman transcriptions of Persian compound verbs that indicate arithmetic operations are indicated in their infinitive form and are prefixed with an asterisk. For example, \gls{sum} (*\jam\ \kardan) `to sum' in passage~[\hyperlink{PEpass1}{1}] on page~\pageref{passage_1_english_persian}. \Vid\ footnote~\ref{compound_action_verb_persian} in the Appendix (on page~\pageref{Appendix_verbs}).
    \item \textbf{Verse numbering in \S~\ref{siddhantasindhu_sanskrit_english}}\quad The numbering of the metrical Sanskrit verses is \textit{different} from the ordering of the passages. I follow the Sanskrit text in placing the verse numbers at the end of the stanzas in my translations. For instance, the verse-number `\tsans{|| 1 ||}' at the end of a stanza in passage [\hyperlink{SpassA}{α}] (on page~\pageref{verse_1_label_sans_example}) also appears as `1' at the end of my English translation (on page~\pageref{verse_1_label_eng_example}).  
    \item \textbf{Poetic meters in \S~\ref{siddhantasindhu_sanskrit_english}}\quad The Sanskrit names of the poetic meters are indicated (in Roman transliteration) in the right margin alongside the Sanskrit verses, \eg verse number `\tsans{|| 2 ||}' (in passage~[\hyperlink{Spass2}{2}]) is in the \textit{pramāṇikā}-meter indicated in the right margin margin (in-line with the verse) on page~\pageref{sanskrit_meter_typography_example}.
    \item \textbf{Folio breaks}
    \begin{itemize}[topsep=0pt]
        \item For the Persian text in \S~\ref{zijshahjahan_persian_english}, I indicate a folio break with `$\rceil$' (in-line with the text) and state the manuscript reference in the left margin. For instance, passage [\hyperlink{Ppass2}{2}] on page~\pageref{folio_break_persian_example} has `\tfarsi{\dots جيب~\te{$\rceil$}بعد\dots}' on its second line and corresponding to it, `f.\thinspace 22a:\thinspace 1~\SjA~$\rceil$' is written in the left margin. This indicates the words `\tfarsi{بعد\dots}' begin on line 1 of f.\thinspace 22a in \SjA.
        \item For the Sanskrit text in \S~\ref{siddhantasindhu_sanskrit_english}, folio breaks are shown with `$\lceil$' (in-line with the text) and manuscript references are in the right margin. For instance, the first line of passage [\hyperlink{SpassB}{β}] on page~\pageref{folio_break_sanskrit_example} has `\dots\tsans{vivara\raisebox{.5ex}{$\lceil$}ga.m dhanu}\dots', and correspondingly, `$\lceil$~ f.\thinspace 20v:\thinspace1~Kh' is in the right margin. This implies the words `\tsans{ga.m dhanu}\dots' begin on line 1 of f.20v in Kh. 
    \end{itemize}
\end{enumerate}

\subsubsection{Critical footnotes in \S\S~\ref{chapter_title_comparision_persian_sanskrit}, \ref{zijshahjahan_persian_english}, \& \ref{siddhantasindhu_sanskrit_english}}\label{critical_footnotes_in_sec_2_3_4}

\begin{enumerate}[topsep=0pt]
    \item Footnotes in \S~\ref{chapter_title_comparision_persian_sanskrit} are numbered 1, 2, 3 \etcp, whereas, those in \S\S~\ref{zijshahjahan_persian_english} \& \ref{siddhantasindhu_sanskrit_english} are numbered [i], [ii], [iii] \etc The numbers are reset at the beginning of each section. 
    \item I use repeated footnote marks for longer footnotes in the Persian text (in \S~\ref{zijshahjahan_persian_english}), \eg footnote~`\ref{wrap_footnote_example}' on page~\pageref{wrap_footnote_example} where `[vi]\==' and `\==[vi]' enclose the commented text in passage [\hyperlink{Ppass9}{9}].   
    \item An edited reading is separated from a variant (attested) reading by a right-square bracket {$\Big]$} like, for instance, \tsans{naama madhyaahnarekheti} {$\Big ]$} \tsans{naama dhyaahnarekheti}~Kh (footnote~\ref{emeneded_attested_sanskrit_example} on page~\pageref{emeneded_attested_sanskrit_example}) or \tfarsi{یا در}~{$\Big ]$}~\tfarsi{یاد در}~\SjB\ (footnote~\ref{emeneded_attested_persian_example} on page~\pageref{emeneded_attested_persian_example}).
    \item I use abbreviated forms of grammatical terms in the critical footnotes. For example, the \acrfull{modifier} \tsans{prathama} in the nominal compound \tsans{prathamaadhyaaye} is indicated as \tsans{prathama}\,$_\text{\acrshort{modifier}}$ in footnote~\ref{grammatical_abbreviation_example} on page~\pageref{grammatical_abbreviation_example}. An expanded list of all grammatical abbreviations used in this study can be found on page~\pageref{acronym}.  
    \item The truncation (of letters) in long Sanskrit words is indicated by the \Nagari\ abbreviation symbol `\tsans{\selip}' (\textit{lāghava-cihna}), comparable in its use to an ellipsis. For example, \tsans{ga.nita\selip} for \tsans{ga.nitasau\-karyaartha.m}  or \tsans{\selip maa.m"saakhya} for \tsans{sphu.taapamaa.m"saakhya}. 
    \item I emphasise the letters in Persian and Sanskrit words by underlying them, \eg \tfarsi{\underline{قسمت}} or \tsans{unna\underline{ta}}. The emphasis is used in the critical footnotes to signify an orthographic feature or a scribal error that is overt.
\end{enumerate}

%%%%%%%%%%%%%%%%%%%%%%%%%%%%%%%%%%%%%%%%%%%%%%%%%%%%%%%%%%%%%%%%%%%%%%%%%%%%%%%%
\subsection{Format of the appendix and the glossary} \label{glossary_appendix_format} 
%%%%%%%%%%%%%%%%%%%%%%%%%%%%%%%%%%%%%%%%%%%%%%%%%%%%%%%%%%%%%%%%%%%%%%%%%%%%%%%%
All Persian and Sanskrit words in the appendix and the glossary are written with Perso-Arabic and \Nagari\ letters respectively, and are accompanied by corresponding Roman transcriptions/transliterations enclosed in parentheses. The grammatical terms are abbreviated in the appendix and the glossary; their expanded forms are listed on page~\pageref{acronym}. 

\subsubsection*{Appendix}\label{format_appendix}
The appendix includes a list of the Persian and Sanskrit verbs seen in the \ZijiShahJahani\ Discourse~II.6 and the \Siddhantasindhu\ Part~II.6 respectively. 
\begin{itemize}[topsep=0pt]
\item The Persian verbs are listed on page~\pageref{persian_verbs} in their infinitive form along with their corresponding present stems. I group together the attested (inflected) forms of these verbs and provide \textit{passage-markers} (in square brackets at the end) to locate them in the text in \S~\ref{zijshahjahan_persian_english}. For example,
{\par\centering\tfarsi{باشد} (\textit{bāshad})  \acrshort{present}-\acrshort{subjunctive}-\acrshort{singular}·\acrshort{third} `[he/she/it] should be' [\hyperlink{Ppass1}{1}, \hyperlink{Ppass3}{3}--\hyperlink{Ppass8}{8}, \hyperlink{Ppass10}{10}, \hyperlink{Ppass11}{11}].\par}
\item The Sanskrit verbs are indicated in their root-form beginning on page~\pageref{sanskrit_verbs}. The verbal roots (of different verb-class numbers) are grouped together based on their common meaning. The attested (inflected forms) are listed under their respective verbal root, and are accompanied by \textit{passage-markers} (in square brackets at the end) to locate them in the text in \S~\ref{siddhantasindhu_sanskrit_english}. For example,
{\par\centering
\tsans{syaat} (\textit{syāt}) \acrshort{optative}-\acrshort{active}-\acrshort{singular}·\acrshort{third} `[he/she/it] should be/exist' [\hyperlink{Spass3}{3}--\hyperlink{Spass7}{7}].
\par}
\end{itemize}

\subsubsection*{Glossary} \label{format_glossary}
Persian and Sanskrit technical expressions are listed in the glossary (beginning on page~\pageref{acronym}). They are derived from  \S\S~\ref{chapter_title_comparision_persian_sanskrit}, \ref{zijshahjahan_persian_english}, and \ref{siddhantasindhu_sanskrit_english} where they appear in bold in corresponding English translations. 
\begin{itemize}[topsep=0pt]
    \item Equivalent Persian and Sanskrit terms are grouped together under their common technical translation in English, separated from each other by a semicolon. Synonyms are separated by commas. For example, 
    {\par\centering\textbf{latitude}\quad \tfarsi{عرض}  (\ard) [\hyperlink{PEpass5}{5}, \hyperlink{PEpass6}{6}, \hyperlink{PEpass10}{10}]; \tsans{"sara} (\textit{śara}) [\hyperlink{SEpass6}{6}], \tsans{baa.na} (\textit{bāṇa}) [\hyperlink{SEpass10}{10}].\par}
    \item At the end of each entry, I provide the \textit{chapter-numbers} and/or the \textit{passage-markers} in square brackets to identify its location in the text. The identifiers refer to apposite passages or chapter-titles corresponding to the language of the entry. For instance, in the example above, \tfarsi{عرض}  (\ard) appears in passage~[\hyperlink{PEpass6}{6}] of \S~\ref{zijshahjahan_persian_english} (page~\pageref{passge_6_ard_glossary_format_example}); whereas, \tsans{"sara} (\textit{śara}) can be found in passage~[\hyperlink{SEpass6}{6}] of \S~\ref{siddhantasindhu_sanskrit_english} (page~\pageref{passge_6_sara_glossary_format_example}).     \item References to multiple chapter-numbers are separated by commas (without repeating `II'). For example,
    {\par\centering
    \textbf{ascendant}\quad \tfarsi{طالع} (\tali), pl.\thinspace \tfarsi{طوالع} (\tawali) \hyperlink{Pii11}{II.11}, \hyperlink{Pii21}{21}, \hyperlink{Pii22}{22}; \tsans{lagna} (\textit{lagna}) \hyperlink{Sii21}{II.21}\par}
    indicates that \tfarsi{طوالع} (\tawali) appears in the chapter-titles of the Persian chapters \hyperlink{Pii11}{II.11}, \hyperlink{Pii21}{II.21}, and \hyperlink{Pii22}{II.22} in \S~\ref{chapter_title_comparision_persian_sanskrit}. 
    \item Successive chapter-numbers or passage-markers are sometimes shown as a range to space space, \eg\enspace \textbf{definition}\quad \tsans{lak.sa.na} (\textit{lakṣaṇa}) \hyperlink{Sii2}{II.2}, \hyperlink{Sii8}{8}--\hyperlink{Sii13}{13}, \hyperlink{Sii17}{17}--\hyperlink{Sii19}{19}, \hyperlink{Sii22}{22}. 
    \item Mutually related technical translations in English are grouped together based on their linguistic or mathematical similarity. For instance, the heading \gls{distance_ce_true_declination} (on page~\pageref{glsentry-distance_ce_true_declination}) includes the expressions $\hookrightarrow$\,\gls{distance}, $\hookrightarrow$\,\gls{distance_celestial_object}, and 
    $\hookrightarrow$\,\gls{true_declination}. 
    \item The glossary entries are arranged following the English alphabetical order.
\end{itemize}



\clearpage

\section[Chapter-titles from the \ZijiShahJahani, Discourse~II and the\linebreak \Siddhantasindhu, Part~II: text and translation]{Chapter-titles from Discourse II (\maqala\idafaconsonant\ \duvum) of the \ZijiShahJahani\ and Part II (\dvitiya\ \kanda) of the \Siddhantasindhu}\label{chapter_title_comparision_persian_sanskrit}

% I include corresponding chapter titles (in Persian) from the second discourse (\maqala\idafaconsonant\ \duvum)\footnote{\,In this study, I follow the EI3 standard of phonetic transcription in Brill's \textit{Encyclopedia of Islam}, third edition, to render Arabic and Persian text into Latin characters, \vid\ \textcite{EIslam}.} of \ZijUlughBeg\ for comparison. These are taken from the Persian edition of the \ZijUlughBeg\ by Louis P.\thinspace E. Amélie Sédillot, \vid\ \textcite[Persian edition on pp.\thinspace 341--386\thinspace(\tfarsi{۳۸۶--۳۴۱})]{Sedillotedition}. The orthography of the Persian text in Sédillot's 1847 edition is modified to match modern spelling conventions. I provide English translations of the Persian titular text for readers unfamiliar with Persian with technical terms.


% Folia 17r--28v of MS~Kh contain Part II (\dvitiya\ \kanda) of \Nityananda's \Siddhantasindhu. The Sanskrit chapter-titles of the twenty-two chapters (\adhyaya) in Part II are presented below (in \Nagari) along with their English translations. At the end of every titular text, the folio and line numbers corresponding to its location in MS~Kh is included in parenthesis. Suspected scribal errors are corrected with explanations in footnotes. Technical terms and phrases are emboldened in the English translation for comparison with similar constructions in the Persian text from \ZijiShahJahani. These are accompanied by corresponding Roman transliterations of Sanskrit expressions in parenthesis.

\setcounter{footnote}{0}
\renewcommand{\thefootnote}{\arabic{footnote}} % for restarting the footnote numbering 
\setlength{\footnotesep}{\baselineskip} % for space between footnotes

\begin{multicols}{2}
\noindent\reversemarginpar\marginnote{\hypertarget{PiiInc}{Discourse~II\\\textit{incipit}}}
\tfarsi{مقاله دوم \hfill}\\
\tfarsi{در معرفت اوقات و طالع هر وقت و آنچه تعلق بدان دارد مشتمل بر بیست و دو باب.%
\te{(f.\thinspace 14a:\thinspace 1~\SjB)}\hfill}
\columnbreak

\noindent\normalmarginpar\marginnote{\hypertarget{SiiInc}{Part~II\\\textit{incipit}}}%
\tsans{nityaanandasvaruupaaya saccidvayamurtaye~||\\ advitiiyaaya vibhave.anantaaya brahma.ne nama.h~||}\\[5pt]
\tsnb{अथ द्वितीयकाण्डे द्वाविंशत्यध्यायैरभिमतसमयस्तात्कालिकलग्नं च तदुपयोगीन्यपि \\ ज्ञायन्ते~॥}%
\enskip (f.\thinspace 17r:\thinspace1--3~Kh)
\end{multicols}%
%
\begin{multicols}{2}
Second Discourse\\
On the \gls{knowledge} (\marifat) of [finding] the \glspl{time} (\avqat) and the \gls{ascendant_each_time} (\tali\ \har\ \vaqt), and whatever belongs to it [\ie all things related to this topic], including twenty-two chapters.
\columnbreak

Obeisance to Brahman (\textit{brahman}) who is the embodiment of eternal bliss (\textit{nitya-ānanda-svarūpa}), who is the form of both existence and thought (\textit{sat-cit-dvaya-mūrti}), who is [the One] without a second (\textit{advitīya}), who is omnipresent (\textit{vibhū}), [and] who is infinite (\textit{ananta}).\\[5pt] 
Now, in the second part, [finding] the \gls{desired_time} (\textit{abhimata-samaya}) and the \gls{ascendant_that_time} (\textit{tātkālika-lagna}), as well as [things] using that, are understood with twenty-two chapters. 
\end{multicols}

\newpage %%%%________________________NEWPAGE

%%%%%%%%%%%%%%%%%%%%%%%%%%%%%%%%%%%
\begin{multicols}{2}
\noindent\reversemarginpar\marginnote{\hypertarget{Pii1}{Discourse~II.1}}
\tfarsi{باب اوّل\\
در بیان معرفت جنس هر یک از حاصل ضرب و خارج قسمت و جذر، یعنی دانستن آنکه حاصل ضرب یا خارج قسمت یا جذر از کدام مرتبه است از مراتب مرفوعات و درج و اجزاء درج مثل دقایق و ثواني و غیر آن.% 
\te{(f.\thinspace 14a:\thinspace 3--1~\SjB)}\hfill}
\columnbreak

\noindent\normalmarginpar\marginnote{\hypertarget{Sii1}{Part~II.1}}%
\tsnb{तत्र प्रथमाध्याये\te{%
\footnote{~\tsans{prathamaadhyaaye} {$\Big ]$} \tsans{prathamadhyaaye}\enskip Kh. A regular sandhi of the words \tsans{prathama}\,$_\text{\acrshort{modifier}}$ + \tsans{adhyaaye}\,$_\text{\acrshort{locative}-\acrshort{singular}}$ generates the compound \tsans{pratha\underline{maa}dhyaaye}. The omission of the long vowel (\textit{ā}-diacritic in \tsans{maa} looks like a scribal oversight.}} %
गुणनभजनफले मूलं च परिवर्तादिस्थानैः कलादिस्थानैर्वा किं\linebreak जातीयं स्यादिति ज्ञायते~॥
तत्र गणितसौकर्यार्थं\te{%
\footnote{~\tsans{ga.nita\selip}{$\Big ]$} \tsans{ga.niita\selip}\enskip Kh. There are no morphological divisions of \tsans{ga.niita} that are syntactically admissible in the larger compound \tsans{ga.niitasaukaryaartha.m}, whereas the word \tsans{ga\underline{.ni}ta} is both grammatically well-formed and contextually apposite. I suspect the use of the long vowel (the \textit{ī}-diacritic in \tsans{.nii}) is either a scribal hypercorrection or an inadvertent misspelling.\label{grammatical_abbreviation_example}}} %
यवनप्रसिद्धप्रकारेणाङ्कस्थानानां संस्कृतशब्दैः संज्ञा कल्प्यते~॥} \\
(f.\thinspace 17r:\thinspace3--5~Kh)
\end{multicols}%
%
\begin{multicols}{2}
First Chapter\\
On the expression of the \gls{knowledge} (\marifat) of each \gls{genus} (\jins) [of digits] from the \gls{result_multiplication} (\hasil\idafaconsonant\ \darb), and the \gls{quotient_division} (\kharij\idafaconsonant\ \qismat), and the \gls{square_root} (\jadr). In other words, to know what is the \gls{position} (\martaba) of [the digits in] the \gls{result_multiplication} (\mbox{\hasil\idafaconsonant} \darb), or the \gls{quotient_division} (\mbox{\kharij\idafaconsonant} \qismat), or the \gls{square_root} (\jadr), from the \glspl{position} (\maratib) of \glspl{elevation_sexagesimal} (\marfuat) [\scl integer number of revolutions], and the \gls{degree} (\daraj), and the \glslink{parts_of_degree}{fractional parts of a degree} (\ajza\idafaconsonant\ \daraj) like \glspl{minute} (\daqaiq) and \glspl{second} (\thawani) and so on.
\columnbreak

There, in the first chapter, with [digits in] the \glspl{position} (\textit{sthāna}) of \gls{revolution} (\textit{parivarta}) \etc or with [digits in] the \glspl{position} (\textit{sthāna}) of \gls{minute} (\textit{kalā}) \etc in the \gls{result_multiplication_division} (\textit{guṇana-bhajana-phala}), and the \gls{square_root} (\textit{mūla}), what [digits] should \gls{belonging_genus} (\textit{jātīya}): this is understood. Therein, for the purpose of facilitating ease in computations (\textit{gaṇita-saukarya-artha}) with the method famous amongst the foreigner (\textit{yavana-prasiddha}), the name of the \glspl{position} (\textit{sthāna}) of \glspl{digit} (\textit{aṅka}) is considered with Sanskrit words.
\end{multicols}
\medskip

\newpage %%%%________________________NEWPAGE

%%%%%%%%%%%%%%%%%%%%%%%%%%%%%%%%%%%
\begin{multicols}{2}
\label{chapter_number_example}
\noindent\reversemarginpar\marginnote{\hypertarget{Pii2}{Discourse~II.2}}%
\tfarsi{باب دوم\\
در عمل تعدیل ما بين السطرين. \\\te{(f.\thinspace 14b\thinspace 6~\SjB)}\hfill}
\columnbreak

\noindent\normalmarginpar\marginnote{\hypertarget{Sii2}{Part~II.2}}%
\tsnb{अथ द्वितीयाध्याये द्विकोष्ठान्तरोत्थफल\-साधनम् ॥} %
\tsans{tasya muula.m trairaa"sikam || atra yavanaa.h parasparasambandhicatuuraa"siinga.nayanti~|| tallak.sa.na.m ca ||}\enskip (f.\thinspace 18r:\thinspace2--4~Kh)
\end{multicols}%
%
\begin{multicols}{2}
Second chapter\\
On the \gls{method_of_interpolation} (\amal\idafaconsonant\ \tadil) 
\gls{between_two_lines} (\ma\ \bayn\ \alsatrayn) [of a table].
\columnbreak

Now, in the second chapter, the \gls{demonstration} (\textit{sādhana}) of the \gls{result} (\textit{phala}) derived from the \gls{difference_between_two_cells} (\textit{dvi-koṣṭha-antara}). The basis of this is the \gls{rule_of_three} (\textit{trai-rāśika}). Here, the foreigner (\textit{yavana}) take into account four \gls{correlated_numbers} (\textit{paraspara-sambandhi-rāśi}). And their \glspl{definition} (\textit{lakṣaṇa}) [are first stated]. 
\end{multicols}
\medskip

%%%%%%%%%%%%%%%%%%%%%%%%%%%%%%%%%%%

\begin{multicols}{2}
\noindent\reversemarginpar\marginnote{\hypertarget{Pii3}{Discourse~II.3}}%
\tfarsi{باب سيوم\\
در معرفت جيب و سهم.%
\te{(f.\thinspace 15a:\thinspace 8~\SjB)}\hfill}
\columnbreak

\noindent\normalmarginpar\marginnote{\hypertarget{Sii3}{Part~II.3}}%
\tsnb{अथ तृतीयाध्याये ज्याशरज्ञानम्~॥}\\
(f.\thinspace 18r:\thinspace23--24~Kh)
\end{multicols}%
%
\begin{multicols}{2}
Third chapter\\
On the \gls{knowledge} (\marifat) of the \gls{Sine} (\jayb) and the \gls{Sagitta} (\sahm) [\ie the Versed Sine].
\columnbreak

Now, in the third chapter, the \gls{knowledge} (\textit{jñāna}) of the \gls{Sine} (\textit{jyā}) and the \gls{Versed_Sine} (\textit{śara}).
\end{multicols}
\medskip

%%%%%%%%%%%%%%%%%%%%%%%%%%%%%%%%%%%

\begin{multicols}{2}
\noindent\reversemarginpar\marginnote{\hypertarget{Pii4}{Discourse~II.4}}%
\tfarsi{باب چهارم\\
در معرفت ظلّ.%
\te{(f.\thinspace 15b:\thinspace 1~\SjB)}\hfill}
\columnbreak

\noindent\normalmarginpar\marginnote{\hypertarget{Sii4}{Part~II.4}}%
\tsnb{अथ चतुर्थाध्याये छायाज्ञानम्~॥}\\ 
(f.\thinspace 19r:\thinspace13~Kh)
\end{multicols}%
%
\begin{multicols}{2}
Fourth chapter\\
On the \gls{knowledge} (\marifat) of the \glslink{shadow}{shadow} (\zill) [of a gnomon].
\columnbreak

Now, in the fourth chapter, the \gls{knowledge} (\textit{jñāna}) of the \glslink{shadow}{shadow} (\textit{chāyā}) [of a gnomon].
\end{multicols}

\newpage %%%%________________________NEWPAGE

%%%%%%%%%%%%%%%%%%%%%%%%%%%%%%%%%%%

\begin{multicols}{2}
\noindent\reversemarginpar\marginnote{\hypertarget{Pii5}{Discourse~II.5}}%
\tfarsi{باب پنجم\\
در معرفت ميل اجزاء فلك البروج از معدّل النهار.%
\te{(f.\thinspace 16a:\thinspace 4~\SjB)}\hfill}
\columnbreak

\noindent\normalmarginpar\marginnote{\hypertarget{Sii5}{Part~II.5}}%
\tsnb{अथ पञ्चमाध्याये क्रान्तिज्ञानम् ॥ \\तत्र तावत्क्रान्तिसूत्रादिसंज्ञोच्यते~॥}\\
(f.\thinspace 19v:\thinspace18--19~Kh)
\end{multicols}%
%
\begin{multicols}{2}
Fifth chapter\\
On the \gls{knowledge} (\marifat) of the \gls{declination_parts_ecliptic} (\mayl\idafaconsonant\ \ajza\idafaconsonant\ \falak\ \alburuj) from the \gls{celestial_equator} (\muaddil\ \alnahar).
\columnbreak

Now, in the fifth chapter, the \gls{knowledge} (\textit{jñāna}) of the \gls{declination} (\textit{krānti}). There, firstly, a [technical] term beginning with \gls{circle_of_declination} (\textit{krānti-sūtra}) is stated.
\end{multicols}
\medskip
%%%%%%%%%%%%%%%%%%%%%%%%%%%%%%%%%%%

\begin{multicols}{2}\label{discourse_part_6_chapter_title_example}
\noindent\reversemarginpar\marginnote{\hypertarget{Pii6}{Discourse~II.6}}%
\tfarsi{باب ششم\\
در معرفت بعد كواكب از معدّل النهار.\\ \te{(f.\thinspace 16a:\thinspace 26--25~\SjB)}\hfill}
\columnbreak

\noindent\normalmarginpar\marginnote{\hypertarget{Sii6}{Part~II.6}}%
\tsnb{अथ षष्ठाध्यायस्पष्टक्रान्तिज्ञानम्~॥}\\
(f.\thinspace 20r:\thinspace16~Kh)
\end{multicols}%
%
\begin{multicols}{2}
On the \gls{knowledge} (\marifat) of the \gls{distance_celestial_object} (\bud\idafaconsonant\ \kawakib\ \az\ \muaddil\ \alnahar).
\columnbreak

Sixth chapter\\
Now, the \gls{knowledge} (\textit{jñāna}) of the \gls{true_declination} (\textit{spaṣṭa-krānti}) in the sixth chapter.
\end{multicols}
\medskip
%%%%%%%%%%%%%%%%%%%%%%%%%%%%%%%%%%%

\begin{multicols}{2}
\noindent\reversemarginpar\marginnote{\hypertarget{Pii7}{Discourse~II.7}}%
\tfarsi{باب هفتم\\
در معرفت غایت ارتفاع و انخفاض كواكب.%
\te{(f.\thinspace 16b:\thinspace 11--10~\SjB)}\hfill}
\columnbreak

\noindent\normalmarginpar\marginnote{\hypertarget{Sii7}{Part~II.7}}%
\tsnb{अथ सप्तमाध्याये ग्रहस्य परमोन्नतांशाना\-मधःस्थपरमभागानां च ज्ञानम्~॥}\\
(f.\thinspace 20v:\thinspace12--13~Kh)
\end{multicols}%
%
\begin{multicols}{2}
Seventh chapter\\
On the \gls{knowledge} (\marifat) of the \gls{maximum_elevation_depression_celestial_object} (\ghayat\idafaconsonant\ \irtifa\ \va\ \inkhifad\idafaconsonant\ \kawakib).
\columnbreak

Now, in the seventh chapter, the \gls{knowledge} (\textit{jñāna}) of the \gls{degrees_maximum_elevation} (\textit{parama-unnata-aṃśa}) and the \gls{degrees_maximum_depression} (\textit{adhaḥstha-parama-bhāga}) of a \gls{celestial_object} (\textit{graha}).
\end{multicols}

\newpage %%%%________________________NEWPAGE

%%%%%%%%%%%%%%%%%%%%%%%%%%%%%%%%%%%

\begin{multicols}{2}
\noindent\reversemarginpar\marginnote{\hypertarget{Pii8}{Discourse~II.8}}%
\tfarsi{باب هشتم\\
در معرفت مطالع خطّ استوا و آنرا مطالع فلك مستقیم نیز گویند.%
\te{(f.\thinspace 16b:\thinspace 17~\SjB)}\hfill}
\columnbreak

\noindent\normalmarginpar\marginnote{\hypertarget{Sii8}{Part~II.8}}%
\tsnb{अथ अष्टमाध्याये व्यक्षोदयांशज्ञानम् ॥\\ तेषां लङ्कोदयांशसंज्ञाप्युच्यते~॥}~%
\tsans{tallak.sa.namaaha~||}\enskip(f.\thinspace 20v:\thinspace24--25~Kh)
\end{multicols}%
%
\begin{multicols}{2}
Eighth chapter\\
On the \gls{knowledge} (\marifat) of \gls{ascension_line_terrestrial_observer} (\matali\ \khatt\idafaconsonant\ \istiva) [\ie the right ascensions of the zodiacal signs]. And that is also called the \gls{ascension_right_sphere} (\matali\ \falak\idafaconsonant\ \mustaqim).
\columnbreak

Now, in the eighth chapter, the \gls{knowledge} (\textit{jñāna}) of the \gls{rising_terrestrial_equator_degrees} (\textit{vyakṣa-udaya-aṃśa}) [\ie the right ascensions of the degrees of the ecliptic]. All of them are also called the \gls{rising_Lanka_degrees} (\textit{laṅkā-udaya-aṃśa}) by name. Their \glspl{definition} (\textit{lakṣaṇa}) state [as follows].
\end{multicols}
\medskip
%%%%%%%%%%%%%%%%%%%%%%%%%%%%%%%%%%%

\begin{multicols}{2}
\noindent\reversemarginpar\marginnote{\hypertarget{Pii9}{Discourse~II.9}}%
\tfarsi{باب نهم\\
در معرفت تعديل النهار و قوس النهار و قوس الليل و ساعات النهار و ساعات الليل.%
\te{(f.\thinspace 16b:\thinspace 25~\SjB)}\hfill}
\columnbreak

\noindent\normalmarginpar\marginnote{\hypertarget{Sii9}{Part~II.9}}%
\tsnb{अथ नवमाध्याये चरदिनरात्रिवामानां दिनरात्रिहोरादीनां च ज्ञानम्~॥} %
\tsans{tatra taavatte.saa.m lak.sa.nam~||}%
\enskip (f.\thinspace 21r:\thinspace5--6~Kh)
\end{multicols}%
%
\begin{multicols}{2}
Ninth chapter\\
On the \gls{knowledge} (\marifat) of the \gls{equation_of_daylight} (\tadil\ \alnahar); and the \gls{arc_of_daylight} (\qaws\ \alnahar) and the \gls{arc_of_night} (\qaws\ \allayl); and the 
\gls{hours_of_daylight} (\saat\ \alnahar) and \gls{hours_of_night} (\saat\ \allayl).
\columnbreak

Now, in the ninth chapter, the \gls{knowledge} (\textit{jñāna}) of the \gls{ascensional_difference} (\textit{cara}) of the \gls{oblique_diurnal_circle} (\textit{dina-rātri-vāma}[-\textit{vrtta}]) and of the \gls{hours_of_day_and_night} (\textit{dina-rātri-horā}) \etcp\ There, firstly, the \glspl{definition} (\textit{lakṣaṇa}) of those [are stated].
\end{multicols}
\medskip
%%%%%%%%%%%%%%%%%%%%%%%%%%%%%%%%%%%

\begin{multicols}{2}
\noindent\reversemarginpar\marginnote{\hypertarget{Pii10}{Discourse~II.10}}%
\tfarsi{باب دهم\\
در معرفت مطالع بلد.\te{(f.\thinspace 17a\thinspace 26~\SjB)}\hfill}
\columnbreak

\noindent\normalmarginpar\marginnote{\hypertarget{Sii10}{Part~II.10}}%
\tsnb{अथ दशमाध्याये निजोदयांशज्ञानम्~॥}\\
\tsans{tallak.sa.na.m ca~||}\enskip (f.\thinspace 21v:\thinspace23~Kh)
\end{multicols}%
%
\begin{multicols}{2}
Tenth chapter\\
On the \gls{knowledge} (\marifat) of the \gls{ascension_ecliptic_locality} (\matali\idafaconsonant\ \balad) [\ie the oblique ascensions of the zodiacal signs].
\columnbreak

Now, in the tenth chapter, the \gls{knowledge} (\textit{jñāna}) of the \gls{rising_location_degree} (\textit{nija-udaya-aṃśa}) [\ie the oblique ascensions of the degrees of the ecliptic]. And their \glspl{definition} (\textit{lakṣaṇa}) [are first stated]. 
\end{multicols}

\newpage %%%%________________________NEWPAGE

%%%%%%%%%%%%%%%%%%%%%%%%%%%%%%%%%%%

\begin{multicols}{2}
\noindent\reversemarginpar\marginnote{\hypertarget{Pii11}{Discourse~II.11}}%
\tfarsi{باب یازدهم \\
در عمل عکس مطالع يعنی معرفت طوالع از مطالع\te{\footnote{~The words \tfarsi{از مطالع} are a marginal addition. They appears in the exterior (left) margin of f.\thinspace 17a~\SjB\ alongside line 5 of the text. The main text has an interlinear insertion mark `$\vee$' at the end of the preceding \tfarsi{طوالع}. The marginal text ends with a terminal \textit{number-like} mark \tfarsi{۴} which, according to \textcite[117]{Gacek}, is an abbreviation for \tfarsi{تمام شد} \textit{tamām shud} `ended/finished' often seen in manuscripts of Indian/Iranian origins.}} بعمل.%
\te{(f.\thinspace 17b:\thinspace 5~\SjB)}\hfill}
\columnbreak

\noindent\normalmarginpar\marginnote{\hypertarget{Sii11}{Part~II.11}}%
\tsnb{अथैकादशाध्याये स्वोदयांशेभ्यो विनैव कोष्ठ\-कैर्विलोमक्रियातो} \mbox{\tsnb{विलग्नांशकज्ञानम्\thinspace॥}}
\tsans{vilomakriyaalak.sa.nam~||}\enskip (f.\thinspace 22r:\thinspace9--11~Kh)
\end{multicols}%
%
\begin{multicols}{2}
Eleventh chapter\\
On the \gls{inverse_method} (\amal\idafaconsonant \aks) [of] \glsuseri{ascension_measure} (\matali); in other words, the knowledge of the [ecliptic degrees of the] \glspl{ascendant} (\tawali) from the [local] \glsuseri{ascension_measure} (\matali) [\ie  from the oblique ascensions of the ascendants] by direct calculation. 
\columnbreak

Now, in the eleventh chapter, the \gls{knowledge} (\textit{jñāna}) of the \glslink{ascendant_ecliptic_degrees}{[ecliptic] degrees of the ascendants} (\textit{vilagna-aṃśaka}) from the \gls{rising_location_degree} (\textit{sva-udaya-aṃśa}) [\ie from the oblique ascensions of the ascendants] without [using] the \glspl{table} (\textit{koṣṭhaka}) [and] by using the \gls{inverse_procedure} (\textit{viloma-kriyā}). The \gls{definition} (\textit{lakṣaṇa}) of the \gls{inverse_procedure} (\textit{viloma-kriyā}) [is first stated].
\end{multicols}
\medskip
%%%%%%%%%%%%%%%%%%%%%%%%%%%%%%%%%%%

\begin{multicols}{2}
\noindent\reversemarginpar\marginnote{\hypertarget{Pii12}{Discourse~II.12}}%
\tfarsi{باب دوازدهم \\
در معرفت مطالع  ممرّ و درجهٔ ممرّ كوكب.%
\te{(f.\thinspace 17b:\thinspace30--29~\SjB)}\hfill}
\columnbreak

\noindent\normalmarginpar\marginnote{\hypertarget{Sii12}{Part~II.12}}%
\tsnb{अथ द्वादशाध्याये नक्षत्रस्य लङ्कायामुदये जाते सति भोदयलग्नव्यक्षोदयांशभोदय-\linebreak लग्नांशयोर्ज्ञानम्~॥} ~\tsans{tallak.sa.nam~||}\\
(f.\thinspace 22v:\thinspace15--17~Kh)
\end{multicols}%
%
\begin{multicols}{2}
Twelfth chapter\\
On the \gls{knowledge} (\marifat) of the \gls{ascension_transit} (\matali\idafaconsonant\ \mamarr) [\ie the right ascension of the zodiacal sign culminating with a celestial object] and the \gls{degree_transit_celestial_object} (\daraji\idafavowel\ \mamarr\idafaconsonant\ \kawkab) [\ie the ecliptic longitude of the zodiacal sign culminating with a celestial object].
\columnbreak

Now, in the twelfth chapter, when a \gls{celestial_object} (\textit{nakṣatra}) \glslink{rising}{rises} (\textit{udaya}) at \Lanka\ (the terrestrial equator), the \gls{knowledge} (\textit{jñāna}) of the \gls{degree_equatorial_ascension} (\textit{bha-udaya-lagna-vyakṣa-udaya-aṃśa}) [\ie the right ascension of the zodiacal sign culminating with the celestial object] and the \gls{degree_meridian_ecliptic_rising_time} (\textit{bha-udaya-lagna-aṃśa}) [\ie the ecliptic longitude of the zodiacal sign culminating with the celestial object]. Their \glspl{definition} (\textit{lakṣaṇa}) [are first stated]. 
\end{multicols}

\newpage %%%%________________________NEWPAGE

%%%%%%%%%%%%%%%%%%%%%%%%%%%%%%%%%%%


\begin{multicols}{2}
\noindent\reversemarginpar\marginnote{\hypertarget{Pii13}{Discourse~II.13}}%
\tfarsi{باب سیزدهم \\
در مطالع طالع و غروب كواكب.%
\te{(f.\thinspace 18a:\thinspace12--11~\SjB)}\hfill}
\columnbreak

\noindent\normalmarginpar\marginnote{\hypertarget{Sii13}{Part~II.13}}%
\tsnb{अथ त्रयोदशाध्याये नक्षत्रस्योदयसमये ऽस्तसमये च निजोदयांशकज्ञानम्~॥}\\
\tsans{tallak.sa.na.m puurvaardhamadhye\te{\footnote{~\tsans{puurvaardhamadhye} {$\Big ]$} \tsans{puurvaadhamaye}\enskip Kh. The compound \tsans{puurvaadhamaye} in Kh can be  segmented as \tsans{puurva}\,$_\text{\acrshort{modifier}}$ + \tsans{adha}\,$_\text{\acrshort{modifier}}$ + \tsans{maye}\,$_\text{\acrshort{locative}-\acrshort{singular}}$; however, this reading is neither syntactically nor contextually coherent with the rest of the sentence. The omission of over-letter \textit{r}-diacritic (\textit{repha}) in \tsans{rdha} and confusing the glyph \tsans{ye} for the ligature \tsans{dhye} are fairly common scribal mistakes.}} %
proktameva~||}\\ (f.\thinspace 23r:\thinspace6--7~Kh)
\end{multicols}%
%
\begin{multicols}{2}
Thirteenth chapter\\
On the [right] \glsuseri{ascension_measure} (\matali) of the \gls{rising} (\tali) and \gls{setting} (\ghurub) of \glspl{celestial_object} (\kawakib).
\columnbreak

Now, in the thirteenth chapter, at the \gls{time_rising} (\textit{udaya-samaya}) and \gls{time_setting} (\textit{asta-samaya}) of a \gls{celestial_object} (\textit{nakṣatra}), the \gls{knowledge} (\textit{jñāna}) of the \gls{rising_location_degree} (\textit{nija-udaya-aṃśaka}) [\ie the oblique ascensions of the degrees of the ecliptic]. The \gls{definition} (\textit{lakṣaṇa}) of that has already been declared in the first half [of Part~II].
\end{multicols}
\medskip
%%%%%%%%%%%%%%%%%%%%%%%%%%%%%%%%%%%

\begin{multicols}{2}
\noindent\reversemarginpar\marginnote{\hypertarget{Pii14}{Discourse~II.14}}%
\tfarsi{باب چهاردهم \\
در معرفت سمت از ارتفاع یا انخفاض.%
\te{(f.\thinspace 18a:\thinspace20~\SjB)}\hfill}
\columnbreak

\noindent\normalmarginpar\marginnote{\hypertarget{Sii14}{Part~II.14}}%
\tsnb{अथ चतुर्दशाध्याये ऽभीप्सितोन्नतांशाधरांशेभ्यः\te{\footnote{~\tsans{.abhiipsitonnataa.m"saadharaa.m"sebhya.h} {$\Big ]$} \tsans{.abhiipsitonnata.m"saadhaaraa"sebhya}\enskip Kh. A regular sandhi of the words \tsans{.abhiipsita}\,$_\text{\acrshort{modifier}}$ + \tsans{unnata}\,$_\text{\acrshort{modifier}}$ + \tsans{a.m"sa}\,$_\text{\acrshort{modifier}}$ + \tsans{adhara}\,$_\text{\acrshort{modifier}}$  + \tsans{a.m"sebhya.h}\,$_\text{\acrshort{dative}/\acrshort{ablative}-\acrshort{plural}}$ generates the contextually apposite compound \tsans{.abhiipsitonna\underline{taa.m}"saa\underline{dha}raa.m"sebhya.h}. Kh attests \tsans{aadhaara}\,$_\text{\acrshort{modifier}}$ `support/base' instead of \tsans{adhara}\,$_\text{\acrshort{modifier}}$ `lower' in the chapter-title, but then uses \tsans{adhara}\,$_\text{\acrshort{modifier}}$ in several other places in this chapter. I suspect the irregular vowel-marks (the \textit{a}-diacritic in \tsans{ta.m} and \textit{ā}-diacritic in \tsans{dhaa}) in \tsans{.abhiipsitonnata.m"saadhaaraa"sebhya} are scribal mistakes (just like the missing \textit{anusvāra} over \tsans{raa}  or the missing \textit{visarga} in \tsans{bhya}).}} %
स्वदिगंशज्ञानम्~॥}\enskip
(f.\thinspace 23r:\thinspace24~Kh)
\end{multicols}%
%
\begin{multicols}{2}
Fourteenth chapter\\
On the \gls{knowledge} (\marifat) of the \gls{azimuth} (\samt) from the \gls{elevation} (\irtifa) or the \gls{depression} (\inkhifad) [of a celestial object].
\columnbreak

Now, in the fourteenth chapter, the \gls{knowledge} (\textit{jñāna}) of the \gls{azimuth_own_location} (\textit{sva-diś-aṃśa}) from the \gls{desired_degree_elevation} (\textit{abhīpsita-unnata-aṃśa}) and the \gls{degrees_depression} (\textit{adharā-aṃśa}) [of a celestial object].
\end{multicols}

\newpage %%%%________________________NEWPAGE

%%%%%%%%%%%%%%%%%%%%%%%%%%%%%%%%%%%

\begin{multicols}{2}
\noindent\reversemarginpar\marginnote{\hypertarget{Pii15}{Discourse~II.15}}%
\tfarsi{باب پانزدهم \\
در معرفت ارتفاع از سمت. \\
\te{(f.\thinspace 18b:\thinspace8--7~\SjB)}\hfill}
\columnbreak

\noindent\normalmarginpar\marginnote{\hypertarget{Sii15}{Part~II.15}}%
\tsnb{अथ पञ्चदशाध्याये दिगंशेभ्यो ऽभीष्टोन्नतांशाधरांशज्ञानम्~॥} %
\tsans{tatraananyatvaprakaaro\-papatti.h~||}
\enskip
(f.\thinspace 23v:\thinspace21--22~Kh)
\end{multicols}%
%
\begin{multicols}{2}
Fifteenth chapter\\
On the \gls{knowledge} (\marifat) of the \gls{elevation} (\irtifa) [of a celestial object] from [its] \gls{azimuth} (\samt).
\columnbreak

Now, in the fifteenth chapter, the \gls{knowledge} (\textit{jñāna}) of the \gls{desired_degree_elevation} (\textit{abhīṣṭa-unnata-aṃśa}) and the \gls{degrees_depression} (\textit{adharā-aṃśa}) [of a celestial object] from the \gls{azimuth_own_location} (\textit{sva-diś-aṃśa}). There, a \gls{demonstration} (\textit{upapatti}) by \gls{method_of_identity} (\textit{ananyatva-prakāra}) [is stated].
\end{multicols}
\medskip

%%%%%%%%%%%%%%%%%%%%%%%%%%%%%%%%%%%

\begin{multicols}{2}
\noindent\reversemarginpar\marginnote{\hypertarget{Pii16}{Discourse~II.16}}%
\tfarsi{باب شانزدهم \\
در معرفت خطّ نصف النهار.\\
\te{(f.\thinspace 18b:\thinspace16--15~\SjB)}\hfill}
\columnbreak

\noindent\normalmarginpar\marginnote{\hypertarget{Sii16}{Part~II.16}}%
\tsnb{अथ षोडशाध्याये याम्योतररेखाज्ञानम्~॥} %
\tsans{tasya eva naama madhyaahnarekheti\te{\footnote{~\tsans{naama madhyaahnarekheti} {$\Big ]$} \tsans{naama\thinspace dhyaahnarekheti}\enskip Kh. The technical word \tsans{madhyaahnarekhaa} is grammatically well-formed and contextually apposite to the discussions in this chapter. The words \tsans{naama} and \tsans{dhyaahnarekheti} occurs across a line break (lines 14 and 15) in Kh. This appears to be a haplography: the scribe inadvertently left out the second \tsans{ma} while copying.\label{emeneded_attested_sanskrit_example}}}~||}\\
(f.\thinspace 24r:\thinspace14--15~Kh)
\end{multicols}%
%
\begin{multicols}{2}
Sixteenth chapter\\
On the \gls{knowledge} (\marifat) of the \gls{line_midday} (\khatt\ \nisf\ \alnahar) [\ie the local meridian line].
\columnbreak

Now, in the sixteenth chapter, the \gls{knowledge} (\textit{jñāna}) of the \gls{meridian_line} (\textit{yāmya-uttara-rekhā}).~It is even called the \gls{line_midday} (\textit{madhyāhna-rekhā}). 
\end{multicols}
\medskip

%%%%%%%%%%%%%%%%%%%%%%%%%%%%%%%%%%%

\begin{multicols}{2}
\noindent\reversemarginpar\marginnote{\hypertarget{Pii17}{Discourse~II.17}}%
\tfarsi{باب هفدهم \\
در معرفت طول و عرض بلد. \\
\te{(f.\thinspace 18b:\thinspace29~SjB)}\hfill}
\columnbreak

\noindent\normalmarginpar\marginnote{\hypertarget{Sii17}{Part~II.17}}%
\tsnb{अथ सप्तदशाध्याये देशान्तराक्षांशज्ञानम्~॥} %
\tsans{tallak.sa.na.m ca~||}\enskip
(f.\thinspace 24v:\thinspace4~Kh)
\end{multicols}%
%
\begin{multicols}{2}
Seventeenth chapter\\
On the \gls{knowledge} (\marifat) of the [terrestrial] \gls{longitude_and_latitude_terrestrial} (\tul\ \va\ \ard\idafaconsonant\ \balad).
\columnbreak

Now, in the seventeenth chapter, the \gls{knowledge} (\textit{jñāna}) of \gls{degrees_terrestrial_long_latitude} (\textit{deśāntara-akṣa-aṃśa}) [in one's own location]. And their \glspl{definition} (\textit{lakṣaṇa}) [are first stated]. 
\end{multicols}

\newpage %%%%________________________NEWPAGE

%%%%%%%%%%%%%%%%%%%%%%%%%%%%%%%%%%%

\begin{multicols}{2}
\noindent\reversemarginpar\marginnote{\hypertarget{Pii18}{Discourse~II.18}}%
\tfarsi{باب هژدهم \\
در معرفت عرض اقلیم رؤیت.\\
\te{(f.\thinspace 19a:\thinspace20--19~SjB)}\hfill}
\columnbreak

\noindent\normalmarginpar\marginnote{\hypertarget{Sii18}{Part~II.18}}%
\tsnb{अथाष्टादशाध्याये दृक्क्षेपदृग्गतिज्ञानम्~॥} %
\tsans{tallak.sa.na.m ca~||}\enskip
(f.\thinspace 25r:\thinspace12--13~Kh)
\end{multicols}%
%
\begin{multicols}{2}
Eighteenth chapter\\
On the \gls{knowledge} (\marifat) of the \gls{latitude_visible_climate}  (\mbox{\ard\idafaconsonant\ \iqlim\idafaconsonant} \ruyat) [\ie the zenith distance of the nonagesimal point].
\columnbreak

Now, in the eighteenth chapter, the \gls{knowledge} (\textit{jñāna}) of \gls{ecliptic_zenith_distance} (\textit{dṛkkṣepa}) and the \gls{ecliptic_pole_zenith_distance} (\textit{drggati}). And their \glspl{definition} (\textit{lakṣaṇa}) [are first stated]. 
\end{multicols}
\medskip

%%%%%%%%%%%%%%%%%%%%%%%%%%%%%%%%%%%

\begin{multicols}{2}
\noindent\reversemarginpar\marginnote{\hypertarget{Pii19}{Discourse~II.19}}%
\tfarsi{باب نوزدهم \\
در استخراج بعد میان دو کوکب. \\
\te{(f.\thinspace 19a:\thinspace28~SjB)}\hfill}
\columnbreak

\noindent\normalmarginpar\marginnote{\hypertarget{Sii19}{Part~II.19}}%
\tsnb{अथैकोनविंशाध्याये\nobreak\te{\footnote{~%
\tsans{athaikonavi.m"saadhyaaye} {$\Big ]$} \tsans{athaikonavi.m"sodhyaaye}\enskip Kh. The locative adverbial phrase \tsans{athaikonavi.m"sodhyaaye} in Kh can be segmented as 
\tsans{atha}\,$_\text{\acrshort{indeclinable}}$ + \tsans{ekonavi.m"sa.h}\,$_\text{\acrshort{nominative}-\acrshort{singular}}$ + \tsans{adhyaaye}\,$_\text{\acrshort{locative}-\acrshort{singular}}$. However, the word \tsans{ekonavi.m"sa.h} `nineteen' is a cardinal number, and if used as an ordinal adjective, it should be in concord with the substantive \tsans{adhyaaye}. The meaning of the phrase `in the nineteenth chapter' is preserved in the compound \tsans{ekonavi.m\underline{"saa}dhyaaye} as well as the words \tsans{ekonavi.m"se .adhyaaye} (with a locative concord). I select the compounded form as it is consistent with the previous chapter-titles in Part II.}} %
द्विनक्षत्रान्तरांशक\-ज्ञानम्~॥}%
\tsans{ tallak.sa.nam~||}\enskip (f.\thinspace 25v:\thinspace6--7~Kh)
\end{multicols}%
%
\begin{multicols}{2}
Nineteenth chapter\\
On the \gls{determination} (\istikhraj) of the \gls{distance_between_celestial_object} (\bud\idafaconsonant\ \miyan\idafaconsonant\ \duvum\idafaconsonant\ \kawkab).
\columnbreak

Now, in the nineteenth chapter, the \gls{knowledge} (\textit{jñāna}) of \gls{degrees_separation_two_celestial_objects} (\textit{dvi-nakṣatra-antara-aṃśaka}). The \gls{definition} (\textit{lakṣaṇa}) of that [is first stated].  
\end{multicols}

\newpage %%%%________________________NEWPAGE

%%%%%%%%%%%%%%%%%%%%%%%%%%%%%%%%%%%

\begin{multicols}{2}
\noindent\reversemarginpar\marginnote{\hypertarget{Pii20}{Discourse~II.20}}%
\tfarsi{باب بیستم \\
در معرفت سمت قبله و انحراف او.\\
\te{(f.\thinspace 19b:\thinspace22--21~SjB)}\hfill}
\columnbreak

\noindent\normalmarginpar\marginnote{\hypertarget{Sii20}{Part~II.20}}%
\tsans{atha vi.m"satime .adhyaaye\te{\footnote{~\tsans{vi.m"satime .adhyaaye}\enskip{$\Big ]$}\enskip\tsans{vi.m"satidhyaaye}\enskip Kh. The attested form is morphologically defective. A regular sandhi of the words \tsans{vi.m"sati}\,$_\text{\acrshort{modifier}}$ + \tsans{adhyaaye}\,$_\text{\acrshort{locative}-\acrshort{singular}}$ generates \tsans{vi.m"sa\underline{tya}dhyaaye}, a locative adverbial phrase meaning  `in twenty chapters'. 
I correct this to \tsans{vi.m"satime .adhyaaye} `in the twentieth chapter' (using the ordinal form \tsans{vi.m"satima} `twentieth' instead of the cardinal number \tsans{vi.m"sati} `twenty') as it is consistent with the next two chapter-titles in Part~II.}} %
 svapure saumyayaamyadigbhyaa.m diga.m"sai.h kaa"sii kvaastiiti%
\te{\footnote{~\tsans{kaa"sii kvaastiiti} {$\Big ]$} \tsans{kaa"siikvaa\char"A8FBmii}$\overset{\text{?}}{\text{/}}$\tsans{sii\char"A8FBti}\,\enskip Kh, (\textit{conjecture}). There are no visible signs of scribal corrections or lacunae, but the writing (in red ink) is partially faded making it difficult to identify the letters with certainty. Nevertheless, there are no combinations of these letters that provide a grammatically valid and contextually apposite reading. I emend the  words to \tsans{kaa"sii kvaa\underline{stii}ti}, \lit the question ``Where is \Kashi?", that serves as the subject of the main sentence. (In other words: ``Where is \Kashi?", \textit{this} is [understood with \dots]). I suspect the scribe unwittingly copied the  glyph \tsans{mii}/\tsans{sii} for the ligature \tsans{stii} as they
often appear very similar in \Nagari\ palaeography.}} %
{\dev{ज्ञा}}yate~||}\enskip(f.\thinspace 26v:\thinspace17--18~Kh)
\end{multicols}%
%
\begin{multicols}{2}
Twentieth chapter\\
On the \gls{knowledge} (\marifat) of the \gls{azimuth_qibla} (\samt\idafaconsonant\ \qibla) and its \glslink{inclination_azimuth_qibla}{inclination} (\inhiraf[\idafaconsonant\ \samt\idafaconsonant\ \qibla]).
\columnbreak

Now, in the twentieth chapter, [the direction of] \Kashi\ is understood with \gls{degree_azimuth} (\textit{diś-aṃśa}) [measured] from both the \glsuseri{north_south_direction} (\textit{saumya-yāmya-diś}) in one's own city.
\end{multicols}
\medskip
%%%%%%%%%%%%%%%%%%%%%%%%%%%%%%%%%%%

\begin{multicols}{2}
\noindent\reversemarginpar\marginnote{\hypertarget{Pii21}{Discourse~II.21}}%
\tfarsi{باب بیست و یکم \\
در معرفت طالع از ارتفاع.%
\te{(f.\thinspace 20a:\thinspace31~SjB)}\hfill}
\columnbreak

\noindent\normalmarginpar\marginnote{\hypertarget{Sii21}{Part~II.21}}%
\tsans{atha ekavi.m"satime .adhyaaye .abhii.s.tonnataa.m\-"sebhyo\nobreak\te{\footnote{~\tsans{.abhii.s.tonnataa.m"sebhyo} {$\Big ]$} \tsans{.abhii.s.tonnavaa.m"sebhyo}\enskip Kh. The compound \tsans{.abhii.s.tonnavaa.m"sebhyo} in Kh can be segmented as \tsans{.abhii.s.ta}\,$_\text{\acrshort{modifier}}$ + \tsans{unnava}\,$_\text{?}$ + \tsans{a.m"sebhyo}\,$_\text{\acrshort{dative}/\acrshort{ablative}-\acrshort{plural}}$; however,  \tsans{unnava} is neither a valid morphophonemic compound nor a standard lexical entry. The word \tsans{unna\underline{ta}} is contextually relevant and also variously attested in this chapter, \eg \tsans{unnatajyaayaa.h} (f.\thinspace 27r:\thinspace27~Kh) or  \tsans{samunnatajyaa} (f.\thinspace 27v:\thinspace14~Kh). I suspect the scribe inattentively copied the glyph \tsans{vaa.m} for \tsans{taa.m} in the chapter-title.}} %
lagna{\dev{ज्ञा}}nam~||}\enskip 
(f.\thinspace 27r:\thinspace26--27~Kh)
\end{multicols}%
%
\begin{multicols}{2}
Twenty-first chapter\\
On the \gls{knowledge} (\marifat) of the \gls{ascendant} (\tali) from the \gls{elevation} (\irtifa). 
\columnbreak

Now, in the twenty-first chapter, the \gls{knowledge} (\textit{jñāna}) of the \gls{ascendant} (\textit{lagna}) from the \gls{desired_degree_elevation} (\textit{abhīṣṭa-unnata-aṃśa}).
\end{multicols}

\newpage %%%%________________________NEWPAGE

%%%%%%%%%%%%%%%%%%%%%%%%%%%%%%%%%%%

\begin{multicols}{2}
\noindent\reversemarginpar\marginnote{\hypertarget{Pii22}{Discourse~II.22}}%
\tfarsi{باب بیست و دوم \\
در معرفت ارتفاع یا انخفاض كواكب از طالع.%
\te{(f.\thinspace 20b:\thinspace25--24~SjB)}\hfill}
\columnbreak

\noindent\normalmarginpar\marginnote{\hypertarget{Sii22}{Part~II.22}}%
\tsnb{अथ द्वाविंशतिमे ऽध्याये खगस्य स्वोदयांशेभ्यो\te{\footnote{~\tsans{svodayaa.m"sebho} {$\Big ]$} \tsans{svodaya.m"sobhyo}\enskip Kh. A regular sandhi of the words \tsans{sva}\,$_\text{\acrshort{modifier}}$ + \tsans{udaya}\,$_\text{\acrshort{modifier}}$ + \tsans{a.m"sebhyo}\,$_\text{\acrshort{dative}/\acrshort{ablative}-\acrshort{plural}}$ generates
\tsans{svodayaa.m\underline{"se}bhyo} where the terminal consonant -\tsans{"sa} (of \tsans{a.m"sa})
changes to -\tsans{"se} (and not -\tsans{"so}) before the dative/ablative case ending -\tsans{bhya.h}. The \textit{o}-diacritic in \tsans{"so} appears to be a scribal hypercorrection.}} %
ऽभीष्टोन्नतांशानामधरांशकानां च ज्ञानम्~॥}%
\tsans{ etallak.sa.na.m puurvame.soktam~||}\enskip (f.\thinspace 28r:\thinspace16--17~Kh)
\end{multicols}%
%
\begin{multicols}{2}
Twenty-second chapter\\
On the \gls{knowledge} (\marifat) of the \gls{elevation} (\irtifa) or \gls{depression} (\inkhifad) of \glspl{celestial_object} (\kawakib) from the \gls{ascendant} (\tali). 
\columnbreak

Now, in the twenty-second chapter, the \gls{knowledge} (\textit{jñāna}) of the \gls{desired_degree_elevation} (\textit{abhīṣṭa-unnata-aṃśa}) and of the \gls{degrees_depression} (\textit{adhara-aṃśaka}) from the \gls{rising_location_degree} (\textit{sva-udaya-aṃśa}) of a \gls{celestial_object} (\textit{khaga}). The \gls{definition} (\textit{lakṣaṇa}) of this has already been declared in the first half [of Part~II].
\end{multicols}

\newpage %%%%________________________NEWPAGE

%%%%%%%%%%%%%%%%%%%%%%%%%%%%%%%%%%%

\begin{multicols}{2}\label{partII_colophon_sanskrit}
The \ZijiShahJahani: Discourse~II does not have a colophon. F~21r:\thinspace  21~\SjB\ ends with the last line of chapter twenty-two: \tfarsi{\dots بهین موامره همین مطلوب حاصل آید.} 
\columnbreak

\noindent\normalmarginpar\marginnote{\hypertarget{SiiCol}{Part~II\\\textit{colophon}}}%
\tsans{ya.h "srii"saahajahaa/\te{\footnote{~\tsans{"srii"saahajahaa/} {$\Big ]$} \tsans{"srii"saahahaa/}\enskip Kh. I suspect the scribe unwittingly left out the letter \tsans{ja} while copying; most other occurrences of \Shahjahan's name in the text read \tsans{"srii"saahajahaa/} (sometimes without the terminal nasal diacritic \textit{candrabindu}), \eg folia 3v:\thinspace1, 5r:\thinspace15, or 6v:\thinspace22 of Kh. Also, this verse is in the nineteen-syllabled \sardulavikridita\ meter which would require a five-syllabled word like \tsans{"srii"saahajahaa/\selip} for metrical completion (\textit{pada-pūrti}).}} %
n.rpaalamuku.taala"nkaara\-cuu.daama.nistasyaa{\dev{ज्ञा}}%
mavalambya dustaramamu.m siddhaanta\-sindhu.m taran~|| nityaananda iti dvijottamak.rpa.h\te{\footnote{~The appositive nominal compound \tsans{dvijottamak.rpa.h} is grammatically irregular. The terminal word \tsans{k.rpa.h}\,$_\text{\acrshort{nominative}-\acrshort{singular}}$ is attested as the sage Kṛpācārya (from the \textit{Mahābhārata}) in Sanskrit lexicons. In this case, however, I suspect \tsans{dvijottamak.rpa.h} is a metrical contraction (\textit{pada-anatireka-karaṇa}) that can be parsed as \tsans{dvijottamaa.naa.m k.rpaapaatra.h} `[the one] worthy of the mercy (\textit{kṛpā}) of the best \Brahmana s'. In the colophon of \Nityananda's \Sarvasiddhantaraja\ (1639), we find a related phrase \tsans{dvijaanaamaa{\tsnb{ज्ञा}}kaarii} %
 `one who executes the commands of the \Brahmana s' (\ie obedient of the authority of \Brahmana s) as an epithet of \Nityananda\ (\vid\ \cite[228]{PetersonCatalogue}; \cite[102]{DvivediGanakaTarangini}).}} %
"sriidevadattaatmajastripra"snapracuroktiyuktisahita.m kaa.n.da.m dvitiiya.m hyagaat~||}%
 \enskip (f.\thinspace 28v:\thinspace15--18~Kh)
\medskip

\Nityananda, who crosses over this unconquerable `Ocean of the \Siddhanta s [\scl composes the \Siddhantasindhu] [by] holding onto the command of Śri \Shahjahan\ who is the crest jewel of the ornamental crown of kings, [the man who is worthy of] the mercy of the best \Brahmana s, the son of Śri \Devadatta, has just finished the second part accompanied by many statements and rationales
on the \triprasna.\footnote{~\Vid\ remark~4 in \S~\ref{comparative_overview_chapters_zij_sindhu}, page~\pageref{triprasna_remark}.}
\end{multicols}

%%%%%%%%%%%%%%%%%%%%%%%%%%%%%%%%%%%
%%%%%%%%%%%%%%%%%%%%%%%%%%%%%%%%%%%



\clearpage

\section{\ZijiShahJahani, Discourse II.6: text and translation}
\label{zijshahjahan_persian_english}
\setcounter{footnote}{0}
\renewcommand{\thefootnote}{[\roman{footnote}]} % for distinct numbering of Persian edition (chapter vi) footnotes (starting at [i])
\setlength{\footnotesep}{\baselineskip} % for space between footnotes
\tfarsi{%
\te{$\rceil$}\tfarsib{باب ششم}\te{\reversemarginpar\marginnote{{\footnotesize f.\thinspace 21b:\thinspace 20~\SjA,}~$\rceil$\\
{\footnotesize f.\thinspace 16a:\thinspace 25~\SjB\phantom{nn}}}}\\
در معرفت بعد كواكب از معدّل النهار.\hfill}
\\[0.5\baselineskip]

\noindent\normalmarginpar\marginnote{\hypertarget{Ppass1}{[1]}}
\tfarsi{%
عرض کوکب و میل ثاني درجه او، %
اگر هر دو در یک جهت باشند
جمع کنیم والا تفاضل بگیریم  و آن را حصّه بعد خوانیم  و جهت حصّه بعد
جهت مجموع یا جهت فضل باشد.\hfill}
\medskip

\noindent\normalmarginpar\marginnote{\hypertarget{Ppass2}{[2]}}\label{folio_break_persian_example}
\tfarsi{%
پس جيب حصّه بعد را در جیب تمام
میل منکوس درجه كوكب منحطّ ضرب کنیم حاصل جيب % 
\te{$\rceil$}%
بعد %
\te{\reversemarginpar\marginnote{{\footnotesize f.\thinspace 22a:\thinspace 1~\SjA}~$\rceil$}}%
بود.\hfill}
\\[\baselineskip]

\noindent\normalmarginpar\marginnote{\hypertarget{Ppass3}{[3]}}
\tfarsi{%
بوجهی\te{\footnote{ %
\tfarsi{بوجهی}~{$\Big ]$}~%
\raisebox{-.6ex}{\includegraphics[scale=.85]{Images/qif_over_text_logograph.eps}}~\SjA. The overlining `\includegraphics[scale=0.85]{Images/qif_arabic_logograph.eps}' over the word \tfarsi{بوجهی} is used to indicate a notable pause in the reading. As \textcite[173]{Gacek} explains, it is most likely a logograph (word-symbol) of the Arabic word \tfarsi{قف} (\qif) `stop' to indicate a pause in reading; or alternatively, the abbreviation \tfarsi{فتــ‬} (\textit{fata--}) of the phrase \tfarsi{فتأملها} (\fatammalha)
`reflect on it'.\label{qif_overlinning}}}
دیگر جيب حصّه بعد را در جیب تمام میل کلّی ضرب کنیم
و حاصل را بر جيب تمام میل\te{\footnote{ % 
\tfarsi{میل}~{$\Big ]$}~\tfarsi{%
میل \cancel{کلّی}}~\SjA, cancellation \textit{intra lineam}.}}
ثاني درجه آن کوکب قسمت کنیم
خارج  
\te{$\rceil$}%
قسمت %  
\te{\reversemarginpar\marginnote{{\footnotesize f.\thinspace 16b:\thinspace 1~\SjB}~$\rceil$}}%
جيب بعد باشد و جهت آن جهت حصّه بعد باشد.\hfill}
\medskip

\noindent\normalmarginpar\marginnote{\hypertarget{Ppass4}{[4]}}
\tfarsi{%
و چون جيب حصّه بعد را در جدول جيب تمام میل کلّی درآرند و حاصل را بر جيب تمام میل ثاني درجه آن کوکب  قسمت کنند خارج قسمت جيب بعد باشد.\hfill}
\medskip

\noindent\normalmarginpar\marginnote{\hypertarget{Ppass5}{[5]}}
\tfarsi{%
و اگر کوکب را عرض نباشد، میل درجه او بعد باشد.\hfill}

%%%%%%%%%%%%%%%%%%%%%%%%%%%%%%%%%
\clearpage\phantomsection
%%%%%%%%%%%%%%%%%%%%%%%%%%%%%%%%%

\textbf{Sixth chapter} \\
On the \gls{knowledge} (\marifat) of the \glspl{distance_celestial_object} (\bud\idafaconsonant\ \kawakib\ \az\ \muaddil\ \alnahar).   
\\[0.5\baselineskip]

\noindent\reversemarginpar\marginnote{\hypertarget{PEpass1}{[1]}}%
{[Given]} the \gls{latitude_celestial_object} (\ard\idafaconsonant\ \kawkab) and the \gls{second_declination_degree} (\mayl\idafaconsonant\ \thani\idafavowel\ \daraji\idafavowel\ \uy), if both should be in \gls{one_direction} (\yik\ \jahat), we \gls{sum} (*\jam\ \kardan) [them]; otherwise, we should take the \gls{difference} (\tafadul). And we call that [result] the \gls{share_of_distance} (\hissi\idafavowel\ \bud). And the \gls{direction_share_of_distance} (\jahat\idafaconsonant\ \hissi\idafavowel\ \bud) should be the \gls{direction_sum} (\jahat\idafaconsonant\ \majmu) or the \gls{direction_residue} (\jahat\idafaconsonant\ \fadla). 
\label{passage_1_english_persian} 
\medskip

\noindent\reversemarginpar\marginnote{\hypertarget{PEpass2}{[2]}}%
Then, we \glslink{low_multiplication}{low-multiply} (*\munhatt\idafaconsonant\ \darb\ \kardan) the \gls{Sine_share_of_distance} (\jayb\idafaconsonant\ \hissi\idafavowel\ \bud) in the \gls{Cosine_inverse_declination_degree_celestial_object} (\jayb\idafaconsonant\ \tamam\idafaconsonant\ \mayl\idafaconsonant\ \mankus\idafaconsonant\ \daraji\idafavowel\ \kawkab). The \gls{result} (\hasil) is the \gls{Sine_distance} (\jayb\idafaconsonant\ \bud). 
\\[\baselineskip]

\noindent\reversemarginpar\marginnote{\hypertarget{PEpass3}{[3]}}%
In another way, we \glslink{multiplication}{multiply} (*\darb\ \kardan) the \gls{Sine_share_of_distance} (\jayb\idafaconsonant\ \hissi\idafavowel\ \bud) in the \gls{Cosine_maximum_declination} (\jayb\idafaconsonant\ \tamam\idafaconsonant\ \mayl\idafaconsonant\ \kulli) [\ie  in the Cosine of the ecliptic obliquity] and we \glslink{division}{divide} (*\qismat\ \kardan) the \gls{result} (\hasil) over the \gls{Cosine_second_declination_degree} (\jayb\idafaconsonant\ \tamam\idafaconsonant\ \mayl\idafaconsonant\ \thani\idafavowel\ \daraji) of that \gls{celestial_object} (\kawkab). The \glslink{quotient_division}{quotient of the division} (\kharij\idafaconsonant\ \qismat) should be the \gls{Sine_distance} (\jayb\idafaconsonant\ \bud), and its \gls{direction} (\jahat)  should be the \gls{direction_share_of_distance} (\jahat\idafaconsonant\ \hissi\idafavowel\ \bud). 
\medskip

\noindent\reversemarginpar\marginnote{\hypertarget{PEpass4}{[4]}}%
And since they \gls{extract} (*\darardan) [the product of the multiplication with] the \gls{Sine_share_of_distance} (\jayb\idafaconsonant\ \hissi\idafavowel\ \bud) from the \gls{table_Cosine_maximum_declination} (\jadval\idafaconsonant\ \jayb\idafaconsonant\ \tamam\idafaconsonant\ \mayl\idafaconsonant\ \kulli) [\ie in the table of the Cosine of the ecliptic obliquity] and they \glslink{division}{divide} (*\qismat\ \kardan) the \gls{result} (\hasil) over the \gls{Cosine_second_declination_degree} (\jayb\idafaconsonant\ \tamam\idafaconsonant\ \mayl\idafaconsonant\ \thani\idafavowel\ \daraji) of that \gls{celestial_object} (\kawkab), the \glslink{quotient_division}{quotient of the division} (\kharij\idafaconsonant\ \qismat) should be the \gls{Sine_distance} (\jayb\idafaconsonant\ \bud).
\medskip

\noindent\reversemarginpar\marginnote{\hypertarget{PEpass5}{[5]}}%
And if a \gls{celestial_object} (\kawkab) should have no \gls{latitude} (\ard), the \gls{declination_degree} (\mayl\idafaconsonant\ \daraji\idafavowel\ \uy) should be the \gls{distance} (\bud). 

%%%%%%%%%%%%%%%%%%%%%%%%%%%%%%%%%
\clearpage\phantomsection
%%%%%%%%%%%%%%%%%%%%%%%%%%%%%%%%%

\noindent\normalmarginpar\marginnote{\hypertarget{Ppass6}{[6]}}
\tfarsi{%
و اگر عرض باشد اما درجه او را میل نباشد، جيب عرض او را در جیب تمام میل كلّی منحطّ ضرب کنیم
یا در\te{\footnote{ %
\tfarsi{یا در}~{$\Big ]$}~%
\tfarsi{یاد در}~\SjB, dittography of the first \tfarsi{د}.\label{emeneded_attested_persian_example}}}   
جدول سابق درآریم
حاصل جيب بعد باشد و جهت او جهت عرض باشد.\hfill}
\medskip

\noindent\normalmarginpar\marginnote{\hypertarget{Ppass7}{[7]}}
\tfarsi{%
و اگر میل درجه او میل كلّی باشد، حصّة البعد بعينه بعد باشد.\hfill}
\\[\baselineskip]

\noindent\normalmarginpar\marginnote{\hypertarget{Ppass8}{[8]}}
\tfarsi{%
و بوجهی\te{\footnote{ %
\tfarsi{و بوجهی}~{$\Big ]$}~%
\raisebox{-.6ex}{\includegraphics[scale=.85]{Images/qif_over_text_logograph_2.eps}}~\SjA, with emphasis. \Vid\ footnote~\ref{qif_overlinning}.}}
دیگر جيب بعد درجه کوکب از انقلاب اقرب در جیب
تمام عرض کوکب منحطّ ضرب کنیم حاصل جیب بعد کوکب از
دایرهٔ ماره باقطاب اربعه %
\wrapfootnote{%
باشد.\hfill\break
\te{\phantom{a}}\break % Required to allow a line break within a single \tfarsi{} group (Hack!)
\te{\noindent\normalmarginpar\marginnote{\hypertarget{Ppass9}{[9]}}}%
پس جیب عرض کوکب را بر جيب تمام بعد از دایرهٔ ماره باقطاب اربعه}%
{ %
\tfarsi{باقطاب اربعه}%
\thinspace\dots\thinspace
\tfarsi{باشد.\@  پس}~\te{$\Big ]$}
\begin{itemize}
    \item \includegraphics[height=.75em]{Images/marginal_insertion_end.eps}~~%
\tfarsi{باقطاب اربعه}%
\thinspace\dots\thinspace
\tfarsi{پس} %
$\stackrel{\text{\tiny\signederenvoi}}{\text{\tfarsi{باشد}}}$~\SjA, %
inserted in the exterior margin by the same hand.
The penultimate word of passage~[8] on f.\thinspace 22a:\thinspace 6  has 
an insertion mark `\signederenvoi' (\textit{signe-de-renvoi}) placed above it, \scl  \dots\thinspace\tfarsi{منحطّ} %
$\stackrel{\text{\tiny\signederenvoi}}{\text{\tfarsi{اربعه} }}$ %
\tfarsi{باقطاب}\thinspace\dots\thinspace .  
The first word of the marginal text also bears the same mark, \textit{supra verbum}, \scl\ \dots\thinspace
\tfarsi{پس} %
$\stackrel{\text{\tiny\signederenvoi}}{\text{\tfarsi{باشد}}}$. The marginal text ends with  `\includegraphics[height=.75em]{Images/marginal_insertion_end.eps}': this could be a calligraphic variant of the abbreviation \tfarsi{هى} (\textit{hā}\Ayn) of the Arabic word \tfarsi{إنتهى} (\intiha) meaning `it is finished', \vid\ \textcite[117]{Gacek};
\item \tfarsi{باقطاب اربعه}%
\thinspace\dots\thinspace
\tfarsi{باشد.\@  پس}~\textit{omitted} \SjB,  per homeoteleuton. The penultimate word \tfarsi{\underline{اربعه}} of passage~[8], %
before the missing text, is identical to the last word of the missing text %
\tfarsi{\underline{اربعه} منحطّ\thinspace\dots}\thinspace %
 in  passage~[9].\\[5pt]
\textbf{\scshape remark}: MS~Or.\thinspace 566 of the \ZijUlughBeg\ from  University Library (Cambridge) is also missing the same amount of text. The missing words \tfarsi{باشد\thinspace\dots\thinspace اربعه} span the length of a line between the end of line 18 (at \tfarsi{اربعه}) %
and the beginning of line 19 (at \tfarsi{منحطّ}) on folio 16r of this manuscript.\label{zij_ulugh_beg_missing_line}
\end{itemize}\vspace{-\topsep}
} %
منحطّ قسمت کنیم %
\wrapfootnote{و به خارج قسمت}%
{\tfarsi{%
و~به~خارج~قسمت
}~\te{$\Big ]$}~\tfarsi{
و به خارج قسمت کنیم  و به خارج قسمت
}~\SjB, dittography of the second \tfarsi{و به خارج قسمت}. %
I suspect a parableptic error as the word \tfarsi{قسمت} %
appears twice on the line in close proximity: 
\dots\thinspace \tfarsi{\underline{قسمت} از}\thinspace\dots\thinspace\tfarsi{\underline{قسمت} کنیم}\thinspace\dots\thinspace.\label{wrap_footnote_example}}
از جدول جيب قوس بگیریم و آن را قوس اوّل خوانیم  و جهت آن جهت عرض کوکب بود.\hfill}

%%%%%%%%%%%%%%%%%%%%%%%%%%%%%%%%%
\clearpage\phantomsection
%%%%%%%%%%%%%%%%%%%%%%%%%%%%%%%%%

\noindent\reversemarginpar\marginnote{\hypertarget{PEpass6}{[6]}}%
And if the \gls{latitude} (\ard) [of a celestial object] should exist but \gls{its_degree} (\daraji\idafavowel\ \uy)  should have no \gls{declination}  (\mayl), we  \glslink{low_multiplication}{low-multiply} (*\munhatt\idafaconsonant\ \darb\ \kardan) the \gls{Sine_latitude} (\jayb\idafaconsonant\ \ard\idafaconsonant\ \uy) in the \gls{Cosine_maximum_declination} (\jayb\idafaconsonant\ \tamam\idafaconsonant\ \mayl\idafaconsonant\ \kulli)  [\ie  in the Cosine of the ecliptic obliquity], or we \gls{extract} (*\darardan) [the product of this multiplication] from the preceding \gls{table} (\jadval) [\ie in the table of the Cosine of the maximum declination]. The \gls{result} (\hasil) should be the \gls{Sine_distance} (\jayb\idafaconsonant\ \bud), and its \gls{direction} (\jahat) should be the \gls{direction_latitude} (\jahat\idafaconsonant\ \ard).\label{passge_6_ard_glossary_format_example}
\medskip

\noindent\reversemarginpar\marginnote{\hypertarget{PEpass7}{[7]}}%
And if the \gls{declination_degree} (\mayl\idafaconsonant\ \daraji\idafavowel\ \uy) should be the \gls{total_declination} (\mayl\idafaconsonant\ \kulli) [\ie the obliquity of the ecliptic], the \gls{share_of_distance} (\hissatalbud) itself should be the \gls{distance} (\bud). 
\\[\baselineskip]

\noindent\reversemarginpar\marginnote{\hypertarget{PEpass8}{[8]}}%
And in another way, we \glslink{low_multiplication}{low-multiply} (*\munhatt\idafaconsonant\ \darb\ \kardan) the \gls{Sine_distance_celestial_object_nearest_solstice} (\jayb\idafaconsonant\ \bud\idafaconsonant\ \daraji\idafavowel\ \kawkab\ \az\ \inqilab\idafaconsonant\ \aqrab) in the \glsuseri{Cosine_latitude_celestial_object} (\jayb\idafaconsonant\ \tamam\idafaconsonant\ \ard\idafaconsonant\ \kawkab). The \gls{result} (\hasil) should be the
\glslink{Sine_distance_celestial_object_solstitial_colure}{Sine of the distance of the celestial object from the `circle passing through the four poles'} [\ie from the `solstitial colure'] (\jayb\idafaconsonant\ \bud\idafaconsonant\ \kawkab\ \az\ \guillemotleft\dayiri\idafavowel\ \marri\ \biaqtab\idafaconsonant\ \arbai\guillemotright). 
\medskip

\noindent\reversemarginpar\marginnote{\hypertarget{PEpass9}{[9]}}%
Then, we \glslink{low_division}{low-divide} (*\munhatt\idafaconsonant\ \qismat\ \kardan) the \gls{Sine_latitude_celestial_object} (\jayb\idafaconsonant\ \ard\idafaconsonant\ \kawkab) over the \glslink{Cosine_distance_solstitial_colure}{Cosine of the distance [of the celestial object] from the `circle passing through the four poles'} (\jayb\idafaconsonant\ \tamam\idafaconsonant\ \bud\ \az\ \guillemotleft\dayiri\idafavowel\ \marri\ \biaqtab\idafaconsonant\ \arbai\guillemotright). And for the \glslink{quotient_division}{quotient of the division} (\kharij\idafaconsonant\ \qismat), we should take the \gls{arc} (\qaws) from the \gls{table_of_Sine} (\jadval\idafaconsonant\ \jayb). And we call it the \gls{first_arc} (\qaws\idafaconsonant\ \avval) and its \gls{direction} (\jahat) is the \gls{direction_latitude_celestial_object} (\jahat\idafaconsonant\ \ard\idafaconsonant\ \kawkab). 

%%%%%%%%%%%%%%%%%%%%%%%%%%%%%%%%%
\clearpage\phantomsection
%%%%%%%%%%%%%%%%%%%%%%%%%%%%%%%%%

\noindent\normalmarginpar\marginnote{\hypertarget{Ppass10}{[10]}}
\tfarsi{%
پس اگر عرض و میل درجه كوكب هر دو در یک جهت باشند، قوس اوّل و میل كلّی را جمع کنیم. و اگر از 
ربع\te{\footnote{ %
\tfarsi{ربع}~{$\Big ]$}  
\begin{itemize}
    \item \tfarsi{بع}\cancel{\tfarsi{ا}}\tfarsi{ر}~\SjA,  erasure
\textit{intra lineam};
    \item \tfarsi{رابع}~\SjB. The Arabic words \tfarsi{ربع} (\rub) and \tfarsi{رابع} (\rabi) are the fractional (`one-fourth') and ordinal (`fourth') forms of the number four \tfarsi{أربعة} (\arbai) respectively. The mathematical context of the passage supports the fractional meaning `one-fourth' or a `quarter'. The reading in \SjB\ could be a semantic mistake by the scribe in copying Arabic loanwords (\muarrab).
\end{itemize}\vspace{-\topsep}}}
دور زیاده شود، تمام مجموع تا نصف دور بگیریم. و اگر در جهت مختلف باشند، تفاضل میان هر دو بگیریم حاصل قوس دوم باشد و جهتش جهت مجموع یا جهت فضل باشد.\hfill}
\medskip

\noindent\normalmarginpar\marginnote{\hypertarget{Ppass11}{[11]}}
\tfarsi{%
پس جیب قوس دوم را در جیب تمام بعد از
دایرهٔ ماره باقطاب اربعه منحطّ ضرب کنیم حاصل جيب بعد کوکب باشد  و جهتش جهت قوس دوم باشد.\hfill}

%%%%%%%%%%%%%%%%%%%%%%%%%%%%%%%%%
\clearpage\phantomsection
%%%%%%%%%%%%%%%%%%%%%%%%%%%%%%%%%

\noindent\reversemarginpar\marginnote{\hypertarget{PEpass10}{[10]}}%
Then, if the \gls{latitude} (\ard) and the \gls{declination_degree_celestial_object} (\mayl\idafaconsonant\ \daraji\idafavowel\ \kawkab) both should be in \gls{one_direction} (\yik\ \jahat), we \gls{sum} (*\jam\ \kardan) the \gls{first_arc} (\qaws\idafaconsonant\ \avval) and the \gls{total_declination} (\mayl\idafaconsonant\ \kulli) [\ie the obliquity of the ecliptic]. And if [the sum] \gls{exceeds} (*\ziyadi\ \shudan) \gls{one_quarter} (\rabi) [\ie is greater than 90\degree], we should take the \gls{whole_sum} (\tamam\idafaconsonant\ \majmu) up to \gls{one_half} (\nisf) [\ie up to 180\degree]. And if they should be in \gls{different_directions} (\jahat\idafaconsonant\ \mukhtalif), we should take the \gls{difference} (\tafadul) between the two; the \gls{result} (\hasil) [in both cases] should be [called] 
the \gls{second_arc} (\qaws\idafaconsonant\ \duvum) and its \gls{direction} (\jahat) should be the \gls{direction_sum} (\jahat\idafaconsonant\ \majmu) or the \gls{direction_residue} (\jahat\idafaconsonant\ \fadla). 
\medskip

\noindent\reversemarginpar\marginnote{\hypertarget{PEpass11}{[11]}}%
Then, we \glslink{low_multiplication}{low-multiply} (*\munhatt\idafaconsonant\ \darb\ \kardan) the \gls{Sine_second_arc} (\jayb\idafaconsonant\ \qaws\idafaconsonant\ \duvum) in the \glslink{Cosine_distance_solstitial_colure}{Cosine of the distance [of a celestial object] from the `circle passing through the four poles'} (\jayb\idafaconsonant\ \tamam\idafaconsonant\ \bud\ \az\ \guillemotleft\dayiri\idafavowel\ \marri\ \biaqtab\idafaconsonant\ \arbai\guillemotright). The result should be the \gls{Sine_distance_celestial_object} (\jayb\idafaconsonant\ \bud\idafaconsonant\ \kawkab) and its \gls{direction} (\jahat) should be the \gls{direction_second_arc}  (\jahat\idafaconsonant\ \qaws\idafaconsonant\ \duvum). 

%%%%%%%%%%%%%%%%%%%%%%%%%%%%%%%%%%%%%%%%%%%%%%%%%%%%%%%%%%%%%%%%%

\clearpage
% \setcounter{footnote}{0}
\renewcommand{\thefootnote}{[\roman{footnote}]}
\setlength{\footnotesep}{\baselineskip}
\tfarsi{%
\te{\raisebox{1.75ex}{\color{darkgray}$\neg$}}\kern-.75ex\tfarsib{باب ششم}\te{\marginnote{\te{\raisebox{1.75ex}{\color{darkgray}$\neg$}}\footnotesize f.~21b:\thinspace 20~\SjA\ \&\\
\hspace{8pt}f.~16a:\thinspace 25~\SjB}}\\
در معرفت بعد كواكب از معدّل النهار.\hfill}
\bigskip


\tfarsi{%
عرض کوکب و میل ثاني درجه او، %
اگر هر دو در یک جهت باشند
جمع کنیم والا تفاضل بگیریم  و آن را حصّه بعد خوانیم  و جهت حصّه بعد
جهت مجموع یا جهت فضل باشد.\te{(1)}\hfill}


\tfarsi{%
پس جيب حصّه بعد را در جیب تمام
میل منکوس درجه كوكب منحطّ ضرب کنیم حاصل جيب % 
\te{\raisebox{1.75ex}{\color{darkgray}$\neg$}}\kern-.75ex
بعد %
\te{\marginnote{\te{\raisebox{1.75ex}{\color{darkgray}$\neg$}}\footnotesize f.~22a:\thinspace 1~\SjA}}%
بود.\te{(2)}\hfill}
\bigskip

\tfarsi{%
بوجهی\te{\footnote{ %
\tfarsi{بوجهی}  {\normalsize ]}  \raisebox{-.6ex}{\includegraphics[scale=.85]{Images/qif_over_text_logograph.eps}}~\SjA. The overlining `\includegraphics[scale=0.85]{Images/qif_arabic_logograph.eps}' over the word \tfarsi{بوجهی} is used to indicate a notable pause in the reading. As \textcite[173]{Gacek} explains, it is most likely a logograph (word-symbol) of the Arabic word \tfarsi{قف} (\qif) `stop' to indicate a pause in reading; or alternatively, the abbreviation \tfarsi{فتــ‬} (\textit{fata--}) of the phrase \tfarsi{فتأملها} (\fatammalha)
`reflect on it'.\label{qif_overlinning}}}
دیگر جيب حصّه بعد را در جیب تمام میل کلّی ضرب کنیم
و حاصل را بر جيب تمام میل\te{\footnote{ % 
\tfarsi{میل} {\normalsize ]} \tfarsi{%
میل \cancel{کلّی}}~\SjA, cancellation \textit{intra lineam}.}}
ثاني درجه آن کوکب قسمت کنیم
خارج  
\te{\raisebox{1.75ex}{\color{darkgray}$\neg$}}\kern-.75ex
قسمت %  
\te{\marginnote{\te{\raisebox{1.75ex}{\color{darkgray}$\neg$}}\footnotesize f.~16b:\thinspace 1~\SjB}}%
جيب بعد باشد و جهت آن جهت حصّه بعد باشد.\te{(3)}\hfill}


\tfarsi{%
و چون جيب حصّه بعد را در جدول جيب تمام میل کلّی درآرند و حاصل را بر جيب تمام میل ثاني درجه آن کوکب  قسمت کنند خارج قسمت جيب بعد باشد.\te{(4)}\hfill}

\tfarsi{%
و اگر کوکب را عرض نباشد، میل درجه او بعد باشد.\te{(5)}\hfill}

\tfarsi{%
و اگر عرض باشد اما درجه او را میل نباشد، جيب عرض او را در جیب تمام میل كلّی منحطّ ضرب کنیم
یا در\te{\footnote{ %
\tfarsi{یا در} {\normalsize ]} 
\tfarsi{یاد در}~\SjB, dittography of the first \tfarsi{د}.}}   
جدول سابق درآریم
حاصل جيب بعد باشد و جهت او جهت عرض باشد.\te{(6)}\hfill}

\tfarsi{%
و اگر میل درجه او میل كلّی باشد، حصّة البعد بعينه بعد باشد.\te{(7)}\hfill}
\newpage

\tfarsi{%
و بوجهی\te{\footnote{ %
\tfarsi{و بوجهی} {\normalsize ]} 
\raisebox{-.6ex}{\includegraphics[scale=.85]{Images/qif_over_text_logograph_2.eps}}~\SjA, with emphasis. \Vid\ footnote~\ref{qif_overlinning}.}}
دیگر جيب بعد درجه کوکب از انقلاب اقرب در جیب
تمام عرض کوکب منحطّ ضرب کنیم حاصل جیب بعد کوکب از
دایرهٔ ماره باقطاب اربعه %
\wrapfootnote{%
باشد.\nobreak\te{(8)}\hfill\break
\te{\phantom{a}}\break % Required to allow a line break within a single \tfarsi{} group (Hack!)
پس جیب عرض کوکب را بر جيب تمام بعد از دایرهٔ ماره باقطاب اربعه}%
{ %
\tfarsi{باقطاب اربعه}%
\thinspace\dots\thinspace
\tfarsi{باشد.\@  پس} 
\te{\normalsize ]}
\begin{itemize}
    \item \includegraphics[height=.75em]{Images/marginal_insertion_end.eps}~~%
\tfarsi{باقطاب اربعه}%
\thinspace\dots\thinspace
\tfarsi{پس} %
$\stackrel{\text{\tiny\signederenvoi}}{\text{\tfarsi{باشد}}}$~\SjA, %
inserted in the exterior margin by the same hand.
The penultimate word of passage~(8) on f.~22a:\thinspace 6  has 
an insertion mark `\signederenvoi' (\textit{signe-de-renvoi}) placed above it, \scl  \dots\thinspace\tfarsi{منحطّ} %
$\stackrel{\text{\tiny\signederenvoi}}{\text{\tfarsi{اربعه} }}$ %
\tfarsi{باقطاب}\thinspace\dots\thinspace .  
The first word of the marginal text also bears the same mark, \textit{supra verbum}, \scl\ \dots\thinspace
\tfarsi{پس} %
$\stackrel{\text{\tiny\signederenvoi}}{\text{\tfarsi{باشد}}}$. The marginal text ends with  `\includegraphics[height=.75em]{Images/marginal_insertion_end.eps}': this could be a calligraphic variant of the abbreviation \tfarsi{هى} (\textit{hā}\Ayn) of the Arabic word \tfarsi{إنتهى} (\intiha) meaning `it is finished', \vid\ \textcite[117]{Gacek};
\item \tfarsi{باقطاب اربعه}%
\thinspace\dots\thinspace
\tfarsi{باشد.\@  پس}~\textit{omitted} \SjB,  per homeoteleuton. The penultimate word \tfarsi{\underline{اربعه}} of passage~(8), %
before the missing text, is identical to the last word of the missing text %
\tfarsi{\underline{اربعه} منحطّ\thinspace\dots}\thinspace %
 in  passage~(9). \textbf{\scshape remark}: MS~UbC of the \ZijUlughBeg\ is also missing the same text, \viz \tfarsi{باشد\thinspace\dots\thinspace اربعه} %
spanning the length of a line between the end of line 18 (at \tfarsi{اربعه}) %
and the beginning of line 19 (at \tfarsi{منحطّ}) on folio 16r.
\end{itemize}\vspace{-\topsep}
} %
منحطّ قسمت کنیم %
\wrapfootnote{و به خارج قسمت}%
{\tfarsi{%
و~به~خارج~قسمت
}\te{\normalsize ]} \tfarsi{
و به خارج قسمت کنیم  و به خارج قسمت
}~\SjB, dittography of the second \tfarsi{و به خارج قسمت}. %
I suspect a parableptic error as the word \tfarsi{قسمت} %
appears twice on the line in close proximity: 
\dots\thinspace \tfarsi{\underline{قسمت} از}\thinspace\dots\thinspace\tfarsi{\underline{قسمت} کنیم}\thinspace\dots\thinspace.}
از جدول جيب قوس بگیریم و آن را قوس اوّل خوانیم  و جهت آن جهت عرض کوکب بود.\te{(9)}\hfill}


\tfarsi{%
پس اگر عرض و میل درجه كوكب هر دو در یک جهت باشند، قوس اوّل و میل كلّی را جمع کنیم. و اگر از 
ربع\te{\footnote{ %
\tfarsi{ربع} {\normalsize ]}  
\begin{itemize}
    \item \tfarsi{بع}\cancel{\tfarsi{ا}}\tfarsi{ر}~\SjA,  erasure
\textit{intra lineam};
    \item \tfarsi{رابع}~\SjB. The Arabic words \tfarsi{ربع} (\rub) and \tfarsi{رابع} (\rabi) are the fractional (`one-fourth') and ordinal (`fourth') forms of the number four \tfarsi{أربعة} (\arbai) respectively. The mathematical context of the passage supports the fractional meaning `one-fourth' or a `quarter'. The reading in \SjB\ could be a semantic mistake by the scribe in copying Arabic loanwords (\moarrab).
\end{itemize}\vspace{-\topsep}}}
دور زیاده شود، تمام مجموع تا نصف دور بگیریم. و اگر در جهت مختلف باشند، تفاضل میان هر دو بگیریم حاصل قوس دوم باشد و جهتش جهت مجموع یا جهت فضل باشد.\te{(10)}\hfill}

\tfarsi{%
پس جیب قوس دوم را در جیب تمام بعد از
دایرهٔ ماره باقطاب اربعه منحطّ ضرب کنیم حاصل جيب بعد کوکب باشد  و جهتش جهت قوس دوم باشد.\te{(11)}\hfill}

% \clearpage
% \textbf{Sixth chapter} \\
On the \gls{knowledge} (\marifat) of the \glspl{distance_celestial_object} (\bud\idafaconsonant\ \kawakib\ \az\ \muaddil\ \alnahar).   
\bigskip

[Given] the \gls{latitude_celestial_object} (\ard\idafaconsonant\ \kawkab) and the \gls{second_declination_degree} (\mayl\idafaconsonant\ \thani\idafavowel\ \daraji\idafavowel\ \uy), if both should be in \gls{one_direction} (\yik\ \jahat), we \gls{sum} (*\jam\ \kardan) [them]; otherwise, we should take the \gls{difference} (\tafadul). And we call that [result] the \gls{argument_of_distance} (\hissi\idafavowel\ \bud) [\lit share of the distance]. And the \gls{direction_argument_of_distance} (\jahat\idafaconsonant\ \hissi\idafavowel\ \bud) should be the \gls{direction_sum} (\jahat\idafaconsonant\ \majmu) or the \gls{direction_residue} (\jahat\idafaconsonant\ \fadla). (1)

Then, we \glslink{low_multiplication}{low multiply} (*\munhatt\idafaconsonant\ \darb\ \kardan) the \gls{Sine_argument_of_distance} (\jayb\idafaconsonant\ \hissi\idafavowel\ \bud) in the \gls{Cosine_inverse_declination_degree_celestial_object} (\jayb\idafaconsonant\ \tamam\idafaconsonant\ \mayl\idafaconsonant\ \mankus\idafaconsonant\ \daraji\idafavowel\ \kawkab). The \gls{result} (\hasil) is the \gls{Sine_distance} (\jayb\idafaconsonant\ \bud). (2)
\bigskip

In another way, we \glslink{multiplication}{multiply} (*\darb\ \kardan) the \gls{Sine_argument_of_distance} (\jayb\idafaconsonant\ \hissi\idafavowel\ \bud) in the \gls{Cosine_maximum_declination} (\jayb\idafaconsonant\ \tamam\idafaconsonant\ \mayl\idafaconsonant\ \kulli) [\ie  in the Cosine of the ecliptic obliquity] and we \glslink{division}{divide} (*\qismat\ \kardan) the \gls{result} (\hasil) over the \gls{Cosine_second_declination_degree} (\jayb\idafaconsonant\ \tamam\idafaconsonant\ \mayl\idafaconsonant\ \thani\idafavowel\ \daraji) of that \gls{celestial_object} (\kawkab). The \glslink{quotient_division}{quotient of the division} (\kharij\idafaconsonant\ \qismat) should be the \gls{Sine_distance} (\jayb\idafaconsonant\ \bud), and its \gls{direction} (\jahat)  should be the \gls{direction_argument_of_distance} (\jahat\idafaconsonant\ \hissi\idafavowel\ \bud). (3)


And since they \gls{extract} (*\darardan) [the product of the multiplication with] the \gls{Sine_argument_of_distance} (\jayb\idafaconsonant\ \hissi\idafavowel\ \bud) from the \gls{table_Cosine_maximum_declination} (\jadval\idafaconsonant\ \jayb\idafaconsonant\ \tamam\idafaconsonant\ \mayl\idafaconsonant\ \kulli) [\ie in the table of the Cosine of the ecliptic obliquity] and they \glslink{division}{divide} (*\qismat\ \kardan) the \gls{result} (\hasil) over the \gls{Cosine_second_declination_degree} (\jayb\idafaconsonant\ \tamam\idafaconsonant\ \mayl\idafaconsonant\ \thani\idafavowel\ \daraji) of that \gls{celestial_object} (\kawkab), the \glslink{quotient_division}{quotient of the division} (\kharij\idafaconsonant\ \qismat) should be the \gls{Sine_distance} (\jayb\idafaconsonant\ \bud). (4)

And if a \gls{celestial_object} (\kawkab) should have no \glslink{latitude_celestial_object}{latitude} (\ard), the \gls{declination_degree} (\mayl\idafaconsonant\ \daraji\idafavowel\ \uy) should be the \glslink{distance_celestial_object}{distance} (\bud). (5)

And if the \glslink{latitude_celestial_object}{latitude} (\ard) [of a celestial object] should exist but \glslink{degree_celestial_object}{its degree} (\daraji\idafavowel\ \uy)  should have no \glslink{declination_degree}{declination}  (\mayl), we  \glslink{low_multiplication}{low multiply} (*\munhatt\idafaconsonant\ \darb\ \kardan) the \gls{Sine_latitude} (\jayb\idafaconsonant\ \ard\idafaconsonant\ \uy) in the \gls{Cosine_maximum_declination} (\jayb\idafaconsonant\ \tamam\idafaconsonant\ \mayl\idafaconsonant\ \kulli)  [\ie  in the Cosine of the ecliptic obliquity], or we \gls{extract} (*\darardan) [the product of this multiplication] from the preceding \gls{table} (\jadval) [\ie in the table of the Cosine of the maximum declination]. The \gls{result} (\hasil) should be the \gls{Sine_distance} (\jayb\idafaconsonant\ \bud), and its \gls{direction} (\jahat) should be the \gls{direction_latitude} (\jahat\idafaconsonant\ \ard). (6) 

And if the \gls{declination_degree} (\mayl\idafaconsonant\ \daraji\idafavowel\ \uy) should be the \gls{maximum_declination} (\mayl\idafaconsonant\ \kulli) [\ie the obliquity of the ecliptic], the \gls{argument_of_distance} (\hissatalbud) itself should be the \glslink{distance_celestial_object}{distance} (\bud). (7)
\bigskip


And in another way, we \glslink{low_multiplication}{low multiply} (*\munhatt\idafaconsonant\ \darb\ \kardan) the \gls{Sine_distance_celestial_object_nearest_solstice} (\jayb\idafaconsonant\ \bud\idafaconsonant\ \daraji\idafavowel\ \kawkab\ \az\ \inqilab\idafaconsonant\ \aqrab) in the \glslink{Cosine_latitude_celestial_object}{Cosine of the latitude of the celestial object} (\jayb\idafaconsonant\ \tamam\idafaconsonant\ \ard\idafaconsonant\ \kawkab). The \gls{result} (\hasil) should be the
\glslink{Sine_distance_celestial_object_solstitial_colure}{Sine of the distance of the celestial object from the `circle passing through the four poles'} [\ie from the `solstitial colure'] (\jayb\idafaconsonant\ \bud\idafaconsonant\ \kawkab\ \az\ \guillemotleft\dayiri\idafavowel\ \marri\ \biaqtab\idafaconsonant\ \arbai\guillemotright). (8)


Then, we \glslink{low_division}{low divide} (*\munhatt\idafaconsonant\ \qismat\ \kardan) the \glslink{Sine_latitude}{Sine of the latitude of a celestial object} (\jayb\idafaconsonant\ \ard\idafaconsonant\ \kawkab) over the \glslink{Cosine_distance_solstitial_colure}{Cosine of the distance [of the celestial object] from the `circle passing through the four poles'} (\jayb\idafaconsonant\ \tamam\idafaconsonant\ \bud\ \az\ \guillemotleft\dayiri\idafavowel\ \marri\ \biaqtab\idafaconsonant\ \arbai\guillemotright). And for the \glslink{quotient_division}{quotient of the division} (\kharij\idafaconsonant\ \qismat), we should take the \gls{arc} (\qaws) from the \gls{table_of_Sine} (\jadval\idafaconsonant\ \jayb). And we call it the \gls{first_arc} (\qaws\idafaconsonant\ \avval) and its \gls{direction} (\jahat) is the \glslink{direction_latitude}{direction of the latitude of the celestial object} (\jahat\idafaconsonant\ \ard\idafaconsonant\ \kawkab). (9)


Then, if the \glslink{latitude_celestial_object}{latitude} (\ard) and the \glslink{declination_degree}{declination of the degree of the celestial object} (\mayl\idafaconsonant\ \daraji\idafavowel\ \kawkab) both should be in \gls{one_direction} (\yik\ \jahat), we \gls{sum} (*\jam\ \kardan) the \gls{first_arc} (\qaws\idafaconsonant\ \avval) and the \gls{maximum_declination} (\mayl\idafaconsonant\ \kulli) [\ie the ecliptic obliquity]. And if [the sum] \gls{exceeds} (*\ziyadi\ \shudan) \gls{one_quarter} (\rabi) [\ie is greater than 90\degree], we should take the \gls{whole_sum} (\tamam\idafaconsonant\ \majmu) up to \gls{one_half} (\nisf) [\ie up to 180\degree]. And if they should be in \gls{different_directions} (\jahat\idafaconsonant\ \mukhtalif), we should take the \gls{difference} (\tafadul) between the two; the \gls{result} (\hasil) [in both cases] should be [called] 
the \gls{second_arc} (\qaws\idafaconsonant\ \duvum) and its \gls{direction} (\jahat) should be the \gls{direction_sum} (\jahat\idafaconsonant\ \majmu) or the \gls{direction_residue} (\jahat\idafaconsonant\ \fadla). (10)

Then, we \glslink{low_multiplication}{low multiply} (*\munhatt\idafaconsonant\ \darb\ \kardan) the \gls{Sine_second_arc} (\jayb\idafaconsonant\ \qaws\idafaconsonant\ \duvum) in the \glslink{Cosine_distance_solstitial_colure}{Cosine of the distance [of a celestial object] from the `circle passing through the four poles'} (\jayb\idafaconsonant\ \tamam\idafaconsonant\ \bud\ \az\ \guillemotleft\dayiri\idafavowel\ \marri\ \biaqtab\idafaconsonant\ \arbai\guillemotright). The result should be the \glslink{Sine_distance}{Sine of the distance of the celestial object} (\jayb\idafaconsonant\ \bud\idafaconsonant\ \kawkab) and its \gls{direction} (\jahat) should be the \gls{direction_second_arc}  (\jahat\idafaconsonant\ \qaws\idafaconsonant\ \duvum). (11)

% \clearpage

\section{\Siddhantasindhu, Part II.6: text and translation}
\label{siddhantasindhu_sanskrit_english}
\setcounter{footnote}{0}
\renewcommand{\thefootnote}{[\roman{footnote}]} % for distinct numbering of Sanskrit edition (chapter vi) footnotes (starting at [i])
\setlength{\footnotesep}{\baselineskip} % for space between footnotes

{\hindifont\selectfont
\OnehalfSpacing
\large
\te{\normalmarginpar\marginnote{$\lceil$~{\footnotesize f.\thinspace 20r:\thinspace16~Kh}}}\raisebox{.5ex}{$\lceil$}{\bfseries || atha .sa.s.thaadhyaayaspa.s.takraanti{\dev{j~naa}}nam\te{\footnote{~Looking at the orthography of the other chapter-titles (in Part~II of Kh, \vid\ \S~\ref{chapter_title_comparision_persian_sanskrit}), 
\tsans{.sa.s.thaadhye spa.s.takraanti{\tsnb{ज्ञा}}nam} %
would be more consistent than \tsans{.sa.s.thaadhyaayaspa.s.takraanti{\tsnb{ज्ञा}}nam}.
Nevertheless, I maintain the reading attested in Kh in the absence of a second manuscript witness. Besides, the locative sense (\textit{adhi\-karaṇa}) of the modifier \tsans{.sa.s.tha-[.a]dhyaaya} in the \textit{tatpuruṣa} compound \tsans{.sa.s.tha-[.a]dhyaaya-spa.s.ta-kraanti-{\tsnb{ज्ञा}}nam}\,$_\text{\acrshort{nominative}-\acrshort{singular}}$ is identical to the use of the
prepositional phrase \tsans{.sa.s.tha-[.a]dhyaaye}\,$_\text{\acrshort{locative}-\acrshort{singular}}$ in 
the sentence \tsans{.sa.s.tha-[.a]dhye spa.s.ta-kraanti-{\tsnb{ज्ञा}}nam}\,$_\text{\acrshort{nominative}-\acrshort{singular}}$.}}||}
\\[.325\baselineskip]

\te{\noindent\reversemarginpar\marginnote{\hypertarget{Spass1}{[1]}}}%
khagasya baa.no.anyataraapama.h\te{\footnote{~\tsans{baa.no.anyatarapama.h}~{$\Big ]$}~\tsans{baa.nenyatarapama.h}\enskip Kh. The conjoined word \tsans{baa.nenyatarapama.h} in Kh can be meaninfully segemnted as \tsans{baa.ne}\,$_\text{\acrshort{locative}-\acrshort{singular}}$  + \tsans{anyatara-apama.h}\,$_\text{\acrshort{nominative}-\acrshort{singular}}$. However, this reading (the second noun \textit{in} the first) is contextually and semantically incoherent. The syntactic structure of the Sanskrit text mimics the syntax of the Persian text in passage~[\hyperlink{Ppass1}{1}], \vid\ \S~\ref{language_content_zij_sindhu_chapter_six}, remark~\ref{subject_fronting_passage_1} on page~\pageref{subject_fronting_passage_1}. This suggests that the emendation \tsans{baa.no}\,$_\text{\acrshort{nominative}-\acrshort{singular}}$ is better suited than \tsans{baa.ne}\,$_\text{\acrshort{locative}-\acrshort{singular}}$ as it agrees with \tsans{apama.h}\,$_\text{\acrshort{nominative}-\acrshort{singular}}$ in the subject-fronted noun phrase \tsans{khagasya-baana.h-.anyatara-apama.h puna.h}.}} puna-\\ ryadaa dvaya.m vaikadi"si sthita.m bhavet || \\
tadaa tayo.h sa.myutiranyathaantara.m\\ sphu.taapamaa.m"saakhya\te{\footnote{~\tsans{\selip maa.m"saakhya}~{$\Big ]$}~\tsans{\selip maa.m"saakhyakhya}\enskip Kh, dittography of the second \tsans{khya}.}}
ihocyate svadik || 1 || \te{\normalmarginpar\marginnote{\footnotesize \textit{vaṃśasthavila}}}
\medskip

\te{\noindent\reversemarginpar\marginnote{\hypertarget{Spass2}{[2]}}}%
sphu.taapamaa"nkasi{\dev{~nji}}nii sabhatrayadyujiivayaa || \\
nihanyate .adhariik.rtaa sphu.taapamajyakaa bhavet || 2 || \te{\normalmarginpar\marginnote{\footnotesize \textit{pramāṇikā}}}
\label{sanskrit_meter_typography_example}
\\[.65\baselineskip]

\te{\noindent\reversemarginpar\marginnote{\hypertarget{Spass3}{[3]}}}%
paramakraantiko.tijyaa sphu.takraantya"nkajiivayaa || \\
hataanya\-kraantiko.tijyaaptaa syaatspa.s.taapamajyakaa || 3 ||
\te{\normalmarginpar\marginnote{\footnotesize \textit{anuṣṭubh}}}
\medskip

\te{\noindent\reversemarginpar\marginnote{\hypertarget{Spass4}{[4]}}}%
ki.mvaa paramakraantijyaako.s.thikebhya.h\te{\footnote{~An apposite reading of the compound \tsans{parama-kraanti-jyaa-ko.s.thikebhya.h} should be \tsans{parama-kraanti-}{\normalsize <}\tsans{ko.ti}{\normalsize >}\tsans{jyaa-ko.s.thakebhya.h}. This would agree with an identical construction in passage~[\hyperlink{Spass6}{6}] construed in the same mathematical context. There is, however, no visible evidence (\eg interlinear lacunae or scribal corrections) on f.\thinspace 20r:\thinspace 20~Kh to suggest an omission. Therefore, I leave the attested reading unaltered in Sanskrit but include an emendation in my English translation.}} sphu.takraantya"nkajyayaa gu.nitaphalamutthaaya \\dvitiiyakraantiko.tijyayaa bhajellabdha.m spa.s.takraantijyaa syaat ||
\medskip

\te{\noindent\reversemarginpar\marginnote{\hypertarget{Spass5}{[5]}}}%
atha ca yadi khagasya  ba.no na syaattadaa 
tasya kraantireva spa.s.takraantirbhavet || 
}

%%%%%%%%%%%%%%%%%%%%%%%%%%%%%%%%%
\clearpage\phantomsection
%%%%%%%%%%%%%%%%%%%%%%%%%%%%%%%%%

\textbf{Now, the \gls{knowledge} (\textit{jñāna}) of the \gls{true_declination} (\textit{spaṣṭa-krānti}) in the sixth chapter.}
\\[.325\baselineskip]

\noindent\reversemarginpar\marginnote{\hypertarget{SEpass1}{[1]}}%
{[Given]} the \gls{latitude_celestial_object} (\textit{khagasya bāṇa}), [and] again, the \gls{other_declination} (\textit{anyatara-apama}) [\ie the second declination]: if indeed both should be situated in \gls{one_direction} (\textit{eka-diś}), then [we take] the \gls{sum} (\textit{saṃyuti}) of both of them; otherwise, [we take their] \gls{difference} (\textit{antara}). [The result] is known as the \gls{share_true_declination} (\textit{sphuṭa-apama-aṃśa}). Here, [it is] said to be [in] its \gls{own_direction} (\textit{sva-diś}). 1 
\medskip

\noindent\reversemarginpar\marginnote{\hypertarget{SEpass2}{[2]}}%
The \gls{Sine_curve_true_declination} (\textit{sphuṭa-apama-aṅka-siñjinī}), \gls{having_been_lowered} (\textit{adharī-kṛtā}), \glslink{multiplication}{is multiplied} (\textit{ni-hanyate}) by the \gls{day_Sine_increased_by_three_signs} (\textit{sa-bha-traya-dyujīvā}) [\ie Cosine of the first declination of the longitude increased by 90\degree]. [The result] will be the \gls{Sine_true_declination} (\textit{sphuṭa-apama-jyakā}). 2
\\[.65\baselineskip]

\noindent\reversemarginpar\marginnote{\hypertarget{SEpass3}{[3]}}%
The \gls{Cosine_maximum_declination} (\textit{parama-krānti-koṭijyā}) [\ie Cosine of the ecliptic obliquity], \glslink{multiplication}{having been multiplied} (\textit{hatā}) by the \gls{Sine_curve_true_declination} (\textit{sphuṭa-krānti-aṅka-jīvā}) [and] \glslink{division}{having been divided} (\textit{āptā}) by the \gls{Cosine_other_declination} (\textit{anya-krānti-koṭijyā}) [\ie Cosine of the second declination], should be the \gls{Sine_true_declination} (\textit{spaṣṭa-apama-jyakā}). 3 
\medskip

\noindent\reversemarginpar\marginnote{\hypertarget{SEpass4}{[4]}}%
Or, \glslink{extract}{having extracted} (\textit{utthāya}) the \glslink{product_multiplication}{product of the multiplication} (\textit{guṇita-phala}) with the \gls{Sine_curve_true_declination} (\textit{sphuṭa-krānti-aṅka-jyā}) from the \glslink{table_Cosine_maximum_declination}{tables of the <Co>sine of the maximum declination} (\textit{parama-krānti-<koṭi>jyā-koṣṭhika}s) [\ie from the tables of the Cosine of the ecliptic obliquity], [one] \glslink{division}{should divide} (\textit{bhajet}) [that product] by the \gls{Cosine_second_declination} (\textit{dvitīya-krānti-koṭijyā}). The [result] \gls{obtained} (\textit{labdha}) [\ie the quotient of the division] should be the \gls{Sine_true_declination} (\textit{spaṣṭa-krānti-jyā}).  
\medskip

\noindent\reversemarginpar\marginnote{\hypertarget{SEpass5}{[5]}}%
And now, if the \glslink{latitude_celestial_object}{latitude of a celestial object} (\textit{khagasya bāṇa}) should not exist, then its [first] \gls{declination} (\textit{krānti}) alone should be the \gls{true_declination} (\textit{spaṣṭa-krānti}). 


%%%%%%%%%%%%%%%%%%%%%%%%%%%%%%%%%
\clearpage\phantomsection
%%%%%%%%%%%%%%%%%%%%%%%%%%%%%%%%%

{\hindifont\selectfont
\OnehalfSpacing
\large
\te{\noindent\reversemarginpar\marginnote{\hypertarget{Spass6}{[6]}}}%
atha kraantiryadi na syaatpuna.h "saro bhavettadaa baa.najyaa  paramakraantiko.tijyayaa sa.mgu.nyaadha.h kuryaat || ki.mvaa paramakraantiko.tijyaako.s.thakebhyo baa.najyayaa gu.nita\-phalamutthaapayetspa.s.takraantijyaa baa.nadigbhavet || 
\medskip

\te{\noindent\reversemarginpar\marginnote{\hypertarget{Spass7}{[7]}}}%
yadi khagasya kraanti.h paramakraantitulyaa syaattadaa spa.s.takraantya"nka eva sphu.takraanti\-rbhavati || 
\bigskip

{\bfseries || atha prakaaraantare.na ||}
\bigskip

\te{\noindent\reversemarginpar\marginnote{\hypertarget{SpassA}{[α]}}}%
kadambavi.savadhruvadvayamupaiti\te{\footnote{~The word \tsans{vi.sava} (as a part of a compound) appears several times in Kh. I suspect this is an irregular (vernacular?) variant of the word \tsans{vi.suva}/\tsans{vi.suvat} that denotes the `equinox/equinoctial point' in Sanskrit astronomical literature. \Vid\ \textcite[p.\thinspace 4934a]{VacaspatyamVI} for the etymology of the word \tsans{vi.suva} (\textit{upapada tatpuṛusa})  or \tsans{vi.suvat} (\textit{matvarthīya taddhitavṛtti} or secondary nominal derivative from \tsans{vi.su}\,$_\text{\acrshort{indeclinable}}$ `in both directions'). In a lager \textit{tatpuruṣa} compound, \tsans{vi.suva}/\tsans{vi.suvat} refers to the equatorial reference frame, \eg in the genitive-compounds \tsans{vi.suva-v.rtta} `circle of the equinox' (\ie the celestial equator) and \tsans{vi.suvat-dhruva} `pole of the equinox' (\ie the celestial pole). The word \tsans{vi.sava} is not an attested form in any Sanskrit lexicon; however, it is consistently and frequently used throughout Kh. Therefore, I maintain \tsans{vi.sava}\,$_\text{\acrshort{irregular}}$ (as attested in Kh) in the \Nagari\ text but transliterate it using \tsans{vi.suva}\,$_\text{\acrshort{regular}}$ in my English translations. Both \tsans{vi.suva} and \tsans{vi.sava} have the same metrical signature ({\msf{\char"23D1}} {\msf{\char"23D1}} {\msf{\char"23D1}}).\label{vishava_deviant}}%
%
\footnote{~\tsans{\selip dvayamupaiti}~{$\Big ]$}~\tsans{\selip dvayamupeti}\enskip Kh. In the (emended) conjoined word \tsans{\selip dvayam-upaiti}, the terminal verb  
\tsans{upaiti}\,$_\text{\acrshort{present}-\acrshort{indicative}-\acrshort{singular}·\acrshort{third}}$  is derived from \tsans{upe}\,$_\text{\acrshort{compound}-\acrshort{verb}}$. A regular sandhi of the words \tsans{upa}\,$_\text{\acrshort{preverb}}$ (indeclinable \textit{upasarga}) + \tsans{eti}\,$_\text{\acrshort{present}-\acrshort{indicative}-\acrshort{singular}·\acrshort{third}}$ produces \tsans{u\underline{pai}ti} and not \tsans{upeti} (\Panini's \Astadhyayi: 6.1.89). The conjugated form \tsans{eti} derives from either √\tsans{i}\,$_{\text{\acrshort{class}\textsubscript{2}}}$ or its intensified version √\tsans{e}\,$_{\text{\acrshort{class}\textsubscript{2}}}$ (\tsans{aa}\,$_\text{\acrshort{preverb}}$ + √\tsans{i}\,$_{\text{\acrshort{class}\textsubscript{2}}}$). Any interpretation of \tsans{upeti} (in Kh), \eg  \tsans{upeti}\,$_\text{\acrshort{noun}}$, is morphosyntactically inadmissible here. The choice of using the \textit{guṇa} vowel (\textit{e}-diacritic in \tsans{pe}) instead of the \textit{vṛddhi} vowel (\textit{ai}-diacritic in \tsans{pai}) is either a grammatical error or a scribal mistake.}} %
%
v.rtta.m ca yat \\ tadaayanamudiirita.m\te{\footnote{~\tsans{\selip yanamudiirita.m}~{$\Big ]$}~\tsans{\selip yanamudiiriita.m}\enskip Kh. In the (emended) conjoined word \tsans{\selip [.a.a]yanam-udiiritam}, the terminal compound verb \tsans{ud}-√\tsans{iir}\,$_\text{\acrshort{class}\textsubscript{2}}$ takes the affix -\tsans{ita} to form \tsans{\selip udii\underline{ri}ta.m}\,$_\text{\acrshort{causative}-\acrshort{past}-\acrshort{passive}-\acrshort{participle}}$ (used as an adjective). The word \tsans{udiiriitam} (in Kh) is grammatically ill-formed; I suspect the \textit{ī}-diacritic in \tsans{rii} is a scribal hypercorrection.}} %
%
dhruvacatu.skayaata.m tathaa || \\ nabhogavi.savadvayopari\te{\footnote{~\tsans{\selip vi.sava\selip}\,$_\text{\acrshort{irregular}}$  is identical to \tsans{\selip  vi.suva\selip}\,$_\text{\acrshort{regular}}$ or \tsans{\selip vi.suvat\selip}\,$_\text{\acrshort{regular}}$, \vid\ footnote~\ref{vishava_deviant}.\label{vishava_deviant_2}}} %
%
patatsuv.rtta.m\te{\footnote{~\tsans{patatsuv.rtta.m}~{$\Big ]$}~\tsans{tatsuv.rtta.m}\enskip Kh. The third quarter \tsans{nabhoga\selip ...\selip ca yad} of the verse in passage~[α] is metrically short by one syllable (hypometric): the verse otherwise follows a regular \textit{atyaṣṭi samavṛtta} metre called \textit{pṛthvī} with seventeen-syllables per quarter. The context of the verse, and its repetition in \Nityananda's \Sarvasiddhantaraja\ (\spastakrantyadhikara:\,verse~4, \textcite{MisraTD}), suggest \tsans{tatsuv.rtta.m} should be \tsans{patatsuv.rtta.m}.}} %
%
ca yad \\
bhacakrasad.r"saahvaya.m taditi kalpayedgolavit || 1 || \te{\normalmarginpar\marginnote{\footnotesize \textit{pṛthvī}}} \label{verse_1_label_sans_example}
}

%%%%%%%%%%%%%%%%%%%%%%%%%%%%%%%%%
\clearpage\phantomsection
%%%%%%%%%%%%%%%%%%%%%%%%%%%%%%%%%

\noindent\reversemarginpar\marginnote{\hypertarget{SEpass6}{[6]}}%
Now, if the [first] \gls{declination} (\textit{krānti}) should not exist but the \gls{latitude} (\textit{śara}) should, then the \gls{Sine_the_latitude} (\textit{bāṇa-jyā}) that
\glslink{multiplication}{must be multiplied} (\textit{saṃ-guṇyā}) by the \gls{Cosine_maximum_declination} (\textit{parama-krānti-koṭijyā}) [\ie by the Cosine of the ecliptic obliquity], \gls{should_be_lowered} (\textit{adhaḥ kuryāt}). Or, [one] \glslink{extract}{may [again] extract} (\textit{utthāpayet}) the \glslink{product_multiplication}{product of the multiplication} (\textit{guṇita-phala}) with the \gls{Sine_the_latitude} (\textit{bāṇa-jyā}) from the \glslink{table_Cosine_maximum_declination}{tables of the Cosine of the maximum declination} (\textit{parama-krānti-koṭijyā-koṣṭhaka}s) [\ie from the tables of Cosine of the ecliptic obliquity]. [The result] should be the \gls{Sine_true_declination} (\textit{spaṣṭa-krānti-jyā}) in the \gls{direction_latitude} (\textit{bāṇa-diś}). \label{passge_6_sara_glossary_format_example}
\medskip

\noindent\reversemarginpar\marginnote{\hypertarget{SEpass7}{[7]}}%
If the \gls{declination_celestial_object} (\textit{khagasya krānti}) should be equal to the \gls{greatest_declination} (\textit{parama-krānti}) [\ie the obliquity of the ecliptic], then the very \gls{curve_true_declination} (\textit{spa\-ṣṭa-krānti-aṅka}) becomes the \gls{true_declination} (\textit{sphuṭa-krānti}).
\bigskip

\textbf{Now, in another way.}
\bigskip

\noindent\reversemarginpar\marginnote{\hypertarget{SEpassA}{[α]}}%
And what \gls{circle} (\textit{vṛtta}) reaches both the \gls{ecliptic_pole} (\textit{kadamba}) and the \gls{celestial_pole} (\textit{viṣuvat}\,$_\text{\acrshort{regular}}$-\textit{dhruva}), that has been stated to be the \glslink{circle_solstice}{solstitial [colure]} (\textit{āyana}[-\textit{vṛtta}]), and also, the \glslink{circle_four_poles}{[circle] passing through the four poles} (\textit{dhruva-catuṣka-yāta}[-\textit{vṛtta}]). And passing over a \gls{celestial_object} (\textit{nabhoga}) and the \gls{pair_equinoctial_points} (\textit{visuvat}\,$_\text{\acrshort{regular}}$-\textit{dvaya}), what [circle] is \gls{well_rounded} (\textit{su-vṛtta}), that the \gls{knower_spheres} (\textit{gola-vid}) should consider as
the \glslink{circle_congruent_ecliptic}{[circle] congruent to the ecliptic} (\textit{bhacakra-sadṛśa}[-\textit{vṛtta}]) by name. 1 
\label{verse_1_label_eng_example}


%%%%%%%%%%%%%%%%%%%%%%%%%%%%%%%%%
\clearpage\phantomsection
%%%%%%%%%%%%%%%%%%%%%%%%%%%%%%%%%

{\hindifont\selectfont
\OnehalfSpacing
\large
\te{\noindent\reversemarginpar\marginnote{\hypertarget{SpassB}{[β]}}}%
vi.savav.rttabhav.rttasad.r"sayorvivara\raisebox{.5ex}{$\lceil$}ga.m\te{\footnote{~\tsans{vi.savav.rtta\selip}\,$_\text{\acrshort{irregular}}$ is identical to \tsans{vi.suvav.rtta\selip}\,$_\text{\acrshort{regular}}$, \vid\ footnote~\ref{vishava_deviant}.\label{vishava_deviant_3}}} %
%
dhanuraayanav.rttajam || \te{\normalmarginpar\marginnote{$\lceil$~\footnotesize f.\thinspace 20v:\thinspace1~Kh}}\\
bhavati yatkathita.h sa parasphu.taapama iti dyucarasya ca samprati || 2 ||\label{folio_break_sanskrit_example}
\te{\normalmarginpar\marginnote{\footnotesize \textit{drutavilambita}}}
\medskip

\te{\noindent\reversemarginpar\marginnote{\hypertarget{SpassC}{[γ]}}}%
bhavanacakrabhacakrasad.rk.sayorvivaraga.m dhanuraayanav.rttajam ||\\
bhavati yatsa pare.surihodito vi.savapaatayuge\te{\footnote{~\tsans{vi.sava\selip}\,$_\text{\acrshort{irregular}}$ is identical to \tsans{vi.suva\selip}\,$_\text{\acrshort{regular}}$, \vid\ footnote~\ref{vishava_deviant}\label{vishava_deviant_4}.}} %
%
sati kalpite || 3 || \te{\normalmarginpar\marginnote{\footnotesize \textit{drutavilambita}}}
\medskip

\te{\noindent\reversemarginpar\marginnote{\hypertarget{SpassD}{[δ]}}}%
vi.savannabhogamadhye\te{\footnote{~The word \tsans{vi.savat} is an attested secondary nominal derivative (from \tsans{vi.sa}\,$_\text{\acrshort{noun}}$ `poison') meaning `poisonous'. However, in the \textit{tatpuruṣa} compound \tsans{vi.savat-nabhoga-madhye}, I believe \tsans{vi.savat\selip}\,$_\text{\acrshort{irregular}}$, like \tsans{vi.sava}\,$_\text{\acrshort{irregular}}$, is identical to \tsans{\selip  vi.suvat\selip}\,$_\text{\acrshort{regular}}$, \vid\ footnote~\ref{vishava_deviant}.\label{vishava_deviant_5}}} %
%
yatkoda.n.da.m bhav.rttasad.r"sasya ||\\
{\dev{j~ne}}ya.h sad.rgbhujo .asau bhaayanavivare sad.rkko.ti.h || 4 || 
\te{\normalmarginpar\marginnote{\footnotesize \textit{āryā}}}
\\[.65\baselineskip]

\te{\noindent\reversemarginpar\marginnote{\hypertarget{Spass8}{[8]}}}%
khagasya ko.tisi{\dev{~nji}}nii svabaa.nako.tijiivayaa ||\\
hataa .adhariik.rtaa\te{\footnote{~\tsans{hataa .adhariik.rtaa}~{$\Big ]$}~\tsans{hataadhariinvataa}\enskip Kh. The conjoined words \tsans{hataadhariinvataa} in Kh can be segmented as 
\tsans{hataa}\,$_\text{\acrshort{past}-\acrshort{passive}-\acrshort{participle}}$ (from √\thinspace\tsans{han}\,$_{\text{\acrshort{class}\textsubscript{2}}}$) +
\tsans{.adhariinvataa}; however, the compound \tsans{.adhariinvataa} is etymologically defective by the rules of Pāṇinian grammar. The adverbial \textit{CvI}-suffixation to \tsans{adhara}\,$_\text{\acrshort{noun}-\acrshort{stem}}$ (making it \tsans{adharii}\,$_\text{\acrshort{preverb}}$) can only occur with terminal verbs √\thinspace\tsans{k.r}\,$_{\text{\acrshort{class}}_8}$, √\thinspace\tsans{bhuu}\,$_{\text{\acrshort{class}}_1}$, and √\thinspace\tsans{as}\,$_{\text{\acrshort{class}\textsubscript{2}}}$ when forming factitive compound verbs like √\thinspace\tsans{adharii-k.r} (\Panini's \Astadhyayi: 5.4.50), \vid\ \textcite[\textbf{1094}, p.\thinspace 357]{Whitney}. \protect\\[0.5\baselineskip]
A regular sandhi of the words \tsans{hataa}\,$_\text{\acrshort{past}-\acrshort{passive}-\acrshort{participle}}$ 
+ \tsans{.adhariik.rtaa}\,$_\text{\acrshort{past}-\acrshort{passive}-\acrshort{participle}}$ (from √\thinspace\tsans{adharii-k.r}\,$_{\text{\acrshort{class}}_8}$) generates \tsans{hataadharii\underline{k.r}taa} that is morphologically correct and contextually apposite. Also, \tsans{.adhariik.r}\,$_\text{\acrshort{class}\textsubscript{8}}$ is variously attested in this chapter, as well as in \Nityananda's  \Sarvasiddhantaraja\ (\textcite{MisraTD}), in the same mathematical context. \Vid\ glossary entry: \glslink{lowering}{lowering}.}} %
%
bhavetsad.rk.sako.tisi{\dev{~nji}}nii || 5 ||
\te{\normalmarginpar\marginnote{\footnotesize \textit{pramāṇikā}}}
\medskip

\te{\noindent\reversemarginpar\marginnote{\hypertarget{Spass9}{[9]}}}%
taddhanurnavatita"scyuta.m yadaa jaayate sad.r"sabaahusa.m{\dev{~nja}}kam ||\\
yaa nabhogavi"sikhasya si{\dev{~nji}}nii bhaajitaa .adharasad.rk.sadorjyayaa || 6 ||
\te{\normalmarginpar\marginnote{\footnotesize \textit{rathoddhatā}}}\\
taddhanu.h para"saraahvayo bhavedvaa---
}


%%%%%%%%%%%%%%%%%%%%%%%%%%%%%%%%%
\clearpage\phantomsection
%%%%%%%%%%%%%%%%%%%%%%%%%%%%%%%%%

\noindent\reversemarginpar\marginnote{\hypertarget{SEpassB}{[β]}}%
What \gls{arc} (\textit{dhanus}) produced on the \gls{circle_solstice} (\textit{āyana-vṛtta}) becomes situated in the  \gls{difference} (\textit{vivara}) between the \gls{celestial_equator} (\textit{viṣuva}\,$_\text{\acrshort{regular}}$-\textit{vṛtta}) and the \glslink{circle_congruent_ecliptic}{[circle] congruent to the ecliptic} (\textit{bhavṛtta-sadṛśa}[-\textit{vṛtta}]), that is the stated [arc of] \gls{maximum_true_declination} (\textit{para-sphuṭa-apama}) of the \gls{celestial_object} (\textit{dyucara}) just at that present moment. 2 
\medskip

\noindent\reversemarginpar\marginnote{\hypertarget{SEpassC}{[γ]}}%
What \gls{arc} (\textit{dhanus}) produced on the \gls{circle_solstice} (\textit{āyana-vṛtta}) becomes situated in the \gls{difference} (\textit{vivara}) between the \gls{ecliptic} (\textit{bhavana-cakra}) and the \glslink{circle_congruent_ecliptic}{[circle] congruent to the ecliptic} (\textit{bhacakra-sadṛkṣa}[-\textit{vṛtta}]), in this case, that is the declared [arc of] \gls{maximum_latitude} (\textit{para-iṣu}) when the \glslink{conjunction_equinox_node}{conjunction of the equinoctial point and the node of the orbit [of the celestial object]} (\textit{viṣuva}\,$_\text{\acrshort{regular}}$-\textit{pāta-yuga}) has been supposed. 3 
\medskip

\noindent\reversemarginpar\marginnote{\hypertarget{SEpassD}{[δ]}}%
What \gls{arc} (\textit{kodaṇḍa}) of the  \glslink{circle_congruent_ecliptic}{[circle] congruent to the ecliptic} (\textit{bhavṛtta-sadṛśa}[-\textit{vṛtta}]) is between the \gls{equinoctial_point} (\textit{viṣuvat}\,$_\text{\acrshort{regular}}$) and the \gls{celestial_object} (\textit{nabhoga}), that [arc] should be known as the \gls{congruent_bhuja} (\textit{sadṛś-bhujā}); [and what is] between the \gls{celestial_object} (\textit{bha}) and the \gls{circle_solstice} (\textit{āyana}[-\textit{vṛtta}]), [that should be known as] the \gls{congruent_koti} (\textit{sadṛś-koṭi}) [\ie complement of \textit{sadṛś-bhujā}].~4
\\[.65\baselineskip]

\noindent\reversemarginpar\marginnote{\hypertarget{SEpass8}{[8]}}%
The \gls{Sine_koti_celestial_object} (\textit{khagasya koṭi-siñjinī}), \glslink{multiplication}{having been multiplied} (\textit{hatā}) by the \gls{Cosine_latitude} (\textit{sva-bāṇa-koṭijīvā}) [and] \gls{having_been_lowered} (\textit{adharī-kṛtā}), should be the \gls{Sine_congruent_koti} (\textit{sadrkṣa-koṭi-siñjinī}) [\ie Sine of the complement of the \textit{sadṛś-bhujā}]. 5 
\medskip

\noindent\reversemarginpar\marginnote{\hypertarget{SEpass9}{[9]}}%
When [the measure of] its \gls{arc} (\textit{dhanus}), \gls{reducing_from_ninety} [degrees] (\textit{navatitaś-cyuta}), is determined, [it] has the name \glslink{congruent_bhuja}{congruent arc} (\textit{sadṛśa-bāhu}). Or, what is the \gls{Sine_latitude_celestial_object} (\textit{nabhoga-viśikhasya siñjinī}), \glslink{division}{having been divided} (\textit{bhājitā}) by the \gls{lowered_Sine_congruent_base} (\textit{adhara-sadrkṣa-dorjyā}), 6\dots\\ \dots its \gls{arc} (\textit{dhanus}) should be [called] the \gls{maximum_latitude} (\textit{para-śara}) by name.---

%%%%%%%%%%%%%%%%%%%%%%%%%%%%%%%%%
\clearpage\phantomsection
%%%%%%%%%%%%%%%%%%%%%%%%%%%%%%%%%

{\hindifont\selectfont
\OnehalfSpacing
\large
\te{\noindent\reversemarginpar\marginnote{\hypertarget{Spass10}{[10]}}}%
---pare.suparamaapamaakhyayo.h ||\\
sa.myutirviyutirasti ca kramaadgolabaa.nasamabhinnadiktayaa || 7 ||
\te{\normalmarginpar\marginnote{\footnotesize \textit{rathoddhatā}}}\\
sa grahasya paramasphu.taapamo jaayate yutiviyogadiksthita.h ||\\
evamabhranavato.adhiko yadaa khaa.s.tabhuu 180 parimitervi"sodhita.h || 8 || 
\te{\normalmarginpar\marginnote{\footnotesize \textit{rathoddhatā}}}
\medskip

\te{\noindent\reversemarginpar\marginnote{\hypertarget{Spass11}{[11]}}}%
parasphu.takraantibhavajyakaa hataa  sad.rk.sabaahujyakayaa .adhariik.rtaa ||  \\
tadiiya caapa.m bhavati sphu.taapamo digasya sa.myogaviyogadiksamaa || 9 ||
\te{\normalmarginpar\marginnote{\footnotesize \textit{vaṃśasthavila}}}% 
}

%%%%%%%%%%%%%%%%%%%%%%%%%%%%%%%%%
\clearpage\phantomsection
%%%%%%%%%%%%%%%%%%%%%%%%%%%%%%%%%

\noindent\reversemarginpar\marginnote{\hypertarget{SEpass10}{[10]}}%
---There is the \gls{sum} (\textit{saṃyuti}) or the \gls{difference} (\textit{viyuti}) of the two [quantities] known as the \gls{maximum_latitude} (\textit{para-iṣu}) and the \gls{greatest_declination} (\textit{parama-apama}) [\ie the obliquity of the ecliptic] with the  \gls{latitude} (\textit{bāṇa}) and the \gls{celestial_hemisphere} (\textit{gola}) [\ie the declination of the celestial object] in the \gls{same_different_directions} (\textit{sama-bhinna-diś}) respectively. 7\\ That [result], being situated in the \gls{conjunction_disjunction_direction} (\textit{yuti-viyoga-diś}), becomes the \gls{maximum_true_declination_celestial_object} (\textit{grahasya parama-sphuṭa-apama}). Thus, when [its measure is] \gls{greater} (\textit{adhika}) than \gls{ninety} [degrees] (\textit{abhra-nava}), [it is] \glslink{subtraction}{made to be subtracted} (\textit{viśodhita}) from a measure of \gls{one_hundred_eighty} [degrees] (\textit{kha-aṣṭa-bhū}). 8 \medskip

\noindent\reversemarginpar\marginnote{\hypertarget{SEpass11}{[11]}}%
The \gls{Sine_maximum_true_declination} (\textit{para-sphuṭa-krānti-bhava-jyakā}), \glslink{multiplication}{having been multiplied} (\textit{hatā}) by the \gls{Sine_congruent_base} (\textit{sadṛkṣa-bāhu-jyakā}) [and] \gls{having_been_lowered} (\textit{adharī-kṛtā}), its \gls{arc} (\textit{cāpa}) becomes the \gls{true_declination} (\textit{sphuṭa-apama}). Its \gls{direction} (\textit{diś}) is the \gls{same} (\textit{sama}) as the \gls{conjunction_disjunction_direction} (\textit{saṃyoga-viyoga-diś}). 9

%%%%%%%%%%%%%%%%%%%%%%%%%%%%%%%%%%%%%%%%%%%%%%%%%%%%%%%%%%%%%%%%%
\clearpage
% {\hindifont\selectfont
\OnehalfSpacing
\large
\setcounter{footnote}{0}
\renewcommand{\thefootnote}{[\roman{footnote}]} % for distinct numbering of Sanskrit edition (chapter vi) footnotes (starting at [i])
\setlength{\footnotesep}{\baselineskip} % for space between footnotes
\normalmarginpar % for direction of margin notes

\te{\marginnote{\raisebox{1ex}{$\ulcorner$}\kern-0.25ex\footnotesize f.~20r:\thinspace16~Kh}}{\raisebox{1.5ex}{$\ulcorner$}\kern-0.5ex}{\bfseries || atha .sa.s.thaadhyaayaspa.s.takraanti{\dev{j~naa}}nam\te{\footnote{\,Looking at the orthography of the other chapter-titles (in Part~II of Kh, \vid\ \S~\ref{chapter_title_comparision_persian_sanskrit}), 
\tsans{.sa.s.thaadhye spa.s.takraanti{\tsnb{ज्ञा}}nam} %
would be more consistent than \tsans{.sa.s.thaadhyaayaspa.s.takraanti{\tsnb{ज्ञा}}nam}.
Nevertheless, I maintain the reading attested in Kh in the absence of a second manuscript witness. Besides, the locative sense (\textit{adhi\-karaṇa}) of the modifier \tsans{.sa.s.tha-[.a]dhyaaya} in the \textit{tatpuruṣa} compound \tsans{.sa.s.tha-[.a]dhyaaya-spa.s.ta-kraanti-{\tsnb{ज्ञा}}nam}\,$_\text{\acrshort{nominative}-\acrshort{singular}}$ is identical to the use of the
prepositional phrase \tsans{.sa.s.tha-[.a]dhyaaye}\,$_\text{\acrshort{locative}-\acrshort{singular}}$ in 
the sentence \tsans{.sa.s.tha-[.a]dhye spa.s.ta-kraanti-{\tsnb{ज्ञा}}nam}\,$_\text{\acrshort{nominative}-\acrshort{singular}}$.}}||}
\bigskip

khagasya baa.no.anyataraapama.h\te{\footnote{\,\tsans{baa.no.anyatarapama.h} {$\normalsize ]$} \tsans{baa.nenyatarapama.h}\enskip Kh. The conjoined word \tsans{baa.nenyatarapama.h} in Kh can be meaninfully segemnted as \tsans{baa.ne}\,$_\text{\acrshort{locative}-\acrshort{singular}}$  + \tsans{anyatara-apama.h}\,$_\text{\acrshort{nominative}-\acrshort{singular}}$. However, this reading (the second noun \textit{in} the first) is contextually and semantically incoherent. The syntactic structure of the Sanskrit text mimics the syntax of the Persian text in passage~(1), \vid\ discussions on <INSERT PAGE \& SECTION REF>. This suggests that the emendation \tsans{baa.no}\,$_\text{\acrshort{nominative}-\acrshort{singular}}$ is better suited than \tsans{baa.ne}\,$_\text{\acrshort{locative}-\acrshort{singular}}$ as it agrees with \tsans{apama.h}\,$_\text{\acrshort{nominative}-\acrshort{singular}}$ in the subject-fronted noun phrase \tsans{khagasya-baana.h-.anyatara-apama.h puna.h}.}} puna-\\ ryadaa dvaya.m vaikadi"si sthita.m bhavet || \\
tadaa tayo.h sa.myutiranyathaantara.m\\ sphu.taapamaa.m"saakaakhya\te{\footnote{\,\tsans{\selip maa.m"saakhya} {\normalsize ]} \tsans{\selip maa.m"saakhyakhya}\enskip Kh, dittography of the second \tsans{khya}.}}
ihocyate svadik || 1 || \marginnote{\te{(1)}} 

sphu.taapamaa"nkasi{\dev{~nji}}nii sabhatrayadyujiivayaa || \\
nihanyate .adhariik.rtaa sphu.taapamajyakaa bhavet || 2 || \marginnote{\te{(2)}}%
\bigskip

paramakraantiko.tijyaa sphu.takraantya"nkajiivayaa || \\
hataanya\-kraantiko.tijyaaptaa syaatspa.s.taapamajyakaa || 3 || \marginnote{\te{(3)}}

ki.mvaa paramakraantijyaako.s.thikebhya.h\te{\footnote{\,An apposite reading of the compound \tsans{parama-kraanti-jyaa-ko.s.thikebhya.h} should be \tsans{parama-kraanti-}{\normalsize <}\tsans{ko.ti}{\normalsize >}\tsans{jyaa-ko.s.thakebhya.h}. This would agree with an identical construction in passage~(6) construed in the same mathematical context. There is, however, no visible evidence (\eg interlinear lacunae or scribal corrections) on f.~20r:\thinspace 20~Kh to suggest an omission. Therefore, I leave the attested reading unaltered in Sanskrit but include an emendation in my English translation.}} sphu.takraantya"nkajyayaa gu.nita-\\phalamutthaaya dvitiiyakraantiko.tijyayaa bhajellabdha.m spa.s.takraantijyaa syaat || \marginnote{\te{(4)}}

atha ca yadi khagasya  ba.no na syaattadaa \\
tasya kraantireva spa.s.takraantirbhavet || \marginnote{\te{(5)}} 
\newpage

atha kraantiryadi na syaatpuna.h "saro bhavettadaa baa.najyaa \\ paramakraantiko.tijyayaa sa.mgu.nyaadha.h kuryaat || \\
ki.mvaa paramakraantiko.tijyaako.s.thakebhyo baa.najyayaa \\ gu.nitaphalamutthaapayetspa.s.takraantijyaa baa.nadigbhavet || \marginnote{\te{(6)}}

yadi khagasya kraanti.h paramakraantitulyaa syaattadaa\\ spa.s.takraantya"nka eva sphu.takraantirbhavati || \marginnote{\te{(7)}}
\bigskip

{\bfseries || atha prakaaraantare.na ||}
\medskip

kadambavi.savadhruvadvayamupaiti\te{\footnote{\,The word \tsans{vi.sava} (as a part of a compound) appears several times in Kh. I suspect this is an irregular (vernacular?) variant of the word \tsans{vi.suva}/\tsans{vi.suvat} that denotes the `equinox/equinoctial point' in Sanskrit astronomical literature. \Vid\ \textcite[p.~4934a]{VacaspatyamVI} for the etymology of the word \tsans{vi.suva} (\textit{upapada tatpuṛusa})  or \tsans{vi.suvat} (\textit{matvarthīya taddhitavṛtti} or secondary nominal derivative from \tsans{vi.su}\,$_\text{\acrshort{indeclinable}}$ `in both directions'). In a lager \textit{tatpuruṣa} compound, \tsans{vi.suva}/\tsans{vi.suvat} refers to the equatorial reference frame, \eg in the genitive-compounds \tsans{vi.suva-v.rtta} `circle of the equinox' (\ie the celestial equator) and \tsans{vi.suvat-dhruva} `pole of the equinox' (\ie the celestial pole). The word \tsans{vi.sava} is not an attested form in any Sanskrit lexicon; however, it is consistently and frequently used throughout Kh. Therefore, I maintain \tsans{vi.sava}\,$_\text{\acrshort{irregular}}$ (as attested in Kh) in the \Nagari\ text but transliterate it using \tsans{vi.suva}\,$_\text{\acrshort{regular}}$ in my English translations. Both \tsans{vi.suva} and \tsans{vi.sava} have the same metrical signature ({\msf{\char"23D1}} {\msf{\char"23D1}} {\msf{\char"23D1}}).\label{vishava_deviant}}%
%
\footnote{\,\tsans{\selip dvayamupaiti} {\normalsize ]} \tsans{\selip dvayamupeti}\enskip Kh. In the (emended) conjoined word \tsans{\selip dvayam-upaiti}, the terminal verb  
\tsans{upaiti}\,$_\text{\acrshort{present}-\acrshort{indicative}-\acrshort{singular}·\acrshort{third}}$  is derived from \tsans{upe}\,$_\text{\acrshort{compound}-\acrshort{verb}}$. A regular sandhi of the words \tsans{upa}\,$_\text{\acrshort{preverb}}$ (indeclinable \textit{upasarga}) + \tsans{eti}\,$_\text{\acrshort{present}-\acrshort{indicative}-\acrshort{singular}·\acrshort{third}}$ produces \tsans{u\underline{pai}ti} and not \tsans{upeti} (\Panini's \Astadhyayi: 6.1.89). The conjugated form \tsans{eti} derives from either √\tsans{i}\,$_{\text{\acrshort{class}\textsubscript{2}}}$ or its intensified version √\tsans{e}\,$_{\text{\acrshort{class}\textsubscript{2}}}$ (\tsans{aa}\,$_\text{\acrshort{preverb}}$ + √\tsans{i}\,$_{\text{\acrshort{class}\textsubscript{2}}}$). Any interpretation of \tsans{upeti} (in Kh), \eg  \tsans{upeti}\,$_\text{\acrshort{noun}}$, is morphosyntactically inadmissible here. The choice of using the \textit{guṇa} vowel (\textit{e}-diacritic in \tsans{pe}) instead of the \textit{vṛddhi} vowel (\textit{ai}-diacritic in \tsans{pai}) is either a grammatical error or a scribal mistake.}} %
%
v.rtta.m ca yat \\ tadaayanamudiirita.m\te{\footnote{\,\tsans{\selip yanamudiirita.m} {\normalsize ]} \tsans{\selip yanamudiiriita.m}\enskip Kh. In the (emended) conjoined word \tsans{\selip [.a.a]yanam-udiiritam}, the terminal compound verb \tsans{ud}-√\tsans{iir}\,$_\text{\acrshort{class}\textsubscript{2}}$ takes the affix -\tsans{ita} to form \tsans{\selip udii\underline{ri}ta.m}\,$_\text{\acrshort{causative}-\acrshort{past}-\acrshort{passive}-\acrshort{participle}}$ (used as an adjective). The word \tsans{udiiriitam} (in Kh) is grammatically ill-formed; I suspect the \textit{ī}-diacritic in \tsans{rii} is a scribal hypercorrection.}} %
%
dhruvacatu.skayaata.m tathaa || \\ nabhogavi.savadvayopari\te{\footnote{\,\tsans{\selip vi.sava\selip}\,$_\text{\acrshort{irregular}}$  is identical to \tsans{\selip  vi.suva\selip}\,$_\text{\acrshort{regular}}$ or \tsans{\selip vi.suvat\selip}\,$_\text{\acrshort{regular}}$, \vid\ footnote~\ref{vishava_deviant}.}} %
%
patatsuv.rtta.m\te{\footnote{\,\tsans{patatsuv.rtta.m} {\normalsize ]} \tsans{tatsuv.rtta.m}\enskip Kh. The third quarter \tsans{nabhoga\selip ...\selip ca yad} of the verse in passage~(α) is metrically short by one syllable (hypometric): the verse otherwise follows a regular \textit{atyaṣṭi samavṛtta} meter called \textit{pṛthvī} with seventeen-syllables per quarter. The context of the verse, and its repetition in \Nityananda's \Sarvasiddhantaraja\ (\spastakrantyadhikara:\,verse~4, \textcite{MisraTD}), suggest \tsans{tatsuv.rtta.m} should be \tsans{patatsuv.rtta.m}.}} %
%
ca yad \\
bhacakrasad.r"saahvaya.m taditi kalpayedgolavit || 1 || \marginnote{\te{(α)}}


vi.savav.rttabhav.rttasad.r"sayorvivara{\raisebox{1.5ex}{$\ulcorner$}\kern-0.9ex}ga.m\te{\footnote{\,\tsans{vi.savav.rtta\selip}\,$_\text{\acrshort{irregular}}$ is identical to \tsans{vi.suvav.rtta\selip}\,$_\text{\acrshort{regular}}$, \vid\ footnote~\ref{vishava_deviant}.}} %
%
dhanuraayanav.rttajam || \te{\marginnote{\raisebox{1ex}{$\ulcorner$}\kern-0.25ex\footnotesize f.~20v:\thinspace1~Kh}}\\
bhavati yatkathita.h sa parasphu.taapama iti dyucarasya ca samprati || 2 || \marginnote{\te{(β)}}

bhavanacakrabhacakrasad.rk.sayorvivaraga.m dhanuraayanav.rttajam ||\\
bhavati yatsa pare.surihodito vi.savapaatayuge\te{\footnote{\,\tsans{vi.sava\selip}\,$_\text{\acrshort{irregular}}$ is identical to \tsans{vi.suva\selip}\,$_\text{\acrshort{regular}}$, \vid\ footnote~\ref{vishava_deviant}.}} %
%
sati kalpite || 3 || \marginnote{\te{(γ)}}

vi.savannabhogamadhye\te{\footnote{\,The word \tsans{vi.savat} is an attested secondary nominal derivative (from \tsans{vi.sa}\,$_\text{\acrshort{noun}}$ `poison') meaning `poisonous'. However, in the \textit{tatpuruṣa} compound \tsans{vi.savat-nabhoga-madhye}, I believe \tsans{vi.savat\selip}\,$_\text{\acrshort{irregular}}$, like \tsans{vi.sava}\,$_\text{\acrshort{irregular}}$, is identical to \tsans{\selip  vi.suvat\selip}\,$_\text{\acrshort{regular}}$, \vid\ footnote~\ref{vishava_deviant}.}} %
%
yatkoda.n.da.m bhav.rttasad.r"sasya ||\\
{\dev{j~ne}}ya.h sad.rgbhujo .asau bhaayanavivare sad.rkko.ti.h || 4 || \marginnote{\te{(δ)}}%
\bigskip

khagasya ko.tisi{\dev{~nji}}nii svabaa.nako.tijiivayaa ||\\
hataa .adhariik.rtaa\te{\footnote{\,\tsans{hataa .adhariik.rtaa} {\normalsize ]} \tsans{hataadhariinvataa}\enskip Kh. The conjoined words \tsans{hataadhariinvataa} in Kh can be segmented as 
\tsans{hataa}\,$_\text{\acrshort{past}-\acrshort{passive}-\acrshort{participle}}$ (from √\thinspace\tsans{han}\,$_{\text{\acrshort{class}\textsubscript{2}}}$) +
\tsans{.adhariinvataa}; however, the compound \tsans{.adhariinvataa} is etymologically defective by the rules of Pāṇinian grammar. The adverbial \textit{CvI}-suffixation to \tsans{adhara}\,$_\text{\acrshort{noun}-\acrshort{stem}}$ (making it \tsans{adharii}\,$_\text{\acrshort{preverb}}$) can only occur with terminal verbs √\thinspace\tsans{k.r}\,$_{\text{\acrshort{class}}_8}$, √\thinspace\tsans{bhuu}\,$_{\text{\acrshort{class}}_1}$, and √\thinspace\tsans{as}\,$_{\text{\acrshort{class}\textsubscript{2}}}$ when forming factitive compound verbs like √\thinspace\tsans{adharii-k.r} (\Panini's \Astadhyayi: 5.4.50), \vid\ \textcite[\textbf{1094}, p.~357]{Whitney}. \protect\\[0.5\baselineskip]
A regular sandhi of the words \tsans{hataa}\,$_\text{\acrshort{past}-\acrshort{passive}-\acrshort{participle}}$ 
+ \tsans{.adhariik.rtaa}\,$_\text{\acrshort{past}-\acrshort{passive}-\acrshort{participle}}$ (from √\thinspace\tsans{adharii-k.r}\,$_{\text{\acrshort{class}}_8}$) generates \tsans{hataadharii\underline{k.r}taa} that is morphologically correct and contextually apposite. Also, \tsans{.adhariik.r}\,$_\text{\acrshort{class}\textsubscript{8}}$ is variously attested in this chapter, as well as in \Nityananda's  \Sarvasiddhantaraja\ (\textcite{MisraTD}), in the same mathematical context. \Vid\ glossary entry: \gls{lowering}.}} %
%
bhavetsad.rk.sako.tisi{\dev{~nji}}nii || 5 ||\marginnote{\te{(8)}}


taddhanurnavatita"scyuta.m yadaa jaayate sad.r"sabaahusa.m{\dev{~nja}}kam ||\\
yaa nabhogavi"sikhasya si{\dev{~nji}}nii bhaajitaa .adharasad.rk.sadorjyayaa || 6 ||

\marginnote{\te{\textsuperscript{\tsans{...\selip hvayo bhavedvaa}\,$^\text{¶}$}(9)}}
taddhanu.h para"saraahvayo bhavedvaa pare.suparamaapamaakhyayo.h ||\\
sa.myutirviyutirasti ca kramaadgolabaa.nasamabhinnadiktayaa || 7 ||

sa grahasya paramasphu.taapamo jaayate yutiviyogadiksthita.h ||\\
evamabhranavato.adhiko yadaa khaa.s.tabhuu 180 parimitervi"sodhita.h || 8 || \marginnote{\te{(10)}}

parasphu.takraantibhavajyakaa hataa  sad.rk.sabaahujyakayaa .adhariik.rtaa ||  \\
tadiiya caapa.m bhavati sphu.taapamo digasya sa.myogaviyogadiksamaa || 9 || \marginnote{\te{(11)}}% 
}
% \clearpage
% \textbf{Now, the \gls{knowledge} (\textit{jñāna}) of the \gls{true_declination} (\textit{spaṣṭa-krānti}) in the sixth chapter.}\medskip

[Given] the \gls{latitude_celestial_object} (\textit{khagasya bāṇa}), [and] again, the \gls{other_declination} (\textit{anyatara-apama}) [\ie the second declination]: if indeed both should be situated in \gls{one_direction} (\textit{eka-diś}), then [we take] the \gls{sum} (\textit{saṃyuti}) of both of them; otherwise, [we take their] \gls{difference} (\textit{antara}). [The result] is known as the \glslink{curve_true_declination}{share of the true declination} (\textit{sphuṭa-apama-aṃśa}). Here, [it is] said to be [in] its \gls{own_direction} (\textit{sva-diś}). 1 \marginnote{(1)} 

The \gls{Sine_curve_true_declination} (\textit{sphuṭa-apama-aṅka-siñjinī}), \glslink{lowering}{having been lowered} (\textit{adharī-kṛtā}), \glslink{multiplication}{is multiplied} (\textit{ni-hanyate}) by the \gls{day_Sine_increased_by_three_signs} (\textit{sa-bha-traya-dyujīvā}) [\ie Cosine of the first declination of the longitude increased by 90\degree]. [The result] should be the \gls{Sine_true_declination} (\textit{sphuṭa-apama-jyakā}). 2 \marginnote{(2)}
\bigskip

The \gls{Cosine_maximum_declination} (\textit{parama-krānti-koṭijyā}) [\ie Cosine of the ecliptic obliquity], \glslink{multiplication}{having been multiplied} (\textit{hatā}) by the \gls{Sine_curve_true_declination} (\textit{sphuṭa-krānti-aṅka-jīvā}) [and] \glslink{division}{having been divided} (\textit{āptā}) by the \gls{Cosine_other_declination} (\textit{anya-krānti-koṭijyā}) [\ie Cosine of the second declination], should be the \gls{Sine_true_declination} (\textit{spaṣṭa-apama-jyakā}). 3 \marginnote{(3)} 


Or, \glslink{extract}{having extracted} (\textit{utthāya}) the \glslink{product_multiplication}{product of the multiplication} (\textit{guṇita-phala}) with the \gls{Sine_curve_true_declination} (\textit{sphuṭa-krānti-aṅka-jyā}) from the \glslink{table_Cosine_maximum_declination}{tables of the <Co>sine of the maximum declination} (\textit{parama-krānti-<koṭi>jyā-koṣṭhika}s) [\ie from the tables of the Cosine of the ecliptic obliquity], [one] \glslink{division}{should divide} (\textit{bhajet}) [that product] by the \glslink{Cosine_second_declination_degree}{Cosine of the second declination} (\textit{dvitīya-krānti-koṭijyā}). The [result] \gls{obtained} (\textit{labdha}) [\ie the quotient of the division] should be the \gls{Sine_true_declination} (\textit{spaṣṭa-krānti-jyā}).  \marginnote{(4)}

And now, if the \glslink{latitude_celestial_object}{latitude of a celestial object} (\textit{khagasya bāṇa}) should not exist, then its [first] \glslink{declination_degree}{declination} (\textit{krānti}) alone should be the \gls{true_declination} (\textit{spaṣṭa-krānti}). \marginnote{(5)}

Now, if the [first] \glslink{declination_degree}{declination} (\textit{krānti}) should not exist but the \glslink{latitude_celestial_object}{latitude} (\textit{śara}) should, then the \glslink{Sine_latitude}{Sine of the latitude} (\textit{bāṇa-jyā}) that
\glslink{multiplication}{must be multiplied} (\textit{saṃ-guṇyā}) by the \gls{Cosine_maximum_declination} (\textit{parama-krānti-koṭijyā}) [\ie by the Cosine of the ecliptic obliquity], \glslink{lowering}{should [be made] lower} (\textit{adhaḥ kuryāt}). Or, [one] \glslink{extract}{may [again] extract} (\textit{utthāpayet}) the \glslink{product_multiplication}{product of the multiplication} (\textit{guṇita-phala}) with the \glslink{Sine_latitude}{Sine of the latitude} (\textit{bāṇa-jyā}) from the \glslink{table_Cosine_maximum_declination}{tables of the Cosine of the maximum declination} (\textit{parama-krānti-koṭijyā-koṣṭhaka}s) [\ie from the tables of Cosine of the ecliptic obliquity]. [The result] should be the \gls{Sine_true_declination} (\textit{spaṣṭa-krānti-jyā}) in the \gls{direction_latitude} (\textit{bāṇa-dik}). \marginnote{(6)}

If the \glslink{declination_degree}{declination of a celestial object} (\textit{khagasya krānti}) should be equal to the \gls{maximum_declination} (\textit{parama-krānti}) [\ie the obliquity of the ecliptic], then the very \gls{curve_true_declination} (\textit{spa\-ṣṭa-krānti-aṅka}) becomes the \gls{true_declination} (\textit{sphuṭa-krānti}). \marginnote{(7)}%
\medskip

\textbf{Now, in another way.}
\medskip

And what \gls{circle} (\textit{vṛtta}) reaches both the \gls{ecliptic_pole} (\textit{kadamba}) and the \gls{celestial_pole} (\textit{viṣuvat}\,$_\text{\acrshort{regular}}$-\textit{dhruva}), that has been stated to be the \glslink{solstitial_colure}{[circle] belonging to the solstice} (\textit{āyana}[-\textit{vṛtta}]), and also, the \glslink{solstitial_colure}{[circle] passing through the four poles} (\textit{dhruva-catuṣka-yāta}[-\textit{vṛtta}]). And passing over a \gls{celestial_object} (\textit{nabhoga}) and the \gls{pair_equinoctial_points} (\textit{visuvat}\,$_\text{\acrshort{regular}}$-\textit{dvaya}), what [circle] is \gls{well_rounded} (\textit{su-vṛtta}), that the \gls{knower_spheres} (\textit{gola-vid}) should consider as
the \glslink{circle_congruent_ecliptic}{[circle] congruent to the ecliptic} (\textit{bhacakra-sadṛśa}[-\textit{vṛtta}]) by name. 1 \marginnote{(α)} 

What \gls{arc} (\textit{dhanus}) produced on the \gls{solstitial_colure} (\textit{āyana-vṛtta}) becomes situated in the  \gls{difference} (\textit{vivara}) between the \gls{celestial_equator} (\textit{viṣuva}\,$_\text{\acrshort{regular}}$-\textit{vṛtta}) and the \glslink{circle_congruent_ecliptic}{[circle] congruent to the ecliptic} (\textit{bhavṛtta-sadṛśa}[-\textit{vṛtta}]), that is the stated [arc of] \gls{maximum_true_declination} (\textit{para-sphuṭa-apama}) of the \gls{celestial_object}(\textit{dyucara}) just at that present moment. 2 \marginnote{(β)}

What \gls{arc} (\textit{dhanus}) produced on the \gls{solstitial_colure} (\textit{āyana-vṛtta}) becomes situated in the \gls{difference} (\textit{vivara}) between the \gls{ecliptic} (\textit{bhavana-cakra}) and the \glslink{circle_congruent_ecliptic}{[circle] congruent to the ecliptic} (\textit{bhacakra-sadṛkṣa}[-\textit{vṛtta}]), in this case, that is the declared [arc of] \gls{maximum_latitude} (\textit{para-iṣu}) when the \glslink{conjunction_equinox_node}{conjunction of the equinoctial point and the node of the orbit of the celestial object} (\textit{viṣuva}\,$_\text{\acrshort{regular}}$-\textit{pāta-yuga}) has been supposed. 3 \marginnote{(γ)}

What \gls{arc} (\textit{kodaṇḍa}) of the  \glslink{circle_congruent_ecliptic}{[circle] congruent to the ecliptic} (\textit{bhavṛtta-sadṛś}[-\textit{vṛtta}]) is between the \gls{equinoctial_point} (\textit{viṣuvat}\,$_\text{\acrshort{regular}}$) and the \gls{celestial_object} (\textit{nabhoga}), that [arc] should be known as the \gls{congruent_bhuja} (\textit{sadṛś-bhujā}); [and what is] between the \gls{celestial_object} (\textit{bha}) and the \gls{solstitial_colure} (\textit{āyana}[-\textit{vṛtta}]), [that should be known as] the \gls{congruent_koti} (\textit{sadṛś-koṭi}) [\ie complement of \textit{sadṛś-bhujā}.]~4~\marginnote{(δ)}
\bigskip

The \gls{Sine_koti_celestial_object} (\textit{khagasya koṭi-siñjinī}), \glslink{multiplication}{having been multiplied} (\textit{hatā}) by the \glslink{Cosine_latitude_celestial_object}{Cosine of its latitude} (\textit{sva-bāṇa-koṭijīvā}) [and] \glslink{lowering}{having been lowered} (\textit{adharī-kṛtā}), should be the \glslink{Sine_distance_celestial_object_solstitial_colure}{Sine of the congruent complementary arc} (\textit{sadrkṣa-koṭi-siñjinī}) [\ie Sine of the complement of the \textit{sadṛś-bhujā}]. 5 \marginnote{(8)}


When [the measure of] its \gls{arc} (\textit{dhanus}), \glslink{reducing_from_ninety}{having been reduced from ninety} [degrees] (\textit{navatitaś-cyuta}), is determined, [it] has the name \glslink{congruent_bhuja}{congruent arc} (\textit{sadṛśa-bāhu}). Or, what is the \glslink{Sine_latitude}{Sine of the latitude of a celestial object} (\textit{nabhoga-viśikhasya siñjinī}), \glslink{division}{having been divided} (\textit{bhājitā}) by the \gls{lowered_Sine_congruent_base} (\textit{adhara-sadrkṣa-dorjyā}), 6$\enskip$\ldots its \gls{arc} (\textit{dhanus}) should be [called] the \gls{maximum_latitude} (\textit{para-śara}) by name. \marginnote{(9)}

There is the \gls{sum} (\textit{saṃyuti}) or the \gls{difference} (\textit{viyuti}) of the two [quantities] known as the \gls{maximum_latitude} (\textit{para-iṣu}) and the \gls{maximum_declination} (\textit{parama-apama}) [\ie the obliquity of the ecliptic] with the  \glslink{latitude_celestial_object}{latitude} (\textit{bāṇa}) and the \gls{celestial_hemisphere} (\textit{gola}) [\ie the declination of the celestial object] in the \gls{same_different_directions} (\textit{sama-bhinna-diś}) respectively. 7$\enskip$That [result], being situated in the \gls{conjunction_disjunction_direction} (\textit{yuti-viyoga-dik}), becomes the \glslink{maximum_true_declination}{maximum true declination of a celestial object} (\textit{grahasya parama-sphuṭa-apama}). Thus, when [its measure is] \gls{greater} (\textit{adhika}) than \gls{ninety} [degrees] (\textit{abhra-nava}), [it is] \glslink{subtraction}{made to be subtracted} (\textit{viśodhita}) from a measure of \gls{one_hundred_eighty} [degrees] (\textit{kha-aṣṭa-bhū}). 8 \marginnote{(10)}

The \gls{Sine_maximum_true_declination} (\textit{para-sphuṭa-krānti-bhava-jyakā}), \glslink{multiplication}{having been multiplied} (\textit{hatā}) by the \gls{Sine_congruent_base} (\textit{sadṛkṣa-bāhu-jyakā}) [and] \glslink{lowering}{having been lowered} (\textit{adharī-kṛtā}), its \gls{arc} (\textit{cāpa}) becomes the \gls{true_declination} (\textit{sphuṭa-apama}). Its \gls{direction} (\textit{diś}) is the \gls{same} (\textit{sama}) as the \gls{conjunction_disjunction_direction} (\textit{saṃyoga-viyoga-dik}). 9 \marginnote{(11)}



%Back Matter
% \renewcommand{\thefootnote}{A\textsubscript{\arabic{footnote}}}
\renewcommand{\thefootnote}{\arabic{footnote}}
\setcounter{footnote}{0}
\setlength{\footnotesep}{\onelineskip} % for space between footnotes
%~~~~~~~~~~~~~~~~~~~~%
\renewcommand{\abstractname}{Acknowledgements}
\begin{abstract}
I thank the Berlin Center for the History of Knowledge hosted at the Max-Planck-Institut Für Wissenschaftsgeschichte in Berlin, the School of Mathematics \& Statistics at University of Canterbury in Christchurch (New Zealand), and the Institut for Tværkulturelle og Regionale Studier at Københavns Universitet for their generous support in hosting me at different stages of this work.

I thank Prof.\@ Kenneth Zysk (University of Copenhagen, Denmark)
for his insightful comments and detailed suggestions in preparing this work. Dr Hassan Amini (University of Tehran) and Dr David Buchta (Brown University) have been immensely helpful in analysing complex (and corrupted) Persian and Sanskrit grammatical forms in the text, and I thank them for their patience and perceptiveness. I also thank Prof.\@ Clemency Montelle, Prof.\@ Kim Plofker, Prof.\@ Benno van Dallen, 
Prof.\@ Erik Reenberg Sand, Dr Jacob Schmidt-Madsen, Dr Robert Middeke-Conlin, Dr Mònica Colominas Aparicio, Dr Flavio D'Abramo, Dr Toke Knudsen, and Razieh-Sadat Mousavi for sharing their time and expertise on several occasions in course of this study.

This document was built and typeset in \XeTeX\ for which I acknowledge the developers and maintainers of the Comprehensive \TeX\ Archive Network. 
\end{abstract}
\clearpage

\printbibliography
\clearpage
%~~~~~~~~~~~~~~~~~~~~%

% Apendices
\appendix
\renewcommand{\thesubsection}{A.\arabic{subsection}}
\phantomsection % Helps the TOC pick up and link to the Appendix properly
\addcontentsline{toc}{chapter}{Appendix: Persian and Sanskrit Verbs}
\OnehalfSpacing
%%%%%%%%%%%%%%%%%%%%%%%%%%%%%%%%%%%%%%%%

\section*{Appendix: Persian and Sanskrit Verbs}\label{Appendix_verbs}

The attested forms of Persian verbs from the \ZijiShahJahani\ Discourse~II.6 (\S~\ref{zijshahjahan_persian_english}) and of Sanskrit verbs from the \Siddhantasindhu\ Part~II.6 (\S~\ref{siddhantasindhu_sanskrit_english}) are listed below separately. At the end of each entry, \textbf{passage-markers} in square brackets point to its location in \S~\ref{zijshahjahan_persian_english} or \S~\ref{siddhantasindhu_sanskrit_english} accordingly.

%%%%%%%%%%%%%%%%%%%%%%%%%%%%%%%%%%%%%%%%
\subsection{Persian verbs in the \ZijiShahJahani, Discourse II.6} \label{persian_verbs}
%%%%%%%%%%%%%%%%%%%%%%%%%%%%%%%%%%%%%%%%

\begin{enumerate}
    \item \tfarsi{بودن} (\textit{budan}) `to be' \hfill \tfarsi{باش} (\textit{bāsh}); \tfarsi{بود} [\acrshort{present}-\acrshort{stem}]
(\textit{buvad})\footnote{\,\Vid\ Chapter XI in \textcite[251--268]{Lenepveu_Hotz} for a diachronic study of the third singular present form of the verb `to be' from the indicative \tfarsi{بود} (\textit{buvad}) to the subjunctive \tfarsi{باشد} (\textit{bāshad}) in the Persian verbal system (10\textsuperscript{th}--16\textsuperscript{th} \ce).}
\begin{itemize}
        \item \tfarsi{باشند} (\textit{bāshand})  \acrshort{present}-\acrshort{subjunctive}-\acrshort{plural}·\acrshort{third} `[they] should be' [\hyperlink{Ppass1}{1}, \hyperlink{Ppass10}{10}]
          \item \tfarsi{باشد} (\textit{bāshad})  \acrshort{present}-\acrshort{subjunctive}-\acrshort{singular}·\acrshort{third} `[he/she/it] should be' [\hyperlink{Ppass1}{1}, \hyperlink{Ppass3}{3}--\hyperlink{Ppass8}{8}, \hyperlink{Ppass10}{10}, \hyperlink{Ppass11}{11}]
          \item \tfarsi{نباشد} (\textit{nabāshad}) \acrshort{present}-\acrshort{subjunctive}-\acrshort{singular}·\acrshort{third} (\acrshort{negation}) `[he/she/it] should not be' [\hyperlink{Ppass5}{5}, \hyperlink{Ppass6}{6}]
          \item \tfarsi{بود} (\textit{buvad}) \acrshort{present}-\acrshort{indicative}-\acrshort{singular}·\acrshort{third} `[he/she/it] is' [\hyperlink{Ppass2}{2}, \hyperlink{Ppass9}{9}]
    \end{itemize}
    \item \tfarsi{خواندن} (\textit{khāndan}) `to recite' \hfill \tfarsi{خوان} (\textit{khān}) [\acrshort{present}-\acrshort{stem}]
    \begin{itemize}
        \item \tfarsi{خوانیم}  (\textit{khānīm}) \acrshort{present}-\acrshort{indicative}-\acrshort{plural}·\acrshort{first} `[we] call' [\hyperlink{Ppass1}{1}, \hyperlink{Ppass9}{9}]
    \end{itemize}
    \item \tfarsi{درآوردن} (\textit{dar āvardan}) `to remove/produce/extract\newline\phantom{a}\hfill \tfarsi{در آور} (\textit{dar āvar}) [\textsc{variant}: \tfarsi{در آر} (\textit{dar ār})] [\acrshort{present}-\acrshort{stem}]
    \begin{itemize}
        \item \tfarsi{درآورند} (\textit{dar āvarand}) [\textsc{variant}: \tfarsi{درآرند} (\textit{dar ārand})] \acrshort{present}-\acrshort{indicative}-\acrshort{plural}·\acrshort{third}\newline `[they] remove/produce/extract' [\hyperlink{Ppass4}{4}]
        \item \tfarsi{درآوریم} (\textit{dar avarīm}) [\textsc{variant}: \tfarsi{درآریم} (\textit{dar ārim})]
        \acrshort{present}-\acrshort{indicative}-\acrshort{plural}·\acrshort{first}\newline `[we] remove/produce/extract' [\hyperlink{Ppass6}{6}]
    \end{itemize}
    \item \tfarsi{کردن} (\textit{kardan}) `to do/make'\footnote{\,The action verb \tfarsi{کردن} (\textit{kardan}) is often used to construct compound verbs, \eg \tfarsi{جمع کردن} (\textit{jama}\Ayn\ \textit{kardan}) `to sum/add' or \tfarsi{ضرب کنیم} (\textit{ḍarb kardan}) `to multiply'\label{compound_action_verb_persian}} \hfill  \tfarsi{کن} (\textit{kun}) [\acrshort{present}-\acrshort{stem}]
    \begin{itemize}
        \item \tfarsi{کنیم} (\textit{kunīm}) \acrshort{present}-\acrshort{indicative}-\acrshort{plural}·\acrshort{first} `[we] do/make' [\hyperlink{Ppass1}{1}--\hyperlink{Ppass3}{3}, \hyperlink{Ppass6}{6}, \hyperlink{Ppass8}{8}--\hyperlink{Ppass11}{11}]
        \item \tfarsi{کنند} (\textit{kunand}) \acrshort{present}-\acrshort{indicative}-\acrshort{plural}·\acrshort{third} `[they] do/make' [\hyperlink{Ppass4}{4}]
    \end{itemize}
    \item \tfarsi{گرفتن} (\textit{giriftan}) `to take/grab' \hfill  \tfarsi{گیر} (\textit{gīr})[\acrshort{present}-\acrshort{stem}]
    \begin{itemize}
        \item \tfarsi{بگیریم} (\textit{bigīrīm}) \acrshort{present}-\acrshort{subjunctive}-\acrshort{plural}·\acrshort{first} `[we] should take' [\hyperlink{Ppass1}{1}, \hyperlink{Ppass9}{9}, \hyperlink{Ppass10}{10}]
    \end{itemize}
\end{enumerate}
\clearpage

%%%%%%%%%%%%%%%%%%%%%%%%%%%%%%%%%%%%%%%%
\subsection{Sanskrit verbal forms in the \Siddhantasindhu, Part II.6} \label{sanskrit_verbs}
%%%%%%%%%%%%%%%%%%%%%%%%%%%%%%%%%%%%%%%%

\begin{enumerate}
\item   \begin{enumerate}
        \item √\tsans{as} (√\textit{as}) \acrshort{class}$_\text{2}$  `to be/exist'
        \begin{itemize}
            \item \tsans{asti} (\textit{asti}) \acrshort{present}-\acrshort{indicative}-\acrshort{active}-\acrshort{singular}·\acrshort{third} `[he/she/it] is' [\hyperlink{Spass10}{10}]
            \item \tsans{sat} (\textit{sat}) \acrshort{present}-\acrshort{active}-\acrshort{participle} `being' [\hyperlink{SpassC}{γ}]        
            \item \tsans{syaat} (\textit{syāt}) \acrshort{optative}-\acrshort{active}-\acrshort{singular}·\acrshort{third} `[he/she/it] should be/exist' [\hyperlink{Spass3}{3}--\hyperlink{Spass7}{7}]
                \begin{itemize}
                \item \tsans{na syaat} (\textit{na syāt}) \acrshort{optative}-\acrshort{active}-\acrshort{singular}·\acrshort{third} (\acrshort{negation}) `[he/she/it] should not be/exist' [\hyperlink{Spass5}{5}, \hyperlink{Spass6}{6}]
                \end{itemize}
        \end{itemize}
        \item √\tsans{jan} (√\textit{jan}) \acrshort{class}$_\text{4}$  `to be born/determined'
        \begin{itemize}
            \item \tsans{jaayate} (\textit{jāyate}) \acrshort{present}-\acrshort{indicative}-\acrshort{middle}-\acrshort{singular}·\acrshort{third} `[he/she/it] comes to be'/`[he/she/it] is determined' [\hyperlink{Spass9}{9}, \hyperlink{Spass10}{10}]
       \end{itemize}
       \item √\tsans{bhuu} (√\textit{bhū}) \acrshort{class}$_\text{1}$  `to be/become'
        \begin{itemize}
            \item \tsans{bhavet} (\textit{bhavet}) \acrshort{optative}-\acrshort{active}-\acrshort{singular}·\acrshort{third} `[he/she/it] should/will be' [\hyperlink{Spass1}{1}, \hyperlink{Spass2}{2}, \hyperlink{Spass5}{5}, \hyperlink{Spass6}{6}, \hyperlink{Spass8}{8}, \hyperlink{Spass9}{9}]
            \item \tsans{bhavati} (\textit{bhavati}) \acrshort{present}-\acrshort{indicative}-\acrshort{active}-\acrshort{singular}·\acrshort{third} `[he/she/it] becomes' [\hyperlink{Spass7}{7}, \hyperlink{SpassB}{β}, \hyperlink{SpassC}{γ}, \hyperlink{Spass11}{11}]
        \end{itemize}
        \item √\tsans{sthaa} (√\textit{sthā}) \acrshort{class}$_\text{1}$  `to stand/situate'
        \begin{itemize}
            \item \tsans{sthita} (\textit{sthita})
            \acrshort{past}-\acrshort{passive}-\acrshort{participle} `being stationed/situated'(\acrshort{active}-sense) [\hyperlink{Spass1}{1}, \hyperlink{Spass10}{10}]
        \end{itemize}
        \end{enumerate}
\item   \begin{enumerate}
        \item √\tsans{udiir} (√\textit{udīr}) \acrshort{class}$_\text{2}$  `to state/utter'
        \begin{itemize}
            \item \tsans{udiirita} (\textit{udīrita})
            \acrshort{causative}-\acrshort{past}-\acrshort{passive}-\acrshort{participle} `has been stated to be' [\hyperlink{SpassA}{α}]
        \end{itemize}
        \item √\tsans{kath} (√\textit{kath}) \acrshort{class}$_\text{10}$  `to declare/tell' 
        \begin{itemize}
            \item \tsans{kathita} (\textit{kathita}) \acrshort{past}-\acrshort{passive}-\acrshort{participle}  `declared/told'  (\acrshort{adjective}-use) [\hyperlink{SpassB}{β}]
        \end{itemize}
        \item √\tsans{vac} (√\textit{vac}) \acrshort{class}$_\text{2}$  `to say'
        \begin{itemize}
            \item \tsans{ucyate} (\textit{ucyate}) \acrshort{present}-\acrshort{passive}-\acrshort{singular}·\acrshort{third} `[he/she/it] is said [to be]' [\hyperlink{Spass1}{1}]
        \end{itemize}
        \item √\tsans{vad} (√\textit{vad}) \acrshort{class}$_\text{1}$  `to declare/state'
        \begin{itemize}
            \item \tsans{udita} (\textit{udita}) \acrshort{past}-\acrshort{passive}-\acrshort{participle}  `declared/stated' (\acrshort{adjective}-use)  [\hyperlink{SpassD}{γ}]
        \end{itemize}
        \end{enumerate} 
\item \begin{enumerate}
        \item √\tsans{aap} (√\textit{āp}) \acrshort{class}$_\text{5}$ `to reach' (in arithmetic, `to divide')
        \begin{itemize}
            \item \tsans{aapta} (\textit{āpta}) \acrshort{past}-\acrshort{passive}-\acrshort{participle} `having been reached/divided' [\hyperlink{Spass3}{3}]
        \end{itemize}
        \item  √\tsans{bhaj} (√\textit{bhaj}) \acrshort{class}$_\text{1}$ `to divide'
        \begin{itemize}
            \item \tsans{bhajet} (\textit{bhajet}) \acrshort{optative}-\acrshort{active}-\acrshort{singular}·\acrshort{third} `[he/she/it] should divide' [\hyperlink{Spass4}{4}]
            \item \tsans{bhaajita} (\textit{bhājita}) \acrshort{causative}-\acrshort{past}-\acrshort{passive}-\acrshort{participle} `having been divided' (in a causal sense) [\hyperlink{Spass9}{9}]
        \end{itemize}
        \end{enumerate}        
\item \begin{enumerate}
        \item √\tsans{ni-han} (√\textit{ni-han}) \acrshort{class}$_\text{2}$  `to strike in' (in arithmetic, `to multiply')
        \begin{itemize}
            \item \tsans{ni-hanyate} (\textit{ni-hanyate}) \acrshort{present}-\acrshort{passive}-\acrshort{singular}·\acrshort{third} `[he/she/it] is struck/multiplied' [\hyperlink{Spass2}{2}]
            \item \tsans{hata} (\textit{hata}) \acrshort{past}-\acrshort{passive}-\acrshort{participle} `having been struck/multiplied' [\hyperlink{Spass3}{3}, \hyperlink{Spass8}{8}, \hyperlink{Spass11}{11}]
        \end{itemize}
        \item √\tsans{sa.m-gu.n} (√\textit{saṃ-guṇ}) \acrshort{class}$_\text{10}$  `to multiply' 
        \begin{itemize}
            \item \tsans{sa.m-gu.nya} (\textit{saṃ-guṇya}) \acrshort{gerundive} `to/must be multiplied' [\hyperlink{Spass6}{6}]
        \end{itemize}
        \end{enumerate}
\item \begin{enumerate}
    \item  √\tsans{cyu} (√\textit{cyu}) \acrshort{class}$_\text{1}$  `to deviate'/`be deprived of' (with \acrshort{ablative}-use)
    \begin{itemize}
        \item {\tsans{cyuta}} (\textit{cyuta}) \acrshort{past}-\acrshort{passive}-\acrshort{participle} `having deviated from'/`having been deprived of' [\hyperlink{Spass9}{9}]
    \end{itemize}
    \item √\tsans{vi-"sudh} (√\textit{vi-śudh}) \acrshort{class}$_\text{1}$  `to purify/subtract' 
    \begin{itemize}
    \item {\tsans{vi-"sodhita}} (\textit{vi-śodhita})  \acrshort{causative}-\acrshort{past}-\acrshort{passive}-\acrshort{participle} `made to be purified/subtracted' [\hyperlink{Spass10}{10}]
    \end{itemize}
    \end{enumerate}        
\item √\tsans{upe} (√\textit{upe}) \acrshort{class}$_\text{2}$ `to reach'
    \begin{itemize}
        \item \tsans{upaiti} (\textit{upaiti}) \acrshort{present}-\acrshort{indicative}-\acrshort{singular}·\acrshort{third} `[he/she/it] reaches' [\hyperlink{SpassA}{α}]
    \end{itemize}
\item √\tsans{utthaa} (√\textit{utthā}) \acrshort{class}$_\text{1}$ `to rise/extract'
    \begin{itemize}
        % \item \tsans{uttha} (\textit{uttha}) \acrshort{primary_nominal_derivative} `derived/extracted' (\acrshort{adjective}-use) 
        \item \tsans{utthaapayet} (\textit{uthayet}) \acrshort{causative}-\acrshort{optative}-\acrshort{active}-\acrshort{singular}·\acrshort{third} `[he/she/it] may raise/extract' [\hyperlink{Spass6}{6}]
        \item \tsans{utthaaya} (\textit{utthāya}) \acrshort{gerund} `having risen/extracted' [\hyperlink{Spass4}{4}]
    \end{itemize}    
\item √\tsans{k.r} (√\textit{kṛ}) \acrshort{class}$_\text{8}$  `to do/make'\footnote{\,The action verb √\tsans{k.r} (√\textit{kr}) is often compounded with an inflected nominal word to form a denominative verb (\acrshort{denominative}, \textit{nāmadhātu}), \eg √\tsans{adharii-k.r} (√\textit{adharī-kṛ}) `to make [something] low', \ie `to lower'.}
        \begin{itemize}
        \item \tsans{kuryaat} (\textit{kuryāt}) \acrshort{optative}-\acrshort{active}-\acrshort{singular}·\acrshort{third} `[he/she/it] should do/make' [\hyperlink{Spass6}{6}]
        \item \tsans{k.rta} (\textit{kṛta}) \acrshort{past}-\acrshort{passive}-\acrshort{participle} `having been done/made' [\hyperlink{Spass2}{2}, \hyperlink{Spass8}{8}, \hyperlink{Spass11}{11}]
        \end{itemize}
\item √\tsans{k.lp} (√\textit{kḷp}) \acrshort{class}$_\text{1}$  `to consider/suppose'
        \begin{itemize}
        \item \tsans{kalpayet} (\textit{kalpayet}) \acrshort{causative}-\acrshort{optative}-\acrshort{active}-\acrshort{singular}·\acrshort{third} `[he/she/it]' should consider/suppose' [\hyperlink{SpassA}{α}]
        \item \tsans{kalpita} (\textit{kalpita})  \acrshort{causative}-\acrshort{past}-\acrshort{passive}-\acrshort{participle} `has been considered/supposed' [\hyperlink{SpassC}{γ}]
        \end{itemize}        
\item √\tsans{pat} (√\textit{pat}) \acrshort{class}$_\text{1}$ `to fall/pass'
    \begin{itemize}
        \item \tsans{patat} (\textit{patat}) \acrshort{present}-
        \acrshort{active}-\acrshort{participle} `falling/passing' [\hyperlink{SpassA}{α}]
    \end{itemize}
\item √{\tsnb{ज्ञा}} (√\textit{jñā}) \acrshort{class}$_\text{9}$  `to know/understand'
    \begin{itemize}
    \item {\tsnb{ज्ञेय}} (\textit{jñeya}) \acrshort{gerundive} `to be known'/`to be understood' [\hyperlink{SpassD}{δ}]
    \end{itemize}
\end{enumerate}
\clearpage

% List of Grammatical Abbreviations
\printglossary[type=\acronymtype, style= mcolindex, title={List of Grammatical Abbreviations}, nonumberlist]
\clearpage

%~~~~~~~~~~~~~~~~~~~~%
% Glossary
\renewcommand{\thefootnote}{\arabic{footnote}}
\setcounter{footnote}{0}
\setlength{\footnotesep}{0.6\baselineskip} % for space between footnotes
% (The tabular line separation in the glossary also uses \\[0.6\baselineskip])


\renewcommand*\glspostdescription{\vspace{\glueexpr0.4\baselineskip}} % Sets the spacing between entries, otherwise the index-style glossary is really compressed. 

\renewcommand\glstreepredesc{\quad} % space between name and description for item
\renewcommand{\glstreechildpredesc}{\quad} % space between name and description for subsequent subitem, subsubitem, etc

\renewcommand{\glstreeitem}{\parindent0pt\par\hangindent0em  
 \everypar{\parindent0pt\hangindent2em}}% Sets the indentation of the description paragraph (item) for the first line of item being \quad separated, and further for the the second line.
 
\renewcommand{\glstreesubitem}{\parindent1em\par\hangindent1em  
 \everypar{\parindent1em\hangindent4em}}% Sets the indentation of the description paragraph (subitem) for the first line of subitem being \quad separated from item, and accordingly for the the second line.
 
\printglossary[style=index, nonumberlist]
%~~~~~~~~~~~~~~~~~~~~%
\end{document}
